% Complete English translation of SGA 4 1/2.
% All book sources, local style definitions, and bibliography entries
% are inlined in this single compilable TeX file.

\documentclass[12pt,oneside,english]{smfbook}
% BEGIN INLINED SGA STYLE
\usepackage[T1]{fontenc}
\usepackage[utf8]{inputenc} % allows French characters
\usepackage[english]{babel}
\usepackage[all]{xy}        % commutative diagrams

\usepackage[margin=2.5cm]{geometry}

\usepackage{
  amsmath,
  amssymb,
  amsthm,
  bm,        % to get bold symbols to work with stmaryrd
 % calligra,  % for sheaf hom
  enumerate, % allows \begin{enumerate}[\indent a)]
  mathrsfs   % allows \mathscr
%  microtype, % improves general appearance
}
\usepackage{newtxtext,newtxmath} % match the text and mathematics fonts of Ng\^o's paper
%\usepackage[T1]{fontenc} % with calligra, but causes everything to crash
\usepackage[kerning=true]{microtype} % improves general appearance
\usepackage[hidelinks,pdfpagelabels,bookmarks=false,hypertexnames=false]{hyperref}
%\setcounter{secnumdepth}{6}

% Use the English SMF front matter and structural names, as in Ng\^o's paper.
\AtBeginDocument{%
  \renewcommand*{\smfbyname}{by}%
  \renewcommand*{\contentsname}{Contents}%
  \renewcommand*{\appendixname}{Appendix}%
  \renewcommand*{\proofname}{Proof}%
  \renewcommand*{\bibname}{References}%
  \renewcommand*{\figurename}{Figure}%
  \renewcommand*{\tablename}{Table}%
}

% operators
  \DeclareMathOperator{\aut}{Aut}
  \DeclareMathOperator{\br}{Br}
  \DeclareMathOperator{\cl}{cl}
  \DeclareMathOperator{\cosp}{cosp}
  \DeclareMathOperator{\dv}{div}
  \DeclareMathOperator{\Div}{Div}
  \DeclareMathOperator{\ext}{Ext}
  \DeclareMathOperator{\Ext}{\mathcal{E}\!xt} % sheaf ext
  \DeclareMathOperator{\gal}{Gal}
  \DeclareMathOperator{\gl}{GL}
  \DeclareMathOperator{\gr}{gr}
  \DeclareMathOperator{\h}{H}
  \DeclareMathOperator{\hh}{\mathsf{H}} % hypercohomology?
  \DeclareMathOperator{\Hom}{\mathcal{H}\!om} % sheaf hom
  \DeclareMathOperator{\im}{im}
  \DeclareMathOperator{\ob}{Ob}
  \DeclareMathOperator{\pgl}{PGL}
  \DeclareMathOperator{\pic}{Pic}
  \DeclareMathOperator{\res}{R} % restriction of scalars
  \DeclareMathOperator{\rHom}{\mathsf{R}\mathcal{H}\!om}
  \DeclareMathOperator{\sh}{Sh}
  \DeclareMathOperator{\spe}{sp} % specialisation
  \DeclareMathOperator{\spec}{Spec}
  \DeclareMathOperator{\swan}{Sw}
  \DeclareMathOperator{\tor}{Tor}
  \DeclareMathOperator{\Tor}{\mathcal{T}\!or} % sheaf tor
  \DeclareMathOperator{\tr}{tr}

% special symbols: a new way of defining them!
  \newcommand{\dmu}{{\bm\mu}}
  \newcounter{char}
  \setcounter{char}{1}
  \loop\ifnum\value{char}<27
    \expandafter\edef\csname c\Alph{char}\endcsname{\noexpand\mathcal{\Alph{char}}}  % mathcal
    \expandafter\edef\csname d\Alph{char}\endcsname{\noexpand\mathbb{\Alph{char}}}   % mathbb
    \expandafter\edef\csname f\Alph{char}\endcsname{\noexpand\mathfrak{\Alph{char}}} % mathfrak
    \expandafter\edef\csname e\Alph{char}\endcsname{\noexpand\mathsf{\Alph{char}}}   % mathsf
    \expandafter\edef\csname s\Alph{char}\endcsname{\noexpand\mathscr{\Alph{char}}}  % mathscr
    \addtocounter{char}{1}
  \repeat
% can't do the previous for lower-case mathfrak because \fi is already defined
% (otherwise, could do \alph instead of \Alph above)
  \def\mydeff#1{\expandafter\def\csname f#1\endcsname{\mathfrak{#1}}}
  \def\mydefallf#1{\ifx#1\mydefallf\else\mydeff#1\expandafter\mydefallf\fi}
  \mydefallf abcdefghjklmnopqrstuvwxyz\mydefallf

% miscellaneous operators
  \providecommand{\llbracket}{\mathopen{[\![}}
  \providecommand{\rrbracket}{\mathclose{]\!]}}
  \newcommand{\an}[1]{{#1}^{\textnormal{an}}}
  \newcommand{\const}[1]{\underline{#1}}
  \newcommand{\et}[1]{{#1}_{\textnormal{et}}}
  \newcommand{\iso}{\xrightarrow\sim}
  \newcommand{\lotimes}{{\overset{\mathsf{L}}{\otimes}}}
  \DeclareMathOperator{\rhom}{\mathsf{R}\hom}
  \renewcommand{\setminus}{\smallsetminus}
  \newcommand{\kloos}{\textnormal{K}} % Kloosterman sum
  \newcommand{\pr}{\mathrm{pr}} % projection map

% theorem types
  \newtheorem{proposition}[subsubsection]{Proposition}
  \newtheorem{corollary}[subsubsection]{Corollary}
  \newtheorem{definition}[subsubsection]{Definition}
  \newtheorem{theorem}[subsubsection]{Theorem}
  \newtheorem{lemma}[subsubsection]{Lemma}

  \newtheorem{lemma_}[subsection]{Lemma}
  \newtheorem{corollary_}[subsection]{Corollary}
  \newtheorem{theorem_}[subsection]{Theorem}
  \newtheorem{proposition_}[subsection]{Proposition}
  \newtheorem{definition_}[subsection]{Definition}
  \newtheorem{prop-def_}[subsection]{Proposition--Definition}
  \setcounter{tocdepth}{1}

  \newtheorem*{theorem*}{Theorem}

%\renewcommand*\thechapter{\Roman{chapter}}
\renewcommand*\thesection{\arabic{section}}
\renewcommand{\theequation}{\arabic{equation}}
\setcounter{secnumdepth}{3}

% allows for large external tensor product symbol
\newcommand{\bigboxtimes}{\mathop{\vphantom{\sum}\mathchoice
  {\vcenter{\hbox{\huge $\boxtimes$}}}
  {\vcenter{\hbox{\Large $\boxtimes$}}}{\boxtimes}{\boxtimes}}\displaylimits}

%\renewcommand{\bullet}{\text{\,\begin{picture}(-1,1)(-1,-2)\circle*{2}\end{picture}\ }}
\let\oldbullet\bullet
\renewcommand{\bullet}{\textnormal{\tiny$\oldbullet$}}

% unmarked footnote
\newcommand\blfootnote[1]{%
  \begingroup
  \renewcommand\thefootnote{}%
  \footnotetext{#1}%
  \endgroup
}
% END INLINED SGA STYLE

\begin{document}
\frontmatter
\pagestyle{plain}
\begin{titlepage}
\thispagestyle{empty}
\centering
\vspace*{0.24\textheight}
{\LARGE\bfseries Étale Cohomology\par}
\vspace{20pt}
{\Large (SGA $4\frac 1 2$)\par}
\vspace{28pt}
{\large P. Deligne\par}
\vspace{12pt}
{\normalsize with the collaboration of J.F. Boutot,\par
A. Grothendieck, L. Illusie, and J.L. Verdier\par}
\vfill
{\normalsize 1977\par}
\end{titlepage}

\clearpage
% BEGIN INLINED FILE: sga4.5-note.tex
\section*{Typesetter's Note}

This is a \LaTeX{} rendition of Deligne's \emph{Cohomologie \'Etale}, Lecture Notes in Mathematics, 569, Springer-Verlag. The typesetting was done by Daniel Miller. The source code may be found at GitHub (\url{https://www.github.com/dkmiller/sga4.5}). It is an essentially unmodified version of the original. Most changes are slight, e.g. using script instead of roman letters for sheaves. The biggest is that $\Hom$ is used to denote ``sheaf hom,'' and $\rHom$ its derived functor, while $\rhom$ is used to denote the derived functor of the regular hom-functor. Any comments or corrections should be sent to \url{dkmiller@math.cornell.edu}.

\medskip
\noindent\textit{Note on this English edition.}
The French prose created by Dan Miller has been translated into English while preserving the mathematical notation, numbering, labels, and bibliographic titles. A small number of evident transcription errors and incorrect internal references in the source have been corrected.
% END INLINED FILE: sga4.5-note.tex
\clearpage
\clearpage
% BEGIN INLINED FILE: sga4.5-intro.tex
\chapter*{Introduction}

This volume is intended to make $\ell$-adic cohomology easier to use for the nonexpert. I hope that it will often enable the reader to avoid consulting the dense Exposés of SGA 4 and SGA 5. It also contains several new results.

The first Exposé, written by J.F. Boutot, surveys SGA 4. It gives the main results---at a minimal level of generality, often insufficient for applications---and an indication of their proofs. For complete statements or detailed proofs, SGA 4 remains indispensable.

The ``\nameref{II}'' contains a complete proof of the trace formula for the Frobenius endomorphism. It is Grothendieck's proof from SGA 5, stripped of all unnecessary detail. This report should allow the reader to dispense with SGA 5, which may then be regarded as a series of digressions, some of them very interesting. Its existence will make it possible to publish SGA 5 shortly in its present form. The report is supplemented by the Exposé ``\nameref{VI},'' which explains how the trace formula can be used to study trigonometric sums and gives examples.

The remaining Exposés are intended for a more restricted audience, and their style reflects this. The Exposé ``\nameref{III}'' is a ``modular'' generalization of the report, based on the study of symmetric powers in SGA 4 XVII 5.5. The Exposé ``\nameref{IV}'' defines the relevant classes in various settings and proves the compatibility of intersections with cup products. The Exposé ``\nameref{V}'' collects several known results for which no reference was available, together with several compatibility statements. The Exposé ``\nameref{VII}'' is new. In particular, it proves finiteness theorems for cohomology without supports analogous to those already known for cohomology with compact supports.

For further details on the Exposés, I refer to their respective introductions.

Finally, I thank J.L. Verdier for allowing me to reproduce his notes ``\nameref{VIII}'' here. I believe that they remain very useful, and they had become unobtainable.

In internal references within this volume, the Exposés are cited by the abbreviated titles given in square brackets in the table of contents.

\vspace{20pt}
Bures-sur-Yvette, September 20, 1976

\vspace{5pt}
Pierre Deligne
% END INLINED FILE: sga4.5-intro.tex
\clearpage
\tableofcontents

\mainmatter
% BEGIN INLINED FILE: sga4.5-1.tex
\chapter{Étale Cohomology: Starting Points}\label{I}
\blfootnote{by P.\ Deligne, notes by J.F.\ Boutot}

\section{Grothendieck Topologies}\label{I:1}

Grothendieck topologies first arose in connection with his theory of descent (cf. SGA 1 VI, VIII); the corresponding cohomology theories came into use later. We follow the same course here. By formalizing the classical notions of localization, local property, and gluing (\ref{I:1-1}, \ref{I:1-2}, \ref{I:1-3}), we arrive at the general concept of a Grothendieck topology (\ref{I:1-6}). To justify its introduction into algebraic geometry, we prove a faithfully flat descent theorem (\ref{I:1-4}), generalizing the classical theorem of Hilbert (\ref{I:1-5}).

The reader will find a more complete, though concise, account of the formalism in Giraud \cite{gi64}. M. Artin's notes, ``Grothendieck Topologies'' \cite{ar62}, Chapters I--III, also remain useful. The 866 pages comprising Exposés I--V of SGA 4 are valuable when one considers exotic topologies, such as the one giving rise to crystalline cohomology. For the use of the étale topology, which is so close to classical intuition, reading them is not indispensable.

\subsection{Sieves}\label{I:1-1}

Let $X$ be a topological space and let $f:X\to \dR$ be a real-valued function on $X$. Continuity of $f$ is a local property: if $f$ is continuous on every sufficiently small open subset of $X$, then it is continuous on all of $X$. To formalize the notion of a ``local property,'' we introduce several definitions.

A collection $\cU$ of open subsets of $X$ is called a \emph{sieve} if $U\in\cU$ and $V\subset U$ imply $V\in\cU$. A sieve is called \emph{covering} if the union of all its members is $X$.

Given a family $\{U_i\}$ of open subsets of $X$, the sieve generated by $\{U_i\}$ is, by definition, the set of open subsets $U$ of $X$ that are contained in one of the $U_i$.

A property $P(U)$ defined for every open subset $U$ of $X$ is called \emph{local} if, for every covering sieve $\cU$ of every open subset $U$ of $X$, the property $P(U)$ holds if and only if $P(V)$ holds for every $V\in\cU$. For example, given $f:X\to\dR$, the property ``$f$ is continuous on $U$'' is local.

\subsection{Sheaves}\label{I:1-2}

We now make precise the notion of a function given locally on $X$.

\subsubsection{The sieve viewpoint}\label{I:1-2-1}

Let $\cU$ be a sieve of open subsets of $X$. A function given $\cU$-locally on $X$ consists of a function $f_U$ on each $U\in\cU$ such that $f_V=f_U|V$ whenever $V\subset U$.

\subsubsection{The Čech viewpoint}\label{I:1-2-2}

If the sieve $\cU$ is generated by a family of open subsets $U_i$ of $X$, giving a function $\cU$-locally amounts to giving a function $f_i$ on each $U_i$ such that $f_i|{U_i\cap U_j}=f_j|{U_i\cap U_j}$.

Equivalently, if $Z=\coprod U_i$, giving a function $\cU$-locally amounts to giving a function on $Z$ that is constant on the fibers of the natural projection $Z\to X$.

\subsubsection{}\label{I:1-2-3}

Continuous functions form a sheaf. This means that, for every covering sieve $\cU$ of an open subset $V$ of $X$ and every function $\{f_U\}$ given $\cU$-locally such that each $f_U$ is continuous on $U$, there is a unique continuous function $f$ on $V$ satisfying $f|U=f_U$ for every $U\in\cU$.

\subsection{Stacks}\label{I:1-3}

We next make precise the notion of a vector bundle given locally on $X$.

\subsubsection{The sieve viewpoint}\label{I:1-3-1}

Let $\cU$ be a sieve of open subsets of $X$. A vector bundle given $\cU$-locally on $X$ consists of
\begin{enumerate}[\indent a)]
  \item a vector bundle $E_U$ on each $U\in\cU$;
  \item for $V\subset U$, an isomorphism $\rho_{U,V}:E_V\iso E_U|V$;
  \item for $W\subset V\subset U$, the commutativity of the diagram
    \[\xymatrix{
      E_W \ar[r]^-{\rho_{U,W}} \ar[dr]_-{\rho_{V,W}}
        & E_U|W \\
      & E_V|W \ar[u]_-{\rho_{U,V}|W}
    }\]
    that is, $\rho_{U,W}=(\rho_{U,V}|W)\circ\rho_{V,W}$.
\end{enumerate}

\subsubsection{The Čech viewpoint}\label{I:1-3-2}

If the sieve $\cU$ is generated by a family of open subsets $U_i$ of $X$, giving a vector bundle $\cU$-locally amounts to giving
\begin{enumerate}[\indent a)]
  \item a vector bundle $E_i$ on each $U_i$;
  \item for $U_{ij}=U_i\cap U_j=U_i\times_XU_j$, an isomorphism $\rho_{ji}:E_i|U_{ij}\iso E_j|U_{ij}$;
  \item for $U_{ijk}=U_i\times_XU_j\times_XU_k$, the commutativity of the diagram
    \[\xymatrix{
      E_i |U_{ijk} \ar[r]^-{\rho_{ki}|U_{ijk}} \ar[dr]_-{\rho_{ji}|U_{ijk}}
        & E_k|U_{ijk} \\
      & E_j|U_{ijk} \ar[u]_-{\rho_{kj}|U_{ijk}}
    }\]
    that is, $\rho_{ki}=\rho_{kj}\circ\rho_{ji}$ on $U_{ijk}$.
\end{enumerate}

Equivalently, if $Z=\coprod U_i$ and $\pi:Z\to X$ is the natural projection, giving a vector bundle $\cU$-locally amounts to giving
\begin{enumerate}[\indent a)]
  \item a vector bundle $E$ on $Z$;
  \item whenever $x$ and $y$ are points of $Z$ with $\pi(x)=\pi(y)$, an isomorphism $\rho_{yx}:E_x\iso E_y$ between the fibers of $E$ at $x$ and $y$, depending continuously on $(x,y)$;
  \item whenever $x,y,z$ are points of $Z$ with $\pi(x)=\pi(y)=\pi(z)$, the equality $\rho_{zx}=\rho_{zy}\circ\rho_{yx}$.
\end{enumerate}

\subsubsection{}\label{I:1-3-3}

A vector bundle $E$ on $X$ defines a vector bundle $E_\cU$ given $\cU$-locally, namely the system of restrictions $E_U$ of $E$ to the members of $\cU$. The fact that the notion of vector bundle is local can be expressed as follows: for every covering sieve $\cU$ of $X$, the functor $E\mapsto E_\cU$ from vector bundles on $X$ to vector bundles given $\cU$-locally is an equivalence of categories.

\subsubsection{}\label{I:1-3-4}

If in \ref{I:1} one replaces ``open subset of $X$'' by ``subset of $X$,'' one obtains the notion of a sieve of subspaces of $X$. Gluing theorems are also available in this setting. For example, let $X$ be a space and let $\cC$ be a sieve of subspaces of $X$ generated by a locally finite closed covering of $X$. Then the functor $E\mapsto E_\cC$ from vector bundles on $X$ to vector bundles given $\cC$-locally is an equivalence of categories.

In algebraic geometry it is also useful to consider ``sieves of spaces over $X$''; this is the subject of the next subsection.

\subsection{Faithfully flat descent}\label{I:1-4}

\subsubsection{}\label{I:1-4-1}

For schemes, the Zariski topology is not fine enough for the study of nonlinear problems, and in the preceding definitions one is led to replace open immersions by more general morphisms. From this viewpoint, descent techniques appear as localization techniques. Thus the following descent statement says that the properties under consideration are local for the faithfully flat topology. Recall that a morphism of schemes is faithfully flat if it is flat and surjective.

\begin{proposition}\label{I:1-4-2}
Let $A$ be a ring and let $B$ be a faithfully flat $A$-algebra. Then:
\begin{enumerate}[(i)]
  \item A sequence $\Sigma=(M'\to M\to M'')$ of $A$-modules is exact whenever the sequence $\Sigma_{(B)}$ obtained from it by extension of scalars to $B$ is exact.
  \item An $A$-module $M$ is finitely generated (respectively, finitely presented, flat, locally free of finite rank, or invertible, that is, locally free of rank one) whenever the $B$-module $M_{(B)}$ has the corresponding property.
\end{enumerate}
\end{proposition}
\begin{proof}
(i) Since the functor $M\mapsto M_{(B)}$ is exact by flatness of $B$, it is enough to show that $N_{(B)}\ne0$ whenever an $A$-module $N$ is nonzero. Such an $N$ contains a nonzero cyclic submodule $A/\fa$. Consequently $N_{(B)}$ contains the cyclic submodule $(A/\fa)_{(B)}=B/\fa B$, which is nonzero because the structural morphism $\varphi:\spec(B)\to\spec(A)$ is surjective: if $V(\fa)$ is nonempty, then $\varphi^{-1}(V(\fa))=V(\fa B)$ is nonempty.

(ii) For every family $(x_i)$ of elements of $M_{(B)}$, there is a finitely generated submodule $M'$ of $M$ such that $M'_{(B)}$ contains all the $x_i$. If $M_{(B)}$ is finitely generated and the $x_i$ generate it, then $M'_{(B)}=M_{(B)}$. Hence $M'=M$ by (i), and $M$ is finitely generated.
\end{proof}

If $M_{(B)}$ is finitely presented, the preceding argument provides a surjection $A^n\to M$. If $N$ is its kernel, then the $B$-module $N_{(B)}$ is finitely generated, and hence so is $N$; thus $M$ is finitely presented. The assertion for flatness follows immediately from (i). To be locally free of finite rank is to be flat and finitely presented, and the rank can be tested after extension of scalars to fields.

\subsubsection{}\label{I:1-4-3}

Let $X$ be a scheme and let $\cS$ be a class of $X$-schemes stable under fiber products over $X$. A subclass $\cU\subset\cS$ is a \emph{sieve} on $X$ relative to $\cS$ if, for every morphism $\varphi:V\to U$ of $X$-schemes with $U,V\in\cS$ and $U\in\cU$, one has $V\in\cU$. The sieve \emph{generated} by a family $\{U_i\}$ of $X$-schemes in $\cS$ is the class of all $V\in\cS$ for which there exists an $X$-morphism from $V$ to one of the $U_i$.

\subsubsection{}\label{I:1-4-4}

Let $\cU$ be a sieve on $X$. A quasi-coherent module given $\cU$-locally on $X$ consists of
\begin{enumerate}[\indent a)]
  \item a quasi-coherent module $E_U$ on each $U\in\cU$;
  \item for every $U\in\cU$ and every morphism $\varphi:V\to U$ of $X$-schemes in $\cS$, an isomorphism $\rho_\varphi:E_V\iso\varphi^*E_U$;
  \item for every morphism $\psi:W\to V$ of $X$-schemes in $\cS$, the commutativity of the diagram
    \[\xymatrix{
      E_W \ar[r]^-{\rho_{\varphi\circ\psi}} \ar[dr]_-{\rho_\psi}
        & \psi^*\varphi^* E_U \\
      & \psi^* E_V \ar[u]_-{\psi^*\rho_\varphi}
    }\]
    that is, $\rho_{\varphi\circ\psi}=(\psi^*\rho_\varphi)\circ\rho_\psi$.
\end{enumerate}

If $E$ is a quasi-coherent module on $X$, we denote by $E_\cU$ the module given $\cU$-locally whose value on $\varphi_U:U\to X$ is $\varphi_U^*E$, and for which the restriction isomorphism $\rho_\psi$ attached to every morphism $\psi:V\to U$ is the canonical isomorphism $E_V=(\varphi_U\circ\psi)^*E\iso\psi^*\varphi_U^*E=\psi^*E_U$.

\begin{theorem}\label{I:1-4-5}
Let $\{U_i\}\subset\cS$ be a finite family of $X$-schemes flat over $X$ such that $X$ is the union of the images of the $U_i$, and let $\cU$ be the sieve generated by $\{U_i\}$. Then the functor $E\mapsto E_\cU$ is an equivalence from the category of quasi-coherent modules on $X$ to the category of quasi-coherent modules given $\cU$-locally.
\end{theorem}
\begin{proof}
We shall treat only the case in which $X$ is affine and $\cU$ is generated by an affine $X$-scheme $U$ faithfully flat over $X$. The reduction to this case is formal. Put $X=\spec(A)$ and $U=\spec(B)$.

If the morphism $U\to X$ has a section, then $X$ belongs to the sieve $\cU$ and the assertion is clear. We shall reduce to this case.

A quasi-coherent module given $\cU$-locally determines modules $M'$, $M''$, and $M'''$ on $U$, $U\times_X U$, and $U\times_X U\times_X U$, together with isomorphisms $\rho:p^*M^\bullet\simeq M^\bullet$ for every projection morphism $p$ between these spaces; this is a \emph{Cartesian diagram}
\[\xymatrix{
  M^* : M' \ar@<2pt>[r] \ar@<-2pt>[r]
    & M'' \ar[r] \ar@<4pt>[r] \ar@<-4pt>[r]
    & M'''
}\]
over
\[\xymatrix{
  U_* : U
    & U\times_X U \ar@<2pt>[l] \ar@<-2pt>[l]
    & U\times_X U\times_X U \ar[l] \ar@<4pt>[l] \ar@<-4pt>[l] \mbox{.}
}\]
Conversely, $M^*$ determines the module given $\cU$-locally: for $V\in\cU$, there is a $\varphi:V\to U$, and one sets $M_V=\varphi^*M'$; for $\varphi_1,\varphi_2:V\to U$, there is a natural identification $\varphi_1^*M'\simeq(\varphi_1\times\varphi_2)^*M''\simeq\varphi_2^*M'$. Using $M'''$, one sees that these identifications are compatible, so the definition is legitimate. In short, giving a module $\cU$-locally is the same as giving a Cartesian diagram $M^*$ over $U_*$.

In algebraic terms, giving $M^*$ amounts to giving a Cartesian diagram of modules
\[\xymatrix{
  M' \ar@<3pt>[r]|-{\partial_0} \ar@<-3pt>[r]|-{\partial_1}
    & M'' \ar@<6pt>[r]|-{\partial_0} \ar[r]|-{\partial_1} \ar@<-6pt>[r]|-{\partial_2}
    & M'''
}\]
over the diagram of rings
\[\xymatrix{
  B \ar@<3pt>[r]|-{\partial_0} \ar@<-3pt>[r]|-{\partial_1}
    & B\otimes_A B \ar@<6pt>[r]|-{\partial_0} \ar[r]|-{\partial_1} \ar@<-6pt>[r]|-{\partial_2}
    & B\otimes_A B\otimes_A B
}\]
(more precisely, $\partial_i(bm)=\partial_i(b)\cdot\partial_i(m)$, the usual identities such as $\partial_0\partial_1=\partial_0\partial_0$ hold, and ``Cartesian'' means that the morphisms $\partial_i:M'\otimes_{B,\partial_i}(B\otimes_A B)\to M''$ and $M''\otimes_{B\otimes_A B,\partial_i}(B\otimes_A B\otimes_A B)\to M'''$ are isomorphisms).

The functor $E\mapsto E_\cU$ becomes the functor that assigns to an $A$-module $M$ the diagram
\[\xymatrix{
  M^* = (M\otimes_A B \ar@<2pt>[r] \ar@<-2pt>[r]
    & M\otimes_A B\otimes_A B \ar@<-4pt>[r] \ar[r] \ar@<4pt>[r]
    & M\otimes_A B\otimes_A B\otimes_A B)\text{.}
}\]
Its right adjoint is the functor
\[\xymatrix{
  (M' \ar@<2pt>[r] \ar@<-2pt>[r]
    & M'' \ar@<-4pt>[r] \ar[r] \ar@<4pt>[r]
    & M''') \ar@{|->}[r]
    & \ker(M' \ar@<2pt>[r] \ar@<-2pt>[r]
    & M''')\text{.}
}\]
We must prove that the adjunction morphisms
\[
  M \to\ker(M\otimes_A B\rightrightarrows M\otimes_A B\otimes_A B)
\]
et
\[
  \ker(M'\rightrightarrows M'')\otimes_A B \to M'
\]
are isomorphisms. By (\ref{I:1-4-2}.i), it is enough to prove this after a faithfully flat base change $A\to A'$ (with $B$ replaced by $B'=B\otimes_A A'$). Taking $A'=B$ reduces us to the case in which $U\to X$ has a section.
\end{proof}

\subsection{A special case: Hilbert's Theorem 90}\label{I:1-5}

\subsubsection{}\label{I:1-5-1}

Let $k$ be a field, let $k'$ be a Galois extension of $k$, and put $G=\gal(k'/k)$. Then the homomorphism
\begin{align*}
  k'\otimes_k k' &\to \oplus_{\sigma\in G} k' \\
  x\otimes y     &\mapsto \{x\cdot \sigma(y)\}_{\sigma\in G}
\end{align*}
is bijective.

It follows that giving a module locally for the sieve generated by $\spec(k')$ over $\spec(k)$ is the same as giving a $k'$-vector space endowed with a semilinear action of $G$, that is:
\begin{enumerate}[\indent a)]
  \item a $k'$-vector space $V'$;
  \item for every $\sigma\in G$, an endomorphism $\varphi_\sigma$ of the underlying group of $V'$ such that $\varphi_\sigma(\lambda v)=\sigma(\lambda)\varphi_\sigma(v)$ for every $\lambda\in k'$ and $v\in V'$, subject to the condition
  \item for every $\sigma,\tau\in G$, one has $\varphi_{\tau\sigma}=\varphi_\tau\circ\varphi_\sigma$.
\end{enumerate}

Let $V={V'}^G$ be the group of invariants under this action of $G$; it is a $k$-vector space, and Theorem \ref{I:1-4-5} gives:

\begin{proposition}\label{I:1-5-2}
The inclusion of $V$ in $V'$ induces an isomorphism $V\otimes_k k' \iso V'$.
\end{proposition}

In particular, if $V'$ has dimension $1$ and $v'\in V'$ is nonzero, $\varphi_\sigma$ is determined by the constant $c(\sigma)\in{k'}^\times$ such that $\varphi_\sigma(v')=c(\sigma)v'$, and condition c) becomes
\[
  c(\tau\sigma) = c(\tau)\cdot \tau(c(\sigma))\text{.}
\]
By the proposition, there is a nonzero invariant vector $v=\mu v'$, with $\mu\in{k'}^\times$. Thus, for every $\sigma\in G$,
\[
  c(\sigma) = \mu\cdot \sigma(\mu^{-1})\text{.}
\]
In other words, every $1$-cocycle of $G$ with values in ${k'}^\times$ is a coboundary:

\begin{corollary}\label{I:1-5-3}
One has $\h^1(G,{k'}^\times)=0$.
\end{corollary}

\subsection{Grothendieck topologies}\label{I:1-6}

We now recast the definitions of the preceding subsections in an abstract framework that encompasses both topological spaces and schemes.

\subsubsection{}\label{I:1-6-1}

Let $\cS$ be a category and $U$ an object of $\cS$. A \emph{sieve} on $U$ is a subset $\cU$ of $\ob(\cS/U)$ such that if $\varphi:V\to U$ belongs to $\cU$ and $\psi:W\to V$ is a morphism in $\cS$, then $\varphi\circ\psi:W\to U$ belongs to $\cU$.

If $\{\varphi_i:U_i\to U\}$ is a family of morphisms, the sieve generated by the $U_i$ is, by definition, the set of morphisms $\varphi:V\to U$ that factor through one of the $\varphi_i$.

If $\cU$ is a sieve on $U$ and $\varphi:V\to U$ is a morphism, the restriction $\cU_V$ of $\cU$ to $V$ is, by definition, the sieve on $V$ consisting of the morphisms $\psi:W\to V$ such that $\varphi\circ\psi:W\to U$ belongs to $\cU$.

\subsubsection{}\label{I:1-6-2}

To give a \emph{Grothendieck topology} on $\cS$ is to give, for every object $U$ of $\cS$, a set $C(U)$ of sieves on $U$, called covering sieves, such that the following axioms hold:
\begin{enumerate}[\indent a)]
  \item The sieve generated by the identity of $U$ is covering.
  \item If $\cU$ is a covering sieve on $U$ and $V\to U$ is a morphism, then the sieve $\cU_V$ is covering.
  \item A locally covering sieve is covering. In other words, if $\cU$ is a covering sieve on $U$ and $\cU'$ is a sieve on $U$ such that, for every $V\to U$ belonging to $\cU$, the sieve $\cU'_V$ is covering, then $\cU'$ is covering.
\end{enumerate}

A category endowed with a Grothendieck topology is called a \emph{site}.

\subsubsection{}\label{I:1-6-3}

Given a site $\cS$, a presheaf on $\cS$ is a contravariant functor $\sF$ from $\cS$ to the category of sets. For every object $U$ of $\cS$, the elements of $\sF(U)$ are called sections of $\sF$ over $U$. For every morphism $V\to U$ and every $s\in\sF(U)$, the image of $s$ in $\sF(V)$ is denoted by $s|V$ ($s$ restricted to $V$).

If $\cU$ is a sieve on $U$, a section given $\cU$-locally consists, for every $V\to U$ belonging to $\cU$, of a section $s_V\in\sF(V)$ such that, for every morphism $W\to V$, one has $s_V|W=s_W$. The presheaf $\sF$ is a \emph{sheaf} if, for every object $U$ of $\cS$, every covering sieve $\cU$ on $U$, and every section $\{s_V\}$ given $\cU$-locally, there is a unique section $s\in\sF(U)$ such that $s|V=s_V$ for every $V\to U$ belonging to $\cU$.

One defines \emph{abelian sheaves} similarly by replacing the category of sets by the category of abelian groups. The category of abelian sheaves on $\cS$ is an abelian category with enough injectives. A sequence $\sF\xrightarrow f\sG\xrightarrow g\sH$ of sheaves is exact if, for every object $U$ of $\cS$ and every $s\in\sG(U)$ such that $g(s)=0$, there is locally a $t$ such that $f(t)=s$; that is, if there is a covering sieve $\cU$ on $U$ and, for every $V\in\cU$, a section $t_V$ of $\sF$ over $V$ such that $f(t_V)=s|V$.

\subsubsection{Examples}\label{I:1-6-4}

We have already seen two examples above.

\begin{enumerate}[\indent a)]
  \item Let $X$ be a topological space and $\cS$ the category whose objects are the open subsets of $X$ and whose morphisms are the natural inclusions. The Grothendieck topology on $\cS$ corresponding to the usual topology of $X$ is the one for which a sieve $\cU$ on an open subset $U$ of $X$ is covering if the union of the open subsets belonging to that sieve is $U$. It is clear that the category of sheaves on $\cS$ is equivalent to the category of sheaves on $X$ in the usual sense.
  \item Let $X$ be a scheme and $\cS$ the category of schemes over $X$. The fpqc (faithfully flat and quasi-compact) topology on $\cS$ is the Grothendieck topology for which a sieve on an $X$-scheme $U$ is covering if it is generated by a finite family of flat morphisms whose images cover $U$.
\end{enumerate}

\subsubsection{Cohomology}\label{I:1-6-5}

We shall always assume that the category $\cS$ has a final object $X$. The group $\sF(X)$ is then called the group of global sections of an abelian sheaf $\sF$ and is denoted $\Gamma\sF$ or $\h^0(X,\sF)$. The functor $\sF\mapsto\Gamma\sF$ is a left exact functor from the category of abelian sheaves on $\cS$ to the category of abelian groups; its derived functors (or satellites) are denoted $\h^i(X,-)$. These cohomology groups represent the obstructions to passing from local to global. By definition, if $0\to\sF\to\sG\to\sH\to0$ is an exact sequence of abelian sheaves, there is a long exact cohomology sequence:
\[\xymatrix{
  0 \ar[r]
    & \h^0(X,\sF) \ar[r]
    & \h^0(X,\sG) \ar[r]
    & \h^0(X,\sH) \ar[r]
    & \h^1(X,\sF) \ar[r]
    & \cdots \\
  \ldots \ar[r]
    & \h^n(X,\sF) \ar[r]
    & \h^n(X,\sG) \ar[r]
    & \h^n(X,\sH) \ar[r]
    & \h^{n+1}(X,\sF) \ar[r]
    & \cdots
}\]

\subsubsection{}\label{I:1-6-6}

Given an abelian sheaf $\sF$ on $\cS$, an \emph{$\sF$-torsor} is a sheaf $\sG$ endowed with an action $\sF\times\sG\to\sG$ of $\sF$ such that locally (after restriction to all the objects of a sieve covering the final object $X$), $\sG$ with its $\sF$-action is isomorphic to $\sF$ with the canonical translation action $\sF\times\sF\to\sF$.

One can show that $\h^1(X,\sF)$ is identified with the set of isomorphism classes of $\sF$-torsors.

\section{The étale topology}\label{I:2}

We specialize the definitions of the preceding section to the étale topology of a scheme $X$ (\ref{I:2-1}, \ref{I:2-2}, \ref{I:2-3}). When $X$ is the spectrum of a field $K$, the resulting cohomology agrees with the Galois cohomology of $K$ (\ref{I:2-4}).

\subsection{The étale topology}\label{I:2-1}

We begin by recalling a few facts about étale morphisms.

\begin{definition}\label{I:2-1-1}
Let $A$ be a (commutative) ring. An $A$-algebra $B$ is said to be étale if $B$ is an $A$-algebra of finite presentation and the following equivalent conditions hold:
\begin{enumerate}[\indent a)]
  \item For every $A$-algebra $C$ and every square-zero ideal $J$ of $C$, the canonical map
    \[
      \hom_{A\textnormal{-alg}}(B,C) \to \hom_{A\textnormal{-alg}}(B,C/J)
    \]
    is a bijection.
  \item $B$ is a flat $A$-module and $\Omega_{B/A}=0$ (where $\Omega_{B/A}$ denotes the module of relative differentials).
  \item Let $B=A[X_1,\dotsc,X_n]/I$ be a presentation of $B$. For every prime ideal $\fp$ of $A[X_1,\dotsc,X_n]$ containing $I$, there are polynomials $P_1,\dotsc,P_n\in I$ such that $I_\fp$ is generated by the images of $P_1,\dotsc,P_n$ and $\det(\partial P_i/\partial X_j)\notin\fp$.
\end{enumerate}
\end{definition}
(cf. \cite[I]{sga1} or \cite[V]{ra70})
%[cf. SGA I, exposé I ou M. {\sc Raynoud}, \emph{Anneaux Locaux Henséliens}, chapitre V].

A morphism of schemes $f:X\to S$ is said to be \emph{étale} if, for every $x\in X$, there is an affine open neighborhood $U=\spec(A)$ of $f(x)$ and an affine open neighborhood $V=\spec(B)$ of $x$ in $X\times_SU$ such that $B$ is an étale $A$-algebra.

\subsubsection{Examples}\label{I:2-1-2}
\begin{enumerate}[\indent a)]
  \item If $A$ is a field, an $A$-algebra $B$ is étale if and only if it is a finite product of separable extensions of $A$.
  \item If $X$ and $S$ are schemes of finite type over $\dC$, a morphism $f:X\to S$ is étale if and only if its analytification $f^{\text{an}}:X^{\text{an}}\to S^{\text{an}}$ is a local isomorphism.
\end{enumerate}

\subsubsection{Elementary properties}\label{I:2-1-3}
\begin{enumerate}[\indent a)]
  \item (Base change) If $f:X\to S$ is an étale morphism, then so is $f_{S'}:X\times_SS'\to S'$ for every morphism $S'\to S$.
  \item If $f:X\to S$ and $g:Y\to S$ are two étale morphisms, every $S$-morphism from $X$ to $Y$ is étale.
  \item (Descent) Let $f:X\to S$ be a morphism. If there is a faithfully flat morphism $S'\to S$ such that $f_{S'}:X\times_SS'\to S'$ is étale, then $f$ is étale.
\end{enumerate}

\subsubsection{}\label{I:2-1-4}

Let $X$ be a scheme, and let $\cS$ be the category of étale $X$-schemes; by (\ref{I:2-1-3}.c), every morphism in $\cS$ is étale. The \emph{étale topology} on $\cS$ is the topology for which a sieve on $U$ is covering if it is generated by a finite family of morphisms $\varphi_i:U_i\to U$ whose images cover $U$. The site defined by $\cS$ with the étale topology is called the \emph{étale site} of $X$ and is denoted $\et X$.

\subsection{Examples of sheaves}\label{I:2-2}

\subsubsection{The constant sheaf}\label{I:2-2-1}

Let $C$ be an abelian group, and, for simplicity, suppose that $X$ is Noetherian. We denote by $\const C_X$ (or simply $C$ when there is no ambiguity) the sheaf defined by $U\mapsto C^{\pi_0(U)}$, where $\pi_0(U)$ is the (finite) set of connected components of $U$. The most important case will be $C=\dZ/n$. Thus, by definition,
\[
  \h^0(X,\dZ/n)=\left(\dZ/n\right)^{\pi_0(X)}\text{.}
\]
Moreover, $\h^1(X,\dZ/n)$ is the set of isomorphism classes of $\dZ/n$-torsors (\ref{I:1-6-6}), in other words, of Galois étale coverings of $X$ with group $\dZ/n$. In particular, if $X$ is connected and $\pi_1(X)$ is its fundamental group for a chosen base point, then
\[
  \h^1(X,\dZ/n) = \hom(\pi_1(X),\dZ/n)\text{.}
\]

\subsubsection{The multiplicative group}\label{I:2-2-2}

We denote by $\dG_{m,X}$ (or $\dG_m$ when there is no ambiguity) the sheaf defined by $U\mapsto\Gamma(U,\sO_U^\times)$; it is indeed a sheaf by the faithfully flat descent theorem (\ref{I:1-4-5}). By definition,
\[
  \h^0(X,\dG_m) = \h^0(X,\sO_X)^\times\text{;}
\]
in particular, if $X$ is reduced, connected, and proper over an algebraically closed field $k$, then
\[
  \h^0(X,\dG_m) = k^\times\text{.}
\]

\begin{proposition}\label{I:2-2-3}
There is an isomorphism
\[
  \h^1(X,\dG_m) = \pic(X)\text{,}
\]
where $\pic(X)$ is the group of classes of invertible sheaves on $X$.
\end{proposition}
\begin{proof}
Let $*$ be the functor that assigns to an invertible sheaf $\sL$ on $X$ the following presheaf $\sL^*$ on $\et X$: for $\varphi:U\to X$ étale,
\[
  \sL^*(U)=\operatorname{Isom}_U(\sO_U,\varphi^*\sL)\text{.}
\]
% the original references here are just to (4.2) and (4.5), but the context
% makes it seem that they refer to the results in chapter 1. Likewise a little
% later
By (\ref{I:1-4-2}.i) and (\ref{I:1-4-5}) (full faithfulness), this presheaf is a sheaf; it is in fact a $\dG_m$-torsor. One immediately checks that
\begin{enumerate}[\indent a)]
  \item the functor $*$ is compatible with (étale) localization;
  \item it induces an equivalence from the category of trivial invertible sheaves (i.e. isomorphic to $\sO_X$) to the category of trivial $\dG_m$-torsors: $\sL$ is trivial if and only if $\sL^*$ is trivial.
\end{enumerate}

Moreover, by (\ref{I:1-4-2}.ii) and (\ref{I:1-4-5}),
\begin{enumerate}[\indent a)]
\setcounter{enumi}{2}
  \item the notion of an invertible sheaf is local for the étale topology.
\end{enumerate}

It follows formally from a), b), and c) that $*$ is an equivalence between the category of invertible sheaves on $X$ and the category of $\dG_m$-torsors on $\et X$; it induces the desired isomorphism. The inverse equivalence is constructed as follows. If $\sT$ is a $\dG_m$-torsor, there is a finite étale covering $\{U_i\}$ of $X$ such that the torsors $\sT|U_i$ are trivial; then $\sT$ is trivial on every $V$ étale over $X$ that belongs to the sieve $\cU\subset\et X$ generated by $\{U_i\}$. On every $V\in\cU$, $\sT|V$ corresponds to an invertible sheaf $\sL_V$ (by b)), and the $\sL_V$ form an invertible sheaf $\sL_\cU$ given $\cU$-locally (by a)). By c), the latter comes from an invertible sheaf $\sL(\sT)$ on $X$, and $\sT\mapsto\sL(\sT)$ is the desired inverse of $*$.
\end{proof}

\subsubsection{Roots of unity}\label{I:2-2-4}

For every integer $n>0$, the kernel of the $n$th-power map on $\dG_m$ is called the sheaf of $n$th roots of unity and is denoted $\dmu_n$. If $X$ is a scheme over a separably closed field $k$ and $n$ is invertible in $k$, the choice of a primitive $n$th root of unity $\zeta\in k$ defines an isomorphism $i\mapsto\zeta^i$ from $\dZ/n$ to $\dmu_n$.

The relation between cohomology with coefficients in $\dmu_n$ and cohomology with coefficients in $\dG_m$ is given by the cohomology exact sequence deduced from the

\begin{theorem}[Kummer theory]\label{I:2-2-5}
If $n$ is invertible on $X$, the $n$th-power map on $\dG_m$ is an epimorphism of sheaves. Thus there is an exact sequence
\[
  0 \to \dmu_n \to \dG_m \to \dG_m \to 0\text{.}
\]
\end{theorem}
\begin{proof}
Let $U\to X$ be an étale morphism and let $a\in\dG_m(U)=\Gamma(U,\sO_U^\times)$. Since $n$ is invertible on $U$, the equation $T^n-a=0$ is separable; in other words, $U'=\spec\left(\sO_U[T]/(T^n-a)\right)$ is étale over $U$. Moreover, $U'\to U$ is surjective and $a$ has an $n$th root on $U'$, which proves the result.
\end{proof}

\subsection{Stalks and direct images}\label{I:2-3}

\subsubsection{}\label{I:2-3-1}

A \emph{geometric point} of $X$ is a morphism $\bar x\to X$, where $\bar x$ is the spectrum of a separably closed field $k(\bar x)$. By abuse of notation we denote it by $\bar x$, with the morphism $\bar x\to X$ understood. If $x$ is the image of $\bar x$ in $X$, one says that $\bar x$ is centered at $x$. If $k(\bar x)$ is an algebraic extension of the residue field $k(x)$, one says that $\bar x$ is an \emph{algebraic} geometric point of $X$.

An \emph{étale neighborhood} of $\bar x$ is a commutative diagram
\[\xymatrix{
  & U \ar[d] \\
  \bar x \ar[ur] \ar[r]
  & X\text{,}
}\]
where $U\to X$ is an étale morphism.

The \emph{strict localization} of $X$ at $\bar x$ is the ring $\sO_{X,\bar x}=\varinjlim\Gamma(U,\sO_U)$, where the direct limit is over the étale neighborhoods of $\bar x$. It is a strictly Henselian local ring whose residue field is the separable closure in $k(\bar x)$ of the residue field $k(x)$ of $X$ at $x$. It plays the role of the local ring for the étale topology.

\subsubsection{}\label{I:2-3-2}

Given a sheaf $\sF$ on $\et X$, the \emph{stalk} of $\sF$ at $\bar x$ is the set (respectively, group, etc.) $\sF_{\bar x}=\varinjlim\sF(U)$, where the direct limit is again taken over the étale neighborhoods of $\bar x$.

A homomorphism of sheaves $\sF\to\sG$ is a monomorphism, epimorphism, or isomorphism if and only if the induced morphisms on stalks $\sF_{\bar x}\to\sG_{\bar x}$ have the corresponding property at every geometric point of $X$. If $X$ is of finite type over an algebraically closed field, it is enough to check this at the rational points of $X$.

\subsubsection{}\label{I:2-3-3}

If $f:X\to Y$ is a morphism of schemes and $\sF$ is a sheaf on $\et X$, the \emph{direct image} $f_*\sF$ of $\sF$ by $f$ is the sheaf on $\et Y$ defined by $f_*\sF(V)=\sF(X\times_YV)$ for every $V$ étale over $Y$. The functor $f_*:\sh(\et X)\to\sh(\et Y)$ is left exact. Its right derived functors $\eR^qf_*$ are called higher direct images. If $\bar y$ is a geometric point of $Y$, then
\[
  (\eR^q f_* \sF)_{\bar y} = \varinjlim \h^q(V\times_Y X,\sF)\text{,}
\]
where the direct limit is taken over the étale neighborhoods $V$ of $\bar y$.

Let $\sO_{Y,\bar y}$ be the strict localization of $Y$ at $\bar y$, put $\widetilde Y=\spec(\sO_{Y,\bar y})$, and put $\widetilde X=X\times_Y\widetilde Y$. One can extend $\sF$ to $\et{\widetilde X}$ (this is a special case of the general notion of inverse image) as follows. Let $\widetilde U$ be a scheme étale over $\widetilde X$. There are an étale neighborhood $V$ of $\bar y$ and a scheme $U$ étale over $X\times_YV$ such that $\widetilde U=U\times_V\widetilde Y$; set
\[
  \sF(\widetilde U)=\varinjlim \sF\left(U\times_V V'\right)\text{,}
\]
where the direct limit is taken over the étale neighborhoods $V'$ of $\bar y$ that dominate $V$. With this definition,
\[
  \left(\eR^q f_* \sF\right)_{\bar y} = \h^q(\widetilde X,\sF)\text{.}
\]

The functor $f_*$ has a left adjoint $f^*$, the ``inverse image'' functor. If $\bar x$ is a geometric point of $X$ and $f(\bar x)$ is its image in $Y$, then $(f^*\sF)_{\bar x}=\sF_{f(\bar x)}$. This formula shows that $f^*$ is exact. Consequently $f_*$ takes injective sheaves to injective sheaves, and the spectral sequence for the composite functor $\Gamma\circ f_*$ (respectively, $g_*f_*$) gives the

\begin{theorem}[Leray spectral sequence]\label{I:2-3-4}
Let $\sF$ be an abelian sheaf on $\et X$ and let $f:X\to Y$ be a morphism of schemes (respectively, let $X\xrightarrow fY\xrightarrow gZ$ be morphisms of schemes). There is a spectral sequence
\begin{align*}
  E_2^{pq} &= \h^p(Y,\eR^q f_* \sF) \Rightarrow \h^{p+q}(X,\sF) \\
  \text{(resp.}\qquad E_2^{pq} &= \eR^p g_* \eR^q f_* \sF \Rightarrow \eR^{p+q}(gf)_* \sF\text{).}
\end{align*}
\end{theorem}

\begin{corollary}\label{I:2-3-5}
If $\eR^qf_*\sF=0$ for every $q>0$, then $\h^p(Y,f_*\sF)=\h^p(X,\sF)$ (respectively, $\eR^pg_*(f_*\sF)=\eR^p(gf)_*\sF$) for every $p\geqslant0$.
\end{corollary}

This applies in particular in the following case.

\begin{proposition}\label{I:2-3-6}
Let $f:X\to Y$ be a finite morphism (or even, by passage to the limit, an integral morphism), and let $\sF$ be an abelian sheaf on $X$. Then $\eR^qf_*\sF=0$ for every $q>0$.
\end{proposition}

Indeed, let $\bar y$ be a geometric point of $Y$, let $\widetilde Y$ be the spectrum of the strict localization of $Y$ at $\bar y$, and put $\widetilde X=X\times_Y\widetilde Y$. By the preceding discussion, it is enough to show that $\h^q(\widetilde X,\sF)=0$ for every $q>0$. But $\widetilde X$ is the spectrum of a product of strictly Henselian local rings (cf. \cite[I]{ra70}), and the functor $\Gamma(\widetilde X,-)$ is exact because every surjective étale $\widetilde X$-scheme has a section. This proves the assertion.

\subsection{Galois cohomology}\label{I:2-4}

For $X=\spec(K)$, the spectrum of a field, we shall see that étale cohomology is identified with Galois cohomology.

\subsubsection{}\label{I:2-4-1}

Let us begin with a topological analogy. If $K$ is the function field of an integral affine algebraic variety $Y=\spec(A)$ over $\dC$, then $K=\varinjlim_{f\in A} A[1/f]$.

In other words, $X=\varprojlim U$, where $U$ ranges over the open subsets of $Y$. It is known that there are arbitrarily small Zariski open subsets that are $K(\pi,1)$'s for the classical topology. It is therefore natural to regard $\spec(K)$ itself as a $K(\pi,1)$, where $\pi$ is the fundamental group of $X$ (in the algebraic sense), namely the Galois group of $\bar K/K$, with $\bar K$ a separable closure of $K$.

\subsubsection{}\label{I:2-4-2}

More precisely, let $K$ be a field, let $\bar K$ be a separable closure of $K$, and let $G=\gal(\bar K/K)$ be the topological Galois group. To every finite étale $K$-algebra $A$ (a finite product of separable extensions of $K$), associate the finite set $\hom_K(A,\bar K)$. The Galois group $G$ acts on this set through a discrete (hence finite) quotient. If $A=K[T]/(F)$, this set is identified with the set of roots of $F$ in $\bar K$. Galois theory, in the form given to it by Grothendieck, says the following.

\begin{proposition}\label{I:2-4-3}
The functor
\[
  \left\{\begin{array}{c}
          \textnormal{finite étale} \\
          \textnormal{$K$-algebras}
        \end{array}\right\}
  \to
  \left\{\begin{array}{c}
          \textnormal{finite sets on which} \\
          \textnormal{$G$ acts continuously}
        \end{array}\right\}
\]
that assigns $\hom_K(A,\bar K)$ to an étale algebra $A$ is an anti-equivalence of categories.
\end{proposition}

One deduces an analogous description of sheaves for the étale topology on $\spec(K)$.

\begin{proposition}\label{I:2-4-4}
The functor
\[
  \left\{\begin{array}{c}
          \textnormal{étale sheaves} \\
          \textnormal{on $\spec(K)$}
        \end{array}\right\}
  \to
  \left\{\begin{array}{c}
          \textnormal{sets on which} \\
          \textnormal{$G$ acts continuously}
        \end{array}\right\}
\]
that assigns to a sheaf $\sF$ its stalk $\sF_{\bar K}$ at the geometric point $\spec(\bar K)$ is an equivalence of categories.
\end{proposition}

One says that $G$ acts continuously on a set $E$ if the stabilizer of every element of $E$ is an open subgroup of $G$. The functor in the opposite direction is described in the evident way: let $A$ be a finite étale $K$-algebra, let $U=\spec(A)$, and let $U(\bar K)=\hom_K(A,\bar K)$ be the corresponding $G$-set; then $\sF(U)= \hom_{G\text{-ens}}(U(\bar K),\sF_{\bar K})$.

In particular, if $X=\spec(K)$, then $\sF(X)=\sF_{\bar K}^G$. Restricting to abelian sheaves and passing to derived functors gives canonical isomorphisms
\[
  \h^q(\et X,\sF) = \h^q(G,\sF_{\bar K})
\]

\subsubsection{Examples}\label{I:2-4-5}

\begin{enumerate}[\indent a)]
  \item The constant sheaf $\dZ/n$ corresponds to $\dZ/n$ with the trivial action of $G$.
  \item The sheaf $\dmu_n$ of $n$th roots of unity corresponds to the group $\dmu_n(\bar K)$ of $n$th roots of unity in $\bar K$, with its natural $G$-action.
  \item The sheaf $\dG_m$ corresponds to the group $\bar K^\times$ with its natural $G$-action.
\end{enumerate}

\section{Cohomology of curves}\label{I:3}

For topological spaces, dévissages using the K\"unneth formula and simplicial decompositions reduce the calculation of cohomology to the interval $I=[0,1]$, for which $\h^0(I,\dZ)=\dZ$ and $\h^q(I,\dZ)=0$ for $q>0$.

In our case the dévissages lead to more complicated objects, namely curves over an algebraically closed field; we shall calculate their cohomology in this section. The situation is more complex than in the topological case, since the cohomology groups vanish only for $q>2$. The essential ingredient in the calculations is the vanishing of the Brauer group of the function field of such a curve (Tsen's theorem, \ref{I:3-2}).

\subsection{The Brauer group}\label{I:3-1}

We first recall its classical definition.

\begin{definition}\label{I:3-1-1}
Let $K$ be a field and $A$ a finite-dimensional $K$-algebra. The algebra $A$ is said to be central simple over $K$ if the following equivalent conditions hold:
\begin{enumerate}[\indent a)]
  \item $A$ has no nontrivial two-sided ideal and its center is $K$.
  \item There is a finite Galois extension $K'/K$ such that $A_{K'}=A\otimes_KK'$ is isomorphic to a full matrix algebra over $K'$.
  \item $A$ is $K$-isomorphic to a full matrix algebra over a division algebra with center $K$.
\end{enumerate}
\end{definition}

Two such algebras are said to be equivalent if the division algebras associated with them in c) are $K$-isomorphic. If the two algebras have the same dimension, this is equivalent to their being $K$-isomorphic. The tensor product induces, after passage to the quotient, an abelian group structure on the set of equivalence classes. This group is classically called the \emph{Brauer group} of $K$ and is denoted $\br(K)$.

\subsubsection{}\label{I:3-1-2}

Let $\br(n,K)$ denote the set of $K$-isomorphism classes of $K$-algebras $A$ for which there is a finite Galois extension $K'$ of $K$ such that $A_{K'}$ is isomorphic to the algebra $M_n(K')$ of $n\times n$ matrices over $K'$. By definition, $\br(K)$ is the union of the subsets $\br(n,K)$ for $n\in\dN$. Let $\bar K$ be an algebraic closure of $K$ and put $G=\gal(\bar K/K)$. The set $\br(n,K)$ is the set of ``forms'' of $M_n(\bar K)$ and is therefore canonically isomorphic to $\h^1\left(G,\aut(M_n(\bar K))\right)$.

Every automorphism of $M_n(\bar K)$ is inner. Consequently $\aut(M_n(\bar K))$ is identified with the projective linear group $\pgl(n,\bar K)$, and there is a canonical bijection
\[
  \theta_n : \br(n,K) \iso \h^1\left(G,\pgl(n,\bar K)\right)\text{.}
\]

On the other hand, the exact sequence
\begin{equation*}\tag{3.1.2.1}\label{I:eq:3-1-2-1}
  1 \to \bar K^\times \to \gl(n,\bar K) \to \pgl(n,\bar K) \to 1\text{,}
\end{equation*}
defines a coboundary map
\[
  \Delta_n : \h^1\left(G,\pgl(n,\bar K)\right) \to \h^2(G,\bar K^\times)\text{.}
\]

Composing $\theta_n$ and $\Delta_n$ gives a map
\[
  \delta_n : \br(n,K) \to \h^2(G,\bar K^\times)\text{.}
\]
It is easy to check that the maps $\delta_n$ are compatible with one another and define a group homomorphism
\[
  \delta : \br(K) \to \h^2(G,\bar K^\times)\text{.}
\]

\begin{proposition}\label{I:3-1-3}
The homomorphism $\delta:\br(K)\to\h^2(G,\bar K^\times)$ is bijective.
\end{proposition}

This follows from the next two lemmas.

\begin{lemma}\label{I:3-1-4}
The map $\Delta_n:\h^1\left(G,\pgl(n,\bar K)\right)\to\h^2(G,\bar K^\times)$ is injective.
\end{lemma}

By \cite[Corollary to Proposition I-44]{se94}, it is enough to check that, whenever the exact sequence (\ref{I:eq:3-1-2-1}) is twisted by an element of $\h^1\left(G,\pgl(n,\bar K)\right)$, the $\h^1$ of the middle group is trivial. This middle group is the group of $\bar K$-points of the multiplicative group of a central simple algebra $A$ of rank $n^2$ over $K$. To prove that $\h^1(G,A_{\bar K}^\times)=0$, interpret $A^\times$ as the automorphism group of the free $A$-module $L$ of rank $1$, and $\h^1$ as the set of ``forms'' of $L$---$A$-modules of rank $n^2$ over $K$, which are automatically free.

\begin{lemma}\label{I:3-1-5}
Let $\alpha\in\h^2(G,\bar K^\times)$, let $K'$ be a finite extension of $K$ contained in $\bar K$, put $n=[K':K]$, and put $G'=\gal(\bar K/K')$. If the image of $\alpha$ in $\h^2(G',\bar K^\times)$ is zero, then $\alpha$ belongs to the image of $\Delta_n$.
\end{lemma}

First observe that
\[
  \h^2(G',\bar K^\times) \simeq \h^2\left(G,(\bar K\otimes_K K')^\times\right)\text{.}
\]
(From a geometric viewpoint, if $x=\spec(K)$, $x'=\spec(K')$, and $\pi:x'\to x$ is the canonical morphism, then $\eR^q\pi_*(\dG_{m,x'})=0$ for $q>0$, and consequently $\h^q(x',\dG_{m,x'})\simeq\h^q(x,\pi_*\dG_{m,x'})$ for $q\geqslant0$.)

Moreover, the choice of a basis of $K'$ as a vector space over $K$ defines a homomorphism
\[
  (\bar K\otimes_K K')^\times \to \gl(n,\bar K)
\]
which assigns to an element $x$ the endomorphism of $\bar K\otimes_KK'$ given by multiplication by $x$. One then has a commutative diagram with exact rows:
\[\xymatrix{
  1 \ar[r]
    & \bar K^\times \ar[r] \ar@{=}[d]
    & (\bar K\otimes_K K')^\times \ar[r] \ar[d]
    & (\bar K\otimes_K K')^\times / \bar K^\times \ar[r] \ar[d]
    & 1 \\
  1 \ar[r]
    & \bar K^\times \ar[r]
    & \gl(n,\bar K) \ar[r]
    & \pgl(n,\bar K) \ar[r]
    & 1
}\]
The lemma follows from the commutative diagram obtained by passing to cohomology:
\[\xymatrix{
  \h^1\left(G,(\bar K\otimes_K K')^\times/\bar K^\times\right) \ar[r] \ar[d]
    & \h^2(G,\bar K^\times) \ar[r] \ar@{=}[d]
    & \h^2\left(G,(\bar K\otimes_K K')^\times\right) \\
  \h^1\left(G,\pgl(n,\bar K)\right) \ar[r]^-{\Delta_n}
    & \h^2(G,\bar K^\times) \text{.}
}\]

\begin{proposition}\label{I:3-1-6}
Let $K$ be a field, let $\bar K$ be an algebraic closure of $K$, and put $G=\gal(\bar K/K)$. Suppose that $\br(K')=0$ for every finite extension $K'$ of $K$. Then:
\begin{enumerate}[\indent i)]
  \item $\h^q(G,\bar K^\times)=0$ for every $q>0$.
  \item $\h^q(G,\sF)=0$ for every torsion $G$-module $\sF$ and every $q\geqslant2$.
\end{enumerate}
\end{proposition}

(For the proof, cf. \cite{se94}.)

\subsection{Tsen's theorem}\label{I:3-2}

\begin{definition}\label{I:3-2-1}
A field $K$ is said to be $C_1$ if every nonconstant homogeneous polynomial $f(x_1,\dotsc,x_n)$ of degree $d<n$ has a nontrivial zero.
\end{definition}

\begin{proposition}\label{I:3-2-2}
If a field $K$ is $C_1$, then $\br(K)=0$.
\end{proposition}

It is enough to show that every division algebra $D$ with center $K$ and finite over $K$ is equal to $K$. Let $r^2$ be the degree of $D$ over $K$ and let $\operatorname{Nrd}:D\to K$ be the reduced norm.

(Locally for the étale topology on $K$, $D$ is isomorphic---noncanonically--- to a matrix algebra $M_r$, and the reduced norm agrees with the determinant. The latter is well defined independently of the chosen isomorphism between $D$ and $M_r$, because every automorphism of $M_r$ is inner and similar matrices have the same determinant. This map, defined locally for the étale topology, descends, by its local uniqueness, to a map $\operatorname{Nrd}:D\to K$.)

The only zero of $\operatorname{Nrd}$ is the zero element of $D$, since, if $x\ne0$, then $\operatorname{Nrd}(x)\cdot\operatorname{Nrd}(x^{-1})=1$. On the other hand, if $\{e_1,\dotsc,e_{r^2}\}$ is a basis of $D$ over $K$ and $x=\sum x_ie_i$, the function $\operatorname{Nrd}(x)$ is a homogeneous polynomial $\operatorname{Nrd}(x_1,\dotsc,x_{r^2})$ of degree $r$ (this is clear locally for the étale topology). Since $K$ is $C_1$, one has $r^2\leqslant r$, hence $r=1$ and $D=K$.

\begin{theorem}[Tsen]\label{I:3-2-3}
Let $k$ be an algebraically closed field and let $K$ be an extension of transcendence degree $1$ over $k$. Then $K$ is $C_1$.
\end{theorem}

First suppose that $K=k(X)$. Let
\[
  f(T) = \sum a_{i_1,\dotsc i_n} T_1^{i_1} \dotsm T_n^{i_n}
\]
be a homogeneous polynomial of degree $d<n$ with coefficients in $k(X)$. After multiplying the coefficients by a common denominator, we may suppose that they belong to $k[X]$. Put $\delta=\sup\deg(a_{i_1\dotsc i_n})$. We seek a nontrivial zero in $k[X]$ by the method of undetermined coefficients, writing each $T_i$ ($i=1,\dotsc,n$) as a polynomial of degree $N$ in $X$. The equation $f(T)=0$ then becomes a system of homogeneous equations in the $n(N+1)$ coefficients of the polynomials $T_i(X)$, expressing the vanishing of the coefficients of the polynomial in $X$ obtained by substituting $T_i(X)$ for $T_i$. This polynomial has degree at most $\delta+Nd$, so there are $\delta+Nd+1$ equations in $n(N+1)$ variables. Since $k$ is algebraically closed, this system has a nontrivial solution if $n(N+1)>Nd+\delta+1$, which holds for sufficiently large $N$ when $d<n$.

To prove the theorem in general, it is clearly enough to treat the case in which $K$ is a finite extension of the purely transcendental extension $k(X)$ of $k$. Let $f(T)=f(T_1,\dotsc,T_n)$ be a homogeneous polynomial of degree $d<n$ with coefficients in $K$. Put $s=[K:k(X)]$, and let $e_1,\dotsc,e_s$ be a basis of $K$ over $k(X)$. Introduce $sn$ new variables $U_{ij}$ such that $T_i=\sum_jU_{ij}e_j$. For $f(T)$ to have a nontrivial zero in $K$, it is enough that the polynomial $g(U_{ij})=N_{K/k(X)}(f(T))$ have a nontrivial zero in $k(X)$. But $g$ is a polynomial of degree $sd$ in $sn$ variables, and the result follows.

\begin{corollary}\label{I:3-2-4}
Let $k$ be an algebraically closed field and let $K$ be an extension of transcendence degree $1$ over $k$. Then the étale cohomology groups $\h^q(\spec(K),\dG_m)$ vanish for every $q>0$.
\end{corollary}

\subsection{Cohomology of smooth curves}

From now on, unless expressly stated otherwise, all cohomology groups are étale cohomology groups.

\begin{proposition}\label{I:3-3-1}
Let $k$ be an algebraically closed field and let $X$ be a connected smooth projective curve over $k$. Then
\begin{align*}
  \h^0(X,\dG_m) &= k^\times \text{,} \\
  \h^1(X,\dG_m) &= \pic(X) \text{,} \\
  \h^q(X,\dG_m) &= 0 \text{ for $q\geqslant 2$.}
\end{align*}
\end{proposition}

Let $\eta$ be the generic point of $X$, let $j:\eta\to X$ be the canonical morphism, and let $\dG_{m,\eta}$ be the multiplicative group of the fraction field $K(X)$. For every closed point $x$ of $X$, let $i_x:x\to X$ be the canonical immersion and let $\dZ_x$ be the constant sheaf with value $\dZ$ on $x$. Thus $j_*\dG_{m,\eta}$ is the sheaf of nonzero meromorphic functions on $X$, and $\bigoplus_{x\in X}i_{x*}\dZ_x$ is the sheaf of divisors. Hence there is an exact sequence of sheaves
\begin{equation*}\tag{3.3.1.1}\label{I:eq:3-3-1-1}
\xymatrix{
  0 \ar[r]
    & \dG_m \ar[r]
    & j_* \dG_{m,\eta} \ar[r]^-{\dv} \ar[r]
    & \displaystyle\bigoplus_{x\in X} i_{x*} \dZ_x \ar[r]
    & 0\text{.}
}
\end{equation*}

\begin{lemma}\label{I:3-3-2}
One has $\eR^qj_*\dG_{m,\eta}=0$ for every $q>0$.
\end{lemma}

It is enough to show that the stalk of this sheaf at every closed point $x$ of $X$ is zero. If $\tilde\sO_{X,x}$ is the strict Henselization of $X$ at $x$ and $K$ is the fraction field of $\tilde\sO_{X,x}$, then
\[
  \spec(K) = \eta\times_X \spec(\tilde\sO_{X,x})\text{,}
\]
and hence $(\eR^qj_*\dG_{m,\eta})_x=\h^q(\spec(K),\dG_m)$.

But $K$ is an algebraic extension of $k(X)$, hence an extension of transcendence degree $1$ over $k$; the lemma follows from (\ref{I:3-2-4}).

\begin{lemma}\label{I:3-3-3}
One has $\h^q(X,j_*\dG_{m,\eta})=0$ for every $q>0$.
\end{lemma}

Indeed, (\ref{I:3-3-2}) and the Leray spectral sequence for $j$ give
\[
  \h^q(X, j_* \dG_{m,\eta}) = \h^q(\eta,\dG_{m,\eta})
\]
for every $q\geqslant0$, and the right-hand side is zero for $q>0$ by (\ref{I:3-2-4}).

\begin{lemma}\label{I:3-3-4}
One has $\h^q\left(X,\bigoplus_{x\in X}i_{x*}\dZ_x\right)=0$ for every $q>0$.
\end{lemma}

Indeed, for every closed point $x$ of $X$, one has $\eR^qi_{x*}\dZ_x=0$ for $q>0$, because $i_x$ is a finite morphism (\ref{I:2-3-6}), and
\[
  \h^q(X,i_{x*}\dZ_x) = \h^q(x,\dZ_x)\text{.}
\]
The right-hand side is zero for every $q>0$, since $x$ is the spectrum of an algebraically closed field. (Thus the lemma holds more generally for every ``skyscraper'' sheaf on $X$.)

The preceding lemmas and the exact sequence (\ref{I:eq:3-3-1-1}) give
\[
  \h^q(X,\dG_m) = 0 \text{ for $q\geqslant 2$,}
\]
and a low-degree exact cohomology sequence
\[
  1 \to \h^0(X,\dG_m) \to \h^0(X, j_*\dG_{m,\eta}) \to \h^0\left(X,\bigoplus_{x\in X} i_{x*}\dZ_x\right) \to \h^1(X,\dG_m) \to 1
\]
which is precisely the exact sequence
\[
  1 \to k^\times \to k(X)^\times\to \Div(X) \to \pic(X) \to 1\text{.}
\]
Proposition \ref{I:3-3-1} yields the expected values for the cohomology groups of $X$ with coefficients in $\dZ/n$, when $n$ is prime to the characteristic of $k$.

\begin{corollary}\label{I:3-3-5}
If $X$ has genus $g$ and $n$ is invertible in $k$, the groups $\h^q(X,\dZ/n)$ vanish for $q>2$ and are free over $\dZ/n$ of ranks $1,2g,1$ for $q=0,1,2$, respectively. Replacing $\dZ/n$ by the isomorphic group $\dmu_n$, there are canonical isomorphisms
\begin{align*}
  \h^0(X,\dmu_n) &= \dmu_n \\
  \h^1(X,\dmu_n) &= \pic^0(X)_n \\
  \h^2(X,\dmu_n) &= \dZ/n \text{.}
\end{align*}
\end{corollary}

Since $k$ is algebraically closed, $\dZ/n$ is isomorphic (noncanonically) to $\dmu_n$. From the Kummer exact sequence
\[
  0 \to \mu_n \to \dG_m \to \dG_m \to 0\text{,}
\]
and Proposition \ref{I:3-3-1}, one obtains
\[
  \h^q(X,\dZ/n) = 0 \text{ for $q>2$,}
\]
and, in low degrees, exact sequences
\[\xymatrix{
  0 \ar[r]
    & \h^0(X,\dmu_n) \ar[r]
    & k^\times \ar[r]^-n
    & k^\times \ar[r]
    & 0
}\]
\[\xymatrix{
  0 \ar[r]
    & \h^1(X,\dmu_n) \ar[r]
    & \pic(X) \ar[r]^-n
    & \h^2(X,\dmu_n) \ar[r]
    & 0\text{.}
}\]
Moreover, there is an exact sequence
\[\xymatrix{
  0 \ar[r]
    & \pic^0(X) \ar[r]
    & \pic(X) \ar[r]^-{\deg}
    & \dZ \ar[r]
    & 0 \text{,}
}\]
and $\pic^0(X)$ is identified with the group of $k$-rational points of an abelian variety of dimension $g$, the Jacobian of $X$. In such a group, multiplication by $n$ is surjective and its kernel is a free $\dZ/n\dZ$-module of rank $2g$ (because $n$ is invertible in $k$); the corollary follows.

A clever dévissage using the ``trace method'' gives, as a consequence, the

\begin{proposition}[{\cite[IX 5.7]{sga4}}]\label{I:3-3-6}
Let $k$ be an algebraically closed field, let $X$ be an algebraic curve over $k$, and let $\sF$ be a torsion sheaf on $X$. Then:
\begin{enumerate}[\indent i)]
  \item One has $\h^q(X,\sF)=0$ for $q>2$.
  \item If $X$ is affine, one even has $\h^q(X,\sF)=0$ for $q>1$.
\end{enumerate}
\end{proposition}

For the proof, as well as for an account of the ``trace method,'' we refer to \cite[IX 5]{sga4}.

\subsection{Dévissages}\label{I:3-4}

To calculate the cohomology of varieties of dimension $>1$, one uses fibrations by curves, thereby reducing to the study of morphisms whose fibers have dimension $\leqslant1$. This principle has several variants; we mention a few of them.

\subsubsection{}\label{I:3-4-1}

Let $A$ be a finitely generated $k$-algebra and let $a_1,\dotsc,a_n$ be generators of $A$. Put $X_0=\spec(k)$, $X_i=\spec(k[a_1,\dotsc,a_i])$, and $X_n=\spec(A)$. The canonical inclusions $k[a_1,\dotsc,a_i]\to k[a_1,\dotsc,a_i,a_{i+1}]$ define morphisms $X_n\to X_{n-1}\to\cdots\to X_1\to X_0$ whose fibers have dimension $\leqslant1$.

\subsubsection{}\label{I:3-4-2}

For a smooth morphism, one can be more precise. An \emph{elementary fibration} is a morphism of schemes $f:X\to S$ that can be embedded in a commutative diagram
\[\xymatrix{
  X \ar@{^{(}->}[r]^-j \ar[dr]_-f
    & \bar X \ar[d]^-{\bar f}
    & \ar@{_{(}->}[l]_-i \ar[dl]^-g Y \\
  & S
}\]
satisfying the following conditions:
\begin{enumerate}[\indent i)]
  \item $j$ is an open immersion dense in every fiber, and $X=\bar X\setminus Y$.
  \item $\bar f$ is smooth and projective, with geometrically irreducible fibers of dimension $1$.
  \item $g$ is an étale covering and none of its fibers is empty.
\end{enumerate}

A \emph{good neighborhood} relative to $S$ is an $S$-scheme $X$ for which there are $S$-schemes $X=X_n,\dotsc,X_0=S$ and elementary fibrations $f_i:X_i\to X_{i-1}$, $i=1,\dotsc,n$. One can show \cite[XI 3.3]{sga4} that if $X$ is a smooth scheme over an algebraically closed field $k$, then every rational point of $X$ has an open neighborhood that is a good neighborhood (relative to $\spec(k)$).

\subsubsection{}\label{I:3-4-3}

A proper morphism $f:X\to S$ can be broken down as follows. By Chow's lemma, there is a commutative diagram
\[\xymatrix{
  X \ar[dr]_-f
    & \ar[l]_-\pi \ar[d]^-{\bar f} \bar X \\
  & S
}\]
where $\pi$ and $\bar f$ are projective morphisms and $\pi$ is moreover an isomorphism over a dense open subset of $X$. Locally on $S$, $\bar X$ is a closed subscheme of a projective space $\dP_S^n$.

The latter is broken down by considering the projection $\varphi:\dP_S^n\dashrightarrow\dP_S^1$ that sends the point with homogeneous coordinates $(x_0,x_1,\dotsc,x_n)$ to $(x_0,x_1)$. This is a rational map defined away from the closed subset $Y\simeq\dP_S^{n-2}$ of $\dP_S^n$ given by the homogeneous equations $x_0=x_1=0$. Let $u:P\to\dP_S^n$ be the blow-up with center $Y$; the fibers of $u$ have dimension $\leqslant1$. Moreover, there is a natural morphism $v:P\to\dP_S^1$ extending the rational map $\varphi$, and $v$ makes $P$ a $\dP_S^1$-scheme locally isomorphic to the standard projective space $\dP^{n-1}$, which can in turn be projected onto a $\dP^1$, and so on.

\subsubsection{}\label{I:3-4-4}

A smooth projective variety $X$ can be swept out by a \emph{Lefschetz pencil}. The blow-up $\widetilde X$ of the intersection of the axis of the pencil with $X$ maps to $\dP^1$, and the fibers of this map are the hyperplane sections of $X$ cut out by the hyperplanes in the pencil.

\section{The base change theorem for a proper morphism}\label{I:4}

\subsection{Introduction}\label{I:4-1}

This section is devoted to the proof and applications of the following theorem.

\begin{theorem}\label{I:4-1-1}
Let $f:X\to S$ be a proper morphism of schemes and let $\sF$ be an abelian torsion sheaf on $X$. For every $q\geqslant0$, the stalk of $\eR^qf_*\sF$ at a geometric point $s$ of $S$ is isomorphic to the cohomology $\h^q(X_s,\sF)$ of the fiber $X_s=X\otimes_S\spec k(s)$ of $f$ at $s$.
\end{theorem}

For a proper separated continuous map $f:X\to S$ between topological spaces (separated means that the diagonal in $X\times_SX$ is closed), and an abelian sheaf $\sF$ on $X$, the analogous result is well known and elementary. Since $f$ is closed, the sets $f^{-1}(V)$, for $V$ a neighborhood of $s$, form a fundamental system of neighborhoods of $X_s$, and one checks that $\h^\bullet(X_s,\sF)=\varinjlim_U\h^\bullet(U,\sF)$, where $U$ ranges over the neighborhoods of $X_s$. In practice, $X_s$ even has a fundamental system $\cU$ of neighborhoods $U$ of which it is a deformation retract, and for constant $\sF$ one thus has $\h^\bullet(X_s,\sF)=\h^\bullet(U,\sF)$. Figuratively speaking, the special fiber swallows the general fiber.

For schemes the proof is more delicate, and it is essential to assume that $\sF$ is torsion (\cite[XII.2]{sga4}). In view of the description of the stalks of $\eR^qf_*\sF$ (\ref{I:2-3-3}), Theorem (\ref{I:4-1-1}) is essentially equivalent to the following theorem.

\begin{theorem}\label{I:4-1-2}
Let $A$ be a strictly Henselian local ring and put $S=\spec(A)$. Let $f:X\to S$ be a proper morphism and let $X_0$ be the closed fiber of $f$. Then, for every abelian torsion sheaf $\sF$ on $X$ and every $q\geqslant0$, one has $\h^q(X,\sF)\iso\h^q(X_0,\sF)$.
\end{theorem}

By passage to the limit, it is enough to prove the theorem when $A$ is the strict Henselization of a finitely generated $\dZ$-algebra at a prime ideal. We first treat the case $q=0$ or $1$ and $\sF=\dZ/n$ (\ref{I:4-2}). An argument based on the notion of a constructible sheaf (\ref{I:4-3}) moreover shows that it is enough to consider the case in which $\sF$ is constant. On the other hand, the dévissage (\ref{I:3-4-3}) allows us to suppose that $X_0$ is a curve; in that case it remains only to prove the theorem for $q=2$ (\ref{I:4-4}).

Among other applications (\ref{I:4-6}), the theorem permits the definition of cohomology with proper support (\ref{I:4-5}).

\subsection{Proof for \texorpdfstring{$q=0$}{q=0} or \texorpdfstring{$1$}{1} and \texorpdfstring{$\sF=\dZ/n$}{\sF=Z/n}}\label{I:4-2}

The result for $q=0$ and constant $\sF$ is equivalent to the following proposition (Zariski's connectedness theorem).

\begin{proposition}\label{I:4-2-1}
Let $A$ be a Noetherian Henselian local ring and put $S=\spec(A)$. Let $f:X\to S$ be a proper morphism and let $X_0$ be its closed fiber. Then the sets of connected components $\pi_0(X)$ and $\pi_0(X_0)$ are in bijection.
\end{proposition}

Equivalently, it is enough to show that the sets $\operatorname{Of}(X)$ and $\operatorname{Of}(X_0)$ of subsets that are both open and closed are in bijection. The set $\operatorname{Of}(X)$ is in bijection with the set of idempotents of $\Gamma(X,\sO_X)$, and similarly $\operatorname{Of}(X_0)$ is in bijection with the set of idempotents of $\Gamma(X_0,\sO_{X_0})$. Thus we must show that the canonical map
\[
  \operatorname{Idem} \Gamma(X,\sO_X) \to \operatorname{Idem}\Gamma(X_0,\sO_{X_0})
\]
is bijective.

Let $\fm$ denote the maximal ideal of $A$, let $\Gamma(X,\sO_X)^\wedge$ be the completion of $\Gamma(X,\sO_X)$ for the $\fm$-adic topology, and, for every integer $n\geqslant0$, put $X_n=X\otimes_AA/\fm^{n+1}$. By the finiteness theorem for proper morphisms \cite[III.3.2]{ega}, $\Gamma(X,\sO_X)$ is a finite $A$-algebra; since $A$ is Henselian, it follows that the canonical map
\[
  \operatorname{Idem} \Gamma(X,\sO_X) \to \operatorname{Idem}\Gamma(X,\sO_X)^\wedge
\]
is bijective. In particular, the canonical map
\[
  \operatorname{Idem} \Gamma(X,\sO_X)^\wedge \to \varprojlim \operatorname{Idem}\Gamma(X_n,\sO_{X_n})
\]
is bijective. But since $X_n$ and $X_0$ have the same underlying topological space, the canonical map
\[
  \operatorname{Idem}\Gamma(X_n,\sO_{X_n}) \to \operatorname{Idem} \Gamma(X_0,\sO_{X_0})
\]
is bijective for every $n$, which completes the proof.

Since $\h^1(X,\dZ/n)$ is in bijection with the set of isomorphism classes of Galois étale coverings of $X$ with group $\dZ/n$, the theorem for $q=1$ and $\sF=\dZ/n$ follows from the next proposition.

\begin{proposition}\label{I:4-2-2}
Let $A$ be a Noetherian Henselian local ring and put $S=\spec(A)$. Let $f:X\to S$ be a proper morphism and let $X_0$ be its closed fiber. Then the restriction functor
\[
  \operatorname{Rev.et}(X) \to \operatorname{Rev.et}(X_0)
\]
is an equivalence of categories.
\end{proposition}

(If $X_0$ is connected and a geometric point of $X_0$ has been chosen as base point, this amounts to saying that the canonical map $\pi_1(X_0)\to\pi_1(X)$ on (profinite) fundamental groups is bijective.)

Proposition (\ref{I:4-2-1}) shows that this functor is fully faithful. Indeed, if $X'$ and $X''$ are two étale coverings of $X$, an $X$-morphism from $X'$ to $X''$ is determined by its graph, which is an open-and-closed subset of $X'\times_XX''$.

It remains to show that every étale covering $X'_0$ of $X_0$ extends to an étale covering of $X$. Étale coverings are insensitive to nilpotent elements \cite[ch.1]{sga1}; consequently $X'_0$ lifts uniquely to an étale covering $X'_n$ of $X_n$ for every $n\geqslant0$, that is, to an étale covering $\fX'$ of the formal scheme $\fX$ obtained by completing $X$ along $X_0$. By Grothendieck's algebraization theorem for formal coherent sheaves (the existence theorem, \cite[III.5]{ega}), $\fX'$ is the formal completion of an étale covering $\bar X'$ of $\bar X=X\otimes_A\hat A$.

By passage to the limit, it is enough to prove the proposition when $A$ is the Henselization of a finitely generated $\dZ$-algebra. We can then apply Artin's approximation theorem to the functor $\sF:(\textnormal{$A$-algebras})\to(\textnormal{sets})$ that assigns to an $A$-algebra $B$ the set of isomorphism classes of étale coverings of $X\otimes_AB$. Indeed, this functor is locally of finite presentation: if $(B_i)$ is a filtered direct system of $A$-algebras and $B=\varinjlim B_i$, then $\sF(B)=\varinjlim\sF(B_i)$. By Artin's theorem, given an element $\bar\xi\in\sF(\hat A)$, here the isomorphism class of $\bar X'$, there is an element $\xi\in\sF(A)$ having the same image as $\bar\xi$ in $\sF(A/\fm)$. In other words, there is an étale covering $X'$ of $X$ whose restriction to $X_0$ is isomorphic to $X'_0$.

\subsection{Constructible sheaves}\label{I:4-3}

In this subsection, $X$ is a \emph{Noetherian} scheme, and by a sheaf on $X$ we mean an \emph{abelian} sheaf on $\et X$.

\begin{definition}\label{I:4-3-1}
A sheaf $\sF$ on $X$ is said to be \emph{locally constant constructible} (abbreviated l.c.c.) if it is represented by an étale covering of $X$.
\end{definition}

\begin{definition}\label{I:4-3-2}
A sheaf $\sF$ on $X$ is said to be \emph{constructible} if it satisfies the following equivalent conditions:
\begin{enumerate}[\indent (i)]
  \item There is a finite surjective family of subschemes $X_i$ of $X$ such that the restriction of $\sF$ to each $X_i$ is l.c.c.
  \item There are a finite family of finite morphisms $p_i:X_i'\to X$, for each $i$ a constructible constant sheaf $C_i$ (that is, one defined by a finite abelian group) on $X_i'$, and a monomorphism $\sF\to\prod p_{i*}C_i$.
\end{enumerate}
\end{definition}

It is easy to check that the category of constructible sheaves on $X$ is an abelian category. Moreover, if $u:\sF\to\sG$ is a homomorphism of sheaves and $\sF$ is constructible, then the sheaf $\im(u)$ is constructible.

\begin{lemma}\label{I:4-3-3}
Every torsion sheaf $\sF$ is a filtered direct limit of constructible sheaves.
\end{lemma}

Indeed, if $j:U\to X$ is an étale scheme of finite type over $X$, an element $\xi\in\sF(U)$ such that $n\xi=0$ defines a homomorphism of sheaves $j_!\const{\dZ/n}\to\sF$ whose image (the smallest subsheaf of $\sF$ of which $\xi$ is a local section) is a constructible subsheaf of $\sF$. It is clear that $\sF$ is the direct limit of such subsheaves.

\begin{definition}\label{I:4-3-4}
Let $\cC$ be an abelian category and let $T$ be a functor from $\cC$ to the category of abelian groups. The functor $T$ is said to be \emph{effaceable} in $\cC$ if, for every object $A$ of $\cC$ and every $\alpha\in T(A)$, there is a monomorphism $u:A\to M$ in $\cC$ such that $T(u)\alpha=0$.
\end{definition}

\begin{lemma}\label{I:4-3-5}
The functors $\h^q(X,-)$ for $q>0$ are effaceable in the category of constructible sheaves on $X$.
\end{lemma}

It is enough to observe that, if $\sF$ is a constructible sheaf, there is necessarily an integer $n>0$ such that $\sF$ is a sheaf of $\dZ/n$-modules. There is then a monomorphism $\sF\hookrightarrow\sG$, where $\sG$ is a sheaf of $\dZ/n$-modules and $\h^q(X,\sG)=0$ for every $q>0$. For example, one may take for $\sG$ the Godement resolution $\prod_{x\in X}i_{x*}\sF_{\bar x}$, where $x$ ranges over the points of $X$ and $i_x:\bar x\to X$ is a geometric point centered at $x$. By (\ref{I:4-3-3}), $\sG$ is a direct limit of constructible sheaves. The lemma follows, since the functors $\h^q(X,-)$ commute with direct limits.

\begin{lemma}\label{I:4-3-6}
Let $\varphi^\bullet:T^\bullet\to{T'}^\bullet$ be a morphism of cohomological functors from an abelian category $\cC$ to the category of abelian groups. Suppose that $T^q$ is effaceable for $q>0$, and let $\cE$ be a collection of objects of $\cC$ such that every object of $\cC$ is contained in an object of $\cE$. Then the following conditions are equivalent:
\begin{enumerate}[\indent (i)]
  \item $\varphi^q(A)$ is bijective for every $q\geqslant0$ and every $A\in\ob\cC$.
  \item $\varphi^0(M)$ is bijective and $\varphi^q(M)$ is surjective for every $q>0$ and every $M\in\cE$.
  \item $\varphi^0(A)$ is bijective for every $A\in\ob\cC$, and ${T'}^q$ is effaceable for every $q>0$.
\end{enumerate}
\end{lemma}

The proof is by induction on $q$ and presents no difficulty.

\begin{proposition}\label{I:4-3-7}
Let $X_0$ be a subscheme of $X$. Suppose that, for every $n\geqslant0$ and every scheme $X'$ finite over $X$, the canonical map
\[
  \h^q(X',\dZ/n) \to \h^q(X_0',\dZ/n)\text{,}
\]
where $X'_0=X'\times_XX_0$, is bijective for $q=0$ and surjective for $q>0$. Then, for every torsion sheaf $\sF$ on $X$ and every $q\geqslant0$, the canonical map
\[
  \h^q(X,\sF) \to \h^q(X_0,\sF)
\]
is bijective.
\end{proposition}

By passage to the limit, it is enough to prove the assertion for constructible $\sF$. Apply Lemma (\ref{I:4-3-6}) with $\cC$ the category of constructible sheaves on $X$, $T^q=\h^q(X,-)$, ${T'}^q=\h^q(X_0,-)$, and $\cE$ the collection of constructible sheaves of the form $\prod p_{i*}C_i$, where $p_i:X'_i\to X$ is a finite morphism and $C_i$ is a finite constant sheaf on $X'_i$.

\subsection{Completion of the proof}\label{I:4-4}

By the method of fibration by curves (\ref{I:3-4-3}), the theorem is reduced to relative dimension $\leqslant1$. By the preceding subsection, it is enough to show the following: if $S$ is the spectrum of a strictly Henselian Noetherian local ring, $f:X\to S$ is a proper morphism whose closed fiber $X_0$ has dimension $\leqslant1$, and $n$ is an integer $\geqslant0$, then the canonical homomorphism
\[
  \h^q(X,\dZ/n) \to \h^q(X_0,\dZ/n)
\]
is bijective for $q=0$ and surjective for $q>0$.

The cases $q=0$ and $1$ were treated above, and $\h^q(X_0,\dZ/n)=0$ for $q\geqslant3$; it is therefore enough to treat the case $q=2$. We may clearly suppose that $n$ is a power of a prime. If $n=p^r$, where $p$ is the characteristic of the residue field of $S$, Artin--Schreier theory shows that $\h^2(X_0,\dZ/p^r)=0$. If $n=\ell^r$, with $\ell\ne p$, Kummer theory gives a commutative diagram
\[\xymatrix{
  \pic(X) \ar[d] \ar[r]^-\alpha
    & \h^2(X,\dZ/\ell^r) \ar[d] \\
  \pic(X_0) \ar[r]^-\beta
    & \h^2(X_0,\dZ/\ell^r)
}\]
where $\beta$ is surjective. (This was shown in Section \ref{I:3} for a smooth curve over an algebraically closed field, but similar arguments apply to any curve over a separably closed field.)

To conclude, it is therefore enough to prove the following proposition.

\begin{proposition}\label{I:4-4-1}
Let $S$ be the spectrum of a Henselian Noetherian local ring, and let $f:X\to S$ be a proper morphism whose closed fiber $X_0$ has dimension $\leqslant1$. Then the canonical restriction map
\[
  \pic(X)\to \pic(X_0)
\]
is surjective. (It is in fact enough for $f$ to be \emph{separated} and of finite type.)
\end{proposition}

To simplify the proof, we shall suppose that $X$ is integral, although this is not necessary. Every invertible sheaf on $X_0$ is associated with a Cartier divisor (because $X_0$ is a curve and hence quasi-projective), so it is enough to show that the canonical map $\Div(X)\to\Div(X_0)$ is surjective.

Every divisor on $X_0$ is a linear combination of divisors whose support is concentrated at a single nonisolated closed point of $X_0$. Let $x$ be such a point, let $t_0\in\sO_{X_0,x}$ be a regular nonunit of $\sO_{X_0,x}$, and let $D_0$ be the divisor concentrated at $x$ with local equation $t_0$. Let $U$ be an open neighborhood of $x$ in $X$ on which there is a section $t\in\Gamma(U,\sO_U)$ lifting $t_0$. Let $Y$ be the closed subset of $U$ defined by $t=0$; after shrinking $U$, we may suppose that $x$ is the only point of $Y\cap X_0$. Then $Y$ is quasi-finite over $S$ at $x$. Since $S$ is the spectrum of a Henselian local ring, it follows that $Y=Y_1\amalg Y_2$, where $Y_1$ is finite over $S$ and $Y_2$ does not meet $X_0$. Moreover, since $X$ is separated over $S$, $Y_1$ is closed in $X$.

After replacing $U$ by a smaller open neighborhood of $x$, we may suppose that $Y=Y_1$, in other words, that $Y$ is closed in $X$. We then define a divisor $D$ on $X$ lifting $D_0$ by setting $D|X\setminus Y=0$ and $D|U=\dv(t)$; this makes sense because $t$ is invertible on $U\setminus Y$.

\subsubsection{Remark}\label{I:4-4-2}

When $f$ is proper, one could also give a proof in the same style as that of Proposition (\ref{I:4-2-2}). Indeed, since $X_0$ is a curve, there is no obstruction to lifting an invertible sheaf on $X_0$ to the infinitesimal neighborhoods $X_n$ of $X_0$, and hence to the formal completion $\fX$ of $X$ along $X_0$. One then concludes by applying successively Grothendieck's existence theorem and Artin's approximation theorem.

\subsection{Cohomology with proper support}\label{I:4-5}

\begin{definition}\label{I:4-5-1}
Let $X$ be a separated scheme of finite type over a field $k$. By a theorem of Nagata, there are a scheme $\bar X$ proper over $k$ and an open immersion $j:X\to\bar X$. For every torsion sheaf $\sF$ on $X$, let $j_!\sF$ denote the extension by zero of $\sF$ to $\bar X$, and define the groups of \emph{cohomology with proper support} by
\[
  \h_c^q(X,\sF)=\h^q(\bar X, j_! \sF) \text{.}
\]
\end{definition}

Let us show that this definition is independent of the chosen compactification $j:X\to\bar X$. Let $j_1:X\to\bar X_1$ and $j_2:X\to\bar X_2$ be two compactifications. The scheme $X$ maps to $\bar X_1\times\bar X_2$ by $x\mapsto(j_1(x),j_2(x))$, and the closure $\bar X_3$ of the image of $X$ under this map is a compactification of $X$. Thus there is a commutative diagram
\[\xymatrix{
  & \bar X_1 \\
  X \ar[ur]^-{j_1} \ar[rr]^-{(j_1,j_2)} \ar[dr]_-{j_2}
    & & \bar X_3 \ar[ul]_-{p_1} \ar[dl]^-{p_2} \\
  & \bar X_2
}\]
where $p_1$ and $p_2$, the restrictions to $\bar X_3$ of the natural projections, are proper morphisms. It is therefore enough to treat the case of a commutative diagram
\[\xymatrix{
  X \ar[r]^-{j_2} \ar[dr]_-{j_1}
    & \bar X_2 \ar[d]^-p \\
  & \bar X_1
}\]
with $p$ proper.

\begin{lemma}\label{I:4-5-2}
One has $p_*(j_{2!}\sF)=j_{1!}\sF$ and $\eR^qp_*(j_{2!}\sF)=0$ for $q>0$.
\end{lemma}

This lemma immediately gives the conclusion. Using the Leray spectral sequence for $p$, one deduces that, for every $q\geqslant0$,
\[
  \h^q(\bar X_2, j_{2!} \sF) = \h^q(\bar X_1, j_{1!} \sF) \text{.}
\]

To prove the lemma, argue fiber by fiber using the base change theorem (\ref{I:4-1-1}) for $p$. The result is immediate: over a point of $X$, $p$ is an isomorphism, while over a point of $\bar X_1\setminus X$, the sheaf $j_{2!}\sF$ is zero on the fiber of $p$.

\subsubsection{}\label{I:4-5-3}

Similarly, if $f:X\to S$ is a separated morphism of finite type between Noetherian schemes, there are a proper morphism $\bar f:\bar X\to S$ and an open immersion $j:X\to\bar X$. For every torsion sheaf $\sF$ on $X$, define the higher direct images with proper support $\eR^qf_!$ by
\[
  \eR^q f_! \sF = \eR^q \bar f_*(j_! \sF) \text{.}
\]
As before, one checks that this definition is independent of the chosen compactification.

\begin{theorem}\label{I:4-5-4}
Let $f:X\to S$ be a separated morphism of finite type between Noetherian schemes, and let $\sF$ be a torsion sheaf on $X$. Then the stalk of $\eR^qf_!\sF$ at a geometric point $s$ of $S$ is isomorphic to the cohomology with proper support $\h_c^q(X_s,\sF)$ of the fiber $X_s$ of $f$ at $s$.
\end{theorem}

This is a straightforward variant of the base change theorem for a proper morphism (\ref{I:4-1-1}). More generally, if
\[\xymatrix{
  X \ar[d]^-f
    & X' \ar[l]_-{g'} \ar[d]^-{f'} \\
  S
    & S' \ar[l]^-g
}\]
is a Cartesian diagram, there is a canonical isomorphism
\begin{equation*}\tag{4.5.4.1}\label{I:eq:4-5-4-1}
  g^* (\eR^q f_! \sF) \simeq \eR^q f_!'({g'}^* \sF) \text{.}
\end{equation*}

\subsection{Applications}\label{I:4-6}

\begin{theorem}[Vanishing theorem]\label{I:4-6-1}
Let $f:X\to S$ be a separated morphism of finite type whose fibers have dimension $\leqslant n$, and let $\sF$ be a torsion sheaf on $X$. Then $\eR^qf_!\sF=0$ for $q>2n$.
\end{theorem}

By the base change theorem, we may suppose that $S$ is the spectrum of a separably closed field. If $\dim X=n$, there is an affine open subset $U$ of $X$ such that $\dim(X\setminus U)<n$; there is then an exact sequence $0\to\sF_U\to\sF\to\sF_{X\setminus U}\to0$, and induction on $n$ reduces the theorem to the case $X=U$ affine. The method of fibration by curves (\ref{I:3-4-1}) and the base change theorem then reduce matters to a curve over a separably closed field, for which the desired result follows from Tsen's theorem (\ref{I:3-3-6}).

\begin{theorem}[Finiteness theorem]\label{I:4-6-2}
Let $f:X\to S$ be a separated morphism of finite type and let $\sF$ be a constructible sheaf on $X$. Then the sheaves $\eR^qf_!\sF$ are constructible.
\end{theorem}

We consider only the case in which $\sF$ is annihilated by an integer invertible on $X$.

To prove the theorem, one reduces to the case in which $\sF$ is the constant sheaf $\const{\dZ/n}$ and $f:X\to S$ is a smooth proper morphism whose fibers are geometrically connected curves of genus $g$. When $n$ is invertible on $X$, the sheaves $\eR^qf_*\sF$ are then locally free of finite rank and vanish for $q>2$ (\ref{I:4-6-1}). Replacing $\dZ/n$ by the sheaf $\dmu_n$, which is locally isomorphic to it (on $S$), one has canonically
\begin{equation*}\tag{4.6.2.1}\label{I:eq:4-6-2-1}
\begin{aligned}
  \eR^0 f_* \dmu_n &= \dmu_n \\
  \eR^1 f_* \dmu_n &= \const\pic(X/S)_n \\
  \eR^2 f_* \dmu_n &= \dZ/n \text{.}
\end{aligned}
\end{equation*}

\begin{theorem}[Comparison with classical cohomology]\label{I:4-6-3}
Let $f:X\to S$ be a separated morphism of schemes of finite type over $\dC$, and let $\sF$ be a torsion sheaf on $X$. Denote by a superscript $\an{(-)}$ the functor of passage to the usual topological spaces, and by $\eR^q\an f_!$ the derived functors of direct image with proper support under $\an f$. Then
\[
  \an{(\eR^q f_! \sF)} \simeq \eR^q \an f_! \an \sF \text{.}
\]
In particular, when $S$ is a point and $\sF$ is the constant sheaf $\dZ/n$,
\[
  \h_c^q(X,\dZ/n) \simeq \h_c^q(\an X,\dZ/n)\text{.}
\]
\end{theorem}

Dévissages using the base change theorem reduce us to the case in which $X$ is a smooth proper curve, $S$ is a point, and $\sF=\dZ/n$. The cohomology groups under consideration then vanish for $q\ne0,1,2$, and one invokes \cite{se55}. Indeed, if $X$ is proper over $\dC$, then $\pi_0(X)=\pi_0(\an X)$ and $\pi_1(X)$ is the profinite completion of $\pi_1(\an X)$, which gives the assertion for $q=0,1$. For $q=2$, one uses the Kummer exact sequence and the fact, again from \cite{se55}, that $\pic(X)=\pic(\an X)$.

\begin{theorem}[Cohomological dimension of affine schemes]\label{I:4-6-4}
Let $X$ be an affine scheme of finite type over a separably closed field and let $\sF$ be a torsion sheaf on $X$. Then $\h^q(X,\sF)=0$ for $q>\dim(X)$.
\end{theorem}

For the very elegant proof, we refer to \cite[XIV, \S\S 2--3]{sga4}.

\subsubsection{Remark}\label{I:4-6-5}

This theorem is, in a sense, a substitute for Morse theory. Indeed, consider the classical case in which $X$ is smooth and affine over $\dC$, embedded in the standard affine space $\dC^N$. For almost every point $p\in\dC^N$, the ``distance from $p$'' function on $X$ is a Morse function, and the indices of its critical points are less than $\dim(X)$. Thus $X$ is obtained by attaching handles of index less than $\dim(X)$, giving the classical analogue of (\ref{I:4-6-4}).

\section{Local acyclicity of smooth morphisms}\label{I:5}

Let $X$ be a complex analytic variety and let $f:X\to D$ be a morphism from $X$ to the disk. Let $[0,t]$ denote the closed line segment in $D$ with endpoints $0$ and $t$, and let $(0,t]$ denote the half-open segment. If $f$ is \emph{smooth}, the inclusion
\[
  j : f^{-1}\left((0,t]\right) \hookrightarrow f^{-1}\left([0,t]\right)
\]
is a homotopy equivalence; one can push the special fiber $X_0=f^{-1}(0)$ inside $f^{-1}([0,t])$.

In practice, for sufficiently small $t$, $f^{-1}((0,t])$ will be a fiber bundle over $(0,t]$, so the inclusion
\[
  X_t = f^{-1}(t) \hookrightarrow f^{-1}\left((0,t]\right)
\]
will also be a homotopy equivalence. The homotopy class of maps
\[
  \cosp : X_0 \hookrightarrow f^{-1}\left([0,t]\right) \xleftarrow\sim f^{-1}\left((0,t]\right) \xleftarrow\sim X_t \text{.}
\]
is then called the \emph{cospecialization morphism}. Figuratively speaking, for a smooth morphism the general fiber swallows the special fiber.

No longer suppose that $f$ is necessarily smooth (but suppose that $f^{-1}((0,t])$ is a fiber bundle over $(0,t]$). A morphism $\cosp^\bullet$ in cohomology can still be defined as soon as $j_*\dZ=\dZ$ and $\eR^qj_*\dZ=0$ for $q>0$. Under these hypotheses, the Leray spectral sequence for $j$ shows that
\[
  \h^\bullet\left(f^{-1}\left([0,t]\right),\dZ\right) \iso \h^\bullet\left(f^{-1}\left((0,t]\right),\dZ\right)
\]
and $\cosp^\bullet$ is the composite morphism
\[
  \cosp^\bullet : \h^\bullet(X_t,\dZ)
    \xleftarrow\sim \h^\bullet\left(f^{-1}\left((0,t]\right),\dZ\right)
    \xleftarrow\sim \h^\bullet\left(f^{-1}\left([0,t]\right),\dZ\right)
    \to \h^\bullet(X_0,\dZ)
\]
The stalk of $\eR^qj_*\dZ$ at a point $x\in X_0$ is calculated as follows. In an ambient space, take a ball $B_\varepsilon$ centered at $x$ with sufficiently small radius $\varepsilon$, and, for sufficiently small $\eta$, put $E=X\cap B_\varepsilon\cap f^{-1}(\eta t)$; this is the \emph{vanishing cycle variety} at $x$. One has
\[
  (\eR^q j_* \dZ)_x \xleftarrow\sim \h^q\left(X\cap B_\varepsilon\cap f^{-1}\left((0,\eta t]\right),\dZ\right) \xrightarrow\sim \h^q(E,\dZ)
\]
and the cospecialization morphism is defined in cohomology whenever the vanishing cycle varieties are acyclic ($\h^0(E,\dZ)=\dZ$ and $\h^q(E,\dZ)=0$ for $q>0$); this is expressed by saying that $f$ is \emph{locally acyclic}.

This section is devoted to the analogue of this situation for a smooth morphism of schemes in étale cohomology. Here, however, it is essential to restrict to torsion coefficients whose orders are \emph{prime to the residue characteristics}. Subsection \ref{I:5-1} gives general facts about locally acyclic morphisms and cospecialization maps. In Subsection \ref{I:5-2} we prove that a smooth morphism is locally acyclic. In Subsection \ref{I:5-3} we combine this with the results of the preceding section to deduce two applications: a specialization theorem for cohomology groups (the cohomology of the geometric fibers of a smooth proper morphism is locally constant) and a base change theorem for a smooth morphism.

Throughout what follows, an integer $n$ is fixed, and ``scheme'' means a ``scheme on which $n$ is invertible.'' A ``geometric point'' will always mean an ``algebraic geometric point'' (\ref{I:2-3-1}) $x:\spec(k)\to X$, with $k$ algebraically closed.

\subsection{Locally acyclic morphisms}\label{I:5-1}

\subsubsection{Notation}\label{I:5-1-1}

Given a scheme $S$ and a geometric point $s$ of $S$, let $\widetilde S^s$ denote the spectrum of the strict localization of $S$ at $s$.

\begin{definition}\label{I:5-1-2}
A geometric point $t$ of $S$ is called a \emph{generization} of $s$ if it is defined by an algebraic closure of the residue field of a point of $\widetilde S^s$. One also says that $s$ is a \emph{specialization} of $t$, and the $S$-morphism $t\to\widetilde S^s$ is called a \emph{specialization map}.
\end{definition}

\begin{definition}\label{I:5-1-3}
Let $f:X\to S$ be a morphism of schemes. Let $s$ be a geometric point of $S$, let $t$ be a generization of $s$, let $x$ be a geometric point of $X$ over $s$, and put $\widetilde X_t^x=\widetilde X^x\times_{\widetilde S^s}t$. The scheme $\widetilde X_t^x$ is called a \emph{vanishing cycle variety} of $f$ at $x$.

The morphism $f$ is said to be \emph{locally acyclic} if the reduced cohomology of every vanishing cycle variety $\widetilde X_t^x$ vanishes:
\begin{equation*}\tag{5.1.3.1}\label{I:eq:5-1-3-1}
  \hat\h^\bullet(\widetilde X_t^x,\dZ/n) = 0\text{,}
\end{equation*}
that is, $\h^0\left(\widetilde X_t^x,\dZ/n\right)=\dZ/n$ and $\h^q(\widetilde X_t^x,\dZ/n)=0$ for $q>0$.
\end{definition}

\begin{lemma}\label{I:5-1-4}
Let $f:X\to S$ be a locally acyclic morphism and let $g:S'\to S$ be a quasi-finite morphism (or an inverse limit of quasi-finite morphisms). Then the morphism $f':X'\to S'$ obtained from $f$ by base change is locally acyclic.
\end{lemma}

Indeed, every vanishing cycle variety of $f'$ is a vanishing cycle variety of $f$.

\begin{lemma}\label{I:5-1-5}
Let $f:X\to S$ be a locally acyclic morphism. For every geometric point $t$ of $S$, giving rise to a Cartesian diagram
\[\xymatrix{
  X_t \ar[r]^-{\varepsilon'} \ar[d]
    & X \ar[d]^-f \\
  t \ar[r]^-\varepsilon
    & S
}\]
one has $\varepsilon'_*\dZ/n=f^*\varepsilon_*\dZ/n$ and $\eR^q\varepsilon'_*\dZ/n=0$ for $q>0$.
\end{lemma}

Let $\bar S$ be the closure of $\varepsilon(t)$, let $S'$ be the normalization of $\bar S$ in $k(t)$, and consider the Cartesian diagram
\[\xymatrix{
  X_t \ar[r]^-{f'} \ar[d]
    & X' \ar[r]^-{\alpha'} \ar[d]^-{f'}
    & X \ar[d]^-f \\
  t \ar[r]^-i
    & S' \ar[r]^-\alpha
    & S
}\]
The local rings of $S'$ are normal with separably closed fraction fields. They are therefore strictly Henselian, and the local acyclicity of $f'$ (\ref{I:5-1-4}) gives $f'_*\dZ/n=\dZ/n$ and $\eR^qi'_*\dZ/n=0$ for $q>0$. Since $\alpha$ is integral, one then has
\[
  \eR^q\varepsilon_*'\dZ/n
    = \alpha_*' \eR^q i_*' \dZ/n
    = \alpha_*' {f'}^* \eR^q i_* \dZ/n
    = f^* \alpha_* \eR^q i_* \dZ/n
    = f^* \eR^q \varepsilon_* \dZ/n \text{,}
\]
which proves the lemma.

\subsubsection{}\label{I:5-1-6}

Given a locally acyclic morphism $f:X\to S$ and a specialization map $t\to\widetilde S^s$, we shall define canonical homomorphisms, called cospecialization maps,
\[
  \cosp^\bullet : \h^\bullet(X_t,\dZ/n) \to \h^\bullet(X_s,\dZ/n)\text{,}
\]
relating the cohomology of the general fiber $X_t=X\times_St$ to that of the special fiber $X_s=X\times_Ss$.

Consider the Cartesian diagram
\[\xymatrix{
  X_t \ar[d] \ar[r]^-{\varepsilon'}
    & \widetilde X \ar[d]^-{f'}
    & \ar[l] X_s \ar[d] \\
  t \ar[r]^-\varepsilon
    & \widetilde S^s
    & \ar[l] s
}\]
obtained from $f$ by base change. By (\ref{I:5-1-4}), $f'$ is again locally acyclic. The definition of local acyclicity immediately shows that the restriction to $X_s$ of the sheaves $\eR^q\varepsilon'_*\dZ/n$ is $\dZ/n$ for $q=0$ and $0$ for $q>0$. By (\ref{I:5-1-5}), one even knows that $\eR^q\varepsilon'_*\dZ/n=0$ for $q>0$. Define $\cosp^\bullet$ to be the composite map
\begin{equation*}\tag{5.1.6.1}\label{I:eq:5-1-6-1}
  \h^\bullet(X_t,\dZ/n) \simeq \h^\bullet(X,\varepsilon_*' \dZ/n) \to \h^\bullet(X_s,\dZ/n) \text{.}
\end{equation*}

\paragraph{Variant}
Let $\bar S$ be the closure of $\varepsilon(t)$ in $\widetilde S^s$, let $S'$ be the normalization of $\bar S$ in $k(t)$, and let $X'/S'$ be obtained from $X/S$ by base change. The map \eqref{I:eq:5-1-6-1} can also be written
\[
  \h^\bullet(X_t,\dZ/n) \simeq \h^\bullet(X',\dZ/n) \to \h^\bullet(X_s,\dZ/n) \text{.}
\]

\begin{theorem}\label{I:5-1-7}
Let $S$ be a locally Noetherian scheme, let $s$ be a geometric point of $S$, and let $f:X\to S$ be a morphism. Suppose that
\begin{enumerate}[\indent a)]
  \item the morphism $f$ is locally acyclic;
  \item for every specialization map $t\to\widetilde S^s$ and every $q\geqslant0$, the cospecialization maps $\h^q(X_t,\dZ/n)\to\h^q(X_s,\dZ/n)$ are bijective.
\end{enumerate}

Then the canonical homomorphism $(\eR^qf_*\dZ/n)_s\to\h^q(X_s,\dZ/n)$ is bijective for every $q\geqslant0$.
\end{theorem}

To prove the theorem, we may clearly suppose that $S=\widetilde S^s$. We shall in fact show that, for every sheaf of $\dZ/n\dZ$-modules $\sF$ on $S$, the canonical homomorphism $\varphi^q(\sF):(\eR^qf_*f^*\sF)_s\to\h^q(X_s,f^*\sF)$ is bijective.

Every sheaf of $\dZ/n\dZ$-modules is a filtered direct limit of constructible sheaves of $\dZ/n\dZ$-modules (\ref{I:4-3-3}). Moreover, every constructible sheaf of $\dZ/n\dZ$-modules embeds in a sheaf of the form $\prod i_{\lambda*}C_\lambda$, where the $i_\lambda:t_\lambda\to S$ form a finite family of generizations of $s$, and $C_\lambda$ is a free $\dZ/n\dZ$-module of finite rank on $t_\lambda$. By the definition of cospecialization maps, condition b) says that the homomorphisms $\varphi^q(\sF)$ are bijective when $\sF$ is of this form.

We conclude using the following variant of Lemma (\ref{I:4-3-6}).

\begin{lemma}\label{I:5-1-8}
Let $\cC$ be an abelian category in which filtered direct limits exist. Let $\varphi^\bullet:T^\bullet\to{T'}^\bullet$ be a morphism of cohomological functors from $\cC$ to the category of abelian groups, both commuting with filtered direct limits. Suppose there are two collections $\cD$ and $\cE$ of objects of $\cC$ such that:
\begin{enumerate}[\indent a)]
  \item every object of $\cC$ is a filtered direct limit of objects belonging to $\cD$;
  \item every object belonging to $\cD$ is contained in an object belonging to $\cE$.
\end{enumerate}

Then the following conditions are equivalent:
\begin{enumerate}[\indent (i)]
  \item $\varphi^q(A)$ is bijective for every $q\geqslant0$ and every $A\in\ob\cC$.
  \item $\varphi^q(M)$ is bijective for every $q\geqslant0$ and every $M\in\cE$.
\end{enumerate}
\end{lemma}

The lemma is proved by passage to the direct limit, induction on $q$, and repeated application of the five lemma to the diagram of exact cohomology sequences arising from an exact sequence $0\to A\to M\to A'\to0$, with $A\in\cD$, $M\in\cE$, and $A'\in\ob\cC$.

\begin{corollary}\label{I:5-1-9}
Let $S$ be the spectrum of a strictly Henselian Noetherian local ring and let $f:X\to S$ be a locally acyclic morphism. Suppose that, for every geometric point $t$ of $S$, one has $\h^0(X_t,\dZ/n)=\dZ/n$ and $\h^q(X_t,\dZ/n)=0$ for $q>0$ (in other words, the geometric fibers of $f$ are acyclic). Then $f_*\dZ/n=\dZ/n$ and $\eR^qf_*\dZ/n=0$ for $q>0$.
\end{corollary}

\begin{corollary}\label{I:5-1-10}
Let $f:X\to Y$ and $g:Y\to Z$ be morphisms of locally Noetherian schemes. If $f$ and $g$ are locally acyclic, then so is $g\circ f$.
\end{corollary}

We may suppose that $X$, $Y$, and $Z$ are strictly local and that $f$ and $g$ are local morphisms. It is then enough to show that, if $z$ is an algebraic geometric point of $Z$, one has $\h^0(X_z,\dZ/n)=\dZ/n$ and $\h^q(X_z,\dZ/n)=0$ for $q>0$.

Since $g$ is locally acyclic, one has $\h^0(Y_z,\dZ/n)=\dZ/n$ and $\h^q(Y_z,\dZ/n)=0$ for $q>0$. Moreover, the morphism $f_z:X_z\to Y_z$ is locally acyclic (\ref{I:5-1-4}), and its geometric fibers are acyclic because they are vanishing cycle varieties of $f$. By (\ref{I:5-1-9}), one therefore has $\eR^q{f_z}_*\dZ/n=0$ for $q>0$. Furthermore, ${f_z}_*\dZ/n$ is constant with fiber $\dZ/n$ on $Y_z$. The conclusion follows from the Leray spectral sequence for $f_z$.

\subsection{Local acyclicity of a smooth morphism}\label{I:5-2}

\begin{theorem}\label{I:5-2-1}
A smooth morphism is locally acyclic.
\end{theorem}

Let $f:X\to S$ be a smooth morphism. The assertion is local for the étale topology on $X$ and $S$, so we may suppose that $X$ is the standard affine space of dimension $d$ over $S$. By passage to the limit, we may suppose that $S$ is Noetherian, and the transitivity of local acyclicity (\ref{I:5-1-10}) shows that it is enough to treat the case $d=1$.

Let $s$ be a geometric point of $S$ and let $x$ be a geometric point of $X$ centered at a closed point of $X_s$. We must show that the geometric fibers of $\widetilde X^x\to\widetilde S^s$ are acyclic. From now on put $S=\widetilde S^s=\spec(A)$ and $X=\widetilde X^x$. One has $X\simeq\spec A\{T\}$, where $A\{T\}$ is the Henselization of $A[T]$ at the point $T=0$ over $s$.

If $t$ is a geometric point of $S$, the fiber $X_t$ is an inverse limit of smooth affine curves over $t$. Thus $\h^q(X_t,\dZ/n)=0$ for $q\geqslant2$, and it is enough to show that $\h^0(X_t,\dZ/n)=\dZ/n$ and $\h^1(X_t,\dZ/n)=0$ when $n$ is prime to the residue characteristic of $S$. This follows from the next two propositions.

\begin{proposition}\label{I:5-2-2}
Let $A$ be a strictly Henselian local ring, put $S=\spec(A)$, and put $X=\spec A\{T\}$. Then the geometric fibers of $X\to S$ are connected.
\end{proposition}

By passage to the limit, we may reduce to the case in which $A$ is the strict Henselization of a finitely generated $\dZ$-algebra.

Let $\bar t$ be a geometric point of $S$ centered at $t$, and let $k'$ be a finite separable extension of $k(t)$ in $k(\bar t)$. Put $t'=\spec(k')$ and $X_{t'}=X\times_S\spec(k')$. We must check that, for all $\bar t$ and $t'$, $X_{t'}$ is connected (by which we mean connected and nonempty). Let $A'$ be the normalization of $A$ in $k'$, that is, the ring of elements of $k'$ that are integral over the image of $A$ in $k(t)$. One has $A\{T\}\otimes_AA'\iso A'\{T\}$: the left-hand side is indeed Henselian local (because $A'$ is finite over $A$ and local) and is a limit of local étale algebras over $A'\{T\}=A\{T\}\otimes_AA'$. Thus $X_{t'}$ is again the fiber at $t'$ of $X'=\spec(A'\{T\})$ over $S'=\spec(A')$. The local scheme $X'$ is normal, hence integral; its localization $X_{t'}$ is still integral and is therefore connected.

\begin{proposition}\label{I:5-2-3}
Let $A$ be a strictly Henselian local ring, put $S=\spec(A)$, and put $X=\spec(A\{T\})$. Let $\bar t$ be a geometric point of $S$ and let $X_{\bar t}$ be the corresponding geometric fiber. Then every Galois étale covering of $X_{\bar t}$ whose order is prime to the characteristic of the residue field of $A$ is trivial.
\end{proposition}

\begin{lemma}[Zariski--Nagata purity theorem in dimension $2$] \label{I:5-2-4} % originally 5.2.3.1
Let $C$ be a regular local ring of dimension $2$, and let $C'$ be a finite normal $C$-algebra that is étale over the open complement of the closed point of $\spec(C)$. Then $C'$ is étale over $C$.
\end{lemma}

Indeed, $C'$ is normal of dimension $2$, hence $\operatorname{depth}(C')=2$. Since $\operatorname{depth}(C')+\operatorname{pd}(C')=\dim(C)=2$, it follows that $C'$ is free over $C$. The set of points of $C$ at which $C'$ is ramified is then defined by one equation, the discriminant; since it contains no point of height $1$, it is empty.

\begin{lemma}[A special case of Abhyankar's lemma]\label{I:5-2-5} % originally 5.2.3.2
Let $S=\spec(V)$ be a trait, let $\pi$ be a uniformizer, let $\eta$ be the generic point of $S$, let $X$ be irreducible, smooth over $S$, and of relative dimension $1$, let $\widetilde X_\eta$ be a Galois étale covering of $X_\eta$ of degree $n$ invertible on $S$, and put $S_1=\spec\left(V[\pi^{1/n}]\right)$. Use a subscript $(-)_1$ to denote base change from $S$ to $S_1$. Then $\widetilde X_{1\eta}$ extends to an étale covering of $X_1$.
\end{lemma}

Let $\widetilde X_1$ be the normalization of $X_1$ in $\widetilde X_{1\eta}$. In view of the structure of the tame inertia groups of the discrete valuation rings obtained by localizing $X$ at the generic points of the special fiber $X_s$, $\widetilde X_1$ is étale over $X_1$ on the generic fiber and at the generic points of the special fiber. By (\ref{I:5-2-4}), it is étale everywhere.

\subsubsection{}\label{I:5-2-6} % originally 5.2.3.3

Let $t$ be the point at which $\bar t$ is centered. We may replace $A$ by its normalization in a finite separable extension $k(t')$ of $k(t)$ contained in $k(\bar t)$ (cf. \ref{I:5-2-2}). This, together with a preliminary passage to the limit, allows us to suppose that
\begin{enumerate}[\indent a)]
  \item $A$ is normal and Noetherian, and $t$ is the generic point of $S$.
  \item The étale covering of $X_{\bar t}$ under consideration comes from an étale covering of $X_t$.
  \item It comes from an étale covering $\widetilde X_U$ of the inverse image $X_U$ of a nonempty open subset $U$ of $S$ (this follows from b), since $t$ is the limit of the $U$).
  \item The complement of $U$ has codimension $\geqslant2$ (after enlarging $k(t)$, by applying (\ref{I:5-2-5}) to the discrete valuation rings obtained by localizing $S$ at the points of $S\setminus U$ of codimension $1$ in $S$; there are only finitely many such points).
  \item The covering $\widetilde X_U$ is trivial over the subscheme $T=0$ (after enlarging $k(t)$ once more).
\end{enumerate}

Under these conditions, we shall show that the covering $\widetilde X_U$ is trivial.

\begin{lemma}\label{I:5-2-7} % originally 5.2.3.4
Let $A$ be a normal, Noetherian, strictly Henselian local ring, let $U$ be an open subset of $\spec(A)$ whose complement has codimension $\geqslant2$, let $V$ be its inverse image in $X=\spec(A\{T\})$, and let $V'$ be an étale covering of $V$. If $V'$ is trivial over $T=0$, then $V'$ is trivial.
\end{lemma}

Let $B=\Gamma(V',\sO)$. Since $V'$ is the inverse image of $V$ in $\spec(B)$, it is enough to see that $B$ is finite étale over $A\{T\}$ (and hence split, since $A\{T\}$ is strictly Henselian). Put $\widehat X=\spec A\llbracket T\rrbracket$, and use $(-)^\wedge$ to denote base change from $X$ to $\widehat X$. The scheme $\widehat X$ is faithfully flat over $X$. Thus $\Gamma(\widehat V',\sO)=B\otimes_{A\{T\}}A\llbracket T\rrbracket$, and it is enough to see that this ring $\widehat B$ is finite étale over $A\llbracket T\rrbracket$.

Let $V_m$ (respectively, $V'_m$) be the subscheme of $\widehat V$ (respectively, $\widehat V'$) defined by $T^{m+1}=0$. By hypothesis, $V'_0$ is a trivial covering of $V_0$, a disjoint union of $n$ copies of $V_0$. The same is true of $V'_m/V_m$, since étale coverings are insensitive to nilpotents. It follows that
\[
  \varphi : \Gamma(\widehat V',\sO) \to \varprojlim_m \Gamma(V_m',\sO) = \left(\varprojlim_m \Gamma(V,\sO)\right)^n \text{.}
\]
By hypothesis, the complement of $U$ has depth $\geqslant2$, so $\Gamma(V_m,\sO)=A[T]/(T^{m+1})$, and $\varphi$ is a homomorphism from $\widehat B$ to $A\llbracket T\rrbracket^n$. Over $U$, it provides $n$ distinct sections of $\widehat V'/\widehat V$; hence $\widehat V'$ is trivial, a disjoint union of $n$ copies of $\widehat V$. Since the complement of $\widehat V$ in $\widehat X$ still has codimension $\geqslant2$ (and hence depth $\geqslant2$), one deduces that $\widehat B=A\llbracket T\rrbracket^n$, proving the lemma.

\subsection{Applications}\label{I:5-3}

\begin{theorem}[Specialization of cohomology groups]\label{I:5-3-1}
Let $f:X\to S$ be a proper locally acyclic morphism, for example a smooth proper morphism. Then the sheaves $\eR^qf_*\dZ/n$ are locally constant constructible, and for every specialization map $t\to\widetilde S^s$, the cospecialization maps $\h^q(X_t,\dZ/n)\to\h^q(X_s,\dZ/n)$ are bijective.
\end{theorem}

This follows immediately from the definition of cospecialization maps and the finiteness and proper base change theorems.

\begin{theorem}[Base change by a smooth morphism]\label{I:5-3-2}
Let
\[\xymatrix{
  X' \ar[r]^-{g'} \ar[d]^-{f'}
    & X \ar[d]^-f \\
  S' \ar[r]^-{g}
    & S
}\]
be a Cartesian diagram with $g$ smooth. For every torsion sheaf $\sF$ on $X$ whose order is prime to the residue characteristics of $S$, one has
\[
  g^* \eR^q f_* \sF \iso \eR^q f_*'({g'}^* \sF) \text{.}
\]
\end{theorem}

Taking an open covering of $X$ reduces us to the case in which $X$ is affine, and passage to the limit then reduces us to the case in which $X$ is of finite type over $S$. The morphism $f$ then factors as an open immersion $j:X\to\bar X$ followed by a proper morphism $\bar f:\bar X\to S$. From the Leray spectral sequence for $\bar f\circ j$ and the proper base change theorem, one deduces that it is enough to prove the theorem when $X\to S$ is an open immersion.

In this case, if $\sF$ has the form $\varepsilon_*C$, where $\varepsilon:t\to X$ is a geometric point of $X$, the theorem follows from (\ref{I:5-1-5}). The general case follows from Lemma (\ref{I:5-1-8}).

\begin{corollary}\label{I:5-3-3}
Let $K/k$ be an extension of separably closed fields, let $X$ be a $k$-scheme, and let $n$ be an integer prime to the characteristic of $k$. Then the canonical map $\h^q(X,\dZ/n)\to\h^q(X_K,\dZ/n)$ is bijective for every $q\geqslant0$.
\end{corollary}

It is enough to observe that $K$ is a direct limit of smooth $k$-algebras.

\begin{theorem}[Relative purity]\label{I:5-3-4}
Let
\begin{equation*}\tag{5.3.4.1}\label{I:eq:5-3-4-1}\xymatrix{
  U \ar@{^{(}->}[r]^-j \ar[dr]
    & X \ar[d]^-f
    & \ar@{_{(}->}[l]_-i Y \ar[dl]^-h \\
  & S
}\end{equation*}
be a commutative diagram in which $f$ is smooth and purely of relative dimension $N$, $h$ is smooth and purely of relative dimension $N-1$, $i$ is a closed immersion, and $U=X\setminus Y$. If $n$ is prime to the residue characteristics of $S$, then
\begin{equation*}
\begin{aligned}
  j_* \dZ/n     &= \dZ/n \\
  \eR^1 j_*\dZ/n &= \dZ/n(-1)_Y \\
  \eR^q j_*\dZ/n &= 0 \qquad \text{for $q\geqslant 2$}
\end{aligned}
\end{equation*}
\end{theorem}

In these formulas, $\dZ/n(-1)$ denotes the $\dZ/n$-dual of $\dmu_n$. If $t$ is a local equation for $Y$, the isomorphism $\eR^1j_*\dZ/n\simeq\dZ/n(-1)_Y$ is defined by the map $a:\dZ/n\to\eR^1j_*\dmu_n$ that sends $1$ to the class of the $\dmu_n$-torsor of $n$th roots of $t$.

The question is local. Thus one may replace $(X,Y)$ by a locally isomorphic pair, for example
\[\xymatrix{
  \dA_T^1 \ar@{^{(}->}[r]^-j \ar[dr]_-g
    & \dP_T^1 \ar[d]^-f
    & \ar@{_{(}->}[l]_-i T \ar[dl] \\
  & T
}\]
with $T=\dA_S^{n-1}$ and $i$ the section at infinity. Corollary (\ref{I:5-1-9}) applies to $g$ and gives $\eR^qg_*\dZ/n=\dZ/n$ for $q=0$ and $0$ for $q>0$. For $f$, on also has \eqref{I:eq:4-6-2-1}
\[
  \eR^q f_* \dZ/n = \dZ/n,0,\dZ/n(-1),0 \text{ for $q>2$.}
\]
It is easy to check that $j_*\dZ/n=\dZ/n$, and that the sheaves $\eR^qj_*\dZ/n$ are supported on $i(T)$ for $q>0$. The Leray spectral sequence $E_2^{pq}=\eR^pf_*\eR^qj_*\dZ/n\Rightarrow\eR^{p+q}g_*\dZ/n$ therefore reduces to
\[\xymatrix{
  i^* \eR^q j_* \dZ/n & 0 & \cdots \\
  i^* \eR^1 j_* \dZ/n \ar[drr]^-{d_2} & 0 & \cdots \\
  \dZ/n              & 0 & \dZ/n(-1) & 0 & \cdots
}\]
$\eR^qj_*\dZ/n=0$ for $q\geqslant2$, and $\eR^1j_*\dZ/n$ is the extension by zero of a locally free sheaf of rank one on $T$ (isomorphic, via $d_2$, to $\dZ/n(-1)$). The map $a$ defined above is injective, as one checks fiber by fiber; hence it is an isomorphism, proving (\ref{I:5-3-4}).

\subsubsection{}\label{I:5-3-5}

We refer to \cite[XVI.4,5]{sga4} for the proof of the following applications of the acyclicity theorem (\ref{I:5-2-1}).

\begin{theorem}\label{I:5-3-6} % originally 5.3.5.1
Let $f:X\to S$ be a morphism of schemes of finite type over $\dC$, and let $\sF$ be a constructible sheaf on $X$. Then
\[
  \an{(\eR^q f_*\sF)} \sim \eR^q \an f_* (\an\sF)
\]
(cf. \ref{I:4-6-3}: for ordinary cohomology it is necessary to suppose that $\sF$ is constructible, and not merely torsion).
\end{theorem}

\begin{theorem}\label{I:5-3-7} % originally 5.3.5.2
Let $f:X\to S$ be a morphism of schemes of finite type over a field $k$ of characteristic $0$, and let $\sF$ be a constructible sheaf on $X$. Then the $\eR^qf_*\sF$ are also constructible.
\end{theorem}

The proof uses resolution of singularities and (\ref{I:5-3-4}). It is generalized to morphisms of finite type between excellent schemes of characteristic $0$ in \cite[XIX.5]{sga4}. Another proof, independent of resolution, is given in this volume (\nameref{VII}, \ref{VII:1-1})
% NOTE: originally (Th. finitude, 1.1),
and applies to a morphism of schemes of finite type over a field or over a Dedekind ring.

\section{Poincaré duality}

\subsection{Introduction}\label{I:6-1}

Let $X$ be an oriented topological manifold purely of dimension $N$, and suppose that $X$ has a finite open covering $\cU=(U_i)_{1\leqslant i\leqslant K}$ such that every nonempty intersection of the open sets $U_i$ is homeomorphic to a ball. For such a manifold, Poincaré duality can be presented as follows.

\begin{enumerate}[\indent A]
  \item The cohomology of $X$ is the Čech cohomology associated with the covering $\cU$. It is the cohomology of the complex
    \begin{equation}\label{I:eq:6-1}
      0 \to \dZ^{A_0} \to \dZ^{A_1} \to \cdots
    \end{equation}
    where $A_k=\{(i_0,\dotsc,i_k):i_0<\cdots<i_k\text{ and } U_{i_0}\cap\cdots\cap U_{i_k}\ne\varnothing\}$.
    \item For $a\in A_k$, $a=(i_0,\dotsc,i_k)$, put $U_a=U_{i_0}\cap\cdots\cap U_{i_k}$, and let $j_a$ be the inclusion of $U_a$ in $X$. The constant sheaf $\dZ$ on $X$ has the (left) resolution
      \begin{equation}\label{I:eq:6-2}
      \xymatrix{
        \cdots \ar[r]
          & \displaystyle\bigoplus_{a\in A_1} j_{a!}\dZ \ar[r]
          & \displaystyle\bigoplus_{a\in A_0} j_{a!} \dZ \ar[r] \ar[d]
          & 0 \\
        & & \dZ \text{.}
      }
      \end{equation}
      The compactly supported cohomology $\h_c^\bullet(X,j_{a!}\dZ)$ is just the cohomology with proper support of the (oriented) ball $U_a$:
        \[
          \h_c^i(X,j_{a!}\dZ) = \begin{cases}
                                  0   & \text{if $i\ne N$} \\
                                  \dZ & \text{if $i = N$}
                                \end{cases}
        \]
\end{enumerate}

The hypercohomology spectral sequence for the complex \eqref{I:eq:6-2}, with cohomology with proper support, therefore says that $\h_c^i(X)$ is the $(i-N)$th cohomology group of the complex
\begin{equation}\label{I:eq:6-3}
  \cdots \to \dZ^{A_1} \to \dZ^{A_0} \to 0 \text{.}
\end{equation}
This complex is dual to the complex \eqref{I:eq:6-1}, which gives Poincaré duality.

The essential points of this construction are
\begin{enumerate}[\indent a)]
  \item the existence of a theory of cohomology with proper support;
  \item the fact that every point $x$ of a manifold $X$ purely of dimension $N$ has a fundamental system of open neighborhoods $U$ for which
    \begin{equation}\label{I:eq:6-4}
      \h_c^i(U) = \begin{cases}
                    0   & \text{for $i\ne N$,} \\
                    \dZ & \text{for $i=N$.}
                  \end{cases}
    \end{equation}
\end{enumerate}

Poincaré duality in étale cohomology can be constructed on this model. Let $X$ be smooth and purely of dimension $N$ over an algebraically closed field $k$, let $n$ be invertible on $X$, and let $x$ be a closed point of $X$. The key point is to calculate the inverse limit over the étale neighborhoods $U$ of $x$:
\begin{equation}\label{I:eq:6-5}
  \varprojlim \h_c^i(U,\dZ/n)
    = \begin{cases}
        0     & \text{if $i\ne 2N$} \\
        \dZ/n & \text{if $i=2N$.}
      \end{cases}
\end{equation}

Just as for topological manifolds one must first treat directly the case of an open ball (or even simply the interval $(0,1)$), here one must first treat the case of curves directly (\ref{I:6-2}). The local acyclicity theorem for smooth morphisms then reduces \eqref{I:eq:6-5} to this special case (\ref{I:6-3}).

The isomorphisms \eqref{I:eq:6-4} and \eqref{I:eq:6-5} are not canonical: they depend on the choice of an \emph{orientation} of $X$. If $n$ is invertible on a scheme $X$, $\dmu_n$ is a sheaf of free $\dZ/n$-modules of rank one. Its $N$th tensor power ($N\in\dZ$) is denoted $\dZ/n(N)$. The intrinsic form of the second line of \eqref{I:eq:6-5} is
\begin{equation}\label{I:eq:6-5'}
  \varprojlim \h_c^{2N} \left(U,\dZ/n(N)\right) = \dZ/n \text{,}
\end{equation}
and $\dZ/n(N)$ is called the \emph{orientation sheaf} of $X$. Since the sheaf $\dZ/n(N)$ is constant and isomorphic to $\dZ/n$, it may be taken outside the inverse-limit sign, and one may instead write
\begin{equation}\label{I:eq:6-5''}
  \varprojlim \h_c^{2N} (U,\dZ/n) = \dZ/n(-N) \text{,}
\end{equation}
and Poincaré duality takes the form of a perfect duality, with values in $\dZ/n(-N)$, between $\h^i(X,\dZ/n)$ and $\h_c^{2N-i}(X,\dZ/n)$.

\subsection{The case of curves}\label{I:6-2}

\subsubsection{}\label{I:6-2-1}

Let $\bar X$ be a smooth projective curve over an algebraically closed field $k$, and let $n$ be invertible on $\bar X$. For connected $\bar X$, the proof of (\ref{I:3-3-5}) gives a canonical isomorphism
\[\xymatrix{
  \h^2(\bar X,\dmu_n) = \pic(\bar X)/n\pic(\bar X) \ar[r]^-{\deg}_-{\sim}
    & \dZ/n
}\]
Let $D$ be a reduced divisor on $\bar X$, and put $X=\bar X\setminus D$:
\[\xymatrix{
  X \ar@{^{(}->}[r]^-j
    & \bar X
    & \ar@{_{(}->}[l]_-i D \text{.}
}\]
The exact sequence $0\to j_!\dmu_n\to\dmu_n\to i_*\dmu_n\to0$ and the fact that $\h^i(\bar X,i_*\dmu_n)=\h^i(D,\dmu_n)=0$ for $i>0$ give an isomorphism
\[\xymatrix{
  \h_c^2(X,\dmu_n) = \h^2(\bar X,j_! \dmu_n) \ar[r]^-\sim
    & \h^2(\bar X,\dmu_n) \ar[r]^-\sim
    & \dZ/n \text{.}
}\]
For disconnected $\bar X$, one similarly has
\[
  \h_c^2(X,\dmu_n) = (\dZ/n)^{\pi_0(X)}
\]
and the \emph{trace morphism} is defined to be the sum
\[\xymatrix{
  \tr : \h_c^2(X,\dmu_n) \simeq (\dZ/n)^{\pi_0(X)} \ar[r]^-\Sigma
    & \dZ/n\text{.}
}\]

\begin{theorem}\label{I:6-2-2}
The pairing $\tr(a\smallsmile b)$ identifies each of the two groups $\h^i(X,\dZ/n)$ and $\h_c^{2-i}(X,\dmu_n)$ with the dual of the other (with values in $\dZ/n$).
\end{theorem}

\emph{Transcendental proof.} Suppose that $\bar X$ is a smooth projective curve over a trait $S$, and that $j:X\hookrightarrow\bar X$ is the inclusion of the complement of a divisor $D$ étale over $S$. The cohomology (respectively, cohomology with proper support) of the special and generic geometric fibers of $X/S$ are ``the same,'' that is, they are stalks of locally constant sheaves on $S$. This follows from the analogous facts for $\bar X$ and $D$, using the exact sequence $0\to j_!\dZ/n\to\dZ/n\to{\dZ/n}_D\to0$ (for cohomology with proper support) and the formulas $j_*\dZ/n=\dZ/n$, $\eR^1j_*\dZ/n\simeq{\dZ/n}_D(-1)$, and $\eR^ij_*\dZ/n=0$ for $i\geqslant2$ (for ordinary cohomology) (\ref{I:5-3-4}).

This specialization principle reduces (\ref{I:6-2-2}) to the case in which $k$ has characteristic $0$. By (\ref{I:5-3-3}), that case reduces to $k=\dC$. Finally, when $k=\dC$, the groups $\h^\bullet(X,\dZ/n)$ and $\h_c^\bullet(X,\dmu_n)$ agree with the groups of the same name calculated for the classical topological space $X_{\textnormal{cl}}$. Under the isomorphism $\dZ/n\to\dmu_n:x\mapsto\exp\left(\frac{2\pi ix}{n}\right)$, the trace morphism is identified with ``integration over the fundamental class.'' Thus (\ref{I:6-2-2}) follows from Poincaré duality for $X_{\textnormal{cl}}$.

\subsubsection{Algebraic proof}\label{I:6-2-3}

For a very economical proof, see (\nameref{V}, \S\ref{V:2}).
% originally (Duality \S 2)
% NOTE: should be hyperlinked V:2
Here is another, related to the self-duality of the Jacobian.

Retain the notation of (\ref{I:6-2-1}). We may and do suppose that $X$ is connected. Since the cases $i=0$ and $i=2$ are trivial, we also suppose that $i=1$. Define $_D\dG_m$ by the exact sequence $0\to{}_D\dG_m\to\dG_m\to i_*\dG_m\to0$ (sections of $\dG_m$ congruent to $1\bmod D$). The group $\h^1(\bar X,{}_D\dG_m)$ classifies invertible sheaves on $\bar X$ trivialized on $D$. It is the group of points of $\pic_D(\bar X)$, an extension of $\dZ$ (the degree) by the group of points of a generalized Jacobian of Rosenlicht (corresponding to conductor $1$ at each point of $D$), $\pic_D^0(\bar X)$; the latter is itself an extension of the abelian variety $\pic^0(\bar X)$ by the torus $\dG_m^D/(\text{diagonal $\dG_m$})$.

\begin{enumerate}[\indent a)]
  \item The exact sequence $0\to j_!\dmu_n\to{}_D\dG_m\xrightarrow{n}{}_D\dG_m\to0$ gives an isomorphism
    \begin{equation*}\tag{6.2.3.1}\label{I:eq:6-2-3-1}
      \h_c^1(X,\dmu_n) = \pic_D^0(\bar X)_n \text{.}
    \end{equation*}
  \item The map that assigns to $x\in X(k)$ the class of the invertible sheaf $\sO(x)$ on $\bar X$, trivialized by $1$ on $D$, comes from a morphism
    \[
      f:X \to \pic_D(\bar X) \text{.}
    \]
\end{enumerate}
For what follows, fix a base point $0$ and put $f_0(x)=f(x)-f(0)$. For every homomorphism $v:\pic_D^0(\bar X)_n\to\dZ/n$, let $\bar v\in\h^1\left(\pic_D^0(\bar X),\dZ/n\right)$ be the image under $v$ of the class in $\h^1\left(\pic_D^0(\bar X),\pic_D^0(\bar X)_n\right)$ of the torsor defined by the extension
\[\xymatrix{
  0 \ar[r]
    & \pic_D^0(\bar X)_n \ar[r]
    & \pic_D^0(\bar X) \ar[r]^-n
    & \pic_D^0(\bar X) \ar[r]
    & 0
}\]
Geometric class field theory (as presented by Serre in \cite{se59}) shows that the map $v\mapsto f_0^*(\bar v)$
\begin{equation*}\tag{6.2.3.2}\label{I:eq:6-2-3-2}
  \hom\left(\pic_D^0(\bar X)_n,\dZ/n\right) \to \h^1(X,\dZ/n)
\end{equation*}
is an isomorphism. To deduce (\ref{I:6-2-2}) from \eqref{I:eq:6-2-3-1} and \eqref{I:eq:6-2-3-2}, it remains to know that
\begin{equation*}\tag{6.2.2.3}\label{I:eq:6-2-3-3}
  \tr\left(u \smallsmile f_0^*(\bar v)\right) = -v(u) \text{.}
\end{equation*}

% originally (Duality, 3.2.4).

\subsection{The general case}\label{I:6-3}

Let $X$ be a smooth algebraic variety purely of dimension $N$ over an algebraically closed field $k$. To state Poincaré duality, one must first define the trace morphism
\[
  \tr : \h_c^{2N}\left(X,\dZ/n(N)\right) \to \dZ/n\text{.}
\]
The definition is a laborious dévissage starting from the case of curves \cite[XVIII, \S 2]{sga4}. One then has the following theorem.

\begin{theorem}\label{I:6-3-1}
The pairing $\tr(a\smallsmile b)$ identifies each of the groups $\h_c^i\left(X,\dZ/n(N)\right)$ and $\h^{2N-i}(X,\dZ/n)$ with the dual of the other.
\end{theorem}

Let $x\in X$ be a closed point and let $X_x$ be the strict localization of $X$ at $x$. As $U$ ranges over the étale neighborhoods of $x$, put
\begin{equation*}\tag{6.3.1.1}\label{I:eq:6-3-1-1}
  \h_c^\bullet(X_x,\dZ/n) = \varprojlim \h_c^\bullet(U,\dZ/n) \text{.}
\end{equation*}
It would be preferable to consider the pro-object $\text{``}\varprojlim\text{''}\h_c^\bullet(U,\dZ/n)$ instead, but since the groups involved are finite, the difference is inessential. As we tried to explain in the introduction, (\ref{I:6-3-1}) follows from the facts that
\begin{equation*}\tag{6.3.1.2}\label{I:eq:6-3-1-2}
\begin{aligned}
  \h_c^i(X_x,\dZ/n) &= 0 \text{ for $i\ne 2N$, and} \\
  \tr : \h_c^{2N}\left(X_x,\dZ/n(N)\right) &\iso \dZ/n \text{ is an isomorphism.}
\end{aligned}
\end{equation*}

The case $N=0$ is trivial. If $N>0$, let $Y_y$ be the strict localization at a closed point of a smooth scheme $Y$ purely of dimension $N-1$, and let $f:X_x\to Y_y$ be an essentially smooth morphism (of relative dimension one). The proof uses the Leray spectral sequence for cohomology with proper support under $f$ to reduce to the case of curves. Since the ``cohomology with proper support'' groups under consideration are defined as limits \eqref{I:eq:6-3-1-1}, the existence of such a spectral sequence raises various passage-to-the-limit issues, treated in excessive detail in \cite[XVIII]{sga4}. For every geometric point $z$ of $Y_y$, one has
\[
  (\eR^if_! \dZ/n)_z = \h_c^i(f^{-1}(z),\dZ/n) \text{.}
\]
The geometric fiber $f^{-1}(z)$ is an inverse limit of smooth curves over an algebraically closed field. It satisfies Poincaré duality. Its ordinary cohomology is given by the local acyclicity theorem for smooth morphisms:
\[
  \h^i(f^{-1}(z),\dZ/n)
    = \begin{cases}
        \dZ/n & \text{for $i=0$} \\
        0     & \text{for $i>0$.}
      \end{cases}
\]
By duality,
\[
  \h_c^i(f^{-1}(z),\dZ/n)
    = \begin{cases}
        \dZ/n(-1) & \text{for $i=2$} \\
        0         & \text{for $i\ne 2$.}
      \end{cases}
\]
and the Leray spectral sequence gives
\[
  \h_c^i\left(X_x,\dZ/n(N)\right) = \h_c^{i-2}\left(Y_y,\dZ/n(N-1)\right) \text{.}
\]
The conclusion follows by induction on $N$.

\subsection{Variants and applications}\label{I:6-4}

In étale cohomology one can construct a ``duality formalism'' (that is, functors $\eR f_!$, $\eR f_*$, and $\eR f^!$ satisfying various compatibility and adjunction formulas) parallel to the one in coherent cohomology. In this language, the results of the preceding subsection are expressed as follows: if $f:X\to S$ is smooth and purely of relative dimension $N$, and $S=\spec(k)$ with $k$ algebraically closed, then
\[
  \eR f^! \dZ/n = \dZ/n [2N](N) \text{.}
\]
This result holds without any hypothesis on $S$. It has the following corollary.

\begin{corollary}\label{I:6-4-1}
If $f$ is smooth and purely of relative dimension $N$, and the sheaves $\eR^if_!\dZ/n$ are locally constant, then the sheaves $\eR^if_*\dZ/n$ are also locally constant, and
\[
  \eR^i f_* \dZ/n = \const\hom\left(\eR^{2 N-i} f_! \dZ/n(N),\dZ/n\right)\text{.}
\]
\end{corollary}

In particular, under the hypotheses of the corollary, the sheaves $\eR^if_*\dZ/n$ are constructible. Starting from this, one can show that, if $S$ is of finite type over the spectrum of a field or of a Dedekind ring, then, for every morphism of finite type $f:X\to S$ and every constructible sheaf $\sF$ on $X$, the sheaves $\eR^if_*\sF$ are constructible (\nameref{VII}, \ref{VII:1-1}).
% originally (Th. finitude, 1.1).
% END INLINED FILE: sga4.5-1.tex
\clearpage
\clearpage
% BEGIN INLINED FILE: sga4.5-2.tex
\chapter{Report on the trace formula}\label{II}

In this text I have tried to present Grothendieck's cohomological theory of $L$-functions as directly as possible. I follow very closely certain lectures given by Grothendieck at the IHES in the spring of 1966. In the spirit of this volume, SGA 5 will not be invoked, apart from two references to passages of Exposé XI, which is otherwise independent of that seminar.

In \S 10 (pp. 81--90) of \cite{de73}, the reader will find a summary of the theory in the crucial case of curves.

\section{Some reminders}\label{II:1}

\subsection{Notation}\label{II:1-1}

$p,q,\dF,\dF_q$: $p$ is a prime number, $q=p^f$ is a power of $p$, and $\dF$ is an algebraic closure of the prime field $\dF_p$. For every power $p^n$ of $p$, let $\dF_{p^n}$ denote the subfield of $\dF$ with $p^n$ elements.

$X_0,X$: An object over $\dF_q$ will often be denoted by a symbol with a subscript $0$. The same symbol without the $0$ then denotes the object obtained by extension of scalars to $\dF$. For example, if $\sF_0$ is a sheaf on a scheme $X_0$ over $\dF_q$, then $\sF$ denotes its inverse image on $X=X_0\otimes_{\dF_q}\dF$.

$F$: If $X_0$ is a scheme over $\dF_q$, the endomorphism of $X_0$ that ``sends the point with coordinates $x$ to the point with coordinates $x^q$'' is called the \emph{Frobenius morphism} and is denoted $F$. It is the identity on the underlying topological space, and, for $x$ a local section of the structure sheaf, $F^*x=x^q$.

The endomorphism of $X$ obtained from it by extension of scalars is also denoted $F$. On the set $X(\dF)=X_0(\dF)$ of rational points of $X$, $F$ acts as the Frobenius substitution $\varphi\in\gal(\dF/\dF_q)$ defined by $\varphi(x)=x^q$. In particular, $X^F=X_0(\dF_q)$.

When there is a risk of ambiguity, we write $F_{(q)}$ instead of $F$.

\subsection{The Frobenius correspondence}\label{II:1-2}

Let $X_0$ be a scheme over $\dF_q$ and let $\sF_0$ be a sheaf on $X_0$. I like to view the \emph{Frobenius correspondence} (an $F$-endomorphism of $\sF_0$) as follows:
\begin{equation*}\tag{1.2.1}\label{II:eq:1-2-1}
  F^* : F^*\sF_0 \iso \sF_0
\end{equation*}

View $\sF_0$ as the sheaf of local sections of a space $[\sF_0]$ étale over $X_0$ ($[\sF_0]$ is an algebraic space in the sense of M. Artin). Functoriality of $F$ gives a commutative diagram
\[\xymatrix{
  [\sF_0] \ar[r]^-F \ar[d]
    & [\sF_0] \ar[d] \\
  X_0 \ar[r]^F
    & X_0
}\]
Because $[\sF_0]$ is étale over $X_0$, this diagram is Cartesian; it gives an isomorphism $[\sF_0]\iso F^*[\sF_0]$ whose inverse is \eqref{II:eq:1-2-1}. We also use $F^*$ to denote the correspondences or morphisms deduced from $F$ by extension of scalars or by functoriality, such as $F^*:\h_c^\bullet(X,\sF)\to\h_c^\bullet(X,\sF)$.

For a formal definition, I refer to \cite[XI.1,2]{sga5}. Only the case $q=p$ is considered there, but the general case is treated in the same way.

\subsection{}\label{II:1-3}

The Frobenius correspondence is functorial in $(X_0,\sF_0)$. Let $f_0:X_0\to Y_0$ be a morphism of $\dF_q$-schemes, let $\sF_0$ be a sheaf on $X_0$, let $\sG_0$ be a sheaf on $Y_0$, and let $u\in\hom(f_0^*\sG_0,\sF_0)$ be an $f_0$-morphism from $\sG_0$ to $\sF_0$. There is a commutative diagram of spaces expressing $f_0F=Ff_0$, and above it a commutative diagram of sheaves:
\[\xymatrix{
  X_0 \ar[r]^-{f_0} \ar[d]^-F
    & Y_0 \ar[d]^-F \\
  X_0 \ar[r]^-{f_0}
    & Y_0
}\qquad\qquad
\xymatrix{
  \sF_0
    & \ar[l]_-u \sG_0 \\
  \sF_0 \ar[u]_-{F^*}
    & \ar[l] \sG_0 \ar[u]^-{F^*}
}\]
In particular, the correspondence on $f_{0*}\sF_0$ deduced from the Frobenius correspondence of $(X_0,\sF_0)$ is the Frobenius correspondence of $(Y_0,f_{0*}\sF_0)$.

\subsection{}\label{II:1-4}

The natural morphism $Y\to Y_0$ is a limit of étale morphisms. It follows that, for every abelian sheaf $\sF_0$ on $X_0$, the inverse images on $Y$ of the $\eR^if_{0*}\sF_0$ are the $\eR^if_*\sF$.

On $\eR^if_*\sF$ one has
\begin{enumerate}[\indent a)]
  \item the correspondence $F^*\eR^if_*\sF\iso\eR^if_*\sF\xrightarrow{F^*}\eR^if_*\sF$ deduced by functoriality from the Frobenius correspondence;
  \item the correspondence $F^*\eR^if_*\sF\to\eR^if_*\sF$ that is the inverse image on $Y$ of the Frobenius correspondence on $\eR^if_{0*}\sF_0$.
\end{enumerate}

It follows from (\ref{II:1-3}) that these correspondences coincide. They induce the initial term of a Frobenius action on the entire spectral sequence
\[
  \h^p(Y,\eR^q f_*\sF) \rightarrow \h^{p+q}(X,\sF) \text{.}
\]
The same holds for cohomology with proper support and for the spectral sequence
\[
  \h_c^p(Y,\eR^q f_* \sF) \rightarrow \h_c^{p+q}(X,\sF)
\]

\subsection{Change of finite field}\label{II:1-5}

Let $(X_0,\sF_0)$ be over $\dF_q$, let $n$ be an integer, and let $(X_1,\sF_1)$ over $\dF_{q^n}$ be obtained from $(X_0,\sF_0)$ by extension of scalars. On the one hand there is the Frobenius correspondence $F_{(q)}:X_0\to X_0$, $F_{(q)}^*:F_{(q)}^*\sF_0\to\sF_0$, and on the other hand there is $F_{(q^n)}:X_1\to X_1$, $F_{(q^n)}^*:F_{(q^n)}^*\sF_1\to\sF_1$. It is easy to check that $F_{(q^n)}$ is obtained from the $n$th power of $F_{(q)}$ by extension of scalars from $\dF_q$ to $\dF_{q^n}$. After extension of scalars to $\dF$, the same identity $F_{(q^n)}=F_{(q)}^n$ holds between endomorphisms of $X$ and correspondences on $X$.

If $x$ is a closed point of $X_0$, with $[k(x):\dF_q]=n$, and $\bar x\in X_0(\dF)=X(\dF)$ is centered at $x$, then $\bar x$ is fixed by $F_{(q^n)}=F_{(q)}^n$. Let $F_x^*$ denote the endomorphism of $\sF_{\bar x}$ induced by $F_{(q^n)}^*$. Up to isomorphism, $(F_x^*,\sF_{\bar x})$ is independent of the choice of $\bar x$. The trace, etc., of $F_x^*$ is denoted by notation such as $\tr(F_x^*,\sF)$.

\subsection{The \texorpdfstring{$L$}{L}-function}\label{II:1-6}

Let $X_0$ be a scheme of finite type over $\dF_q$. Let $|X_0|$ denote its set of closed points and, for $x\in|X_0|$, put $\deg(x)=[k(x):\dF_p]$. If $\sF_0$ is a $\dQ_\ell$-sheaf on $X_0$ (for this notion, see below), or a constructible flat sheaf of modules over a commutative Noetherian ring $A$, let $L(X_0,\sF_0)$ denote the formal power series with constant term $1$ in $\dQ_\ell\llbracket t\rrbracket$ or $A\llbracket t\rrbracket$ given by
\[
  L(X_0,\sF_0) = \prod_{x\in|X_0|} \det\left(1-F_x^* t^{\deg(x)},\sF\right)^{-1} \text{.}
\]

\subsection{Remark}\label{II:1-7}

In Definition (\ref{II:1-6}), $\deg(x)=[k(x):\dF_p]$ is the absolute degree, not $\deg(x)=[k(x):\dF_q]$, the degree relative to $\dF_q$. Consequently $L(X_0,\sF_0)$ depends only on the scheme $X_0$ and the sheaf $\sF_0$, not on $q$ or on the structural morphism $X_0\to\spec(\dF_q)$. The indeterminate $t$ is an avatar of $p^{-s}$.

\subsection{The Galois viewpoint}\label{II:1-8}

This subsection will not be used in the rest of the report. If $\sF_0$ is an abelian sheaf on $X_0$, the Galois group $\gal(\dF/\dF_q)$ acts on $\dF$ and, by transport of structure, on $\h_c^i(X,\sF)$. In particular, the Frobenius substitution $\varphi$ (\ref{II:1-1}) induces an automorphism, still denoted $\varphi$, of $\h_c^i(X,\sF)$. The Frobenius correspondence, on the other hand, defines an endomorphism $F^*$ of $\h_c^i(X,\sF)$. One has
\begin{equation*}\tag{1.8.1}\label{II:eq:1-8-1}
  {F^*}^{-1} = \varphi \qquad \text{(on $\h_c^i(X,\sF)$).}
\end{equation*}

To see this, apply (\ref{II:1-4}) when $Y_0$ is a point, $Y_0=\spec(\dF_q)$. One finds that $\h_c^i(X,\sF)$ is the stalk at the geometric point $\spec(\dF)$ of $Y_0$ of $\eR^if_!\sF_0$, and that $F^*$ is the stalk of the Frobenius correspondence of $\eR^if_!\sF_0$. Let $[\eR^if_!\sF_0]$ have the meaning of (\ref{II:1-2}). Then $\h_c^i(X,\sF) = [\eR^i f_!\sF_0](\dF)$, and
\begin{enumerate}[\indent a)]
  \item $\varphi$ acts by transport of structure through its action on $\dF$;
  \item $F^*$ is the inverse of $F:[\eR^i f_!\sF_0](\dF) \to [\eR^i f_!\sF_0](\dF)$.
\end{enumerate}

The conclusion follows from the identity $F=\varphi$ of (\ref{II:1-1}).

One advantage of the Galois viewpoint is that it permits arguments by transport of structure.

\section{\texorpdfstring{$\dQ_\ell$}{Ql}-sheaves}\label{II:2}

The notion of a $\dQ_\ell$-sheaf is developed in detail in \cite[V,VI]{sga5}. We shall use only the definitions and results below.

\begin{definition_}\label{II:2-1}
Let $X$ be a Noetherian scheme. A $\dZ_\ell$-sheaf $\sF$ on $X$ is an inverse system of sheaves $\sF_n$ ($n\geqslant0$), where $\sF_n$ is a constructible sheaf of $\dZ/\ell^{n+1}$-modules \cite[IX.2]{sga4}, such that the transition morphism $\sF_n\to\sF_{n-1}$ factors through an isomorphism
\[
  \sF_n\otimes_{\dZ/\ell^{n+1}} \dZ/\ell^n \iso \sF_{n-1}\text{.}
\]
The sheaf $\sF$ is said to be \emph{lisse} if the $\sF_n$ are locally constant.
\end{definition_}

The terminology of \cite{sga5} is different: it uses ``constructible $\dZ_\ell$-sheaf'' for ``$\dZ_\ell$-sheaf,'' and ``constructible twisted constant'' for ``lisse.'' The same applies below to $\dQ_\ell$-sheaves.

\subsection{}\label{II:2-2}

If $k$ is a separably closed field, a sheaf of $\dZ/\ell^n$-modules on $\spec(k)$ is identified with a $\dZ/\ell^n$-module, and the functor $\sF\mapsto\varprojlim\sF_n$ identifies the $\dZ_\ell$-sheaves on $\spec(k)$ with finitely generated $\dZ_\ell$-modules. This justifies the following definition: the \emph{stalk} of a $\dZ_\ell$-sheaf $\sF$ at a geometric point $x$ of $X$ is the $\dZ_\ell$-module $\varprojlim(\sF_n)_x$.

\subsection{}\label{II:2-3}

Suppose that $X$ is connected, and let $x$ be a geometric point of $X$. The ``stalk at $x$'' functor is an equivalence from the category of locally constant constructible sheaves of $\dZ/\ell^n$-modules to the category of finitely generated $\dZ/\ell^n$-modules endowed with a continuous action of $\pi_1(X,x)$. Passing to the limit gives the following proposition.

\begin{proposition_}\label{II:2-4}
Under the hypotheses above, the ``stalk at $x$'' functor is an equivalence from the category of lisse $\dZ_\ell$-sheaves on $X$ to the category of finitely generated $\dZ_\ell$-modules endowed with a continuous action of $\pi_1(X,x)$.
\end{proposition_}

\begin{proposition_}\label{II:2-5}
Let $\sF$ be a $\dZ_\ell$-sheaf on a Noetherian scheme $X$. There is a finite partition of $X$ into locally closed subsets $X_i$ such that the $\sF|X_i$ are lisse.
\end{proposition_}

Put $\gr_\ell^n\sF=\ell^n\sF_m/\ell^{m+1}\sF_m$ for $m\geqslant n$, and put $\gr_\ell\sF=\bigoplus\gr_\ell^n\sF$. This is a sheaf of modules over the Noetherian ring $\gr_\ell\dZ=\dZ/\ell[t]$, the associated graded ring of $\dZ$ for the $\ell$-adic filtration. It is a quotient of $\gr_\ell^0\sF\otimes_{\dZ/\ell}\dZ/\ell[t]$, hence a constructible sheaf of $\gr_\ell\dZ$-modules \cite[IX.2]{sga4}. There is therefore a finite partition of $X$ into locally closed subsets $X_i$ such that $\gr_\ell\sF|X_i$ is locally constant. Each sheaf $\gr_\ell^n\sF|X_i$ is then locally constant, as are the $\sF_m|X_i$, which are successive extensions of them.

\subsection{}\label{II:2-6}

If $\sF$ and $\sG$ are two $\dZ_\ell$-sheaves on $X$, put $\hom(\sF,\sG)=\varprojlim\hom(\sF_n,\sG_n)$; this is a $\dZ_\ell$-module. Since $\hom(\sF_n,\sG_n)=\hom(\sF_m,\sG_n)$ for $m\geqslant n$, one also has
\[
  \hom(\sF,\sG) = \varprojlim_n \varinjlim_m \hom(\sF_m,\sG_n) \text{,}
\]
so the functor $\sF\mapsto\text{``}\varprojlim\text{''}\sF_n$ is fully faithful from the category of $\dZ_\ell$-sheaves to that of pro-sheaves on $X$ (that is, pro-objects in the category of sheaves on $X$). We shall use this to identify $\dZ_\ell$-sheaves with certain pro-sheaves, namely those for which every $\sF_n=\sF/\ell^{n+1}\sF$ is a sheaf and $\sF\iso\text{``}\varprojlim\text{''}\sF_n$. It is easy to check that, if $\sF=\text{``}\varprojlim\text{''}\sF_n$, with $\sF_n$ a constructible sheaf of $\dZ/\ell^n$-modules, then $\sF$ is a $\dZ_\ell$-sheaf if and only if the inverse systems (in $n$) $\sF_n/\ell^{k+1}\sF_n$ are essentially constant. The inverse system of $\sF'_k=\varprojlim_n\sF_n/\ell^{k+1}\sF_n$ then satisfies the conditions of (\ref{II:2-1}), and $\sF=\text{``}\varprojlim\text{''}\sF'_k$.

In what follows, we abandon the notation $\sF_n$ of (\ref{II:2-1}). If $\sF$ is a $\dZ_\ell$-sheaf, the sheaf $\sF_n$ of (\ref{II:2-1}) will be denoted $\sF\otimes\dZ/\ell^{n+1}$.

\subsection{}\label{II:2-7}

In the abelian category of pro-sheaves, the subcategory of $\dZ_\ell$-sheaves is stable under kernels, cokernels, and extensions. This is clear for cokernels. For kernels, if $f:\sF\to\sG$ is a morphism of $\dZ_\ell$-sheaves, it follows readily from (\ref{II:2-4}, \ref{II:2-5}) and Artin--Rees that the inverse system of $\ker(f:\sF\otimes\dZ/\ell^n\to\sG\otimes\dZ/\ell^n)$ satisfies the criterion given in (\ref{II:2-6}) above. There is even an integer $r$ such that, for every $n$,
\[
  \ker(f)\otimes\dZ/\ell^n = \ker(f:\sF\otimes\dZ/\ell^{n+r}\to\sG\otimes\dZ/\ell^{n+r})\otimes\dZ/\ell^n \text{.}
\]
The case of extensions is left to the reader.

\subsection{}\label{II:2-8}

A $\dZ_\ell$-sheaf $\sF$ is \emph{torsion-free} if $\ker(\ell:\sF\to\sF)=0$. Each $\sF\otimes\dZ/\ell^n$ is then flat over $\dZ/\ell^n$ (that is, its stalks are finitely generated free $\dZ/\ell^n$-modules). It follows from (\ref{II:2-4}, \ref{II:2-5}) that, for every $\dZ_\ell$-sheaf $\sF$, there is an integer $n$ such that $\sF/\ker(\ell^n)$ is torsion-free.

\subsection{}\label{II:2-9}

The abelian category of \emph{$\dQ_\ell$-sheaves} on $X$ is the quotient of the category of $\dZ_\ell$-sheaves by the subcategory of torsion $\dZ_\ell$-sheaves. Equivalently, its objects are $\dZ_\ell$-sheaves and, if $\sF\otimes\dQ_\ell$ denotes the $\dZ_\ell$-sheaf $\sF$ viewed as an object of the category of $\dQ_\ell$-sheaves, then
\[
  \hom(\sF\otimes\dQ_\ell,\sG\otimes\dQ_\ell) = \hom(\sF,\sG)\otimes_{\dZ_\ell}\dQ_\ell \text{.}
\]

The \emph{stalk} of a $\dQ_\ell$-sheaf $\sF\otimes\dQ_\ell$ at a geometric point $x$ of $X$ is the finite-dimensional $\dQ_\ell$-vector space $\sF_x\otimes_{\dZ_\ell}\dQ_\ell$.

\subsection{}\label{II:2-10}

Let $X$ be a separated scheme of finite type over an algebraically closed field $k$, and let $\sF$ be a constructible $\dQ_\ell$-sheaf on $X$. Let $\sF'$ be a constructible $\dZ_\ell$-sheaf such that $\sF\sim\sF'\otimes_{\dZ_\ell}\dQ_\ell$. We shall see in (\ref{II:4-11}) that
\[
  \left(\varprojlim \h_c^q(X,\sF'\otimes \dZ/\ell^n)\right)\otimes_{\dZ_\ell}\dQ_\ell
\]
is a finite-dimensional $\dQ_\ell$-vector space. It depends only on $\sF$ and is denoted $\h_c^q(X,\sF)$. The proof does not use the assumption that $\ell$ is prime to the characteristic, but only that case concerns us.

\subsection{}\label{II:2-11}

In this exposé, all important calculations will be made at the level of $\dZ/\ell^n$-sheaves, and passage to $\dQ_\ell$-sheaves will occur only at the very end. In practice, however, it is often convenient to work systematically with $\dQ_\ell$-sheaves. In Exposés V and VI of \cite{sga5}, Jouanolou gives the formalism that makes this possible. One essential result (VI 2.2) is that, if $f:X\to Y$ is a separated morphism of finite type between Noetherian schemes and $\sF$ is a $\dZ_\ell$-sheaf on $X$, then, for every $q$, the pro-sheaf $\text{``}\varprojlim\text{''}\eR^qf_!(\sF\otimes\dZ/\ell^n)$ is a $\dZ_\ell$-sheaf on $Y$. He proves an even stronger regularity property (of Artin--Rees type) for the inverse system of the $\eR^qf_!(\sF\otimes\dZ/\ell^n)$.

Jouanolou in fact works in a more general framework that includes the very useful case of $E_\lambda$-sheaves, where $E_\lambda$ is a finite extension of $\dQ_\ell$. They are defined as $\dQ_\ell$-sheaves were, with $\dQ_\ell$ replaced by $E_\lambda$, $\dZ_\ell$ by the valuation ring $\sO_\lambda$, and $\dZ/\ell^n$ by $\sO_\lambda/\pi^n$, where $\pi$ is a uniformizer. Equivalently, an $E_\lambda$-sheaf can be defined as a $\dQ_\ell$-sheaf $\sF$ endowed with a homomorphism $E_\lambda\to\operatorname{End}(\sF)$.

\section{The trace formula and \texorpdfstring{$L$}{L}-functions}\label{II:3}

\emph{We use the notation of \S\ref{II:1}, and $\ell$ is a prime number different from $p$.}

Let $X_0$ be a separated scheme of finite type over $\dF_q$ ($q=p^f$), and let $\sF_0$ be a constructible $\dQ_\ell$-sheaf on $X_0$. Grothendieck's cohomological interpretation of $L$-functions is the following theorem.

\begin{theorem_}\label{II:3-1}
$L(X_0,\sF_0) = \prod_i \det\left(1-F^* t^f, \h_c^i(X,\sF)\right)^{(-1)^{i+1}}$.
\end{theorem_}

It is obtained as a consequence of the trace formula.

\begin{theorem_}\label{II:3-2}
For every integer $n$, $\sum_{x\in X^{F^n}} \tr\left({F^n}^*,\sF_x\right) = \sum_i (-1)^i \tr\left({F^n}^*,\h_c^i(X,\sF)\right)$.
\end{theorem_}

Let us briefly recall the classical method for deducing (\ref{II:3-1}) from (\ref{II:3-2}). Put $T=t^f$; both sides are formal power series in $T$ with constant term $1$. Since $\dQ_\ell$ has characteristic $0$, it is enough to prove that the logarithmic derivatives $T\frac{d}{dT}\log$ of the two sides are equal.

To calculate these logarithmic derivatives, we use the formula
\[
  T \frac{d}{dT}\log\det(1-f T)^{-1} = \sum_{n>0} \tr(f^n) T^n \text{,}
\]
% NOTE: the following lines are not typesetting correctly
whose proof (in a slightly more general setting) is recalled below. For the first side one finds
\begin{align*}
  T\frac{d}{dT}\log L(X_0,\sF_0)
    &= \sum_{x\in |X_0|} \sum_n [k(x):\dF_q] \tr\left(F_x^n,\sF\right) T^{n\cdot[k(x):\dF_q]} \\
    &= \sum_{n>0} T^n \sum_{x\in X^{F^n}} \tr\left({F^n}^*,\sF_x\right)
\end{align*}
and for the second side
\[
  \sum_{n>0} T^n \sum (-1)^i \tr\left({F^n}^*,\h_c^i(X,\sF)\right) \text{ ; }
\]
the terms are compared one by one.

\begin{proposition_}\label{II:3-3}
Let $A$ be a commutative ring, let $M$ be a finitely generated projective $A$-module, and let $f$ be an endomorphism of $M$. For every invertible formal power series $Q$, put $\frac{d}{dT}\log Q=Q^{-1}\frac{d}{dT}Q$. Then
\begin{equation*}\tag{3.3.1}\label{II:eq:3-3-1}
  T\frac{d}{dT}\log\det(1-f T)^{-1} = \sum_{n>0} \tr(f^n) T^n \text{.}
\end{equation*}
\end{proposition_}

Let $M'$ be an $A$-module such that $M\oplus M'$ is free of finite rank. Replacing $(M,f)$ by $(M\oplus M',f\oplus0)$ does not change either side of (\ref{II:3-3}), and reduces us to the case in which $M$ is free. For $M=A^d$, \eqref{II:eq:3-3-1} is an algebraic identity in the coefficients $f_i^j$ of $f$. By the principle of extension of algebraic identities, it is enough to verify it when $A$ is an algebraically closed field of characteristic $0$. Since both sides are additive in $M$ in an extension $0\to M'\to M\to M''\to0$, it is even enough to treat the case in which $M$ has rank $1$. If $f$ is multiplication by $a$, \eqref{II:eq:3-3-1} reduces to the formula
\[
  T\frac{d}{dT}\log\left((1-a T)^{-1}\right) = \sum_{n>0} a^n T^n \text{.}
\]

\begin{corollary_}\label{II:3-4}
Let $n$ be an integer, and let $f^{(n)}$ be the endomorphism $(x_1,\dotsc,x_n)\mapsto \left(f(x_n),x_1,\dotsc,x_{n-1}\right)$ of $M^n$. Then
\begin{equation*}\tag{3.4.1}\label{II:eq:3-4-1}
  \det(1-f T^n) = \det(1-f^{(n)} T) \text{.}
\end{equation*}
\end{corollary_}

As in (\ref{II:3-3}), one reduces to the case in which $A$ has characteristic $0$. It is then enough to check that the logarithmic derivatives $-T\frac{d}{dT}\log$ of the two sides are equal.
\begin{align*}
  -T\frac{d}{dT}\log\det(1-f T^n) &= - n T^n \frac{d}{dT^n} \log\det(1-f T^n) = n \sum_{m>0} \tr(f^m) T^{n m} \text{,} \\
  -T \frac{d}{dT}\log\det\left(1-f^{(n)} T\right) &= \sum_{m>0} \tr\left(f^{(n) m}\right) T^m \text{,}
\end{align*}
and one observes that
\begin{align*}
  \tr(f^{(n) m}) &= 0 \qquad \text{if $n\nmid m$} \\
  \tr(f^{(n) n m}) &= n \tr(f^m) \text{.}
\end{align*}

\subsection{Remark}\label{II:3-5}

For $A=\dQ_\ell$, \eqref{II:eq:3-4-1} is the special case $X_0=\spec(\dF_{q^n})$ of (\ref{II:3-1}). Using this identity, it is easy to check directly that the right-hand side of (\ref{II:3-1}) is independent of $q$ and of the structural morphism $X_0\to\spec(\dF_q)$.

\subsection{Remark}\label{II:3-6}

If $j:X_0\hookrightarrow\bar X_0$ is a compactification of $X_0$, then $L(X_0,\sF_0)=L(\bar X_0,j_!\sF_0)$ and $\h_c^i(X,\sF)=\h^i(\bar X,j_!\sF)$. This reduces (\ref{II:3-1}) to the proper case and explains the use of cohomology with proper support. When $X$ is proper, (\ref{II:3-2}) has the form of a Lefschetz trace formula. Notice that the local terms $\tr(F^{n*},\sF_x)$ ($x\in X^{F^n}$) depend only on the stalk of $\sF$ at $x$ and carry no multiplicity. When $X$ and $\sF$ are smooth and lisse, respectively, this corresponds to the fact that the graph of $F$ is transverse to the diagonal (since $dF=0$).

\subsection{Remark}\label{II:3-7}

Assuming the formalism of $\dQ_\ell$-sheaves (cf. \ref{II:2-11}; one needs the Leray spectral sequence and the base change theorem for the higher direct images $\eR^qf_!$), it is easy to reduce the proof of (\ref{II:3-1}, \ref{II:3-2}) to the case in which $X_0$ is a smooth curve and $\sF_0$ is lisse. This is explained clearly in \cite[\S 5]{gr95} (for \ref{II:3-1}; \ref{II:3-2} is treated in the same way).

\subsection{}\label{II:3-8}

There are two methods for proving (\ref{II:3-2}).

\paragraph{Lefschetz-Verdier}
If $X_0$ is proper (cf. \ref{II:3-6}), the general Lefschetz--Verdier trace formula expresses the right-hand side of (\ref{II:3-2}) as a sum of local terms, one for each point of $X^{F^n}$. In the original version of \cite{sga5}, this formula was proved only assuming resolution of singularities. The reader will find an unconditional proof in the definitive version. In the case of curves, to which one can reduce (\ref{II:3-7}), all the ingredients of the proof were in fact already available.

To deduce (\ref{II:3-2}) from the Lefschetz--Verdier formula, one must be able to calculate its local terms. For a curve and the Frobenius endomorphism, this had been done by Artin and Verdier \cite{ve67}.

\paragraph{Nielsen--Wecken}
A method inspired by the work of Nielsen and Wecken reduces (\ref{II:3-2}) to a special case proved by Weil; this is explained in the following sections.

\section{Reduction to theorems in \texorpdfstring{$\dZ/\ell^n$}{Z/l n}-cohomology}\label{II:4}

We shall prove (\ref{II:3-2}) by passage to the limit from an analogous theorem in $\dZ/\ell^n$-cohomology. Stating it presents the following difficulty. Let $X_0$ be a separated scheme of finite type over $\dF_q$, and let $\sF_0$ be a flat constructible sheaf of $\dZ/\ell^n$-modules. Its stalks are finitely generated free $\dZ/\ell^n$-modules, so the left-hand side of (\ref{II:3-2}) makes sense (it is an element of $\dZ/\ell^n$). On the other hand, the $\h_c^i(X,\sF)$ will not in general be free, so the right-hand side is not defined. Overcoming this difficulty is one of Grothendieck's essential uses of the theory of derived categories.

\subsection{}\label{II:4-1}

For the definition of the derived category $\eD(\cA)$ of an abelian category $\cA$, and of the subcategories $\eD^+(\cA)$, $\eD^-(\cA)$, and $\eD^b(\cA)$, I refer to Verdier's text in this volume. For a ring $\Lambda$, we write $\eD(\Lambda)$ instead of $\eD(\text{category of left $\Lambda$-modules})$; for a site $X$, we write $\eD(X,\Lambda)$ instead of $\eD(\text{category of sheaves of left $\Lambda$-modules})$. Likewise for $\eD^-$, etc.

The symbol $\eR$ (respectively, $\mathsf{L}$) indicates a right (respectively, left) derived functor. For example, $\lotimes$ denotes the derived tensor product. For $K$ in $\eD^-(\Lambda^\circ)$ and $L$ in $\eD^-(\Lambda)$ (where $\Lambda^\circ$ is the opposite ring), $K\lotimes L$ is calculated as follows: take quasi-isomorphisms $K'\to K$ and $L'\to L$, with $K'$ and $L'$ bounded above and one of them having flat components, and form the total complex $s(K'\otimes_\Lambda L')$ associated with the double complex $K'\otimes_\Lambda L'$.

\subsection{}\label{II:4-2}

Even to treat only $\dZ/\ell^n$-cohomology, we shall need the theory of traces of endomorphisms of modules over rings that need not be commutative (finite-group algebras $\dZ/\ell^n[G]$). These traces were introduced by Stallings and Hattori \cite{ba76}.

Let $\Lambda$ be a ring (with identity, but not necessarily commutative). Let $\Lambda^\natural$ be the quotient of the additive group of $\Lambda$ by the subgroup generated by the commutators $ab-ba$. If $f$ is an endomorphism of the free left $\Lambda$-module $\Lambda^n$, with matrix $f_i^j$, let $\tr(f)$ be the image of $\sum_i f_i^i$ in $\Lambda^\natural$. If $f$ and $g$ are homomorphisms $f:\Lambda^n\leftrightarrows\Lambda^m:g$, one immediately checks that $\tr(fg)=\tr(gf)$.

Let $f$ be an endomorphism of a finitely generated projective $\Lambda$-module $P$. Write $P$ as a direct summand of a finite free module $\Lambda^n$: this means choosing a module $P'$ and an isomorphism $\alpha:P\oplus P'\iso\Lambda^n$, or equivalently choosing a homomorphism $a:P\to\Lambda^n$ and a retraction $b:\Lambda^n\to P$ (with $P'=\ker(b)$). Put $f'=\alpha(f\oplus0)\alpha^{-1}=afb$, and define $\tr(f)=\tr(f')$. This element of $\Lambda^\natural$ depends only on $f$: if another choice $c:P\leftrightarrows\Lambda^m:d$ is made, then $a=a(dc)(ba)$, and hence
\[
 \tr(a f b) = \tr(adcbafb) = \tr(cbafbad) = \tr(cfd) \text{.}
\]

If $f$ is an endomorphism of a $\dZ/2$-graded finitely generated projective $\Lambda$-module, with components $f_i^j:P^j\to P^i$, put
\[
  \tr(f) = \tr(f_0^0) - \tr(f_1^1) \text{.}
\]
When there is a risk of ambiguity, we write $\tr(f;P)$. If the homomorphisms $f:P\leftrightarrows Q:g$ are homogeneous, reduction to the free case immediately gives
\begin{equation*}\tag{4.2.1}\label{II:eq:4-2-1}
  \tr(f g) = \tr\left((-1)^{\deg f \deg g} g f\right) \text{.}
\end{equation*}

If $f$ is an endomorphism of a filtered $\dZ/2$-graded $\Lambda$-module $(P,F)$, with a finite filtration compatible with the grading, and the $\gr_F^i(P)$ are finitely generated projective, then $P$ is finitely generated projective and
\begin{equation*}\tag{4.2.2}\label{II:eq:4-2-2}
  \tr(f;P) = \sum \tr\left(f; \gr_F^i(P) \right)
\end{equation*}
(split the filtration to reduce to the case of a direct sum).

\subsection{}\label{II:4-3}

A $\dZ$-graded module will be regarded as $\dZ/2$-graded by parity of degree. In particular, if $f$ is an endomorphism of a bounded complex of finitely generated projective $\Lambda$-modules, put
\begin{equation*}\tag{4.3.1}\label{II:eq:4-3-1}
  \tr(f) = \sum (-1)^i \tr(f^i) \text{.}
\end{equation*}

By \eqref{II:eq:4-2-1}, if $f$ is homotopic to zero, then $\tr(f)=0$:
\begin{equation*}\tag{4.3.2}\label{II:eq:4-3-2}
  f = dH + Hd \rightarrow \tr(f) = 0 \text{.}
\end{equation*}

Let $\mathsf{K}_\text{parf}(\Lambda)$ be the category whose objects are bounded complexes of finitely generated projective $\Lambda$-modules and whose morphisms are complex morphisms modulo homotopy. The functor $\mathsf{K}_\text{parf}(\Lambda)\to\eD(\Lambda)$ is fully faithful. Denote its essential image by $\mathsf{D}_\text{parf}(\Lambda)$. Formulas \eqref{II:eq:4-3-1} and \eqref{II:eq:4-3-2} define $\tr(f)$ for an endomorphism $f$ of $K\in\ob\mathsf{D}_\text{parf}(\Lambda)$, and this trace depends only on the isomorphism class of $(K,f)$.

\subsection{}\label{II:4-4}

Recall that the filtered derived category $\mathsf{DF}(\cA)$ of an abelian category $\cA$ is obtained from the category of complexes of finitely filtered objects of $\cA$ by inverting filtered quasi-isomorphisms (that is, morphisms of filtered complexes inducing a quasi-isomorphism on the associated graded). There are functors ``underlying complex'' $\mathsf{DF}(\cA)\to\eD(\cA):(K,F)\mapsto K$ and ``$p$th graded component'' $\mathsf{DF}(\cA)\to\eD(\cA):(K,F)\mapsto\gr_F^p(K)$.

Let $\mathsf{KF}_\text{parf}(\Lambda)$ be the category whose objects are bounded complexes $(K,F)$ with finite filtrations and with the $\gr_F^p(K^q)$ finitely generated projective, and whose morphisms are filtration-preserving complex morphisms modulo filtration-preserving homotopy. The functor $\mathsf{KF}_\text{parf}(\Lambda)\to\mathsf{DF}(\Lambda)$ is fully faithful; denote its essential image by $\mathsf{DF}_\text{parf}(\Lambda)$. If $(K,F)$ belongs to $\mathsf{DF}_\text{parf}(\Lambda)$, then $K$ belongs to $\mathsf{D}_\text{parf}(\Lambda)$.

By \eqref{II:eq:4-2-2}, if $f$ is an endomorphism of $(K,F)\in\ob\mathsf{DF}_\text{parf}(\Lambda)$, then
\begin{equation*}\tag{4.4.1}\label{II:eq:4-4-1}
 \tr(f,K) = \sum_p \tr(f,\gr_F^p(K)) \text{.}
\end{equation*}

\subsection{}\label{II:4-5}

Let $\Lambda$ be a ring and let $K\in\ob\eD^-(\Lambda)$. The complex $K$ is said to have tor-dimension $\leqslant r$ if, for every right $\Lambda$-module $N$, the hyper-Tor groups $\h^i(N\lotimes_\Lambda K)$ vanish for $i<-r$. It is said to have finite tor-dimension if this condition holds for some $r$.

The same terminology is used for a complex of sheaves $K\in\ob\eD^-(X,\Lambda)$.

\begin{lemma}\label{II:4-5-1}
Let $\Lambda$ be a left Noetherian ring and let $K\in\ob\eD^-(\Lambda)$. Then $K$ belongs to $\mathsf{D}_\text{parf}(\Lambda)$ if and only if it has finite tor-dimension and the $\h^i(K)$ are finitely generated.
\end{lemma}

After replacing $K$ by an isomorphic complex (in $\eD^-(\Lambda)$), we may suppose that $K$ is bounded above and has finitely generated free components (cf. \ref{II:4-7} below). If $K$ has tor-dimension $\leqslant r$, then the $\h^i(K)$ vanish for $i<-r$, and $K^{-r}/\im(d)$ has the free resolution
\[
  (\cdots \to K^{-r-1} \to K^{-r}) \to K^{-r}/\im(d) \text{,}
\]
and, for every right $\Lambda$-module $N$, $\h^{-r-k}(N\lotimes_\Lambda K) = \h^{-r-k}(N\otimes_\Lambda K) = \tor_k(N,K^{-r}/\im d)$ for $k\geqslant1$.

By hypothesis, these groups vanish; hence $K^{-r}/\im d$ is flat and of finite presentation, that is, finitely generated projective. The quotient
\[
  0 \to K^{-r}/\im(d) \to K^{-r+1} \to K^{-r+2}\to \cdots
\]
of $K$ is quasi-isomorphic to $K$, which shows that $K\in\ob\mathsf{D}_\text{parf}(\Lambda)$.

\begin{prop-def_}\label{II:4-6}
Let $X$ be a Noetherian scheme and let $\Lambda$ be a left Noetherian ring. Denote by $\eD_\textnormal{ctf}^b(X,\Lambda)$ the subcategory of $\eD^-(X,\Lambda)$ consisting of the complexes $\sK$ satisfying the following equivalent conditions:
\begin{enumerate}[\indent i)]
  \item $\sK$ is isomorphic (in $\eD^-(X,\Lambda)$) to a bounded complex of flat constructible sheaves of $\Lambda$-modules.
  \item $\sK$ has finite tor-dimension and the sheaves $\h^i(\sK)$ are constructible.
\end{enumerate}
\end{prop-def_}

The same argument as in (\ref{II:4-5}) applies, starting from the following lemma.

\begin{lemma_}\label{II:4-7}
If a complex $\sK$ of sheaves of $\Lambda$-modules has constructible $\h^i(\sK)$ that vanish for sufficiently large $i$, there is a quasi-isomorphism $\sK'\to\sK$, with $\sK'$ bounded above and having flat constructible components.
\end{lemma_}

Construct $\sK'$ from right to left. For an $m$ such that the $\h^i(\sK)$ vanish for $i\geqslant m$, take ${\sK'}^i=0$ for $i\geqslant m$. Suppose next that the ${\sK'}^i$ have already been constructed for $i>n$:
\[\xymatrix{
  \cdots \ar[r]
    & \sK^{n-1} \ar[r]
    & \sK^n \ar[r]
    & \sK^{n+1} \ar[r]
    & \cdots \\
  & & & {\sK'}^{n+1} \ar[u] \ar[r]
    & \cdots \text{,}
}\]
with $\h^i(\sK')\iso\h^i(\sK)$ for $i>n+1$, and with $\ker(d)\to\h^{n+1}(\sK)$ surjective in degree $n+1$. Define sheaves $\sA$ and $\sB$ by the Cartesian diagram
\[\xymatrix{
  \sK^n \ar[r]
    & \sK^n/\im(d) \ar[r]
    & \ker(d) \\
  \sA \ar[u] \ar[r]^u
    & \sB \ar[u]\ar[r]
    & \ker(d) \ar[u] \text{.}
}\]

Then $\sB$ is constructible and $u$ is surjective. To advance one step, it is enough to construct a flat constructible $K^n$ and a map $v:K^n\to\sA$ such that $uv$ is surjective. The next lemma guarantees that this is possible.

\begin{lemma_}\label{II:4-8}
Let $v:\sA\to\sB$ be an epimorphism of sheaves of $\Lambda$-modules, with $\sB$ constructible. Then there is a map $u:\sK\to\sA$, with $\sK$ flat and constructible, such that $vu$ is an epimorphism.
\end{lemma_}

The sheaf $\sA$ is a quotient of a sum of sheaves of the form $\varphi_!\Lambda$, for $\varphi:U\to X$ étale. Since $\sB$ is Noetherian (\cite[IX.2.10]{sga4}), a finite sum of these already maps onto $\sB$, which proves the lemma.

\begin{theorem_}\label{II:4-9}
Let $f:X\to Y$ be a separated morphism of finite type between Noetherian schemes. If $\sK\in\ob\eD_\textnormal{ctf}^b(X,\Lambda)$, then $\eR f_!\sK\in\ob\eD_\textnormal{ctf}^b(Y,\Lambda)$.
\end{theorem_}

Recall that if $j:X\hookrightarrow\bar X\xrightarrow{\bar f}Y$ is a relative compactification of $X/Y$, then by definition $\eR f_!=\eR\bar f_*\circ j_!$, where $\eR\bar f_*$ is the derived functor of $\bar f_*$. This formula reduces the proof of (\ref{II:4-9}) to the case in which $f$ is proper. When $f$ is proper, $f_*$ has finite cohomological dimension; this permits the definition of $\eR f_*$ on the entire derived category, and also as a functor
\[
  \eR f_* : \eD^-(X,\Lambda) \to \eD^-(Y,\Lambda) \text{.}
\]
By composition, one similarly defines in general
\[
  \eR f_! : \eD^-(X,\Lambda) \to \eD^-(Y,\Lambda) \text{.}
\]

For every right $\Lambda$-module $N$, one has (the proof is given below)
\begin{equation*}\tag{4.9.1}\label{II:eq:4-9-1}
  N\lotimes \eR f_! \sK \iso \eR f_! (N\lotimes \sK) \text{.}
\end{equation*}

Let us deduce (\ref{II:4-9}) from \eqref{II:eq:4-9-1}.
\begin{enumerate}[\indent a)]
  \item The spectral sequence
    \[
      E_2^{p q} = \eR^p f_! \h^q(\sK) \rightarrow \h^{p+q} \eR f_! \sK \text{.}
    \]
    and the finiteness theorem for $\eR f_!$ show that $\eR f_!\sK\in\ob\eD^b(Y,\Lambda)$ and has constructible cohomology. It remains to prove finite tor-dimension (\ref{II:4-6}).
  \item If $\sK$ has finite tor-dimension, then so does $\eR f_!\sK$: one has $\h^i(N\lotimes\eR f_!\sK)=\h^i(\eR f_!(N\lotimes\sK))$, and the spectral sequence above for $N\lotimes\sK$, together with the finite cohomological dimension of $f_!$, gives a uniform lower bound for the degrees in which these groups can be nonzero.
\end{enumerate}

We prove \eqref{II:eq:4-9-1} at the level of complexes.
\begin{enumerate}[\indent a)]
  \item Using the defining formula $\eR f_!=\eR\bar f_*\circ j_!$, we reduce to the case in which $f$ is proper.
  \item Represent $\sK$ by a bounded-above complex whose components are acyclic for $f_*$: $\eR^pf_*\sK^q=0$ for $p>0$. Then $\eR f_*\sK\sim f_*\sK$. Working in this way with bounded-above complexes is possible because $f_*$ has finite cohomological dimension.
  \item If $L$ is a free right $\Lambda$-module, then $\eR^pf_*(L\otimes\sK^q)=L\otimes\eR^pf_*\sK^q$: this is immediate when $L$ is free of finite rank, and the general case follows by passage to the limit. The sheaf $L\otimes\sK^q$ is acyclic for $f_*$, and $f_*(L\otimes\sK^q)=L\otimes f_*\sK^q$.
  \item Let $N_\bullet$ be a free resolution of $N$. One has $N\lotimes\sK\sim s(N_\bullet\otimes\sK)$, the total complex associated with the double complex $N_\bullet\otimes\sK$. By c),
    \[
      \eR f_*(N\lotimes \sK) \sim \eR f_* (s(N_\bullet\otimes \sK)) \sim f_*(s(N_\bullet\otimes \sK)) \sim s(N_\bullet\otimes f_* \sK) \sim N\lotimes \eR f_* \sK \text{,}
    \]
    which is the asserted formula.
\end{enumerate}

When $Y$ is the spectrum of a separably closed field, the functor $\eR f_!$ is identified with a functor denoted $\eR \Gamma_c : \eD_\text{ctf}^b(X,\Lambda)\to\eD_\text{parf}(\Lambda)$.

\begin{theorem_}\label{II:4-10}
With the notation of \S\ref{II:1}, let $X_0$ be a separated scheme of finite type over $\dF_q$, and let $\Lambda$ be a Noetherian torsion ring annihilated by an integer prime to $p$. Let $\sK_0\in\ob\eD_\textnormal{ctf}^b(X_0,\Lambda)$. With the definition of traces in (\ref{II:4-3}), one has, for every $n$,
\begin{equation*}\label{II:eq:4-10}
  \sum_{x\in X^{F^n}} \tr\left({F^n}^*,\sK_x\right) = \tr\left({F^n}^*,\eR\Gamma_c(X,\sK)\right) \text{.}
\end{equation*}
\end{theorem_}

In the remainder of this section, we shall deduce \ref{II:3-2} from \ref{II:4-10} and prove \ref{II:2-10}.

\subsection{Proof of \texorpdfstring{\ref{II:2-10}}{2.10}}\label{II:4-11}

Let $X$ be a separated scheme of finite type over an algebraically closed field $k$, and let $\sG$ be a $\dQ_\ell$-sheaf on $X$. By \ref{II:2-8}, it can be written $\sG=\sF\otimes\dQ_\ell$, with $\sF$ a torsion-free $\dZ_\ell$-sheaf. Put
\[
  K_n = \eR \Gamma_c(X,\sF\otimes\dZ/\ell^{n+1})\in\eD(\dZ/\ell^{n+1}) \text{.}
\]
By \ref{II:4-9}, one has $K_n\in\ob\eD_\text{parf}(\dZ/\ell^{n+1})$, and it follows from \eqref{II:eq:4-9-1} that the natural map (for $n\geqslant1$)
\begin{equation*}\tag{4.11.1}\label{II:eq:4-11-1}
\xymatrix{
  K_n\lotimes_{\dZ/\ell^{n+1}} \dZ/\ell^n \ar[r] & K_{n-1} \text{,}
}
\end{equation*}
is an isomorphism. The issue of defining this map is discussed in \ref{II:4-12}.

We can now invoke \cite[XI.3.3]{sga5} (p. 31, lines 1--2 from the bottom; this passage is independent of the rest of \cite{sga5}) to represent each $K_n$ by a bounded complex of finitely generated free $\dZ/\ell^{n+1}$-modules, and the maps \eqref{II:eq:4-11-1} by \emph{isomorphisms of complexes}
\[\xymatrix{
  K_n\otimes \dZ/\ell^n \ar[r]^-\sim
    & K_{n-1} \text{.}
}\]
Let $K$ be the complex $\varprojlim K_n$. It is a bounded complex of free $\dZ_\ell$-modules, and $K_n\simeq K\otimes_{\dZ_\ell}\dZ/\ell^{n+1}$. One has
\[
  \h^i(K) = \varprojlim \h^i(K_n) \text{,}
\]
(indeed, every inverse system of finite groups satisfies the Mittag--Leffler condition), and $\h^i(K)$ is a finitely generated $\dZ_\ell$-module. Thus
\[
  \h_c^i(X,\sG) = \left(\varprojlim \h^i(K_n) \right)\otimes_{\dZ_\ell} \dQ_\ell \text{:}
\]
is a finite-dimensional $\dQ_\ell$-vector space.

\subsection{The map \texorpdfstring{\eqref{II:eq:4-11-1}}{(4.11.1)}}\label{II:4-12}

To \emph{define} the map \eqref{II:eq:4-11-1}, one cannot proceed as in \ref{II:4-9}, that is, resolve $\dZ/\ell^n$ by a complex of free $\dZ/\ell^{n+1}$-modules, since one wants to construct an isomorphism in $\eD_\text{parf}(\dZ/\ell^n)$. The desired map is a special case of the following extension-of-scalars map.

(*) Let $\Lambda\to\Lambda'$ be a homomorphism of Noetherian torsion rings. For $X$ a separated scheme of finite type over an algebraically closed field $k$, and $\sK\in\ob\eD_\text{ctf}(X,\Lambda)$, there is an isomorphism
\[\xymatrix{
  \eR\Gamma_c(X,\sK)\lotimes_\Lambda \Lambda' \ar[r]^-\sim & \eR \Gamma_c(X,\sK\lotimes_\Lambda \Lambda') \text{.}
}\]

More general maps are constructed in \cite[XVII 4.2.12]{sga4}. The method is as follows.
\begin{enumerate}[\indent a)]
  \item The definition of $\eR\Gamma_c$ and the formula $j_!(\sK\otimes_\Lambda\Lambda')=(j_!\sK)\otimes_\Lambda\Lambda'$ for an open immersion reduce us to the case in which $X$ is proper.
  \item For proper $X$, one replaces $\sK$ by a bounded complex whose components are acyclic for $\Gamma$ and whose stalk at every geometric point (or merely every closed point) is homotopy equivalent to a complex of flat $\Lambda$-modules. For such a complex, $\Gamma(X,\sK)\sim\eR\Gamma(X,\sK)$ and $\sK\otimes_\Lambda\Lambda'\sim\sK\lotimes_\Lambda\Lambda'$; the map \eqref{II:eq:4-11-1} is defined as the composite
    \[
      \eR\Gamma(X,\sK)\lotimes_\Lambda \Lambda' \to \Gamma(X,\sK)\otimes_\Lambda\Lambda' \to \Gamma(X,\sK\otimes_\Lambda\Lambda') \xleftarrow\sim \eR\Gamma(X,\sK\lotimes_\Lambda\Lambda') \text{.}
    \]
  \item It remains to show that $\sK$ can be replaced by a complex of the desired kind \cite[XVII.4.2.10]{sga4}. First replace $\sK$ by a bounded complex with flat components (\ref{II:4-6}), then take a truncated canonical flasque resolution. This does not change the homotopy type of the stalk complexes, and one truncates far enough for the acyclicity condition for $\Gamma$ to hold, that is, beyond the cohomological dimension of $\Gamma$.
\end{enumerate}

\subsection{Proof of \texorpdfstring{(\ref{II:3-2})}{(3.2)}}\label{II:4-13}

To simplify notation, we treat only the case $n=1$ of (\ref{II:3-2}). The general case follows by replacing $F$ below by $F^n$.

Let $X_0$ therefore be separated and of finite type over $\dF_q$, and let $\sG_0$ be a $\dQ_\ell$-sheaf on $X_0$. One has $\sG_0=\sF_0\otimes\dQ_\ell$ for a suitable torsion-free $\dZ_\ell$-sheaf $\sF_0$. Put $K_n=\eR\Gamma_c(X,\sF\otimes\dZ/\ell^{n+1})$. The Frobenius morphisms induce morphisms
\[
  F^* : K_n\to K_n \qquad \text{(in $\eD_\text{parf}(\dZ/\ell^{n+1})$),}
\]
which are obtained from one another via the isomorphisms \eqref{II:eq:4-11-1}.

As in (\ref{II:4-11}), represent the $K_n$ and the isomorphisms \eqref{II:eq:4-11-1} at the level of complexes. The already cited passage \cite[XI.3.3]{sga5} permits the morphisms $F^*$ to be represented by endomorphisms of complexes, still denoted $F^*$, that reduce to one another. Their inverse limit and the induced endomorphisms are again denoted $F^*$. Since $\h_c^i(X,\sG)=\h^i(K)\otimes\dQ_\ell=\h^i(K\otimes\dQ_\ell)$, one has, with the sign convention of \eqref{II:eq:4-3-1},
\[
  \tr\left(F^*, \h_c^\bullet(X,\sG)\right) = \tr\left(F^*,K^\bullet\otimes\dQ_\ell\right) = \varprojlim_n \tr\left(F^*,K_n^\bullet\right) = \varprojlim_n \tr\left(F^*,\eR\Gamma_c(X,\sF\otimes\dZ/\ell^{n+1})\right) \text{.}
\]
Apply (\ref{II:4-10}) to $\sF_0\otimes\dZ/\ell^{n+1}$, viewed as a complex concentrated in degree $0$, with $\Lambda=\dZ/\ell^{n+1}$. One finds
\begin{align*}
  \tr\left(F^*,\eR\Gamma_c(X,\sF\otimes\dZ/\ell^{n+1})\right)
    &= \sum_{x\in X^F} \tr\left(F^*,\sF_x\otimes\dZ/\ell^{n+1}\right) \\
    &= \sum_{x\in X^F} \tr\left(F^*,\sG_x\right) \mod \ell^{n+1} \text{,}
\end{align*}
and (\ref{II:3-2}) follows by passage to the limit.

\section{The Nielsen--Wecken method}\label{II:5}

In this section we prove two special cases of (\ref{II:4-10}). The general case will be considered in the next section.

\begin{lemma_}\label{II:5-1}
Theorem (\ref{II:4-10}) holds when $X_0$ has dimension $0$.
\end{lemma_}

In dimension $0$, formula \eqref{II:eq:4-10} holds for every correspondence, not only for iterates of a Frobenius. It is a special case of the following statement (applied to $X(\dF)$, $K$, and the ${F^*}^n$).

(*) Let $X$ be a finite set, let $F:X\to X$, let $\sK\in\ob\eD_\text{parf}(X,\Lambda)$, and let $F^*:F^*\sK\to\sK$. Then
\[
  \sum_{X^F} \tr\left(F^*,\sK_x\right) = \tr\left(F^*,\Gamma(X,K)\right) \text{.}
\]

One has $\Gamma(X,\sK)=\bigoplus_{x\in X}\sK_x$, and the formula follows from the fact that, for every morphism $u:\bigoplus_{x\in X}\sK_x\to\bigoplus_{x\in X}\sK_x$ with matrix $u_x^y$, one has $\tr(u)=\sum_{x\in X}\tr(u_x^x)$.

\begin{lemma_}\label{II:5-2}
Let $\Lambda$ be a Noetherian torsion ring annihilated by a power of the prime $\ell\ne p$, let $X_0$ be a smooth connected projective curve over $\dF_q$, let $j:U_0\hookrightarrow X_0$ be a dense open subset, and let $\sG_0$ be a locally constant sheaf of projective $\Lambda$-modules on $U_0$. Then
\[
  \sum_{x\in U^F} \tr\left(F^*,\sG_x\right) = \tr\left(F^*,\eR\Gamma_c(U,\sG)\right) \text{.}
\]
\end{lemma_}

More precisely, we shall deduce (\ref{II:5-2}) from the following theorem.

\begin{theorem_}\label{II:5-3}
Let $X$ be a smooth projective curve over an algebraically closed field $k$, let $f$ be an endomorphism of $X$ with isolated fixed points, and let $v(f)$ be the number of fixed points of $f$, each counted with its multiplicity; $v(f)$ is the intersection number of the graph of $f$ with the diagonal of $X\times X$. If $\ell$ is prime to the characteristic exponent of $k$, then
\[
  v(f) = \sum (-1)^i \tr\left(f^*,\h^i(X,\dQ_\ell)\right) \text{.}
\]
\end{theorem_}

This theorem is due to Weil (\cite{we48}, nos. 43, 66, and 68), and Weil's proof is reproduced in Lang \cite{la83} (VI \S3 combined with VII \S2, Theorem 3). It can also be deduced from the formalism of the cohomology class associated with a cycle (\nameref{IV}, \ref{IV:3-7}).

\begin{corollary_}\label{II:5-4}
Let $j:U\hookrightarrow X$ be a dense open subset such that $f^{-1}(U)=U$. If the fixed points of $f$ in $X\setminus U$ have multiplicity one, then the number $v_U(f)$ of fixed points of $f$ in $U$ (each counted with its multiplicity) is
\[
  v_U(f) = \sum (-1)^i \tr\left(f^*,\h_c^i(U,\dQ_\ell)\right) \text{.}
\]
\end{corollary_}

Let $F$ be the finite set $X\setminus U$. The cohomology exact sequence defined by the short exact sequence $0\to j_! \dQ_\ell\to\dQ_\ell\to{\dQ_\ell}_F\to 0$:
\[\xymatrix{
  \ar[r]^-\partial
    & \h_c^i(U,\dQ_\ell) \ar[r]
    & \h^i(X,\dQ_\ell) \ar[r]
    & \h^i(F,\dQ_\ell) \ar[r]^-\partial
    &
}\]
gives, with the sign convention of (\ref{II:3-1}),
\[
  \tr\left(f^*,\h_c^\bullet(U,\dQ_\ell)\right)
    = \tr\left(f^*,\h^\bullet(X,\dQ_\ell)\right) - \tr\left(f^*,\h^\bullet(F,\dQ_\ell)\right)
    = v(f) - \tr(f^*,\dQ_\ell^F) \text{.}
\]
The trace of $f^*$ on $\dQ_\ell^F$ is the number $v_F(f)$ of fixed points of $f$ in $F$; the multiplicity-one hypothesis therefore ensures that
\[
  \tr\left(f^*,\h_c^\bullet(U,\dQ_\ell)\right)
    = v_U(f) \text{,}
\]
as asserted.

The proof of (\ref{II:5-2}) uses a few lemmas, which we now develop, about a special case of noncommutative traces (\ref{II:4-2}). \emph{For brevity, we shall say projective to mean finitely generated projective.}

\subsection{}\label{II:5-5}

Let $\Lambda$ be a ring and let $H$ be a finite group. The map from the group algebra $\Lambda[H]$ to $\Lambda^\natural$ that sends $\sum a_hh$ to the class of $a_e$ vanishes on the commutators $ab-ba$ of $\Lambda[H]$. It therefore factors through
\[
  \varepsilon : \Lambda[H]^\natural \to \Lambda^\natural \text{.}
\]
If $F$ is an endomorphism of a projective $\Lambda[H]$-module $P$, put
\[
  \tr_\Lambda^H(F) = \varepsilon \tr_{\Lambda[H]} (F)
\]
where the trace on the right is the one from (\ref{II:4-2}). To avoid ambiguity, we shall sometimes write this trace as $\tr_\Lambda^H(F,P)$.

\begin{proposition_}\label{II:5-6}
$\tr_\Lambda(F) = |H|\tr_\Lambda^H(F)$.
\end{proposition_}

The formal properties of traces reduce us to the case $P=\Lambda[H]$. The endomorphism $F$ is then right multiplication by an element $\sum a_hh$ of $\Lambda[H]$, and $\tr_\Lambda(F)=|H|a_e$, while $\tr_\Lambda^H(F)=a_e$.

\subsection{}\label{II:5-7}

Let $A$ be a commutative ring and let $\Lambda$ be an $A$-algebra. Multiplication by an element of $A$ descends to the quotient and defines an $A$-module structure on $\Lambda^\natural$. If $H$ is a finite group, $P$ is a projective $A[H]$-module, and $M$ is a $\Lambda[H]$-module projective as a $\Lambda$-module, then the $\Lambda$-module $P\otimes_AM$, endowed with the diagonal action of $H$, is a projective $\Lambda[H]$-module. Indeed:
\begin{enumerate}[\indent a)]
  \item It is enough to check this for $P=A[H]$.
  \item The $\Lambda$-module $P\otimes_AM$, endowed with the $H$-action on the first factor, is clearly a projective $\Lambda[H]$-module; denote it by $(P\otimes_AM)'$.
  \item The map $h\otimes m\mapsto h\otimes h^{-1}m$ induces an isomorphism
    \begin{equation*}\tag{5.7.1}\label{II:eq:5-7-1}
      A[H]\otimes_A M \iso (A[H]\otimes_A M)' \text{.}
    \end{equation*}
\end{enumerate}

\begin{proposition_}\label{II:5-8}
With the preceding notation, if $u$ is an endomorphism of $P$ and $v$ an endomorphism of $M$, then
\[
  \tr_\Lambda^H(u\otimes v) = \tr_A^H(u)\tr_\Lambda(v) \text{.}
\]
\end{proposition_}

One reduces to $P=A[H]$; the endomorphism $u$ is then right multiplication $m_x$ by an element $x=\sum a_hh$ of $A[H]$. The isomorphism \eqref{II:eq:5-7-1} transforms $u\otimes v$ into $\sum a_hm_h\otimes h^{-1}v$, whose trace is $a_e\cdot\tr_\Lambda(v)$.

\subsection{}\label{II:5-9}

Let $G$ be an extension of $\dZ$ by a finite group $H$,
\[
  1 \to H \to G \to \dZ \to 1 \text{,}
\]
and let $G^+$ be the inverse image of $\dN$. Let $\Lambda$ be a ring and let $P$ be a $\Lambda[G^+]$-module projective as a $\Lambda[H]$-module. If $g\in G^+$, let $Z_g$ be the centralizer of $g$ in $H$; multiplication by $g$ is an endomorphism of $P$, viewed as a $\Lambda[Z_g]$-module. Since $P$ is projective over $\Lambda[Z_g]$, a trace $\tr_\Lambda^{Z_g}(g)$ is defined. By \ref{II:5-6},
\begin{equation*}\tag{5.9.1}\label{II:eq:5-9-1}
  \tr_\Lambda(g) = |Z_g| \cdot \tr_\Lambda^{Z_g}(g) \text{.}
\end{equation*}

Suppose that $\Lambda$ is an algebra over a commutative ring $A$. If $P$ is an $A[G^+]$-module projective as an $A[H]$-module and $M$ is a $\Lambda[G^+]$-module projective as a $\Lambda$-module, then the $\Lambda[G^+]$-module $P\otimes_AM$, with the diagonal $G^+$-action, is projective as a $\Lambda[H]$-module and, by \ref{II:5-8}, for every $g\in G^+$ one has
\begin{equation*}\tag{5.9.2}\label{II:eq:5-9-2}
  \tr_\Lambda^{Z_g}(g,P\otimes_A M) = \tr_A^{Z_g}(g,P) \cdot \tr_\Lambda(g,M) \text{.}
\end{equation*}

\subsection{}\label{II:5-10}

Let $P$ be a $\Lambda[G^+]$-module projective as a $\Lambda[H]$-module, and let $P_H$ be the $H$-coinvariants of $P$. This is a projective $\Lambda$-module. The action of $g\in G^+$ descends to the quotient and defines an endomorphism of $P_H$ depending only on the image of $g$ in $\dN$. Denote by $F$ the action of the elements $g\in G^+$ with image $1$.

\begin{proposition_}\label{II:5-11}
$\tr_\Lambda(F,P_H)=\sum_{g\mapsto1}'\tr_\Lambda^{Z_g}(g,P)$, where $\sum_{g\mapsto1}'$ denotes a sum over the $H$-conjugacy classes of elements of $G^+$ with image $1$ in $\dN$.
\end{proposition_}

After multiplication by $|H|$, this formula can be verified as follows:
\[
  |H|\cdot \tr_\Lambda(F,P_H)
    = \tr_\Lambda\left(\sum_{g\mapsto 1} g,P\right)
    = \sum_{g\mapsto 1}' \frac{|H|}{|Z_g|}\tr_\Lambda(g,P)
    = \sum_{g\mapsto 1}' |H|\cdot \tr_\Lambda^{Z_g}(g,P) \text{.}
\]

Choose an element $g\in G^+$ with image $1$ in $\dN$, and let $\sigma$ be the automorphism of $\Lambda[H]$ induced by the automorphism $\operatorname{ad}g$ of $H$. Giving a $\Lambda[G^+]$-module is equivalent to giving a $\Lambda[H]$-module endowed with a $\sigma$-linear endomorphism $\gamma$: let $g^nh$ act as $\gamma^nh$. In this language, formula \ref{II:5-11} takes the following form: for $P$ a projective $\Lambda[H]$-module and $\gamma$ a $\sigma$-linear endomorphism of $P$,
\[
  \tr_\Lambda(\gamma,P_H) = \sum_h' \tr_\Lambda^{Z_g}(h\gamma,P)\text{,}
\]
where $\sum'$ is a sum over the $\operatorname{ad}g$-conjugacy classes in $H$ (the orbits of $H$ acting on itself by $x\mapsto hx\operatorname{ad}g(h)^{-1}$).

In this form, the properties of traces immediately reduce us to the case $P=\Lambda[H]$. The endomorphism $\gamma$ then has the form $x\mapsto\sigma(x)a$, with $a=\sum a_hh$ in $\Lambda[H]$. It is enough to treat the universal case in which the $a_h$ are indeterminates. In that case, $\Lambda=\dZ\langle a_h\rangle_{h\in H}$ (noncommutative polynomials) and $\Lambda^\natural$ is torsion-free: \ref{II:5-11} follows from the preceding argument by division by $|H|$. One can also calculate directly.

\subsubsection{Remark}\label{II:5-11-1}

Everything said above applies equally to $\dZ/2$-graded modules and to complexes of modules (cf. \ref{II:4-2}, \ref{II:4-3}).

\subsection{}\label{II:5-12}

Retain the notation of \ref{II:5-2}, and let $A=\dZ/\ell^n$, with $n$ large enough that $\Lambda$ is an $A$-algebra. Let $f:V_0\to U_0$ be a connected finite étale covering of $U_0$, Galois with group $H$, such that the sheaf $f^*\sG_0$ is constant. Let $M=\h^0(V_0,f^*\sG_0)$ be its constant value; this is a $\Lambda[H]$-module projective as a $\Lambda$-module.

The sheaf $f_*A$ on $U_0$, endowed with the natural action of $H$, is a locally free sheaf of rank one of $A[H]$-modules. Thus the complex $\eR\Gamma_c(U,f_*A)$ is an object of $\eD_\text{parf}(A[H])$. It is endowed with a Frobenius endomorphism $F^*$.

For the natural action of $H$ on $f_*f^*\sG_0$, the trace morphism $f_*f^*\sG_0\to\sG_0$ factors through an isomorphism
\[
  \left(f_* f^*\sG_0\right)_H \to \sG_0 \text{.}
\]
Moreover, $f_*f^*\sG_0=f_*A\otimes_AM$, with diagonal $H$-action. The isomorphism
\[\xymatrix{
  \sG_0
    & \ar[l]_-\sim (f_* A\otimes_A M)_H = (f_* A\otimes_A M)\otimes_{\Lambda[H]} \Lambda
}\]
derives (\ref{II:4-12}) to an isomorphism
\[
  \eR\Gamma_c(U,\sG) \sim \left(\eR\Gamma_c(U,f_* A)\otimes_A M\right)\otimes_{\Lambda[H]} \Lambda \text{.}
\]

The endomorphism $F^*$ of the left-hand side is deduced from the endomorphism $F^*$ of $\eR\Gamma_c(U,f_*A)$. The complex of $\Lambda[H]$-modules $\eR\Gamma_c(U,f_*A)$, endowed with $F^*$, can be viewed as a complex of $\Lambda[G^+]$-modules, for the extension $G^+=H\rtimes\dN$. Formulas \ref{II:5-10} and \eqref{II:eq:5-9-2}, extended by \ref{II:5-11-1}, therefore give
\[
  \tr\left(F^*,\eR\Gamma_c(U,\sG)\right) = \sum_h' \tr_A^{Z_g}\left(h F^*,\eR\Gamma_c(U,f_* A)\right)\cdot \tr_\Lambda(h,M)
\]
where $\sum'$ is a sum over the conjugacy classes in $H$. Moreover,
\[
  |Z_h|\cdot \tr_A^{Z_h}\left(h F^*,\eR\Gamma_c(U,f_*A)\right) = \tr_A\left((F h^{-1})^*,\eR\Gamma_c(V,A)\right) \text{.}
\]
This formula remains true for $A=\dZ/\ell^m$ as $m$ increases. Passing to the limit, one finds that $\tr_A^{Z_h}\left(hF^*,\eR\Gamma_c(U,f_*A)\right)$ is
\[
  \frac{1}{|Z_h|} \sum (-1)^i \tr\left((F h^{-1})^*,\h_c^i(V,\dQ_\ell)\right)
\]
(an element of $\dZ_\ell$), and
\begin{equation*}\tag{5.12.1}\label{II:eq:5-12-1}
\tr\left(F^*,\eR\Gamma_c(U,\sG)\right) = \sum_h' \frac{1}{|Z_h|}
  \left(\sum_i(-1)^i\tr\left((F h^{-1})^*,\h_c^i(V,\dQ_\ell)\right)\right)
  \cdot \tr_\Lambda(h,M) \text{.}
\end{equation*}

\subsection{}\label{II:5-13}

It remains to apply \ref{II:5-4}. Notice that, if $\bar V$ is the smooth projective completion of $V$, all fixed points of $Fh^{-1}$ have multiplicity one. This can be done without any calculation:
\begin{enumerate}[\indent A.]
  \item If $U^F=\varnothing$, then $\tr\left(F^*,\eR\Gamma_c(U,\sG)\right)=0$.
\end{enumerate}

Indeed, $Fh^{-1}$ has no fixed point on $V$; \ref{II:5-4} shows that $\sum(-1)^i\tr\left((Fh^{-1})^*,\h_c^i(V,\dQ_\ell)\right)=0$, and one applies \eqref{II:eq:5-12-1}.
\begin{enumerate}[\indent A.]
\setcounter{enumi}{1}
  \item In general, let $j:U'=U\setminus U^F\hookrightarrow U$. The filtration $j_!j^*\sG_0\subset\sG_0$ of $\sG_0$ induces a filtration of $\eR\Gamma_c(U,\sG)$ with successive quotients $\eR\Gamma_c(U',\sG)$ and $\eR\Gamma_c(U^F,\sG)$. By \eqref{II:eq:4-4-1} and \ref{II:5-1},
    \[
      \tr\left(F^*,\eR\Gamma_c(U,\sG)\right) = \tr\left(F^*,\eR\Gamma_c(U',\sG)\right) + \sum_{x\in U^F} \tr(F^*,\sG_x) \text{,}
    \]
    and, since ${U'}^F=\varnothing$, the first term on the right is $0$.
\end{enumerate}

\section{Completion of the proof of \ref{II:4-10}}\label{II:6}

We prove \ref{II:4-10} under the additional hypotheses that the ring $\Lambda$ is annihilated by a power of a prime $\ell\ne p$ and that $n=1$. The general case follows: decompose $\Lambda$ into its $\ell$-primary components (as $\ell$ ranges over the primes), and replace $\dF_q$ by $\dF_{q^n}$ and $X_0$ by $X_0\otimes_{\dF_q}\dF_{q^n}$.

For $X_0$ separated and of finite type over $\dF_q$, and $\sK_0\in\ob\eD_\text{ctf}(X_0,\Lambda)$, put
\[
  T'(X_0,\sK_0) = \sum_{x\in X^F} \tr(F^*,\sK_x)
\]
and
\[
  T''(X_0,\sK_0) = \tr(F^*,\eR\Gamma_c(X,\sK)) \text{.}
\]
We must prove the identity
\[
  (1;X_0,\sK_0) \qquad T'(X_0,\sK_0) = T''(X_0,\sK_0) \text{.}
\]
We begin by observing that $T'$ and $T''$ (denoted generically by $T$) obey the same formalism.

(A) If $j:U_0\hookrightarrow X_0$ is an open subset with closed complement $Y_0$, then
\[
  T(X_0,\sK_0) = T(U_0,\sK_0|U_0) + T(Y_0,\sK_0|Y_0) \text{.}
\]
This is clear for $T'$. For $T''$, observe that the filtration $0\subset j_!j^*\sK_0\subset\sK_0$ of $\sK_0$ induces a filtration of $\eR\Gamma_c(X,\sK)$ with successive quotients $\eR\Gamma_c(U,\sK)$ and $\eR\Gamma_c(Y,\sK)$, and invoke \eqref{II:eq:4-4-1}. More precisely, $\sK_0$, endowed with this filtration, defines an object of $\mathsf{DF}(X,\Lambda)$; applying $\eR\Gamma_c$ gives an object of $\mathsf{DF}_\text{parf}(\Lambda)$ whose underlying object is $\eR\Gamma_c(X,\sK)$ and whose associated graded is as stated.

(B) If $\sK_0$ is a bounded complex with flat constructible components, then
\[
  T(X_0,\sK_0) = \sum (-1)^i T(X_0,\sK^i) \text{.}
\]
This is clear for $T'$. For $T''$, apply the preceding argument to the stupid filtration of $\sK_0$ to obtain
\[
  T''(X_0,\sK_0) = \sum T''(X_0,\sK_0^i[-i]) = \sum (-1)^i T''(X_0,\sK_0^i) \text{.}
\]

(C) For a morphism $f:X_0\to Y_0$, one has
\[
  T''(X_0,\sK_0) = T''(Y_0,\eR f_! \sK_0)
\]
and the same identity holds for $T'$ if the fibers of $f$ at the rational points of $Y_0$ satisfy $\left(1;f^{-1}(y),\sK_0|f^{-1}(y)\right)$; one uses
\[
  T'(X_0,\sK_0) = \sum_{y\in Y^F} T'\left(f^{-1}(y), \sK_0|f^{-1}(y)\right)
\]
and
\[
  \eR\Gamma_c(Y, \eR f_! \sK) = \eR\Gamma_c(X,\sK) \text{.}
\]

Using this formalism and \ref{II:4-7}, the proof of \ref{II:4-10} reduces to the two special cases \ref{II:5-1} and \ref{II:5-2}.
% END INLINED FILE: sga4.5-2.tex
\clearpage
\clearpage
% BEGIN INLINED FILE: sga4.5-3.tex
\chapter{\texorpdfstring{$L$}{L}-functions modulo \texorpdfstring{$\ell^n$}{l n} and modulo \texorpdfstring{$p$}{p}}\label{III}

In this exposé, I prove a formula analogous to the theorem in the report on the trace formula (in this volume; cited below as ``Report'') for a coefficient ring $A$ that is either of torsion prime to $p$, or a field of characteristic $p$. We shall make essential use of \cite[XVII 5.5]{sga4}.

\section{Symmetric tensors}\label{III:1}

\subsection{}\label{III:1-1}

Let $A$ be a commutative ring. For a list of properties of the functors $\Gamma^n:(\text{$A$-modules})\to(\text{$A$-modules})$, I refer to \cite[XVII 5.5.1 and 2]{sga4}. We recall only that, for $M$ flat (in fact, we shall only need the case in which $M$ is finitely generated projective), one has
\[
  \Gamma^n M = \left(M^{\otimes n}\right)^{S_n} \qquad
  \text{(symmetric tensors of degree $n$).}
\]

\subsection{}\label{III:1-2}

Let $X$ be an $S$-scheme, let $X^n/S$ be the $n$-fold fiber power of $X$, and let $\pi$ be the projection from $X^n/S$ onto $(X^n/S)/S_n=\operatorname{Sym}_S^n(X)$ (assumed to exist). If $\sG$ is a sheaf of $A$-modules on $X$, the external tensor product of $n$ copies of $\sG$, denoted $\sG^{\boxtimes n}$, is an $S_n$-equivariant sheaf on $X^n/S$. We shall use the \emph{external} functor $\Gamma^n$ of \cite[XVII 5.5.7--9]{sga4}. We recall only that, for $\sG$ a sheaf of flat $A$-modules (in fact, only the constructible flat case will concern us), one has
\[
  \Gamma_\text{ext}^n(\sG) = \left(\pi_*\sG^{\boxtimes n}\right)^{S_n} \text{.}
\]

\begin{lemma_}\label{III:1-3}
Let $S=\spec(k)$, where $k$ is algebraically closed, let $X$ be a finite disjoint union of copies of $S$, and let $\sG$ be a flat sheaf of $A$-modules on $X$. Then
\[
  \Gamma\left(\operatorname{Sym}_S^n(X),\Gamma_\text{ext}^n(\sG)\right) = \Gamma^n \Gamma(X,\sG) \text{.}
\]
\end{lemma_}

Indeed,
\begin{align*}
  \Gamma\left(\operatorname{Sym}_S^n(X),\Gamma_\text{ext}^n(\sG)\right)
    &= \Gamma\left(\operatorname{Sym}_S^n(X),\pi_*\sG^{\boxtimes n}\right)^{S_n} \\
    &= \Gamma\left(X^n/S,\sG^{\boxtimes n}\right)^{S_n}
    = \left(\Gamma(X,\sG)^{\otimes n}\right)^{S_n}
    = \Gamma^n \Gamma(X,\sG) \text{.}
\end{align*}

\subsection{}\label{III:1-4}

We shall need the derived functors of the nonadditive functors $\Gamma^n$ and $\Gamma_\text{ext}^n$.

The theory of such derived functors is due to Dold and Puppe \cite{do61}. For a fuller review than the one given below, I refer to \cite[XVII 5.5.3]{sga4}.

Let $\cA$ be an abelian category (or even an additive category in which every idempotent endomorphism is the projection onto a direct summand). Let $N$ denote the normalization functor
\[
  N:(\text{cosimplicial complexes of objects of $\cA$})
  \longrightarrow
  (\text{cochain complexes of objects of $\cA$, with $K^i=0$ for $i<0$}).
\]
It is an equivalence of categories. Both it and its inverse carry homotopic morphisms to homotopic morphisms. If $T$ is a functor $\cA\to\cB$, and if $K$ is a complex of objects of $\cA$ concentrated in nonnegative degrees ($K^i=0$ for $i<0$), put $T K=N T N^{-1}K$.

\subsection{Example}\label{III:1-5}

If $K$ is concentrated in degree $0$, with value $M$, then $TK$ is concentrated in degree $0$, with value $TM$.

\subsection{}\label{III:1-6}

If $K$ is a bounded complex in nonnegative degrees of finitely generated projective $A$-modules, the complex of $A$-modules $\Gamma^nK$ has the same properties.

Let $\eD_\text{parf}^{\geqslant 0}(A)$ denote the subcategory of $\eD_\text{parf}(A)$ (Report, \ref{II:4-3}) consisting of complexes of Tor-dimension $\leqslant0$. One sees that $\Gamma^n$ induces a functor
\[
  \eL\Gamma^n : \eD_\text{parf}^{\geqslant 0}(A) \to \eD_\text{parf}^{\geqslant 0}(A) \text{.}
\]

For the analogous, but more complicated, definition of $\eL\Gamma_\text{ext}^n$, I refer to \cite[XVII 5.5.14]{sga4}.

\subsection{}\label{III:1-7}

Let $K$ be a bounded complex of finitely generated projective $A$-modules, and let $u$ be an endomorphism of $K$. We denote by $\det(1-ut,K)$ the formal power series with constant term $1$
\[
  \det(1-u t,K) = \prod_i \det(1-u t,K^i)^{(-1)^i} \text{.}
\]

We leave verification of the following properties to the reader.

\subsubsection{}\label{III:1-7-1}

Let $F$ be a finite filtration stable under $u$, such that the $\gr_F^p(K^q)$ are finitely generated projective. Then
\[
  \det(1-u t,K) = \prod_p \det(1-u t,\gr_F^p(K)) \text{.}
\]

\subsubsection{}\label{III:1-7-2}

If $K$ is shifted $q$ places to the left, then
\[
  \det(1-u t,K[q]) = \det(1-u t,K)^{(-1)^q} \text{.}
\]

\begin{proposition_}\label{III:1-8}
Let $u$ be an endomorphism of a bounded complex in nonnegative degrees of finitely generated projective $A$-modules. Then
\begin{equation*}\tag{1.8.1}\label{III:eq:1-8-1}
  \det(1-u t,K)^{-1} = \sum_n \tr(u,\Gamma^n K) t^n \text{.}
\end{equation*}
\end{proposition_}

Let $D(u,K)$ denote the right-hand side of \eqref{III:eq:1-8-1}.

\begin{lemma_}\label{III:1-9}
Let $u$ be an endomorphism of a short exact sequence of bounded complexes in degrees $\geqslant0$ of finitely generated projective $A$-modules
\[
  0 \to K' \to K \to K'' \to 0 \text{.}
\]
Then $D(u,K)=D(u,K')\cdot D(u,K'')$.
\end{lemma_}

If $0\to M'\to M\to M''\to0$ is a split short exact sequence of $A$-modules, $\Gamma^nM$ has a natural filtration whose successive quotients are the $\Gamma^iM'\otimes\Gamma^jM''$ ($i+j=n$). Consequently, $\Gamma^nK$ has a natural filtration whose successive quotients are the complexes $N(N^{-1}\Gamma^iK'\otimes N^{-1}\Gamma^jK'')$ ($i+j=n$), canonically homotopic to the total complexes $s(\Gamma^iK'\otimes\Gamma^jK'')$ associated with the double complexes $\Gamma^iK'\otimes\Gamma^jK''$, and
\[
  \tr(u,\Gamma^n K) = \sum_{i+j=n} \tr(u,\Gamma^i K')\cdot \tr(u,\Gamma^j K'') \text{,}
\]
which proves the lemma. By induction, \ref{III:1-9} gives the following.

\begin{lemma_}\label{III:1-10}
Let $F$ be a finite $u$-stable filtration of $K$, such that the $\gr_F^p(K^q)$ are finitely generated projective. Then
\[
  D(u,K) = \prod_p D\left(u,\gr_F^p(K)\right) \text{.}
\]
\end{lemma_}

\begin{lemma_}\label{III:1-11}
Let $u$ be an endomorphism of a finitely generated projective $A$-module $M$, and let $M[-q]$ be the complex concentrated in degree $q$, with value $M$. Then
\[
  D\left(u,M[-q]\right) = \left(\sum \tr(u,\Gamma^n M) t^n\right)^{(-1)^q} \qquad \text{($q\geqslant 0$.)}
\]
\end{lemma_}

By \ref{III:1-5}, this is true for $q=0$. Proceeding by induction, it remains to show that
\[
  D\left(u,M[-q]\right) \cdot D\left(u,M[-(q+1)]\right) = 1 \text{.}
\]
This formula follows from \ref{III:1-9}, applied to the complex $K$ equal to $M$ in degrees $q$ and $q+1$ (with $d=\operatorname{id}$), and to its stupid filtration:
\[
  (\cdots 0 \to M \cdots) \to (\cdots M \to M) \to (\cdots M \to 0 \to 0 \cdots) \text{.}
\]
Indeed, $K$ is homotopic to $0$. Thus, in the derived category, $\Gamma^n(K)=\Gamma^n(0)$: this is $0$ for $n>0$ and $A$ for $n=0$; hence $D(u,K)=1$.

\subsection{}\label{III:1-12}

Let us prove \ref{III:1-8}. By \ref{III:1-10}, applied to the ``stupid'' filtration of $K$ by the subcomplexes
\[
  0 \to \cdots \to 0 \to K^q \to K^{q+1} \to \cdots \text{,}
\]
one has
\[
  D(u,K) = \prod_q D\left(u,K^q[-q]\right) =_{\text{(\ref{III:1-10})}} \prod_q \left(\sum \tr\left(u,\Gamma^n(K^q)\right) t^n \right)^{(-1)^q} \text{.}
\]
It remains to prove that, for $u$ an endomorphism of a finitely generated projective module $M$, one has
\begin{equation*}\tag{1.12.1}\label{III:eq:1-12-1}
  \det\left(1-u t,M\right)^{-1} = \sum \tr(u,\Gamma^n M) t^n \text{.}
\end{equation*}
If $u=0$, both sides equal $1$. By adjoining to $M$ a module $N$ on which $u=0$, and applying \ref{III:1-9}, we reduce to the case in which $M$ is free. Choose a basis. It is then enough to prove algebraic identities among the coordinates of $u$. The principle of extension of algebraic identities reduces us to the case in which $A$ is an algebraically closed field. The module $M$ then has a $u$-stable filtration with successive quotients of rank $1$, and \ref{III:1-10} reduces us to the case in which $M$ has rank $1$. Formula \eqref{III:eq:1-12-1} then reduces to the identity
\[
  \frac{1}{1-a t} = \sum a^n t^n \text{.}
\]

\begin{corollary_}\label{III:1-13}
Let $u$ and $u'$ be endomorphisms of bounded complexes $K$ and $K'$ of finitely generated projective $A$-modules. If, in the derived category, $K$, equipped with $u$, is isomorphic to $K'$, equipped with $u'$, then
\[
  \det(1-u t,K) = \det(1-u' t,K') \text{.}
\]
\end{corollary_}

Applying \ref{III:1-7-2}, we reduce to the case in which $K$ and $K'$ are concentrated in nonnegative degrees. The right-hand side of \ref{III:1-8} then plainly has the required invariance property.

\subsection{}\label{III:1-14}

This corollary makes it possible to define $\det(1-ut,K)$ for $u$ an endomorphism of $K\in\ob\eD_\text{parf}(A)$. This construction is ad hoc; in fact, for $v$ an automorphism of $L\in\ob\eD_\text{parf}(\Lambda)$, one can define $\det(v)\in\Lambda^\times$, and $\det(1-ut,K)$ is obtained by taking $v=1-ut$, $\Lambda=A\llbracket t\rrbracket$, and $L=K\otimes_A\Lambda$.

\section{The theorem}\label{III:2}

\subsection{}\label{III:2-1}

We retain the notation of (Report, \S\ref{II:1}). For $A$ a commutative Noetherian torsion ring, $X_0$ a separated scheme of finite type over $\dF_q$, and $\sK_0\in\ob\eD_\text{ctf}(X_0;A)$, put
\[
  L(X_0,\sK_0) = \prod_{x\in |X_0|} \det\left(1-F_x^* t^{\deg(x)}, \sK\right)^{-1}
\]
(\ref{III:1-7}, \ref{III:1-13}, and Report, \ref{II:1-5}, \ref{II:1-6}).

\begin{theorem_}\label{III:2-2}
In each of the following two cases:
\begin{enumerate}[\indent a)]
  \item $A$ is of torsion prime to $p$;
  \item $A$ is reduced and of characteristic $p$;
\end{enumerate}
one has
\begin{equation*}\tag{2.2.1}\label{III:eq:2-2-1}
  L(X_0,\sK_0) = \det\left(1-F^* t, \eR\Gamma_c(X,\sK)\right)^{-1} \text{.}
\end{equation*}
\end{theorem_}

\subsection{Remark}\label{III:2-3}

In case a), the method of Report, \S\ref{II:3}, applied to the trace formula (Report, \ref{II:4-10}), shows that the logarithmic derivatives of the two sides of \eqref{III:eq:2-2-1} are equal. Since the ring $A$ is a torsion ring, this does not suffice to prove \ref{III:2-2}.

\subsection{Remark}\label{III:2-4}

If $A$ is a field, then
\[
  L(X_0,\sK_0) = \prod_i L\left(X_0,\underline\h^i(\sK_0)\right)^{(-1)^i}
\]
and the spectral sequence $E_2^{p q}=\h_c^p(X,\underline\h^q(\sK))\Rightarrow \h_c^{p+q}(X,\sK)$ shows that
\begin{align*}
  \det\left(1-F^* t, \eR\Gamma_c(X,\sK)\right)
    &= \prod_n \det\left(1-F^* t,\h_c^n(X,\sK)\right)^{(-1)^n} \\
    &= \prod_{p,q} \det\left(1-F^* t,\h_c^p(X,\underline\h^q(\sK))\right)^{(-1)^{p+1}} \text{.}
\end{align*}
This reduces \ref{III:2-2} (for $A$ a field) to the following special case.

\subsubsection{}\label{III:2-4-1}

If $\sG_0$ is a constructible sheaf of $A$-vector spaces, then
\[
  L(X_0,\sG_0) = \prod_n \det\left(1-F^* t,\h_c^n(X,\sG)\right)^{(-1)^{n+1}} \text{.}
\]

\subsection{Remark}\label{III:2-5}

It is easy to reduce case b) to the case in which $A$ is a field (or even a finite field) of characteristic $p$; cf. \ref{III:4-2}.

\subsection{Remark}\label{III:2-6}

The special case of b) in which $A=\dZ/p$ and $\sK$ is the constant sheaf $\dZ/p$, concentrated in degree $0$,
\[
  Z(X_0,t) = \prod_{x\in |X_0|} \left(1-t^{\deg(x)}\right)^{-1}
      \equiv \prod_n \det\left(1-F^* t^f,\h_c^n(X,\dZ/p)\right)^{(-1)^{n+1}} \pmod p
\]
was proved by N. Katz \cite[XXII 3.1]{sga7}. His method, entirely different from the one followed here, starts from Dwork's $p$-adic theory of $Z(X_0,t)$.

\subsection{Remark}\label{III:2-7}

Contrary to what I had rashly told some friends, in b) one cannot omit the hypothesis ``$A$ reduced.'' For a counterexample, see \ref{III:4-5}.

\subsection{Outline of the proof of \ref{III:2-2}}\label{III:2-8}

A shift in degrees \ref{III:1-7-2} reduces us to the case in which $\sK_0$ is a bounded complex, in nonnegative degrees, of constructible flat sheaves of $A$-modules. Then $\eR\Gamma_c(X,\sK)\in\ob\eD_\text{parf}^{\geqslant0}(A)$ (Report, proof of \ref{II:4-9}, b). We shall use \ref{III:1-8} to expand both sides of \eqref{III:eq:2-2-1} as power series in $T=t$. Equality of the coefficients of $T^n$ will follow from \cite[XVII 5.5.21]{sga4} and a trace formula on $\operatorname{Sym}^n(X)$: in case a), the formula of Report, \ref{II:4-10}, and in case b), a formula to be proved in the next section. First, one must reduce to the case in which the symmetric powers of $X_0$ exist (for example, to the case in which $X$ is affine), using the multiplicativity in $X_0$ of both sides of \eqref{III:eq:2-2-1}: if $U_0$ is an open subscheme of $X_0$ with closed complement $Y_0$, then $L(X_0,\sK_0)=L(U_0,\sK_0)\cdot L(Y_0,\sK_0)$, and the analogous formula for the right-hand side follows from \ref{III:1-7-1} by the method of (Report, \S\ref{II:6}A).

\subsection{The case in which $X_0$ is finite}\label{III:2-9}

The multiplicativity in $X_0$ of both sides of \eqref{III:eq:2-2-1} reduces this case to $X_0=\spec(\dF_{q^k})$. Returning to the definitions, one sees that \eqref{III:eq:2-2-1} is then equivalent to Report, \ref{II:3-4}.

\subsection{The left-hand side}\label{III:2-10}

By \ref{III:1-8}, one has
\[
  L(X_0,\sK_0) = \prod_{x\in |X_0|} \sum_n
  \tr\left(F_x^*,\Gamma^n\sK_x\right)t^{n\deg(x)} \text{.}
\]
The coefficient of $T^n$ is therefore the sum, over all formal linear combinations $\sum_{x\in|X_0|}n_x\cdot x$ of elements of $|X_0|$ ($n_x\geqslant0$, with almost all $n_x$ zero) such that $\sum n_x[k(x):\dF_q]=n$, of the products
\begin{equation*}\tag{2.10.1}\label{III:eq:2-10-1}
  \prod_x \tr\left(F_x^*,\Gamma^{n_x} \sK\right) \text{.}
\end{equation*}

For $x\in|X_0|$, let $(x)$ denote the $0$-cycle consisting of the $[k(x):\dF_q]$ points of $X$ above $x$ (an orbit of $F$). The map $\sum n_x\cdot x\mapsto\sum n_x\cdot(x)$ identifies the formal linear combinations above, with $\sum n_x[k(x):\dF_q]=n$, with the effective $0$-cycles of degree $n$ on $X$ fixed by $F$, that is, with the fixed points of $F$ in $\operatorname{Sym}^n(X)$. Thus the coefficient of $T^n$ in $L(X_0,\sK_0)$ appears as a sum of the products \eqref{III:eq:2-10-1}, indexed by $\operatorname{Sym}^n(X)^F$. More precisely:

\begin{lemma_}\label{III:2-11}
 $L(X_0,\sK_0) = \sum_n \sum_{x\in\operatorname{Sym}^n(X)^F} \tr\left(F_x^*,\eL\Gamma_\textnormal{ext}^n \sK\right) T^n$.
\end{lemma_}

As the finite subscheme $Y_0$ of $X_0$ increases, one has
\[
  L(X_0,\sK_0) = \lim L(Y_0,\sK_0) \text{,}
\]
that is, for every $m$ there is a finite subscheme $Z_0\subset X_0$ such that, if $Y_0\supset Z_0$, then $L(X_0,\sK_0)$ and $L(Y_0,\sK_0|Y_0)$ are congruent modulo $T^m$. It suffices that $Z_0$ contain the closed points of degree $\leqslant m$. The same assertion holds for the right-hand side, so it is enough to prove \ref{III:2-11} for finite $X_0$. In that case,
\begin{align*}
  L(X_0,\sK_0)
    &=_\text{\ref{III:2-9}} \det\left(1-F^* T,\Gamma(X,\sK)\right)^{-1} \\
    &=_\text{\ref{III:1-8}} \sum_n \tr\left(F^*,\Gamma^n \Gamma(X,\sK)\right) T^n \\
    &=_\text{\ref{III:1-3}} \sum_n \tr\left(F^*,\Gamma(\operatorname{Sym}^n(X),\Gamma_\text{ext}^n(\sK)\right) T^n \\
    &= \sum_n \sum_{x\in \operatorname{Sym}^n(X)^F} \tr\left(F_x^*,\Gamma_\text{ext}^n(\sK)\right) T^n
\end{align*}
(by the trace formula in the trivial case of a finite set).

\subsection{The right-hand side}\label{III:2-12}

By \ref{III:1-8} and \cite[XVII 5.5.12]{sga4}, one has
\begin{align*}
  \det\left(1-F^* T,\eR\Gamma_c(X,\sK)\right)^{-1}
    &= \sum_n \tr\left(F^*,\eL\Gamma^n\eR\Gamma_c(X,\sK)\right) T^n \\
    &= \sum_n \tr\left(F^*,\eR\Gamma_c(\operatorname{Sym}^n(X),\Gamma_\text{ext}^n\sK)\right) T^n \text{.}
\end{align*}

\subsection{}\label{III:2-13}

The theorem is therefore equivalent to the trace formulas
\[
  \tr\left(F^*,\eR\Gamma_c(\operatorname{Sym}^n(X),\Gamma_\text{ext}^n \sK)\right)
  = \sum_{x\in\operatorname{Sym}^n(X)^F}
  \tr\left(F_x^*,\Gamma_\text{ext}^n\sK\right) \text{.}
\]
In case a), they follow from Report, \ref{II:4-10}. The analogous formula in case b) will be proved in \S\ref{III:4}.

\section{Artin--Schreier theory}\label{III:3}

\subsection{}\label{III:3-1}

In what follows, a coherent sheaf on a scheme $S$ will always be regarded as a sheaf on $S_\text{ét}$. Let $S$ be a scheme of characteristic $p>0$. If $G$ is a locally constant sheaf of finite-dimensional $\dZ/p$-vector spaces, then $\sG=G\otimes_{\dZ/p}\sO$ is a locally free coherent sheaf on $S$. We use $\Phi$ to denote both the endomorphism $f\mapsto f^p$ of $\sO$ and its tensor product with the identity of $G$, $\Phi:\sG\to\sG$. Artin--Schreier theory \cite[IX 3.5]{sga4} asserts that the sequence
\[\xymatrix{
  0 \ar[r]
    & \dZ/p \ar[r]
    & \sO \ar[r]^-{\Phi-1}
    & \sO \ar[r]
    & 0
}\]
is exact. Tensoring with $G$ gives an exact sequence
\[\xymatrix{
  0 \ar[r]
    & G \ar[r]
    & \sG \ar[r]^-{\Phi-1}
    & \sG \ar[r]
    & 0
}\]

\subsection{}\label{III:3-2}

Suppose that $S$ is an open subscheme of a Noetherian scheme $\bar S$. Let $j:S\hookrightarrow\bar S$ be the inclusion, and let $I$ be a sheaf of ideals defining the closed complement $\bar S\setminus S$. We again denote by $\Phi$ the functorially induced extension of $\Phi$ to $j_*\sG$.

\begin{lemma_}\label{III:3-3}
There exist coherent extensions $\bar\sG\subset j_*\sG$ of $\sG$ such that
\begin{equation*}\tag{3.3.1}\label{III:eq:3-3-1}
  \Phi(\bar\sG)\subset I\cdot \bar\sG\text{.}
\end{equation*}
If $\bar\sG$ is such an extension, the sequence
\begin{equation*}\tag{3.3.2}\label{III:eq:3-3-2}
\xymatrix{
  0 \ar[r]
    & j_! G \ar[r]
    & \bar\sG \ar[r]^-{\Phi-1}
    & \bar\sG \ar[r]
    & 0
}
\end{equation*}
is exact.
\end{lemma_}
\begin{enumerate}[a)]
  \item \emph{Existence}: Let $\sG'\subset j_*\sG$ be a coherent extension of $\sG$ to $\bar S$. For every $n$, one has $\Phi(I^n\sG')\subset I^{pn}\Phi(\sG')$. Choose $n$ large enough that $I^{n(p-1)-1}\Phi(\sG')\subset\sG'$, and put $\bar\sG=I^n\sG'$. Then $\Phi(\bar\sG)\subset I^{pn}\Phi(\sG') =I\cdot I^n\cdot I^{n(p-1)-1}\Phi(\sG')\subset I\cdot I^n\sG' =I\bar\sG$.
  \item \emph{Kernel of $\Phi-1$}: On $S$, the assertion follows from \ref{III:3-1}; it remains to check that a section $x$ of $\ker(\Phi-1)$ over $U$ étale over $\bar S$ vanishes in a neighborhood of the inverse image of $\bar S\setminus S$. Since $x=\Phi(x)$, \eqref{III:eq:3-3-1} indeed gives $x\in I\bar\sG$; moreover, $x\in I^n\bar\sG\Rightarrow x=\Phi(x)\in\Phi(I^n\bar\sG)\subset I^{np}\bar\sG$. Thus $x$ lies in every $I^n\bar\sG$, and hence vanishes near $\bar S\setminus S$.
  \item \emph{Surjectivity of $\Phi-1$}: We may assume that $\bar S$ is affine and write $\sG$ as a quotient of $\sO^n$; $\Phi$ then lifts to a $p$-linear map $\widetilde\Phi:\sO^n\to\sO^n$, and it suffices to check the surjectivity of $\widetilde\Phi-1$:
    \[\xymatrix{
      \sO^n \ar[r]^-{\widetilde \Phi-1} \ar[d]
        & \sO^n \ar[d] \\
      \sG \ar[r]^-{\Phi-1}
        & \sG
    }\]
    The sheaf $\sO^n$ is the sheaf of local sections of the group scheme $\dG_a^n$, and $\widetilde\Phi-1$ is induced by a homomorphism $\dG_a^n\to\dG_a^n$. This homomorphism is étale. To check that it is surjective, it suffices to do so after every base change $\spec(k)\to\bar S$ with $k$ algebraically closed; over an algebraically closed field, every étale homomorphism between connected groups of the same dimension is surjective. Finally, a surjective étale morphism of $\bar S$-schemes induces an epimorphism between the corresponding sheaves on $\et{\bar S}$.
\end{enumerate}

\subsection{Frobenius and Frobenius}\label{III:3-4}

The absolute Frobenius morphism $F_\text{abs}:S\to S$ is the identity on the underlying topological space of $S$, and is $f\mapsto f^p$ on the structure sheaf. For $U/S$ étale, there is a canonical isomorphism $F_\text{abs}^*U=U$: the diagram
\[\xymatrix{
  U \ar[r]^-{F_\text{abs}} \ar[d]
    & S \ar[d] \\
  S \ar[r]^-{F_\text{abs}}
    & S
}\]
is Cartesian. The endomorphism of ringed sites $(\et{\bar S},\sO)\to(\et{\bar S},\sO)$ defined by $F_\text{abs}$ can therefore be described as follows:
\begin{center}
  $F_\text{abs}^*$ is the identity on $\et S$; the homomorphism $F_\text{abs}^*\sO=\sO\to\sO$ is $\Phi$.
\end{center}
With this description, for every étale sheaf $\sG$ on $S$, the Frobenius correspondence $F_\text{abs}^*\sG\to\sG$ is the identity.

Giving a $p$-linear morphism of quasi-coherent sheaves $u:\sM\to\sN$ on $S$ is equivalent to giving a linear morphism $u':F_\text{abs}^*\sM\to\sN$. In particular, with the notation of \ref{III:3-5}, the homomorphism $\Phi:\sG\to\sG$ corresponds to $\Phi':F_\text{abs}^*\sG\to\sG$.

\begin{lemma_}\label{III:3-5}
With the notation of \ref{III:3-4}, the diagram
\[\xymatrix{
  F_\textnormal{abs}^* G \ar[r] \ar[d]
    & G \ar[d] \\
  F_\textnormal{abs}^*\sG \ar[r]^-{\Phi'}
    & \sG
}\]
is commutative.
\end{lemma_}

It suffices to check this locally. We may therefore assume that $G$ is constant, in which case the lemma is clear.

The same theory applies to powers of $F_\text{abs}$.

\subsection{}\label{III:3-6}

Now suppose that $\bar S$ is a scheme over $\dF_q$ ($q=p^f$), and, in accordance with the conventions of Report, \S\ref{II:1}, write $S_0,\bar S_0,G_0,\sG_0,\bar\sG_0$ in place of $S,\dotsc$. Lemma \ref{III:3-5}, applied to $F_\text{abs}^f:S_0\to S_0$, gives a commutative diagram
\[\xymatrix{
  F^* G_0 \ar[r] \ar[d]
    & G_0 \ar[d] \\
  F^*\sG_0 \ar[r]^-{(\Phi^f)'}
    & \sG_0
}\qquad\text{and then}\qquad
\xymatrix{
  F^*(j_! G_0) \ar[r] \ar[d]
    & j_!G_0 \ar[d] \\
  F^*\bar\sG_0 \ar[r]^-{(\Phi^f)'}
    & \sG_0
}\text{,}
\]
Then, after extending scalars from $\dF_q$ to $\dF$ and passing to cohomology, one obtains
\[\xymatrix{
  F^*(j_! G) \ar[r]^-{F^*} \ar[d]
    & j_!G \ar[d] \\
  F^*\bar\sG \ar[r]
    & \bar\sG
}\qquad\text{and}\qquad
\xymatrix{
  \h^\bullet(\bar S,j_! G) \ar[r]^-{F^*} \ar[d]
    & \h^\bullet(\bar S,j_!G) \ar[d] \\
  \h^\bullet(\bar S,\bar\sG) \ar[r]
    & \h^\bullet(\bar S,\bar\sG)
}\text{,}
\]
where the second row is obtained by extending scalars from $\dF_q$ to $\dF$ in the $\dF_q$-linear endomorphism $\Phi^f$ of $\h^\bullet(\bar S_0,\bar\sG_0)$.

Artin--Schreier theory gives a long exact sequence
\[\xymatrix{
  \cdots \ar[r]
    & \h^i(\bar S,j_! G) \ar[r]
    & \h^i(\bar S,\bar\sG) \ar[r]^-{\Phi-1}
    & \h^i(\bar S,\bar\sG) \ar[r]
    & \cdots
}\]
where $\Phi$ is the $p$-linear extension to $\h^\bullet(\bar S,\bar\sG) =\h^\bullet(\bar S_0,\bar\sG_0)\otimes_{\dF_q}\dF$ of the $p$-linear endomorphism $\Phi$ of $\h^\bullet(\bar S_0,\bar\sG_0)$.

\subsection{}\label{III:3-7}

If $\bar S_0$ is proper over $\dF_q$, the $\h^i(\bar S,\bar\sG)$ are finite-dimensional and $\Phi-1$ is surjective (cf. \ref{III:3-3}c). Hence
\[
  \h_c^i(S,G) = \ker\left(\Phi-1 :
  \h^i(\bar S,\bar\sG)\to\h^i(\bar S,\bar\sG)\right) \text{,}
\]
and $F^*$ is induced by the $\dF$-linear extension to $\h^i(\bar S,\bar\sG) =\h^i(\bar S_0,\bar\sG_0)\otimes_{\dF_q}\dF$ of the $\dF_q$-linear endomorphism $\Phi^f$ of $\h^i(\bar S_0,\bar\sG_0)$.

Applying Lemma \ref{III:3-8} below, we obtain the identity
\begin{equation*}\tag{3.7.1}\label{III:eq:3-7-1}
  \tr\left(F^*,\h_c^i(S,G)\right) = \tr\left(\Phi^f,\h^i(\bar S_0,\bar\sG_0)\right) \text{.}
\end{equation*}

\begin{lemma_}\label{III:3-8}
Let $\Phi$ be a $p$-linear endomorphism of a vector space $V$ over $\dF_q$: $\Phi(\lambda x)=\lambda^p\Phi(x)$. Put $F=\Phi^f$ (so $F$ is linear), and again denote by $\Phi$ and $F$ their respective $p$-linear and linear extensions to $V\otimes_{\dF_q}\dF$. Then $F$ stabilizes $\ker(\Phi-1)\subset V\otimes_{\dF_q}\dF$, and
\[
  \tr\left(F,\ker(\Phi-1)\right) = \tr(F,V) = \tr\left(F,V\otimes_{\dF_q}\dF\right) \text{.}
\]
\end{lemma_}

On $V\otimes_{\dF_q}\dF$, one has $F\Phi=\Phi F$: the $p$-linear endomorphism $F\Phi-\Phi F$ vanishes on $V$, hence everywhere. Write $V=V'\oplus V''$, with $F$ invertible on $V'$ and nilpotent on $V''$. The theory of $p$-linear endomorphisms ensures that
\[
  \ker(\Phi-1)\otimes_{\dF_q}\dF \iso V'\otimes_{\dF_q}\dF \text{.}
\]
Since $\tr(F,V'')=0$, it follows that
\begin{align*}
  \tr\left(F,\ker(\Phi-1)\right)
    &= \tr\left(F,\ker(\Phi-1)\otimes_{\dF_q}\dF\right) = \tr\left(F,V'\otimes_{\dF_q}\dF\right) \\
    &= \tr\left(F,V\otimes_{\dF_q}\dF\right) = \tr(F,V) \text{.}
\end{align*}

\section{The trace formula modulo \texorpdfstring{$p$}{p}}\label{III:4}

\begin{theorem_}\label{III:4-1}
Let $X_0$ be a separated scheme of finite type over $\dF_q$, let $A$ be a reduced commutative Noetherian ring of characteristic $p$, and let $\sK_0\in\ob\eD_\textnormal{ctf}(X_0,A)$. Then
\begin{equation*}\tag{4.1.1}\label{III:eq:4-1-1}
  \sum_{x\in X^F} \tr\left(F_x^*,\sK_x\right) = \tr\left(F^*,\eR\Gamma_c(X,\sK)\right) \text{.}
\end{equation*}
\end{theorem_}

\subsection{Reduction to the case in which \texorpdfstring{$A$}{A} is a finite field}\label{III:4-2}

% NOTE: in this paragraph and others, (3.1.1) in the printed source appears
% to refer to (4.1.1).
A standard limit argument reduces us to the case in which $A$ is of finite type over $\dZ/p$. For each maximal ideal $\fm$ of $A$, the image of either side of \eqref{III:eq:4-1-1} in $A/\fm$ is given by the same formula, with $\sK_0$ replaced by $\sK_0\lotimes_A A/\fm\in\ob\eD_\text{ctf}(X_0,A/\fm)$ (for the right-hand side, cf. Report, \ref{II:4-12}). Since the $A/\fm$ are finite, and the intersection of the maximal ideals of $A$ is zero, this gives the stated reduction.

\subsection{Reduction to the case \texorpdfstring{$A=\dZ/p$}{A=Z/p}}\label{III:4-3}

Suppose that $A$ is a finite field, and denote the two sides of \eqref{III:eq:4-1-1} by $T_A'(\sK_0)$ and $T_A''(\sK_0)$. If, by restriction of scalars, we regard $\sK_0$ as a complex of sheaves of $\dZ/p$-modules, then
\[
  T_{\dZ/p}'(\sK_0) = \tr_{A/(\dZ/p)} T_A'(\sK_0)
\]
and similarly for $T''$. If \eqref{III:eq:4-1-1} holds for $A=\dZ/p$, then
\begin{equation*}\tag{4.3.1}\label{III:eq:4-3-1}
  \tr_{A/(\dZ/p)} T_A'(\sK_0) = \tr_{A/(\dZ/p)} T_A''(\sK_0) \text{.}
\end{equation*}
I claim that, for every $\lambda\in A$, one even has
\begin{equation*}\tag{4.3.2}\label{III:eq:4-3-2}
  \tr_{A/(\dZ/p)}\left(\lambda\cdot T_A'(\sK_0)\right) = \tr_{A/(\dZ/p)}\left(\lambda\cdot T_A''(\sK_0)\right) \text{,}
\end{equation*}
and hence $T_A'(\sK_0)=T_A''(\sK_0)$. This is clear for $\lambda=0$. For $\lambda$ invertible, let $A(\lambda)$ be the inverse image on $X$ of the free rank-one sheaf of $A$-modules on $\spec(\dF_q)$ for which $F^*=\lambda$. Formula \eqref{III:eq:4-3-2} is simply \eqref{III:eq:4-3-1} applied to $A(\lambda)\otimes_A\sK_0$.

% originally theorem 4.1bis (i.e. theorem 4.1 again)
\begin{theorem_}\label{III:4-1bis}
Let $\bar S_0$ be a smooth projective curve, let $j:S_0\hookrightarrow\bar S_0$ be a dense open subscheme, and let $G_0$ be a locally constant sheaf of finite-dimensional $\dZ/p$-vector spaces on $S_0$. Then
\begin{equation*}\tag{4.4.1}\label{III:eq:4-4-1}
  \sum_{x\in S^F} \tr\left(F_x^*,G_x\right) = \sum_i (-1)^i \tr\left(F^*,\h_c^i(S,G)\right) \text{.}
\end{equation*}
\end{theorem_}
\begin{proof}
Apply \ref{III:3-3} to $(S_0,\bar S_0,G_0)$. The coherent extension $\bar\sG_0$ of $\sG_0$ whose existence is guaranteed by \ref{III:3-3} is locally free: it is torsion-free (because $\bar\sG_0\subset j_*\sG_0$), and $\bar S_0$ is regular of dimension one. The $q$-linear endomorphism $\Phi^f$ of $\bar\sG_0$ may be interpreted as a morphism
\[
  u : F^*\bar\sG_0 \to \bar\sG_0 \text{.}
\]
It factors through $I\bar\sG_0$, where $I$ defines the closed subset $\bar S_0\setminus S_0$. Let $v:j_!G_0\to\bar\sG_0$ be the unique extension of the natural map $G_0\to\sG_0=G_0\otimes_{\dZ/p}\sO$. The diagram
\begin{equation*}\tag{4.4.2}\label{III:eq:4-4-2}
\xymatrix{
  F^* j_! G_0 \ar[r]^-{F^*} \ar[d]^-{F^* v}
    & j_! G_0 \ar[d]^-v \\
  F^*\bar \sG_0 \ar[r]^-u
    & \bar\sG_0
}
\end{equation*}
is commutative. Finally, denote by $uF^*$ the composite map
\[\xymatrix{
  \h^\bullet(\bar S,\bar\sG) \ar[r]^-{F^*}
    & \h^\bullet(\bar S,F^*\bar\sG) \ar[r]^-u
    & \h^\bullet(\bar S,\bar\sG) \text{,}
}\]
Then, by \eqref{III:eq:3-7-1},
\begin{equation*}\tag{4.4.3}\label{III:eq:4-4-3}
  \tr\left(F^*,\h_c^i(S,G)\right) = \tr\left(u F^*,\h^i(\bar S,\bar\sG)\right) \text{.}
\end{equation*}

% NOTE: trace the reference to the Woods Hole seminar.
For each closed point $i_x:x\hookrightarrow\bar S$ of $\bar S$, put $\bar\sG_x=i_x^*\bar\sG$. The Lefschetz trace formula in coherent cohomology (the ``Woods Hole fixed-point formula''), applied to $F$ and $u$, says that
\begin{equation*}\tag{4.4.4}\label{III:eq:4-4-4}
  \sum_i (-1)^i \tr\left(u F^*, \h^i(\bar S,\bar\sG)\right)
  = \sum_{x\in \bar S^F}
  \frac{\tr(u_x,\bar\sG_x)}{\det(1-d F_x)} \text{.}
\end{equation*}
On the right-hand side:
\begin{enumerate}[\indent a)]
  \item since $dF=0$, the denominators equal $1$;
  \item for $x\in\bar S^F\setminus S^F$, the endomorphism $u_x$ of $\bar\sG_x$ is zero, and so is $\tr(u_x,\bar\sG_x)$;
  \item for $x\in S^F$, \eqref{III:eq:4-4-4} gives the contribution
    \[
      \tr(u_x,\bar\sG_x) = \tr(F_x^*,G_x) \text{.}
    \]
    Thus \eqref{III:eq:4-4-1} follows from \eqref{III:eq:4-4-3} and \eqref{III:eq:4-4-4}.
\end{enumerate}
\end{proof}

\subsection{A counterexample}\label{III:4-5}

The following example shows that the hypothesis that $A$ be reduced cannot be omitted from \ref{III:4-1}.

Let $Y_0$ be the elliptic curve over $\dF_2$ for which the eigenvalues of Frobenius are $\pm\sqrt{-2}$. It has $3$ rational points and is supersingular. The involution $\sigma:x\mapsto-x$ therefore has only one fixed point, namely $0$. The quotient of $Y_0$ by $\sigma$ is a projective line; the quotient of $Y_0\setminus0$ by $\sigma$ is consequently an affine line $\dA_0^1$, and $Y_0\setminus0$ is an unramified double covering of $\dA_0^1$:
\[
  \pi:Y_0\setminus 0 \to \dA_0^1 \qquad \text{, $\pi=\pi\sigma$.}
\]

The sheaf $\pi_*\dZ/2$ is a free module of rank one over the group algebra $A$ of the two-element group $\{1,\sigma\}$ over $\dZ/2$. This algebra is isomorphic to the algebra of dual numbers over $\dZ/2$.

The sheaf of $A$-modules $\pi_*\dZ/2$ is our example:
\begin{enumerate}[\indent a)]
  \item One has $\h_c^\bullet(\dA^1,\pi_*\dZ/2) = \h_c^\bullet(Y\setminus 0,\dZ/2)$. Since $Y_0$ is supersingular, $\h^i(Y,\dZ/2)=0$ for $i>0$, and the long exact sequence
    \[
      \cdots \to \h_c^i(Y\setminus 0,\dZ/2) \to \h^i(Y,\dZ/2) \to \h^i(\{0\},\dZ/2) \to \cdots
    \]
    shows that $\h_c^\bullet(Y\setminus0,\dZ/2)=0$. Hence $\h_c^\bullet(\dA^1,\pi_*\dZ/2)=0$, and
    \[
      \tr_A\left(F^*,\eR\Gamma_c(\dA^1,\pi_*\dZ/2)\right) = 0 \text{.}
    \]
  \item Since $Y_0\setminus0$ has two rational points, the Frobenius substitution is the identity at one of the rational points of $\dA_0^1$, and is $\sigma$ at the other. Thus
    \[
      \sum_{x\in {\dA^1}^F} \tr\left(F^*,\pi_*\dZ/2\right) = 1+\sigma \text{,}
    \]
    and $1+\sigma\ne0$.
\end{enumerate}

The equation of the covering $\pi$ is $y^2-y=x^3$.
% END INLINED FILE: sga4.5-3.tex
\clearpage
\clearpage
% BEGIN INLINED FILE: sga4.5-4.tex
\chapter{The cohomology class associated with a cycle}\label{IV}
\blfootnote{by A.\ Grothendieck, written up by P.\ Deligne}

This exposé is inspired by notes of Grothendieck that formed a preliminary version of \cite[IV]{sga5}. We define the cohomology class of a cycle in a separated smooth scheme of finite type over a field, and prove that intersection corresponds to the cup product.

Chapter \ref{IV:1} contains some general formalities. In Chapter \ref{IV:2}, the class of a cycle is defined in several situations more general than the one just mentioned. The principal compatibility not considered is that between the direct image of a cycle and the trace morphism in cohomology. In Chapter \ref{IV:3}, this formalism is used to deduce the Lefschetz trace formula for an endomorphism with isolated fixed points of a proper smooth scheme over an algebraically closed field, and for a Frobenius endomorphism of a curve.

We impose the following conventions:
\begin{enumerate}[\indent 1)]
  \item ``scheme'' means separated Noetherian scheme (this is largely a hypothesis of convenience).
  \item In Chapter \ref{IV:2}, an integer $n$ is fixed, and $n$ is invertible on every scheme considered.
  \item In Chapter \ref{IV:3}, a prime number $\ell$ is fixed, and $\ell$ is invertible on every scheme considered. The cohomology used is always $\ell$-adic cohomology:
    \[
      \h^\bullet(X) = \dQ_\ell\otimes_{\dZ_\ell} \varprojlim \h^\bullet(X,\dZ/\ell^n) \text{.}
    \]
\end{enumerate}
Cohomology classes of cycles, trace morphisms, and so forth are defined by passage to the limit from the case of finite coefficients $\dZ/\ell^n$ (cf. \cite[VI]{sga5}).

\section{Cohomology with support and cup products}\label{IV:1}

This section contains reminders from general topology, which the reader is invited to consult only as the need arises.

\subsection{\texorpdfstring{$H^1$}{H1} and torsors}\label{IV:1-1}

\subsubsection{}\label{IV:1-1-1}

Let $\sF$ be an abelian sheaf on a site $X$. It is known that $\h^1(X,\sF)$ classifies the $\sF$-torsors on $X$. We normalize (that is, choose the sign of) the isomorphism
\[
  (\text{set of isomorphism classes of $\sF$-torsors})
  \longrightarrow\h^1(X,\sF)
\]
so that, for every exact sequence $0\to\sF\xrightarrow\alpha\sG\xrightarrow\beta\sH\to0$ and every $h\in\h^0(X,\sH)$, the $\sF$-torsor $\beta^{-1}(h)\subset\sG$, on which $\sF$ acts by $(f,x)\mapsto\alpha(f)+x$, has class $\partial h$.

\subsubsection{}\label{IV:1-1-2}

Let $\sP$ be an $\sF$-torsor. If $(U_i)$ is an open covering of $X$, and $p_i$ is a section of $\sP$ over $U_i$, the \v{C}ech cocycle associated with $\sP$ is
\[
  p_{i j} = p_j - p_i \qquad \text{($p_{i j}\in \h^0(U_i\times U_j,\sF)$).}
\]
If, according to the usual convention, the map $\operatorname{\check H}^\bullet(X,\sF)\to\h^\bullet(X,\sF)$ is defined so as to be a morphism of $\delta$-functors, the image of $(p_{ij})\in\operatorname{\check H}^1(X,\sF)$ in $\h^1(X,\sF)$ is the class of $\sP$, normalized as in \ref{IV:1-1-1}.

\subsubsection{}\label{IV:1-1-3}

The definition of $\h^\bullet(X,\sF)$ is as follows: for a resolution $\sF^\bullet$ of $\sF$ with acyclic terms, $\h^i(X,\sF)=\h^i\Gamma(X,\sF^\bullet)$. The $\delta$-functor structure is obtained by associating with a short exact sequence of sheaves a short exact sequence of resolutions that remains exact after applying $\Gamma$. If $\sF^\bullet$ is a resolution of $\sF$, the connecting homomorphism $\partial$ associated with
\[\xymatrix{
  0 \ar[r]
    & \sF \ar[r]
    & \sF^0 \ar[r]^-d
    & \ker(d) \ar[r]
    & 0
}\]
induces the negative of the defining isomorphism $\h^1(X,\sF) = \Gamma\left(X,\ker(d)\right)/d \Gamma(X,\sF^0)$.

\subsubsection{}\label{IV:1-1-4}

Let $U$ be an open part of $X$ (a subsheaf of the final sheaf), and let $D$ be the ``closed complement.'' It is known that $\h_D^1(X,\sF)$ classifies the $\sF$-torsors on $X$ trivialized on $U$. For every short exact sequence $0\to\sF\xrightarrow\alpha\sG\xrightarrow\beta\sH\to0$, and every section $h\in\h_D^0(X,\sH)$ supported in $D$, the torsor $\beta^{-1}(h)$, trivialized on $U$ by the section $0$, has class $\partial h$.

The long exact sequence of cohomology with support
\[\xymatrix{
  \cdots \ar[r]^-\partial
    & \h_D^i(X,\sF) \ar[r]
    & \h^i(X,\sF) \ar[r]
    & \h^i(U,\sF) \ar[r]^-\partial
    & \cdots
}\]
is defined from the sequence of functors
\[\xymatrix{
  0 \ar[r]
    & \Gamma_D \ar[r]
    & \Gamma \ar[r]
    & \Gamma(U,-) \ar[r]
    & 0
}\]
(which is exact on injective sheaves). For every section $f\in\h^0(U,\sF)$, $\partial f\in\h_D^1(X,\sF)$ is the class of the trivial torsor $\sF$, trivialized on $U$ by the section $f$.

\subsubsection{}\label{IV:1-1-5}

Let $j:U\hookrightarrow X$. If $\sF$ injects into $j_*j^*\sF$, the exact sequence
\[\xymatrix{
  0 \ar[r]
  & \sF \ar[r]
  & j_* j^* \sF \ar[r]
  & j_* j^* \sF / \sF \ar[r]
  & 0
}\]
gives $\partial:\h^0(j_* j^*\sF/\sF) = \h_D^0(j_* j^* \sF/\sF) \to \h_D^1(X,\sF)$. The composite map
\[\xymatrix{
  \h^0(U,\sF) = \h^0(X,j_* j^*\sF) \ar[r]
    & \h^0(X,j_* j^*\sF / \sF) \ar[r]^-\partial
    & \h_D^1(X,\sF)
}\]
is the negative of the map considered in \ref{IV:1-1-4}.

\subsubsection{}\label{IV:1-1-6}

Recall that if $\sL$ is an invertible sheaf on a scheme $X$, the corresponding $\dG_m$-torsor is the sheaf $\operatorname{Isom}(\sO,\sL)$, on which $\dG_m$ acts by $(\lambda,f)\mapsto f\circ(\lambda\cdot)=\lambda f$. Recall also that if $D$ is a Cartier divisor on $X$, and $j:U\hookrightarrow X$ is the inclusion of the complementary open subscheme, the invertible sheaf $\sO(D)$ is the subsheaf of $j_*\sO_U$ consisting of local sections $s$ such that $sf$ lies in $\sO_X$, where $f$ is a local equation for $D$.

\subsection{Cup products}\label{IV:1-2}

In this subsection we develop some observations on cup products in cohomology with support that will be used again in [\nameref{V}] to relate Poincaré duality for curves to the autoduality of the Jacobian.

\subsubsection{}\label{IV:1-2-1}

Let $X$ be a site, $Y$ a closed part of $X$, and $\sF,\sG$ two abelian sheaves on $X$. For example, $X$ may be a scheme, $Y$ a closed subscheme, and $\sF,\sG$ sheaves on $\et X$. Define a product
\begin{equation*}\tag{1.2.1.1}\label{IV:eq:1-2-1-1}
  \Gamma_Y(X,\sF)\otimes \Gamma(Y,\sG) \to \Gamma_Y(X,\sF\otimes\sG) \text{.}
\end{equation*}
The product of a section $s$ of $\sF$ supported in $Y$ with a section $t$ of $\sG$ on $Y$ is obtained as follows: locally on $X$, $t$ is the restriction to $Y$ of a section $t'$ of $\sG$ on $X$, and one forms $s\otimes t'$. It is supported in $Y$ and does not depend on the choice of $t'$, which makes the definition legitimate and allows it to be globalized.

Let $i:Y\hookrightarrow X$ be the inclusion. The local analogue of \eqref{IV:eq:1-2-1-1} is the product
\begin{equation*}\tag{1.2.1.2}\label{IV:eq:1-2-1-2}
  i^!\sF\otimes i^*\sG \to i^!(\sF\otimes\sG) \text{,}
\end{equation*}
from which \eqref{IV:eq:1-2-1-1} follows by applying $\Gamma(Y,-)$.

\subsubsection{}\label{IV:1-2-2}

Let us derive these maps. Let $\sK$, $\sL$, and $\sM$ be in the derived category, and let $\sK\lotimes\sL\to\sM$ be a bilinear map. Deriving \eqref{IV:eq:1-2-1-1} gives
\begin{equation*}\tag{1.2.2.1}\label{IV:eq:1-2-2-1}
  \eR\Gamma_Y(X,\sK) \lotimes \eR\Gamma(Y,\sL) \to \eR\Gamma_Y(X,\sM)
\end{equation*}
which induces
\begin{equation*}\tag{1.2.2.2}\label{IV:eq:1-2-2-2}
  \h_Y^i(X,\sK)\otimes \h^j(Y,\sL) \to \h_Y^{i+j}(X,\sM)
\end{equation*}
(the cup product). The local map \eqref{IV:eq:1-2-1-2} gives
\begin{equation*}\tag{1.2.2.3}\label{IV:eq:1-2-2-3}
  \eR i^! \sK\lotimes i^*\sL \to \eR i^!\sM \text{,}
\end{equation*}
from which \eqref{IV:eq:1-2-2-1} follows by applying $\eR\Gamma(Y,-)$.

\subsubsection{}\label{IV:1-2-3}

Above, I have passed over the conflicts caused by using right-derived functors ($\eR\Gamma$) and left-derived functors ($\lotimes$) in the same formula.
\begin{enumerate}[\indent a)]
  \item In the next subsection, all the right-derived functors considered will have finite cohomological dimension. This makes it possible to work systematically in the derived categories $\eD^-$. To obtain complexes that are both flat and flasque, one uses the canonical flasque resolutions as in \cite[XVII]{sga4}.
  \item For a more general theory, one must cease to interpret bilinear maps from $\sK$ and $\sL$ to $\sM$ as morphisms $\sK\lotimes\sL\to\sM$. For example, $\eR\Gamma_Y(X,\sK)\lotimes\eR\Gamma(Y,\sL)$ is not defined if both factors lie in $\eD^+$. One solution is to work in $\eD^+$ and define
    \[
      \operatorname{Bil}(\sK,\sL;\sM) = \varinjlim \hom(\sK'\otimes\sL',\sM') \text{,}
    \]
    where the limit is taken over the quasi-isomorphisms $\sK'\iso\sK$, $\sL'\iso\sL$, $\sM\iso\sM'$ (with $\sK'$, $\sL'$, $\sM'$ bounded below), and where $\hom$ means ``morphism of complexes, up to homotopy''; one then systematically uses such bilinear maps without ever mentioning $\otimes$.
\end{enumerate}

\subsubsection{Second theme}\label{IV:1-2-4}

Let $U$ be an open part of $X$, and let $j:U\hookrightarrow X$ be the inclusion. For $\sK$ in the derived category on $U$, put $\eR\Gamma_!(U,\sK)=\eR\Gamma(X,j_!\sK)$. Given $\sK$, $\sL$, and $\sM$ on $U$, and a bilinear map $\sK\lotimes\sL\to\sM$, we wish to define
\begin{equation*}\tag{1.2.4.1}\label{IV:eq:1-2-4-1}
  \eR\Gamma(U,\sK)\lotimes\eR\Gamma_!(U,\sL)\to\eR\Gamma_!(U,\sM) \text{.}
\end{equation*}
At the level of sheaves and their global sections, such a product follows from the isomorphism $j_*\sF\otimes j_!\sG\xleftarrow\sim j_!(\sF\otimes\sG)$, but one must remember that $\eR\Gamma_!$ is not, in general, the derived functor of $\Gamma(X,j_!-)$.

One begins by defining
\[\xymatrix{
  \eR j_* \sK\lotimes j_! \sL
    & \ar[l]_-\sim j_!(\sK\lotimes \sL) \ar[r]
    & j_!\sM \text{.}
}\]
Applying $\eR\Gamma(X,-)$ gives
\[\xymatrix{
  \eR\Gamma(U,\sK)\lotimes\eR\Gamma_!(U,\sL) = \eR\Gamma(X,\eR j_*\sK)\lotimes \eR\Gamma(X,j_!\sL) \ar[r]
    & \eR\Gamma(X,j_! \sM) \text{.}
}\]

Here again, the conflict between left and right will be resolved in the next subsection by working in $\eD^-$.

\subsubsection{Coda}\label{IV:1-2-5}

Let $j:U\hookrightarrow X$ be an open part, and let $i:Y\hookrightarrow U$ be a closed part of $U$. Let $\bar Y$ be a closed part of $X$ such that $\bar Y\cap U=Y$ (for example, the complement of $U\setminus Y$), and let $\sK$ be on $U$. Put $\eR\Gamma_{Y!}(U,\sK)=\eR\Gamma_!(Y,\eR i^!\sK)$, where $\eR\Gamma_!$ is relative to the inclusion of $Y$ in $\bar Y$. For every open $V$ of $U$ containing $Y$, there is a map $\eR\Gamma_{Y!}(U,\sK)\to\eR\Gamma_!(V,\sK)$: if $i$ also denotes the inclusion of $\bar Y$ in $X$, $j$ that of $Y$ in $\bar Y$, and $k$ that of $V$ in $X$, then $\eR i^!\sK=j^*\eR i^!k_!(k^*\sK)$, whence a morphism $j_!\eR i^!\sK\to\eR i^!k_!(k^*\sK)$. Applying $\eR\Gamma(\bar Y,-)$ gives
\[
  \eR\Gamma_{Y!}(U,\sK)\to\eR\Gamma_{\bar Y}(k_!k^*\sK)
  \to\eR\Gamma(k_!k^*\sK)=\eR\Gamma_!(V,\sK) \text{.}
\]

For $\sK,\sL,\sM$ on $U$, and a bilinear map $\sK\lotimes\sL\to\sM$, we wish to define
\begin{equation*}\tag{1.2.5.1}\label{IV:eq:1-2-5-1}
  \h^n(U,\sK)\otimes \h_!^m(Y,\sL) \to \h_{Y!}^{n+m}(U,\sM)
\end{equation*}
(the latter group itself mapping to $\h_!^{n+m}(V,\sM)$). In the derived category, this amounts to defining
\begin{equation*}\tag{1.2.5.2}\label{IV:eq:1-2-5-2}
  \eR\Gamma_Y(U,\sK) \lotimes \eR\Gamma_!(Y,\sL)
  \to \eR\Gamma_{Y!}(U,\sM) \text{.}
\end{equation*}
We identify $\eR\Gamma_Y$ with $\eR\Gamma(Y,\eR i^!\sK)$. The desired product is then of type \eqref{IV:eq:1-2-4-1}, relative to the local product \eqref{IV:eq:1-2-2-3} on $Y$: $\eR i^!\sK\lotimes\sL\to\eR i^!\sM$.

\subsubsection{}\label{IV:1-2-6}

Above we have unfolded the absolute formalism. There is a parallel relative formalism, with $\Gamma$ replaced by $f_*$ for a morphism $f:X\to S$.

\subsection{The Koszul rule}\label{IV:1-3}

Let $A$ be a commutative ring, and let $(V_i)_{i\in I}$ be a finite family of graded (or $\dZ/2$-graded) $A$-modules. Recall the definition of the graded tensor product $\bigotimes_{i\in I}V_i$ according to the Koszul rule (cf. \cite[XVII 1.1]{sga4}). For each total ordering $a$ on $I$, define a module $V(a)$. We also define a transitive system of isomorphisms $\varphi_{ab}:V(b)\iso V(a)$, and the graded tensor product of the $V_i$ will be the ``common value'' $\varprojlim V(a)$ of this system of modules. Set
\begin{enumerate}[\indent a)]
  \item $V(a)=\bigotimes_{i\in I}|V_i|$ (the ordinary tensor product of the ungraded modules underlying the $V_i$);
  \item if the $x_i\in V_i$ are homogeneous, set $\varphi_{ab}(\otimes x_i)=(-1)^N\otimes x_i$, where $N$ is the sum of the $\deg(x_i)\deg(x_j)$ over the pairs $(i,j)$ such that $i<_a j$ and $i>_b j$.
\end{enumerate}

\subsubsection*{Example 1}

Take $I=\{1,2\}$, let $a$ be the ordering in which $1<2$, and let $b$ be the ordering in which $1>2$. Let $v_i\in V_i$ be homogeneous, and denote by $v_1\otimes v_2$ (respectively, $v_2\otimes v_1$) the image in the graded tensor product of the product of the $v_i$ in $V(a)$ (respectively, $V(b)$). Then
\[
  v_1\otimes v_2 = (-1)^{\deg(v_1)\deg(v_2)} v_2\otimes v_1
\]
(the Koszul rule).

\subsubsection*{Example 2}

If all the $V_i$ are in degree $1$, there is a canonical isomorphism, relating the module underlying the graded tensor product to the ordinary tensor product of the modules $|V_i|$ underlying the $V_i$:
\[
  \left|\bigotimes V_i\right| \simeq \bigotimes |V_i|\otimes_{\dZ}\bigwedge^{|I|}\dZ^I \text{.}
\]

\subsubsection*{Example 3}

For $A$ a field and $(X_i)$ a finite family of spaces, the K\"unneth formula reads $\h^\bullet\left(\prod X_i,A\right) = \bigotimes \h^\bullet(X_i,A)$.

Let $V^\vee$ be the graded dual of $V$. The canonical map
\begin{equation*}\tag{1.3.1}\label{IV:eq:1-3-1}
  V^\vee\otimes V = V\otimes V^\vee \to A
\end{equation*}
is $v'\otimes v\mapsto v'(v)$.

Now suppose that $A$ is a field, and consider only finite-dimensional vector spaces. The canonical isomorphism
\begin{equation*}\tag{1.3.2}\label{IV:eq:1-3-2}
  W\otimes V^\vee \to \hom(V,W)
\end{equation*}
is $w\otimes v'\mapsto(v\mapsto w\cdot v'(v))$. Via this isomorphism, the composition $\hom(Y,Z)\otimes\hom(X,Y)\to\hom(X,Z)$ is identified with the morphism induced by \eqref{IV:eq:1-3-1}: $Z\otimes Y^\vee\otimes Y\otimes X^\vee \to Z\otimes X^\vee$.

The trace of $f:V\to V$ (zero when $f$ is homogeneous of degree $\ne0$) is the image of $f$ under
\begin{equation*}\tag{1.3.4}\label{IV:eq:1-3-4}
  \tr:\hom(V,V) {\xleftarrow\sim}_\text{\eqref{IV:eq:1-3-2}} V\otimes V^\vee = V^\vee \otimes V \to_\text{\eqref{IV:eq:1-3-1}} A .
\end{equation*}

It is readily checked that, for $f$ of degree $0$,
\begin{equation*}\tag{1.3.5}\label{IV:eq:1-3-5}
  \tr(f,V) = \sum (-1)^i \tr(f,V^i) \text{.}
\end{equation*}

Expressing the fact that the two composites of maps \eqref{IV:eq:1-3-1}, $V^\vee\otimes V\otimes W^\vee\otimes W\to k$, agree, one finds that, for homogeneous $f:V\to W$ and $g:W\to V$,
\begin{equation*}\tag{1.3.6}\label{IV:eq:1-3-6}
  \tr(f g) = (-1)^{\deg(f)\deg(g)} \tr(g f) \text{.}
\end{equation*}

\section{The cohomology class associated with a cycle}\label{IV:2}

\subsection{The class of a divisor}\label{IV:2-1}

\subsubsection{}\label{IV:2-1-1}

Let $D$ be a Cartier divisor in a scheme $X$. Away from $D$, the invertible sheaf $\sO(D)$ is trivialized by the section $1$. The class $\cl(D)$ of $D$ in $\h_D^1(X,\dG_m)$ is the class of the corresponding $\dG_m$-torsor trivialized on $X\setminus D$ (\ref{IV:1-1-6} and \ref{IV:1-1-4}).

Let $\partial:\h^0(X\setminus D,\dG_m)\to\h_D^1(X,\dG_m)$ be the map of \ref{IV:1-1-4}. If $D$ has a global equation $f$, multiplication by $f$ is an isomorphism from $\sO(D)$, trivialized by $1$ on $X\setminus D$, to $\sO$, trivialized by $f$ on $X\setminus D$. By \ref{IV:1-1-4}, therefore,
\begin{equation*}\tag{2.1.1}\label{IV:eq:2-1-1}
  \cl(D) = \partial f \text{.}
\end{equation*}

For every morphism $u:X'\to X$ such that $u^*D$ is again a Cartier divisor (that is, $u^{-1}D$ is disjoint from $\operatorname{Ass}(X')$), one has $\cl(u^*D)=u^*\cl(D)$. To obtain such functoriality for every morphism $u$, one would have to consider, more generally, invertible sheaves equipped with a section rather than merely Cartier divisors.

Recall that the integer $n$ is henceforth assumed invertible on all schemes considered. Let $\partial:\h_D^i(X,\dG_m)\to\h_D^{i+1}(X,\dmu_n)$ be the connecting map for the Kummer exact sequence $0\to\dmu_n\to\dG_m\xrightarrow{(\cdot)^n}\dG_m\to0$.

\begin{definition}\label{IV:2-1-2}
The \emph{class $\cl_n(D)$} of $D$ in $\h_D^2(X,\dmu_n)$ is $\partial\cl(D)$.
\end{definition}

When there is no risk of confusion, the subscript $n$ will be omitted.

\subsubsection{}\label{IV:2-1-3}

The diagram
\[\xymatrix{
  \h^0(X\setminus D,\dG_m) \ar[r]^-\partial \ar[d]^-\partial
    & \h_D^1(X,\dG_m) \ar[d]^-\partial \\
  \h^1(X\setminus D,\dmu_n) \ar[r]^-\partial
    & \h_D^2(X,\dmu_n)
}\]
is anticommutative. If $D$ has a global equation $f$, then $\cl_n(D)$ is the negative of the image under $\partial$ of the class in $\h^1(X\setminus D,\dmu_n)$ of the $\dmu_n$-torsor of $n$th roots of $f$.

\begin{proposition}\label{IV:2-1-4}
Let $i$ be the inclusion of $D$ in $X$. If $D$ and $X$ are regular, the cohomology sheaves with support $\eR^p i^!\dmu_n$ vanish for $p=0,1$, and $\eR^2i^!\dmu_n=\underline{\dZ/n}$, generated by $\cl_n(D)$.
\end{proposition}

It is enough to prove that, when $X$ is strictly local and $D$ is defined by a regular parameter, $\h_D^p(X,\dmu_n)=0$ for $p=0,1$, and $\h_D^2(X,\dmu_n)=\dZ/n$, generated by $\cl(D)$. Denoting reduced cohomology by a tilde, one has $\widetilde\h^{p-1}(X\setminus D,\dmu_n)\iso\h_D^p(X,\dmu_n)$. For $p=0,1$ the assertion says that $D$ does not disconnect $X$, and for $p=2$ it follows, via \ref{IV:2-1-3}, from Abhyankar's lemma.

This is a partial analogue of the relative theorem (\hyperref[I]{Étale cohomology}, \ref{I:5-3-4}). Grothendieck conjectures that the $\eR^pi^!\dmu_n$ vanish for $p\ne2$, at least when $X$ is excellent (the purity conjecture), but this is known only in characteristic $0$ (\cite[XIX]{sga4}).

\begin{theorem}[Fundamental compatibility]\label{IV:2-1-5}
Let $X$ be a smooth curve over an algebraically closed field $k$, let $P$ be a closed point of $X$, and let $\tr$ be the composite $\h_P^2(X,\dmu_n)\to\h_c^2(X,\dmu_n)\xrightarrow{\tr}\dZ/n$. Then
\[
  \tr \cl(P) = 1 \text{.}
\]
\end{theorem}

Let $\bar X$ be the smooth projective curve completing $X$. The formula expresses the fact that the invertible sheaf $\sO(P)$ on $\bar X$ has degree $1$.

\subsection{The cohomological method}\label{IV:2-2}

\subsubsection{}\label{IV:2-2-1}

Let $X$ be a (Noetherian) scheme. Recall that a subscheme $Y$ of $X$ is called a local complete intersection of codimension $c$ if, locally on $Y$, it is defined in $X$ by a regular sequence of $c$ equations. If $X$ is the spectrum of a local ring $A$ with maximal ideal $\fm$, and $Y$ has ideal $\fa$, this means that $\operatorname{Ext}^i(A/\fa,A)=0$ for $i<\dim(\fa/\fm\fa)=c$, and every sequence of elements of $\fa$ whose image in $\fa/\fm\fa$ is a basis is a regular sequence of equations for $Y$.

\subsubsection{}\label{IV:2-2-2}

Let $i:Y\hookrightarrow X$ be a local complete intersection of codimension $c$. We wish to define a local fundamental class $\cl(Y)$ that is a global section of the cohomology sheaf with support $\eR^{2c}i^!\dZ/n(c)$ (recall that $\dZ/n(c)=\dmu_n^{\otimes c}$).

Locally, $Y$ is the intersection of a sequence of $c$ divisors $D_i$, and we define $\cl(Y)$ as the cup product of the $\cl(D_i)$. Each $\cl(D_i)$ is supported in $D_i$, and their product is supported in $Y$. The fact that, locally, this product does not depend on the choice of the $D_i$ follows from \ref{IV:2-2-3} below and the following invariance properties:
\begin{enumerate}[\indent a)]
  \item compatibility with localization;
  \item independence of the order of the $D_i$ (the $\cl(D_i)$ have the even degree $2$, so the cup product is commutative);
  \item the product depends only on the ``flag'' $D_1\supset D_1\cap D_2\supset\cdots\supset Y$.
\end{enumerate}
To prove c), note the existence of a product (a variant of \ref{IV:1-2-1})
\[
  \h_{D_1}^\bullet(X)\otimes \h_{D_1\cap D_2}^\bullet(D_1)\otimes \cdots \otimes \h_Y^\bullet(D_1\cap \cdots \cap D_{c-1}) \to \h_Y^\bullet(X) \text{;}
\]
the product of the $\cl(D_i)$ is also the product of $\cl(D_i)$ restricted to $D_1\cap\cdots\cap D_{i-1}$, that is, of $\cl(D_1\cap\cdots\cap D_i)$ in $D_1\cap\cdots\cap D_{i-1}$.

\begin{lemma}\label{IV:2-2-3}
Let $A$ be a local ring with maximal ideal $\fm$, and let $u=(u_1,\dotsc,u_c)$ and $v=(v_1,\dotsc,v_c)$ be two regular sequences generating the same ideal $\fa$. There is then a sequence $w_i$ ($1\leqslant i\leqslant N$) of such sequences joining them, such that $w_{i+1}$ is obtained from $w_i$ by a permutation or by changing only the last element.
\end{lemma}

Because permutations are available, it is equivalent to allow the change of any one element rather than only the last. By \ref{IV:2-2-1}, it remains to check that in the vector space $\fa/\fm\fa$ one can pass from one basis to another by a sequence of permutations and basis changes modifying only one vector. Indeed, the linear group is generated by diagonal, elementary, and permutation matrices.

\subsubsection{}\label{IV:2-2-4}

The same methods define a local fundamental class for every $Y\subset X$ locally definable by $c$ equations (where fewer equations suffice, the class is zero).

\subsubsection{}\label{IV:2-2-5}

One can pass from such a local fundamental class in $\h^0(Y,\eR^{2c}i^!\dZ/n(c))$ to a global fundamental class in $\h_Y^{2c}(X,\dZ/n(c))$ when semipurity results are available:

\begin{proposition}\label{IV:2-2-6}
\begin{enumerate}[(i)]
  \item If $\eR^pi^!\dZ/n=0$ for $p<2c$, then
    \[
      \h_Y^{2 c}\left(X,\dZ/n(c)\right) \iso \h^0\left(Y,\eR^{2 c} i^! \dZ/n(c)\right) \text{.}
    \]
  \item Let $Z$ be a closed part of $Y$, with complement $V$ in $Y$, and let $k$ be the inclusion of $Z$ in $X$. If $\eR^pi^!\dZ/n=0$ for $p\leqslant2c$, then
    \[
      \h_Y^{2 c}\left(X,\dZ/n(c)\right) \hookrightarrow \h_V^{2 c}\left(X,\dZ/n(c)\right) \text{.}
    \]
    If $\eR^pi^!\dZ/n=0$ for $p\leqslant2c+1$, this map is an isomorphism.
\end{enumerate}
\end{proposition}

The hypothesis $\eR^pi^!\dZ/n=0$ is equivalent to $\eR^pi^!\dZ/n(c)=0$. With this understood, (i) follows from the spectral sequence $\h^p(Y,\eR^qi^!)\Rightarrow\h_Y^{p+q}(X,-)$. If $k_1$ is the inclusion of $Z$ in $Y$, then $k=ik_1$, whence $\eR k^!\dZ/n=\eR k_1^!\eR i^!\dZ/n$, and the long exact cohomology sequence for $Z\subset Y$ gives a long exact sequence
\[\xymatrix{
  \cdots \ar[r]
    & \h^i\left(Z,\eR k^! \dZ/n(c)\right) \ar[r]
    & \h_Y^i\left(X,\dZ/n(c)\right) \ar[r]
    & \h_V^i\left(X,\dZ/n(c)\right) \ar[r]
    & \cdots
}\]
from which (ii) follows.

\subsubsection{Generalization}\label{IV:2-2-7}

Proposition \ref{IV:2-2-6} remains valid for any separated morphism $i$ of finite type, and for any integer $2c$ (positive or negative), provided $\h_Y^{2c}(X,\dZ/n(c))$ is replaced by $\h^{2c}(Y,\eR i^!\dZ/n(c))$, and similarly for $\h_V$.

The following semipurity results follow from \cite[1.8, 1.10, 1.15]{sga2}. We recall their proof.

\begin{theorem}\label{IV:2-2-8}
Let there be a morphism of finite-type $S$-schemes
\[\xymatrix{
  Y \ar[r]^-i \ar[dr]_-f
    & X \ar[d]^-g \\
  & S
}\]
Assume that $X$ is smooth and purely of relative dimension $N$, and that $Y$ has fiber dimension $\leqslant d$. Put $c=N-d$.
\begin{enumerate}[\indent (i)]
  \item One has $\eR^pi^!\dZ/n=0$ for $p<2c$; the same holds with $\dZ/n$ replaced by a sheaf $g^*\sF$.
  \item If, over a dense open $U$ of $S$, $Y$ has fiber dimension $<d$, then $\eR^pi^!\dZ/n=0$ for $p\leqslant2c$. If, moreover, the complement $Z$ of $U$ does not locally disconnect $S$, then $\eR^pi^!\dZ/n=0$ for $p\leqslant2c+1$.
\end{enumerate}
\end{theorem}

Since $\eR g^!\sF=g^*\sF(N)[2N]$, the transitivity formula $\eR i^!\eR g^!=\eR f^!$ shows that $\eR^{2c+q}i^!(g^*\sF)=0\Leftrightarrow\eR^{-2d+q}f^!(\sF)=0$. This reduces us to studying $f$, that is, to assuming that $X=S$. Assertion (i) is then \cite[XVIII 3.17]{sga4}.

Let $Y'$ be the inverse image of $Z$ in $Y$:
\[\xymatrix{
  Y' \ar[r]^-v \ar[d]^-f
    & Y \ar[d]^-f \\
  Z \ar[r]^-u
    & S
}\]
By (i), the $\eR^pf^!\dZ/n$ are supported in $Y'$ for $p\leqslant-2d+1$; and the spectral sequence $\eR^av^!\eR^bf^!\Rightarrow\eR^{a+b}(fv)^!$ shows that they coincide with the $\eR^p(fv)^!\dZ/n=\eR^p(uf)^!\dZ/n$. Applying (i) to $Y'/Z$ and to the spectral sequence $\eR^af^!\eR^bu^!\Rightarrow\eR^{a+b}(uf)^!$, one finds that (ii) follows from the vanishing of $\eR^bu^!\dZ/n$ for $b=0$, or for $b=0$ and $1$, respectively.

\subsubsection{}\label{IV:2-2-9}

Grothendieck conjectures the following absolute analogue of \ref{IV:2-2-8} (the semipurity conjecture, a consequence of the purity conjecture): if $Y$ has codimension $\geqslant c$ in a regular $X$, then $\h_Y^i(X,\dZ/n)=0$ for $i<2c$, at least when $X$ is excellent.

\subsubsection{}\label{IV:2-2-10}

We are now in a position to define the class of a cycle $Y$ of codimension $c$ in $X$, where $X$ is smooth over a field $k$. Write $Y=\sum d_iY_i$, with the $Y_i$ reduced and irreducible. An open $U_i\subset Y_i$ whose complement has codimension $>c$ in $X$ is then a local complete intersection in $X$. This defines the local fundamental class of $U_i$. By \ref{IV:2-2-6} and \ref{IV:2-2-8}, it comes from a unique fundamental class $\cl(Y_i)\in\h_{Y_i}^{2c}(X,\dZ/n(c))$, and we put
\[
  \cl(Y) = \sum d_i \cl(Y_i) \in \h_{|Y|}^{2 c}\left(X,\dZ/n(c)\right) \text{.}
\]

\subsubsection{}\label{IV:2-2-11}

In complex analytic geometry, a construction of Baum, Fulton, and MacPherson \cite{bfm75} defines without restriction the cohomology class of a local complete intersection $Y\subset X$. Suppose that $Y$ is purely of codimension $c$, and let $\sN$ be the normal vector bundle of $Y$ in $X$. Its local sections are those of $(\sI/\sI^2)^\vee$, where $\sI$ is the ideal sheaf of $Y$. Recall that the sheaf of $C^\infty$ functions on $X$ is defined locally, in terms of local embeddings of $X$ in $\dC^n$, as the restriction to $X$ of the quotient of the sheaf of $C^\infty$ functions on $\dC^n$ by the ideal generated by the real and imaginary parts of the equations defining $X$. The bundle $\sN$ extends to a complex $C^\infty$ vector bundle $N$ on a neighborhood $U$ of $Y$ in $X$, and, for $U$ sufficiently small, there are sections $f$ of $N$ whose zero locus is $Y$ and such that $df:\sN\to N|_Y$ is the identity. Two choices of $N$ and $f$ are homotopic after shrinking $U$ about $Y$.

The cohomology class $\cl(U)$ of the zero section of $N$ (denoted $z:U\to N$) is defined: $U$ is a local complete intersection in $N$, and $\eR^iz^!\dZ=0$ for $i\ne2c$, while $\eR^{2c}z^!\dZ=\dZ$. There is a map $f^*:\h_U^\bullet(N,\dZ)\to\h_Y^\bullet(X,\dZ)$, and we put
\[
  \cl(Y) = f^*\cl(U) \text{.}
\]

The same construction works whenever an analytic subspace $Y$ of $X$ is equipped with the following normal structure: a locally free sheaf $\sC$ of rank $c$ on $Y_\text{red}$, and an epimorphism $\sC\to\sI/\sI\cdot\sI_\text{red}$.

\subsection{The homological method}\label{IV:2-3}

\subsubsection{}\label{IV:2-3-1}

Let $f:Y\to X$ be a flat morphism of finite type whose fibers have dimension $\leqslant d$. In \cite[XVIII 2.9]{sga4}, we defined a trace morphism
\begin{equation*}\tag{2.3.1.1}\label{IV:eq:2-3-1-1}
  \tr_f:\eR^{2 d} f_! \dZ/n(d) \to \dZ/n \text{.}
\end{equation*}
One has $\eR^if_!\dZ/n(d)=0$ for $i>2d$, and $\eR^if^!\dZ/n=0$ for $i<-2d$. These facts, together with the adjunction between $\eR f_!$ and $\eR f^!$, give isomorphisms
\begin{equation*}\tag{2.1.3.2}\label{IV:eq:2-1-3-2}
\begin{aligned}
  \hom\left(\eR^{2 d}f_!\dZ/n(d),\dZ/n\right)
    &= \hom\left(\eR f_!\dZ/n(d),\dZ/n[-2 d]\right) \\
    &= \hom\left(\dZ/n,\eR f^!\dZ/n(-d)[-2 d]\right) \\
    &= \h^0\left(Y,\eR^{-2 d}f^!\dZ/n\right) \text{.}
\end{aligned}
\end{equation*}
At least when $S$ is a point, the image of $\tr_f$ in the last two groups deserves the name fundamental class of $Y$ (in homology).

We again denote by $\tr_f$ the image of $\tr_f$ in the second group, and the morphisms obtained from it by functoriality. For example, when $Y/S$ is proper, one obtains the morphism
\begin{equation*}\tag{2.3.1.3}\label{IV:eq:2-3-1-3}
\xymatrix{
  \h^i\left(Y,\dZ/n(d)\right) = \h^i\left(S,\eR f_!\dZ/n(d)\right) \ar[r]
    & \h^{i-2 d}(S,\dZ/n) \text{.}
}
\end{equation*}

Suppose that $Y$ is contained in $X$, with $X$ smooth over $S$ and purely of relative dimension $N$, and put $c=N-d$.
\[\xymatrix{
  Y \ar@{^{(}->}[r]^-i \ar[dr]_-f
    & X \ar[d] \\
  & S
}\]
One has $\eR f^!=\eR i^!\eR g^!$, and $\eR g^!\dZ/n=\dZ/n(N)[2N]$. The fundamental class of $Y$ is identified with an element of $\h^0(Y,\eR i^!\dZ/n(c)[2c])=\h_Y^{2c}(X,\dZ/n(c))$, the class $\cl(Y)$ of $Y$ in $X$. We shall see below that it depends only on $Y\subset X$, not on the projection of $X$ to $S$, and that, when $Y$ is a local complete intersection, it induces the local class of the preceding subsection.

Let us make explicit the isomorphism
\[
  \h_Y^{2 c}\left(X,\dZ/n(c)\right) \iso \hom\left(\eR^{2 d}f_!\dZ/n(d),\dZ/n\right)
\]
that carries $\cl(Y)$ to $\tr_f$. Via the isomorphisms
\begin{align*}
  \h_Y^{2 c}\left(X,\dZ/n(c)\right)
    &= \h_Y^{2 c}\left(\eR i^!\dZ/n(c)\right) \\
    &= \hom_Y\left(\dZ/n(d)[2 d],\eR i^!\dZ/n(N)[2 N]\right) \\
    &= \hom_X\left(\dZ/n(d)[2 d]_Y,\dZ/n(N)[2 N]\right) \text{,}
\end{align*}
it is the top row of
\[\xymatrix{
  \h_Y^{2 c} \ar[r] \ar[dr]_-{(3)}
    & \hom\left(\eR f_!\dZ/n(d)[2 d],\eR f_! \eR i^!\dZ/n(N)[2 n]\right) \ar[r]^-{(1)} \ar[d]^-{\tr_i} \ar@{}[dr]|-{(2)}
    & \hom\left(\eR f_!\dZ/n(d)[2d],\dZ/n\right) \ar@{=}[d] \\
  & \hom\left(\eR g_!\dZ/n(d)[2 d]_Y,\eR g_!\dZ/n(N)[2 N]\right) \ar[r]^-{\tr_g}
    & \hom\left(\eR g_!\dZ/n(d)[2 d]_Y,\dZ/n\right) \text{,}
}\]
where (1) is the adjunction map $\eR f_!\eR f^!\to\operatorname{id}$. The commutativity of (2) expresses that the isomorphism $\eR f^!=\eR i^!\eR g^!$ is defined by adjunction. Finally, the map deduced from (3),
\[
  \h_Y^{2c}\left(X,\dZ/n(c)\right) \otimes \eR^{2 d} g_!\dZ/n(d) \to \eR^{2 N} g_!\dZ/n(N)
\]
is interpreted as a cup product (cf. \ref{IV:1-2}).

\begin{definition}\label{IV:2-3-2}
The class $\cl(Y)\in\h_Y^{2c}(X,\dZ/n(c))$ is characterized by the fact that, for every local section $u$ of $\eR^{2d}f_!\dZ/n(d)$,
\[
  \tr_f(u) = \tr_g\left(\cl(Y)\smallsmile u\right) \text{.}
\]
\end{definition}

Because $\cl(Y)$ already has an image $\tr_f$ in $\hom(\eR f_!\dZ/n(d)[2d],\dZ/n)$, formula \ref{IV:2-3-2} holds for the maps deduced from $\tr_f$ by functoriality. For example, when $X$ is proper over $S$ and $u\in\h^\bullet(Y,\dZ/n)$, the formula holds in $\h^\bullet(S,\dZ/n(-d))$. For $v\in\h^\bullet(X,\dZ/n)$, it gives
\[
  \tr_f(i^* v) = \tr_g\left(\cl(Y)\smallsmile v\right) \text{,}
\]
where the cup product may be computed in $\h^\bullet(X,\dZ/n)$.

\subsubsection{}\label{IV:2-3-3}

We shall define the trace morphisms, and hence the class $\cl(Y)$, under more general hypotheses. Consider a diagram
\[\xymatrix{
  Y \ar@{^{(}->}[r] \ar[dr]_-f
    & X \ar[d]^-g \\
  & S
}\]
where $g$ is smooth and purely of relative dimension $N$, and $Y$ is closed in $X$, with fiber dimension $\leqslant d$. In addition, equip $Y$ with a ``weight'' $\sK$ of the following kind: $\sK\in\eD_{\text{parf},X}^b$ is a bounded complex of sheaves of $\sO$-modules on $X$, of finite Tor-dimension over $S$ (equivalently, over $X$), whose cohomology sheaves are coherent and supported in $Y$. We wish to define a morphism $\tr_{f,\sK}:\eR^{2d}f_!\dZ/n(d)\to\dZ/n$. When $Y$ is flat over $S$ and $\sK=\sO_Y$, this will be the preceding trace morphism. In general, it depends only on the lengths of $\sK$ at the generic points $y$ of the irreducible components of $Y$: $\operatorname{lg}_y(\sK)=\sum(-1)^i\operatorname{lg}\h^i(\sK)_y$.
% NOTE: notation is not clear, is it a calligraphic H?
Finally, denote by $\cl(Y,\sK)$ the class supported in $Y$ such that
\begin{equation*}\tag{2.3.3.1}\label{IV:eq:2-3-3-1}
  \tr_g\left(\cl(Y,\sK)\smallsmile u\right) = \tr_{f,\sK}(u) \text{.}
\end{equation*}

\subsubsection{}\label{IV:2-3-4}

The construction of $\tr_{f,\sK}$ parallels that of \cite[XVIII.2]{sga4}; we shall give only the main outlines.

\paragraph{A}
$d=0$ ($f$ quasi-finite). Let $x$ be a geometric point of $Y$, let $s$ be its image in $S$, let $\sK_{(x)}$ be the inverse image of $\sK$ on the strict localization of $X$ at $x$, and let $\sK_{(x)s}=\sK_{(x)}\lotimes_{\sO_{S,s}^h}k(s)$ be its inverse image on the geometric fiber at $s$. Put $n(x)=\sum(-1)^i\dim_{k(s)}\h^i(\sK_{(x)s})$. The function $x\mapsto n(x)$ is a weighting of $f$, and we take the corresponding trace morphism \cite[XVII.6.2.5]{sga4}. Formation of the weighting $n(x)$ and of $\tr_{f,\sK}$ is compatible with every base change.

\paragraph{B}
The general case. If $u:X\to\dA_S^d$ is such that $ui$ is quasi-finite, put $\tr_{f,\sK}=\tr_{\dA_S^d}\circ\tr_{ui,\sK}$. This trace morphism is plainly compatible with every base change $S'/S$. To prove that it does not depend on $u$, we may therefore assume that $S$ is the spectrum of an algebraically closed field $k$. Let $Y_i$ be the irreducible components of $Y$, and let $\operatorname{lg}_i$ be the length of $\sK$ at the generic point of $Y_i$. If $\alpha_i$ is the inclusion of $(Y_i)_\text{red}$ in $Y$, then
\begin{equation*}\tag{2.3.4}\label{IV:eq:2-3-4}
  \tr_{f,\sK}
  = \sum \operatorname{lg}_i\,\tr_{Y_{i,\text{red}}/S}\,\alpha_i^\ast
\end{equation*}
(this follows from the analogous and elementary formula for $\tr_{ui}$).

This defines $\tr_{f,\sK}$ locally on $Y$; one then proceeds as in \cite[XVIII.2.9]{sga4}.

\begin{lemma}\label{IV:2-3-5}
Suppose that $g=g'g'':X\xrightarrow{g''}S'\xrightarrow{g'}S$, where $g'$ and $g''$ are smooth and purely of relative dimensions $N'$ and $N''$, and that over $S'$ the fiberwise codimension of $Y$ is still $\geqslant c$. Then the class $\operatorname{cl}(Y,\sK)$ is the same whether computed in terms of $g$ or of $g''$.
\end{lemma}

Put $f'=g''i$ and $d'=d-N'$. The formula $\tr_g=\tr_{g'}\tr_{g''}$ (\cite[XVIII.2.9]{sga4}, Var.\ 3) ensures the commutativity of
\[\xymatrix{
  \h_Y^{2 c}\left(X,\dZ/n(c)\right) \ar[d] \\
  \hom_{S'}\left(\eR f_!\dZ/n(d')[2 d'],\eR g_!''\dZ/n(N'')[2 N'']\right) \ar[r] \ar[d]^-{\eR g_!'(N')}
    & \hom_{S'}\left(\eR^{2 d'} f_!' \dZ/n,\dZ/n\right) \ar[d]^-{\tr_{g'}\circ \eR^{2 N'} g_!'} \\
  \hom_S\left(\eR f_!\dZ/n(d)[2d],\eR g_!'\dZ/n(N)[2 N]\right) \ar[r]
    & \hom_S\left(\eR^{2 d} f_!\dZ/n,\dZ/n\right)
}\]
while the analogous and easily checked formula $\tr_{f,\sK}=\tr_{g'}\tr_{f',\sK}$ ensures that the image of $\tr_{f',\sK}$ is $\tr_{f,\sK}$. The assertion follows.

\begin{lemma}\label{IV:2-3-6}
If $Y$ is a divisor in $X$, the class of \ref{IV:2-3-2} coincides with the class of \ref{IV:2-1-2}.
\end{lemma}

By semipurity (\ref{IV:2-2-6}(i) and \ref{IV:2-2-8}(i)), the question is local on $Y$; using \ref{IV:2-3-5} to replace $S$ by a suitable $S'$ allows us to assume that $N=1$. Using semipurity again (\ref{IV:2-2-6}(ii) and \ref{IV:2-2-8}(ii)), it is enough to prove \ref{IV:2-3-6} over the generic points of $S$. An étale localization then reduces us to the case in which $S$ is the spectrum of an algebraically closed field $k$. Since the classes of \ref{IV:2-3-2} and \ref{IV:2-1-2} are each additive in $Y$, this reduces us to the case in which $Y$ is a closed point on a smooth curve over $k$. The assertion that the characterizing property \ref{IV:2-3-2} holds is then precisely the fundamental compatibility \ref{IV:2-1-5}.

\begin{lemma}\label{IV:2-3-7}
Formation of $\cl(Y,\sK)$ is compatible with every base change $S'/S$.
\end{lemma}

This follows from the same assertion for the trace morphisms.

\begin{theorem}\label{IV:2-3-8}
(i) The class $\cl(Y,\sK)$ depends only on $Y\subset X$ and on the lengths of $\sK$ at the generic points of $Y$. It is additive in $\sK$. For $\sK=\sO_Y$ and $Y$ a local complete intersection, it induces the local class of \ref{IV:2-2-2}.

(ii) Let there be a commutative diagram
\[\xymatrix{
  X' \ar[r]^-u \ar[d]
    & X \ar[d] \\
  S' \ar[r]
    & S
}\]
with $X'$ smooth over $S'$. If $u^{-1}(Y)$ still has fiberwise codimension $\geqslant c$, then
\[
  \cl\left(u^{-1}(Y),\eL u^\ast\sK\right) = u^\ast \cl(Y,\sK) \text{.}
\]

(iii) Let $Y'\subset X$ have fiberwise codimension $\geqslant c'$, let $\sK'$ be a weight on $Y'$, and suppose that $Y\cap Y'$ has fiberwise codimension $\geqslant c+c'$. Then
\[
\cl\left(Y\cap Y',\sK\lotimes \sK'\right) = \cl(Y,\sK) \smallsmile \cl(Y',\sK') \text{.}
\]
\end{theorem}

% The following paragraphs were originally numbered

% (2.3.8.1)
\paragraph{(A)}
Given $Y\subset X/S$ as above, semipurity (\ref{IV:2-2-6}, \ref{IV:2-2-8}) shows that, to check that two classes in $\h_Y^{2c}(X,\dZ/n(c))$ coincide, it is enough to check this locally at the generic points of $Y$, or even after any base change $s\to S$, with $s$ a geometric generic point of $S$, at the generic points of $Y$.

% (2.3.8.2)
\paragraph{(B) Proof of (ii) for $S=S'$ and $u$ a closed immersion}
By localization (A), we may assume that $X'$ is the inverse image of the zero section under a smooth morphism $v:X\to\dA_S^{N'}$:
\[\xymatrix{
  X' \ar[r]^-u \ar[d]
    & X \ar[d]^-v \\
  S \ar[r]^-0
    & \dA_S^{N'} \ar[r]^-\pi
    & S
}\]
Apply \ref{IV:2-3-5} to $g=\pi v$ and \ref{IV:2-3-7} to the base change $S\to\dA_S^{N'}$.

% (2.3.8.3)
\paragraph{(C) In a product situation:}
If $X=X'\times_SX''$, $Y=Y'\times_SY''$, and $\sK=\sK'\lotimes_S\sK''$, then $\cl(Y,\sK)=\cl(Y',\sK')\smallsmile\cl(Y'',\sK'')$ (the external cup product, that is, $\operatorname{pr}_1^\ast\cl(Y',\sK')\smallsmile \operatorname{pr}_2^\ast\cl(Y'',\sK'')$).

The functorial properties of $\tr_g$ \cite[XVIII.2.12]{sga4} imply the commutativity of \small
\[\xymatrix@=0.5cm{
  \h_{X'}^{2c'}\otimes \h_{X''}^{2 c''}\ar[r] \ar[d]
    & \hom_S\left(\eR f_!' \dZ/n(d')[2 d'],\eR g_!'\dZ/n(N')[2 N']\right)\otimes \cdots \ar[r] \ar[d]
    & \hom_S\left(\eR^{2 d'} f_!' \dZ/n,\dZ/n\right) \otimes \cdots \ar[d] \\
  \h_X^{2 c} \ar[r]
    & \hom_S\left(\eR f_!\dZ/n(d)[2 d],\eR g_!\dZ/n(N)[2 N]\right) \ar[r]
    & \hom_S\left(\eR^{2 d}f_!\dZ/n,\dZ/n\right)
}\]
\normalsize and one checks that $\tr_f$ is the image of $\tr_{f'}\otimes\tr_{f''}$, which proves the assertion.

\paragraph{Proof of (ii)}
By \ref{IV:2-3-7} and a preliminary base change, we may assume that $S=S'$. Factor $u$ through its graph:
\[\xymatrix{
  X' \ar[r]^-{(\text{id},u)}
    & X'\times X \ar[r]^-{\operatorname{pr}_2}
    & X \text{.}
}\]
This leads us to treat $\operatorname{pr}_2$ and $(\text{id},u)$ separately. The first is covered by (C), taking $Y'=X'$ and $\sK'=\sO_{X'}$ (with $\cl(Y')=1\in\h^0(X',\dZ/n)$), and the second by (B).

\paragraph{Proof of (iii)}
Let $\delta:X\to X\times_SX$ be the diagonal immersion. The formula $\delta^\ast(Y_1\times_SY_2)=Y_1\cap Y_2$ reduces the assertion to (B) and (C).

\paragraph{Proof of (i)}
The base change $S_\text{red}\to S$ replaces $X$ by $X_\text{red}$; this reduces us to the case in which $X$ and $S$ are reduced. Localization near the generic points of $Y$ then reduces us to the case in which $Y_\text{red}$ is irreducible and is the complete intersection of $c$ divisors $D_i$. The same localization, together with the definition of $\tr$ in the quasi-finite case, then shows that $\cl(Y,\sK)=\operatorname{lg}\cl(Y_\text{red},\sO_{Y_\text{red}})$, where $\operatorname{lg}$ is the length of $\sK$ at the generic point of $Y$. By (iii), therefore,
\[
  \cl(Y,\sK) = \operatorname{lg} \prod_i \cl(D_i,\sO_{D_i})
\]
and the conclusion follows from \ref{IV:2-3-6}.

\subsubsection{Remark}\label{IV:2-3-9}

When $S$ is the spectrum of a field and the class of a cycle is defined as in \ref{IV:2-2-10}, assertion (iii) says that, if two cycles $Y'$ and $Y''$ intersect in the expected dimension, then
\[
  \cl(Y'\cap Y'') = \cl(Y')\smallsmile \cl(Y'')\text{,}
\]
provided the multiplicities of the components of $Y'\cap Y''$ are computed as alternating sums of Tor.

\subsubsection{Remark}\label{IV:2-3-10}

Let $X/\spec(k)$, with $k$ algebraically closed. If two cycles of codimension $c$ in $X$ are algebraically equivalent, the existence of the relative theory above shows that the images of their classes in $\h^{2c}(X,\dZ/n(c))$ are the fibers, at two points of a suitable connected $S$, of a section over $S$ of the constant sheaf $\h^{2c}(X,\dZ/n(c))$ (the sheaf $\eR^{2c}f_*\dZ/n(c)$ for $f=\operatorname{pr}_2:X\times S\to S$). Thus algebraically equivalent cycles have the same class in $\h^{2c}(X,\dZ/n(c))$.

\section{Application: the Lefschetz trace formula in the proper smooth case}\label{IV:3}

\subsection{}\label{IV:3-1}

Let $X$ and $Y$ be proper smooth algebraic varieties over an algebraically closed field $k$, purely of dimensions $N$ and $M$. Via the K\"unneth isomorphism $\h^\bullet(X\times Y)=\h^\bullet(X)\otimes\h^\bullet(Y)$, the trace morphism $\h^\bullet(X\times Y)(M)\to\h^\bullet(X)$ is simply $\h^\bullet(X)$ tensored with the trace morphism $\h^\bullet(Y)(M)\to\dQ_\ell$. Since the morphism is homogeneous of even degree, there is no sign issue. We denote it by $\int_Y$.

\subsection{}\label{IV:3-2}

A class $\eta\in\h^{2q}(X\times Y)(q)$ defines a morphism of degree $2(q-N)$, $\eta_{XY}:\h^\bullet(Y)(N-q)\to\h^\bullet(X)$, by the formula
\[
  \beta \mapsto \int_Y \eta\cdot \operatorname{pr}_2^\ast \beta \text{.}
\]

\begin{proposition_}\label{IV:3-3}
Let $p$ and $q$ be integers with $p+q=N+M$, and let $\varepsilon\in\h^{2p}(X\times Y)(p)$ and $\eta\in\h^{2q}(X\times Y)(q)$. Interpreting the trace of $\eta_{XY}\varepsilon_{YX}:\h^\bullet(X)\to\h^\bullet(X)$ in the sense of \ref{IV:1-3}, one has
\[
  \tr_{X\times Y}(\eta\cdot \varepsilon) = \tr\left(\eta_{X Y}\varepsilon_{Y X},\h^\bullet(X)\right) \text{.}
\]
\end{proposition_}

To exorcise the signs, it is useful to retain only the underlying $\dZ/2$-grading of $\h^\bullet$. To avoid carrying Tate twists throughout, also fix an isomorphism $\dQ_\ell(1)=\dQ_\ell$. Let $\alpha_Y$ be the homomorphism $\h^\bullet(Y)\to\h^\bullet(Y)^\vee$ (an isomorphism) making the diagram
\[\xymatrix{
  \h^\bullet(Y)\otimes \h^\bullet(Y) \ar[r] \ar[d]^-{\alpha_Y\otimes 1}
    & \h^\bullet(Y) \ar[r]^-{\tr}
    & \dQ_\ell \ar@{=}[d] \\
  \h^\bullet(Y)^\vee \otimes \h^\bullet(Y) \ar[rr]^-{\text{(\ref{IV:eq:1-3-1})}}
    & & \dQ_\ell \text{.}
}\]
commute. The map $\eta_{XY}$ is the image of $\eta$ under
\[\xymatrix{
  \h^\bullet(X\times Y) = \h^\bullet(X)\otimes \h^\bullet(Y) \ar[r]^-{\alpha_Y}
    & \h^\bullet(X) \otimes \h^\bullet(Y)^\vee \ar@{=}[r]^-{\text{(\ref{IV:eq:1-3-2})}}
    & \hom\left(\h^\bullet(Y),\h^\bullet(X)\right) \text{,}
}\]
and similarly for $\varepsilon$. The definition \ref{IV:eq:1-3-4} of the trace then reduces \ref{IV:3-3} to the formula $\tr_{X\times Y}=\tr_X\otimes\tr_Y$.

\subsection{Remark}\label{IV:3-4}

The proof of \ref{IV:3-3} does not use the fact that $\alpha_X$ and $\alpha_Y$ are isomorphisms. It remains valid without assuming $X$ and $Y$ smooth. Outside the smooth case, however, it is difficult to find \emph{cohomology} classes to which one would want to apply \ref{IV:3-3}.

\subsection{Remark}\label{IV:3-5}

If $\varepsilon,\eta$ are the cohomology classes of algebraic cycles $A$ and $B$ on $X\times Y$, of codimensions $p$ and $q$, and if $A\cap B$ has dimension $0$, then $\tr_{X\times Y}(\eta\varepsilon)$ is the intersection multiplicity $A\cdot B$, computed as an alternating sum of Tor (\ref{IV:2-3-9}).

\begin{proposition_}\label{IV:3-6}
For $q=M$, with $\eta=\cl(B)$ the class of a subscheme $B\subset X\times Y$ finite over $X$, the map $\eta_{XY}$ is the composite
\[\xymatrix{
  \eta_{X Y} : \h^\bullet(Y) \ar[r]^-{\operatorname{pr}_2^\ast}
    & \h^\bullet(B) \ar[r]^-{\tr_{B/X}}
    & \h^\bullet(X) \text{.}
}\]
\end{proposition_}

Indeed, by \ref{IV:3-2} and \eqref{IV:eq:2-3-3-1} for $\sK=\sO_B$,
\[
  \eta_{X Y}(\beta) = \tr_{X\times Y/X}\left(\cl(B)\cdot \operatorname{pr}_2^\ast\beta\right) = \tr_{B/X}\left(\operatorname{pr}_2^\ast \beta\right) \text{.}
\]

The most important case is that in which $B$ is the graph of a map $f:X\to Y$. In this case, \ref{IV:3-6} shows that $\eta_{XY}=f^\ast$.

\begin{corollary_}\label{IV:3-7}
Let $f$ be an endomorphism of $X$, where $X$ is proper and smooth over an algebraically closed field $k$. Suppose that the fixed points of $f$ are isolated. Then the trace $\sum(-1)^i\tr(f^\ast,\h^i(X))$ is the number of fixed points of $f$, each counted with its multiplicity.
\end{corollary_}

In \ref{IV:3-5}, take $A$ to be the graph of the identity and $B$ the graph of $f$, and apply \ref{IV:3-3} and \ref{IV:3-6}.

\begin{corollary_}\label{IV:3-8}
Let $X$ be a curve over an algebraically closed field $k$, obtained from $X_0$ over $\dF_q$ by extension of scalars, and let $f$ be the Frobenius morphism. Then $\sum(-1)^i\tr(f^\ast,\h_c^i(X))$ is the number of fixed points of $f$.
\end{corollary_}

We may replace $X$ by $X_\text{red}$, and thus assume that $X$ is reduced. Let $U$ be the open subset of $X$ on which $X$ is smooth of dimension $1$, and let $\bar U$ be the smooth projective completion of $U$. The long exact cohomology sequences for $(X,U)$ and $(\bar U,U)$ give
\[
  \tr\left(f^\ast,\h_c^\bullet(X)\right) = \tr\left(f^\ast,\h_c^\bullet(U)\right) + \tr\left(f^\ast,\h_c^\bullet(X\setminus U)\right)
\]
and
\[
  \tr\left(f^\ast,\h_c^\bullet(\bar U)\right) = \tr\left(f^\ast,\h_c^\bullet(U)\right) + \tr\left(f^\ast,\h_c^\bullet(\bar U\setminus U)\right) \text{.}
\]
The same formulas hold for the numbers of fixed points. Since \ref{IV:3-8} is clear for a zero-dimensional scheme, this reduces the proof of \ref{IV:3-8} to $\bar U$. That case follows from \ref{IV:3-7}. The fact that $df=0$ ensures that all fixed points have multiplicity one.

\subsection{Remark}\label{IV:3-9}

For $X$ of dimension $\leqslant1$ over an algebraically closed field $k$, and $f:X\to X$ such that $df=0$, it is not difficult to check that \ref{IV:3-8} remains valid.
% END INLINED FILE: sga4.5-4.tex
\clearpage
\clearpage
% BEGIN INLINED FILE: sga4.5-5.tex
\chapter{Duality}\label{V}

This exposé contains several theorems and compatibility results, all concerning Poincaré duality. Section \ref{V:1} gives the local biduality theorem in dimension $1$ \cite[I.5.1]{sga5} and some calculations of duals. Section \ref{V:2} gives a very economical proof of Poincaré duality for curves that I learned from M.\ Artin. Section \ref{V:3} establishes a compatibility linking two definitions of the pairing that gives Poincaré duality for curves: one by cup product, the other by autoduality of the Jacobian. Finally, Section \ref{V:4} proves the compatibility named in its title.

Throughout the exposé, schemes are Noetherian and separated, and $n$ is an integer invertible on every scheme considered.

\section{Local biduality in dimension \texorpdfstring{$1$}{1}}\label{V:1}

\subsection{}\label{V:1-1}

Let $S$ be a regular scheme purely of dimension $1$. We shall show that the complex consisting of $\dZ/n$ in degree $0$ is dualizing: namely, for $\sK\in\ob\eD_c^b(S,\dZ/n)$ (where the subscript $c$ means constructible), if $D\sK=\rHom(\sK,\dZ/n)$, then $D\sK$ is again constructible with bounded cohomology, and the canonical morphism $\alpha:\sK\to DD\sK$ is an isomorphism.

For $\sK$ in $\eD^-$ and arbitrary $\sL$, one has
\[
  \hom(\sK\lotimes \sL,\dZ/n) = \hom(\sL,\Hom(\sK,\dZ/n)) \text{.}
\]
The morphism $\sK\to DD\sK$ is defined assuming that $D\sK$ lies in $\eD^-$; if $\beta:D\sK\lotimes\sK\to\dZ/n$ is the canonical pairing, it is defined by the pairing $\sK\lotimes D\sK=D\sK\lotimes \sK \xrightarrow\beta \dZ/n$.

\subsection{}\label{V:1-2}

Let $f:X\to S$ be a separated morphism of finite type. Put $\sK_X=\eR f^!\dZ/n$. For $\sK\in\ob\eD^-(X,\dZ/n)$, put $D\sK=\rHom(\sK,\sK_X)$. The adjunction between $\eR f_!$ and $\eR f^!$ gives
\[\xymatrix{
  \eR f_\ast \sK\lotimes \eR f_\ast D\sK \ar[r]
    & \eR f_! (\sK\lotimes D\sK) \ar[r]
    & \eR f_! \eR f^! \dZ/n \ar[r]
    & \dZ/n \text{.}
}\]
The symmetry of this description shows that, for $f$ proper, the diagram
\begin{equation*}\tag{1.2.1}\label{V:eq:1-2-1}
\xymatrix{
  \eR f_\ast \sK \ar[r] \ar[d]
    & D D \eR f_\ast \sK \ar[d]^-\wr \\
  \eR f_\ast D D \sK \ar[r]^\sim
    & D\eR f_\ast D\sK
}
\end{equation*}
is commutative.

If $f$ is the inclusion of a closed point, then $\eR f^!\dZ/n=\dZ/n(-1)[-2]$, that is, up to twist and shift, $\dZ/n$. Over the spectrum of a field this complex is dualizing (Pontryagin duality for $\dZ/n$-modules). By \eqref{V:eq:1-2-1}, therefore, $\sK\iso DD\sK$ for $\sK$ of the form $\eR f_!\sL$, and hence whenever the support of $\underline\h^\bullet(\sK)$ is finite.

\begin{theorem_}\label{V:1-3}
Let $j:U\hookrightarrow S$ be a dense open subscheme, and let $\sF$ be a constructible locally constant sheaf of $\dZ/n$-modules on $U$. Then $D j_*\sF=j_*D\sF$; that is, $\Hom(j_*\sF,\dZ/n)=j_*\Hom(\sF,\dZ/n)$ and $\Ext^i(j_*\sF,\dZ/n)=0$ for $i>0$.
\end{theorem_}

On $U$, $\Ext^i(j_*\sF,\dZ/n)=0$ for $i>0$, because $\sF$ is locally constant (and $\dZ/n$ is injective as a $\dZ/n$-module), while for $i=0$ it is the dual $\sF^\vee$ of $\sF$. One checks that $\hom(j_*\sF,\dZ/n)=j_*\sF^\vee$, and it remains to verify the vanishing of the $\Ext^i$ ($i>0$) at the points of $S\setminus U$. The question is local at these points. We may thus assume that $S$ is a strictly local trait and that $U$ consists of its generic point $\eta$. Let $I=\gal(\bar\eta/\eta)$. The sheaf $\sF$ is identified with the Galois module $\sF_{\bar\eta}$, and the special stalk of $j_*\sF$ with $\sF_{\bar\eta}^I$.

Let $i$ be the inclusion of the closed point $s$, and apply $D$ to the exact sequences
\[\xymatrix{
  0 \ar[r]
    & j_! \sF \ar[r]
    & j_\ast \sF \ar[r]
    & i_\ast \sF_{\bar\eta}^I \ar[r]
    & 0 \\
  0 \ar[r]
    & j_! \sF_{\bar\eta}^I \ar[r] \ar[u]
    & \sF_{\bar\eta}^I \ar[r] \ar[u]
    & i_\ast \sF_{\bar\eta}^I \ar[r] \ar[u]
    & 0 \text{.}
}\]

This gives a morphism of triangles
\[\xymatrix{
  i_\ast\left(\sF_{\bar\eta}^I\right)^\vee (-1)[-2] \ar[r] \ar@{=}[d]
    & D j_\ast\sF \ar[r] \ar[d]
    & \eR j_\ast \sF^\vee \ar[d] \\
  i_\ast\left(\sF_{\bar\eta}^I\right)^\vee(-1)[-2] \ar[r]
    & \left(\sF_{\bar\eta}^I\right)^\vee \ar[r]
    & \eR j_\ast \sF_{\bar\eta}^I \text{.}
}\]
The long exact sequence deduced from the first row gives the vanishing of $\Ext^i$ for $i>2$, and, taking the stalk at $s$, gives
\[\xymatrix@=0.7cm{
  0 \ar[r]
    & \left(\underline\h^1(D j_\ast\sF\right)_s \ar[r]
    & \h^1\left(I,\sF_{\bar\eta}^\vee\right) \ar[r]^-\partial \ar[d]
    & \left(\sF_{\bar\eta}^I\right)^\vee(-1) \ar[r] \ar@{=}[d]
    & \left(\underline\h^2 D j_\ast\sF\right)_s \ar[r]
    & 0 \\
  & 0 \ar[r]
    & \h^1\left(I,(\sF_{\bar\eta}^I)^\vee\right) \ar[r]^-\partial
    & \left(\sF_{\bar\eta}^I\right)^\vee(-1) \ar[r]
    & 0 \text{.}
}\]
Since $\left(\sF_{\bar\eta}^I\right)^\vee =\left(\sF_{\bar\eta}^\vee\right)_I$, putting $M=\sF_{\bar\eta}^\vee$, it remains to verify that
\[\xymatrix{
  \h^1(I,M) \ar[r]^-\sim
    & \h^1(I,M_I) \text{.}
}\]
If $p$ is the residual characteristic exponent, $I$ is an extension of a group isomorphic to $\widehat\dZ_{p'}=\varprojlim_{(m,p)=1}\dZ/m$ by a $p$-group $P$. Since the order of $M$ is prime to $p$, one has $\h^i(P,M)=0$ for $i>0$ and $(M_I)^P=(M^P)_I$. This permits us to replace $M$ by $M^P$ and $I$ by $I/P$. Finally, there is a functorial isomorphism
\[
  \h^1(\widehat\dZ_{p'},M)\iso
  (\text{coinvariants of $\widehat\dZ_{p'}$ in $M$}),
\]
which proves the theorem.

\begin{theorem_}\label{V:1-4}
For $\sK\in\ob\eD_c^b(S,\dZ/n)$, one has $\sK\iso DD\sK$.
\end{theorem_}

By dévissage, we reduce to the case in which $\sK$ consists of a constructible sheaf $\sF$ in degree $0$, and $\sF$ either has finite support (\ref{V:1-2}) or is of the form $j_*\sF_1$ as in \ref{V:1-3}. In the latter case, \ref{V:1-3} reduces us to local biduality for $\sF$ locally constant on $U$.

\section{Poincaré duality for curves, after M.\ Artin}\label{V:2}

\subsection{}\label{V:2-1}

Let $X$ be a smooth projective curve over an algebraically closed field $k$. Put $\sK_X=\dZ/n(1)[2]$ and, for $\sK\in\ob\eD_c^b(X,\dZ/n)$, put $D\sK=\rHom(\sK,\sK_X)$. For a $\dZ/n$-module $M$, also put $DM=\hom(M,\dZ/n)$, and similarly for complexes of modules. The trace morphism $\h^0(X,\sK_X)\to\dZ/n$, or $\eR\Gamma(X,\sK_X)\to\dZ/n$, defines a pairing $\eR\Gamma(X,\sK)\lotimes\eR\Gamma(X,D\sK)\to\dZ/n$. Poincaré duality between the cohomology and the cohomology with proper support of an open subscheme $j:U\hookrightarrow X$ says that, for $\sK=j_!\dZ/n$, this pairing identifies each factor with the dual of the other.

I give below a proof, communicated to me by M.\ Artin, that for every $\sK\in\ob\eD_c^b(X,\dZ/n)$ this pairing is perfect, that is, defines an isomorphism
\begin{equation*}\tag{2.1.1}\label{V:eq:2-1-1}
\xymatrix{
  \eR\Gamma(X,\sK) \ar[r]^-\sim
    & D \eR\Gamma(X,D\sK) \text{.}
}
\end{equation*}

For every constructible sheaf $\sF$, put
\[
  '\h^i(X,\sF) = \text{the $\dZ/n$-dual of $\h^{-i}(X,D\sF)$.}
\]
The assertion that \eqref{V:eq:2-1-1} is an isomorphism is equivalent to the following theorem.

\begin{theorem_}\label{V:2-2}
For every constructible sheaf $\sF$ of $\dZ/n$-modules, one has
\begin{equation*}\tag{2.2.1}\label{V:eq:2-2-1}
\xymatrix{
  \h^i(X,\sF) \ar[r]^-\sim
    & '\h^i(X,\sF) \text{.}
}
\end{equation*}
\end{theorem_}

\begin{lemma_}\label{V:2-3}
Let $f:X\to Y$ be a generically étale morphism between smooth curves over $k$. Then $\sK_X=f^\ast\sK_Y=\eR f^!\sK_Y$, with $\tr_f$ as the adjunction map $\eR f_!\eR f^!\sK_Y\to\sK_Y$.
\end{lemma_}

% NOTE: \rhom or \rHom ? originally \rhom
It is enough to check that, for $f$ finite, $f_*\eR f^!\sK_Y=\rHom(f_*\dZ/n,\sK_Y)$ is $f_*\dZ/n$, with $\tr_f$ as the adjunction map. The vanishing of the $\Ext^i$ for $i>0$ follows from \ref{V:1-3}, and the assertion follows.

More explicitly, $\eR f^!$ is the derived functor of $f^!:f^!\sF(U)=\hom(f_!\dZ_U,\sF)$, and \ref{V:1-3} implies the vanishing of $\eR^if^!\dZ/n$ for $i>0$.

\begin{corollary_}\label{V:2-4}
Let $f:X'\to X$ be a generically étale morphism of smooth projective curves over $k$. Then $'\h^i(X,f_*\sF)='\h^i(X',\sF)$; this isomorphism is functorial, and the diagram
\[\xymatrix{
  \h^i(X,f_\ast\sF) \ar[r] \ar@{=}[d]
    & '\h^i(X,f_\ast\sF) \ar@{=}[d] \\
  \h^i(X',\sF) \ar[r]
    & '\h^i(X',\sF)
}\]
is commutative.
\end{corollary_}

The isomorphism is given by \ref{V:1-2}: $Df_*\sF=f_*D\sF$.

\subsection{}\label{V:2-5}

Let us prove \ref{V:2-2}. If $\sF$ is $\dZ/n$ at one point and extended by zero, then $\h^\bullet(D\sF)=\ext^\bullet(\sF,\sK_X)=\dZ/n$ in degree $0$, and in the pairing $\h^\bullet(\sF)\otimes\h^\bullet(D\sF)\to\dZ/n$ one has $1\otimes1\mapsto1$ (\hyperref[IV]{Cycles}, \ref{IV:2-1-5}); thus the duality is perfect. The case in which $\sF$ has finite support is treated in the same way.

For arbitrary constructible $\sF$, let $\sA$ be the subsheaf of sections with finite support, and put $\sG=\sF/\sA$. The exact sequence
\begin{equation*}\label{V:eq:2-5-1}\tag{2.5.1}
\xymatrix{
  0 \ar[r]
    & \sA \ar[r]
    & \sF \ar[r]
    & \sG \ar[r]
    & 0
}
\end{equation*}
gives a diagram
\begin{equation*}\tag{2.5.2}\label{V:eq:2-5-2}
\xymatrix@=.5cm{
  & & 0 \ar[r]
    & \h^0(X,\sA) \ar[r] \ar[d]^-\wr
    & \h^0(X,\sF) \ar[r] \ar[d]
    & \h^0(X,\sG) \ar[r] \ar[d]
    & 0 \\
  0 \ar[r]
  & '\h^{-1}(X,\sF) \ar[r]
  & '\h^{-1}(X,\sF) \ar[r]
  & '\h^0(X,\sA) \ar[r]
  & '\h^0(X,\sF) \ar[r]
  & '\h^0(X,\sF) \ar[r]
  & 0
}
\end{equation*}
and compatible isomorphisms $\h^i(X,\sF)\iso\h^i(X,\sG)$ and $'\h^i(X,\sF)\iso'\h^i(X,\sG)$ for $i\ne0,-1$. If $j:U\to X$ is an open subscheme on which $\sF$ is locally constant, there is an exact sequence
\begin{equation*}\tag{2.5.3}\label{V:eq:2-5-3}
\xymatrix{
  0 \ar[r]
    & \sG \ar[r]
    & j_\ast j^\ast \sF \ar[r]
    & \sB \ar[r]
    & 0
}
\end{equation*}
with $\sB$ of finite support. By \ref{V:1-3}, $'\h^i(X,j_*j^*\sF)=D\h^{2-i}(X,j_*(j^*\sF)^\vee(1))$; hence this $'\h^i$ vanishes for $i<0$, and \eqref{V:eq:2-5-2}, together with the long exact sequence defined by \eqref{V:eq:2-5-3}, shows that $'\h^i$ always vanishes for $i<0$.

If $\sF=\dZ/n$, the theorem is true for $i=0$ and $2$. This expresses the fact that, for connected $X$, the trace morphism $\h^2(X,\dmu_n)\iso\dZ/n$ is an isomorphism. By \ref{V:2-4}, the theorem remains true for $i=0$ and $\sF=f_*\dZ/n$.

For arbitrary $\sF$ again, and $\sG$ as above, there is an $f:X'\to X$ as in \ref{V:2-4} such that $\sG$ embeds in $f_*\dZ/n$:
\[\xymatrix{
  0 \ar[r]
    & \sG \ar[r]
    & f_\ast\dZ/n \ar[r]
    & \sQ \ar[r]
    & 0 \text{.}
}\]
Replacing $X'$ by a suitable $X''$ over $X'$, we may assume that $\h^i(X,\sF)\to\h^i(X,f_*\dZ/n)=\h^i(X',\dZ/n)$ is zero for $i>0$. Consider then \small
\[\xymatrix@=.3cm{
  0 \ar[r]
    & \h^0(X,\sG) \ar[r] \ar[d]
    & \h^0(X',\dZ/n) \ar[r] \ar[d]^-\wr
    & \h^0(X,\sQ) \ar[r] \ar[d]
    & \h^1(X,\sG) \ar[r]^-0 \ar[d]
    & \h^1(X',\dZ/n) \ar[r] \ar[d]
    & \h^1(X,\sQ) \ar[r] \ar[d]
    & \cdots \\
  0 \ar[r]
    & '\h^0(X,\sG) \ar[r]
    & '\h^0(X',\dZ/n) \ar[r]
    & '\h^0(X,\sQ) \ar[r]
    & '\h^1(X,\sG) \ar[r]
    & '\h^1(X',\dZ/n) \ar[r]
    & '\h^1(X,\sQ) \ar[r]
    & \cdots
}\]
\normalsize

The isomorphism at $\h^0(X')$ gives $\h^0(X,\sG)\hookrightarrow'\h^0(X,\sG)$, and \eqref{V:eq:2-5-2} gives the same injectivity for $\sF$. Applied to $\sQ$, this gives $\h^1(X,\sG)\hookrightarrow'\h^1(X,\sG)$. Apply the same argument to $\h^1(X',\dZ/n)$. Since this group has the same order as $'\h^1(X',\dZ/n)=D(\h^1(X',\dZ/n)(1))$, the injection is an isomorphism, and consequently $\h^1(X,\sG)\iso'\h^1(X,\sG)$. The same holds for $\sF$. For $i=2$, one already knows that $\h^2(X',\dZ/n)\iso'\h^2(X',\dZ/n)$, and one argues in the same way; there are no further cases.

\section{Poincaré duality for curves}\label{V:3}

In this section, we prove the compatibility (\hyperref[I]{Arcata} \ref{I:eq:6-2-3-3}).

\subsection{}\label{V:3-1}

We use the notation of (\hyperref[I]{Arcata} \ref{I:6-2-3}): $\bar X$ is a smooth connected projective curve over an algebraically closed field $k$, $D$ is a reduced divisor on $\bar X$, and $0$ is a point of $X=\bar X\setminus D$:
\[\xymatrix{
  X \ar@{^{(}->}[r]^-j
    & \bar X
    & \ar[l]_-i D \text{,}
}\]
and $\pic_D(\bar X)=\h^1(\bar X,{}_{D}\dG_m)$, where ${}_{D}\dG_m$ is the sheaf of sections of $\dG_m$ congruent to $1$ modulo $D$. Recall that $\h^1(X,\dmu_n)=\pic_D^0(\bar X)_n$, and that $x\mapsto\sO(x)$ defines a canonical map $f:X\to\pic_D(\bar X)$. Put $f_0(x)=f(x)-f(0):X\to\pic_D^0(\bar X)$.

Again denote by $j$ the inclusion of $X\times X$ in $X\times\bar X$. The diagonal $\Delta$ of $X\times X$ is closed in $X\times\bar X$; it defines a class in $\h_\Delta^2(X\times X,\dmu_n) =\h_\Delta^2(X\times\bar X,j_!\dmu_n)$. Denote its image in $\h^2(X\times\bar X,j_!\dmu_n)$ by $c$.

The K\"unneth formula gives
\[
  \h^\bullet(X\times \bar X,j_!\dmu_n) = \h^\bullet(X,\dZ/n)\otimes \h^\bullet(\bar X,j_!\dmu_n) = \h^\bullet(X,\dZ/n)\otimes \h_c^\bullet(X,\dmu_n) \text{.}
\]
We shall compute the $(1,1)$-component of $c$:
\[
  c^{1,1} \in \h^1(X,\dZ/n)\otimes \h_c^1(X,\dmu_n) = \h^1(X,\h_c^1(X,\dmu_n)) = \h^1\left(X,\pic_D^0(\bar X)_n\right) \text{.}
\]

The exact sequence
\[\xymatrix{
  0 \ar[r]
    & \pic_D^0(\bar X)_n \ar[r]
    & \pic_D^0(\bar X) \ar[r]^n
    & \pic_D^0(\bar X) \ar[r]
    & 0
}\]
makes $\pic_D^0(\bar X)$ a $\pic_D^0(\bar X)_n$-torsor over $\pic_D^0(\bar X)$. Let $u$ denote, in $\h^1(X,\h_c^1(X,\dmu_n)) =\h^1(X,\pic_D^0(\bar X)_n)$, the class of its inverse image under $f_0$.

\begin{proposition_}\label{V:3-2}
With the preceding notation, $c^{1,1}=-u$.
\end{proposition_}

Let $e$ be the difference between the images in $\h^2(X\times\bar X,j_!\dmu_n)$ of the cohomology classes of $\Delta$ and $X\times\{0\}$. Since the latter is of K\"unneth type $(0,2)$, one has $c^{1,1}=e^{1,1}$. Consider the Leray spectral sequence for $\pr_1:X\times\bar X\to X$ and the sheaf $j_!\dmu_n$:
\[
  E_2^{pq} = \h^p(X,\eR^q {\pr_1}_\ast j_!\dmu_n) = \h^p(X,\h_c^q(X,\dmu_n)) \Rightarrow \h^{p+q}(X\times \bar X,j_!\dmu_n) \text{.}
\]
It degenerates: it is the tensor product of $\h^\bullet(X)$ with the (trivial) Leray spectral sequence for $\bar X\to\spec(k)$ and the sheaf $j_!\dmu_n$. The divisor $\Delta-X\times\{0\}$ has degree zero on every fiber, so the image of $e$ in $E^{02}$ is zero. We may therefore speak of its image $\bar e$ in $E^{1,1}$, and must prove that $\bar e=-u$.

Let $\sG$ on $X\times\bar X$ be the subsheaf of $\dG_m$ consisting of local sections whose restriction to the subscheme $X\times D$ is $1$. There is again an exact sequence
\begin{equation*}\tag{3.2.1}\label{V:eq:3-2-1}
\xymatrix{
  0 \ar[r]
    & j_!\dmu_n \ar[r]
    & \sG \ar[r]
    & \sG \ar[r]
    & 0
}
\end{equation*}
and $e$ is the image under $\partial$ of the class $e_1\in\h^1(X\times\bar X,\sG)$ of the invertible sheaf $\sO(\Delta-X\times\{0\})$, trivialized by $1$ on $X\times D$.

The direct image under $\eR\pr_{1*}$ of the distinguished triangle defined by \eqref{V:eq:3-2-1} is a distinguished triangle
\begin{equation*}\tag{3.2.2}\label{V:eq:3-2-2}
\xymatrix{
  & \ar[r]^-\partial
    & \eR \pr_{1\ast} j_!\dmu_n \ar[r]
    & \eR \pr_{1\ast}\sG \ar[r]
    & \eR \pr_{1\ast}\sG \ar[r]^-\partial
    &
}
\end{equation*}
and the diagram
\[\xymatrix{
  \h^1(X\times \bar X,\sG) \ar[r]^-\partial \ar@{=}[d]^-{(1)}
    & \h^2(X\times \bar X,j_!\dmu_n) \ar@{=}[d] \\
  \h^1(X,\eR\pr_{1\ast} \sG) \ar[r]^-\partial
    & \h^2(X,\eR\pr_{1\ast} j_!\dmu_n)
}\]
is commutative. Our problem is now to identify the image in $E^{11}=\h^1(X,\h_c^1(X,\dmu_n))$ (the image in $E^{02}$ being zero) of the image under $\partial_{\text{\eqref{V:eq:3-2-2}}}:\h^1(X,\eR\pr_{1\ast}\sG) \to \h^2(X,\eR\pr_{1\ast} j_!\dmu_n)$ of the class, again denoted $e_1$, corresponding to $e_1$ under (1).

The sheaf $\eR^1\pr_{1*}\sG$ is represented by the group scheme over $X$ obtained by pulling back the algebraic group $\pic_D(\bar X)$ over $\spec(k)$, and the image of $e_1$ in $\h^0(X,\eR^1\pr_{1*}\sG)=\hom(X,\pic_D(\bar X))$ is $f_0$.

Represent the triangle \eqref{V:eq:3-2-2} by a short exact sequence of complexes of sheaves:
\[\xymatrix{
  0 \ar[r]
    & \sA \ar[r]
    & \sB \ar[r]
    & \sC \ar[r]
    & 0 \text{.}
}\]
Let $\sB'$ be the subcomplex of
\[
  \tau_{\leqslant 1} \sB : \cdots \sB^0 \to \ker(d) \to 0 \cdots
\]
obtained by replacing ${\sB'}^1=\ker(d)$ by the inverse image of $\pic_D^0(\bar X)$ under $\ker(d)\to\underline\h^1(\sB)=\pic_D(\bar X)$ over $X$. Let $\sA'=\sA\cap\sB'$ and let $\sC'$ be the image of $\sB'$. One has $\sA'=\tau_{\leqslant1}(\sA)$, and $\sC'$ is obtained from $\sC$ in the same way as $\sB'$ from $\sB$, giving a commutative diagram of short exact sequences
\begin{equation*}\tag{3.2.3}\label{V:eq:3-2-3}
\xymatrix{
  0 \ar[r]
    & \sA \ar[r]
    & \sB \ar[r]
    & \sC \ar[r]
    & 0 \\
  0 \ar[r]
    & \tau_{\leqslant 1} \sA \ar[r] \ar[u] \ar[d]
    & \sB' \ar[r] \ar[u] \ar[d]
    & \sC' \ar[r] \ar[u] \ar[d]
    & 0 \\
  0 \ar[r]
    & \h_c^1(X,\dmu_n)[-1] \ar[r]
    & \pic_D^0(\bar X)[-1] \ar[r]
    & \pic_D^0(\bar X)[-1] \ar[r]
    & 0
}
\end{equation*}
where the last row is to be viewed as an exact sequence of complexes of sheaves concentrated in degree $1$ on $X$.
\begin{enumerate}[a)]
  \item The element $e\in\hh^2(X,\sA)$ comes from $\widetilde e\in\hh^2(X,\tau_{\leqslant1}\sA)$ (because its image in $E^{02}$ is zero). We seek the image $\bar e$ of $\widetilde e$ in $\hh^2(X,\h_c^1(X,\dmu_n)[-1]) = \h^1(X,\h_c^1(X,\dmu_n))$.
  \item The element $e_1\in\hh^1(X,\sC)$ comes from $\widetilde e_1\in\hh^1(X,\sC')$, because its image $f_0$ in $\h^0(X,\pic_D(\bar X))$ lies in $\h^0(X,\pic_D^0(\bar X))$. We may take $\widetilde e=\partial\widetilde e_1$.
  \item Thus $\bar e=\partial f_0$, where $\partial$ is defined by \eqref{V:eq:3-2-2}, and where the usual isomorphisms $\hh^2(X,\sF[-1])=\h^1(X,\sF)$ are used. These isomorphisms anticommute with $\partial$, so that $\bar e=-\partial f_0$, where $\partial$ is defined by the exact sequence
    \[\xymatrix{
      0 \ar[r]
        & \h_c^1(X,\dmu_n) \ar[r]
        & \pic_D^0(\bar X) \ar[r]
        & \pic_D^0(\bar X) \ar[r]
        & 0
    }\]
    The conclusion follows from (\hyperref[IV]{Cycles}, \ref{IV:1-1-1}).
\end{enumerate}

\subsection{}\label{V:3-3}

For every homomorphism $\varphi:\h_c^1(X,\dmu_n)\to\dZ/n$, let $\varphi(u)\in\h^1(X,\dZ/n)$ be the image of $u$ under $\varphi:\h^1(X,\h_c^1(X,\dmu_n))\to\h^1(X,\dZ/n)$. If $x\in\h_c^1(X,\dmu_n)$, the cup product $\varphi(u)\smallsmile x$ lies in $\h_c^2(X,\dmu_n)$. We shall prove:

\begin{proposition_}\label{V:3-4}
$\tr(\varphi(u)\smallsmile x) = \varphi(x)$.
\end{proposition_}

This identity follows, on applying $\varphi$, from
\begin{equation*}\tag{3.4.1}\label{V:eq:3-4-1}
  \tr(u\smallsmile x) = x
\end{equation*}
where this time $\tr$ is the map
\[\xymatrix{
  \tr:\h_c^2(X,\h_c^1(X,\dmu_n)\otimes\dmu_n) \ar[r]
    & \h_c^1(X,\dmu_n) \text{,}
}\]
so that $x\mapsto\tr(u\smallsmile x)$ is the composite map \small
\begin{equation*}\tag{3.4.2}\label{V:eq:3-4-2}
\begin{aligned}
&\xymatrix{
  \h_c^1(X,\dmu_n) \ar[r]^-{u\otimes x}
    & \h^1(X,\h_c^1(X,\dmu_n)) \otimes \h_c^1(X,\dmu_n)
    & \ar[l]_-\sim \h^1(X,\dZ/n)\otimes \h_c^1(X,\dmu_n) \otimes \h_c^1(X,\dmu_n)} \\ &\xymatrix{
    \ar[r]^-{(1)}
      & \h_c^2(X,\dmu_n) \otimes \h_c^1(X,\dmu_n) \ar[r]^-{\tr\otimes\operatorname{id}}
      & \h_c^1(X,\dmu_n)
}
\end{aligned}
\end{equation*}
\normalsize where (1) is given by $a\otimes b\otimes c\mapsto (a\smallsmile c)\otimes b$.

Express this morphism in terms of the trace morphism $\tr:\h_c^3(X\times X,\dmu_n^{\otimes2})\to\h_c^2(X,\dmu_n)$ relative to the second projection. For $x\in\h_c^3$, $\tr(x)$ is computed as follows: retain only the K\"unneth component of type $(2,1)$ and apply $\tr:\h_c^2(X,\dmu_n)\to\dZ/n$.

Let $u'$ be the image of $u$ under the map
\begin{align*}
&\xymatrix{
  \h^1(X,\h_c^1(X,\dmu_n))
    & \ar[l]_-\sim \h^1(X,\dZ/n) \otimes \h_c^1(X,\dmu_n) = \h^1(X,\dZ/n)\otimes \h^1(\bar X,j_!\dmu_n)
} \\
&\xymatrix{
  \ar[r]
    & \h^2(X\times \bar X,\dZ/n\boxtimes j_!\dmu_n) \text{.}
}
\end{align*}
By the relation between K\"unneth and cup product,
\[
  \tr(u\smallsmile x) = -\tr(u'\smallsmile \pr_1^\ast x) \text{.}
\]
The second cup product is that of (\hyperref[IV]{Cycles}, \ref{IV:1-2-4}):
\[\xymatrix{
  \h^2(X\times \bar X,\dZ/n\boxtimes j_!\dmu_n) \otimes \h^1(\bar X\times \bar X,j_!\dmu_n\boxtimes \dZ/n) \ar[r]
    & \h^3(\bar X\times \bar X,j_!\dmu_n \boxtimes j_!\dmu_n) \text{.}
}\]
The minus sign comes from interchanging two symbols of degree $1$ in (1). Since $c\smallsmile\pr_1^\ast x$ and $c^{1,1}\smallsmile\pr_1^\ast x$ have the same $(2,1)$ K\"unneth component, and since $u'=-c^{1,1}$ by \ref{V:3-2}, one has
\[
  \tr(u\smallsmile x) = \tr(c\smallsmile \pr_1^\ast x) \text{.}
\]
The cup product on the right may now be interpreted as the cup product (\hyperref[IV]{Cycles}, \ref{IV:1-2-5}) of $\cl(\Delta)\in\h_\Delta^2(X\times X,\dmu_n)$ with $\pr_1^\ast x\in\h_c^1(\Delta,\dmu_n)$. On $\Delta$ one has $\pr_1=\pr_2$, whence
\[
  \tr(c\smallsmile \pr_1^\ast x) = \tr(\cl(\Delta)\smallsmile \pr_1^\ast x) = \tr(\cl(\Delta)\smallsmile \pr_2^\ast x) = \tr(\cl(\Delta))\smallsmile x = x \text{.}
\]

\section{Compatibility with SGA 4, XVIII 3.1.10.3}\label{V:4}

In \cite[XVIII]{sga4}, overcome by discouragement, I did not verify the compatibility named in the title. Although it turns out to be no more than smoke, I feel obliged to repair the omission here. For a separated morphism of finite type $f:X\to S$, the point was to verify that the adjunction morphism $\eR f_!\eR f^!\sL\to\sL$ is compatible with every étale localization $k:V\to S$. Denoting several arrows by $k$ and $f$, as in the diagrams below, we must verify commutativity:
\[\xymatrix{
  X_V \ar[r]^-k \ar[d]^-f
    & X \ar[d]^-f \\
  V \ar[r]^-k
    & S
} \qquad
\xymatrix{
  k^\ast \eR f_! \eR f^!\sL \ar[r]^-{k^\ast(\text{adj})} \ar[d]^-\sim
    & k^\ast \sL \ar@{=}[d] \\
  \eR f_! \eR f^! k^\ast \sL \ar[r]^-{\text{adj}}
    & k^\ast \sL \text{.}
}\]
At the level of complexes, the required commutativity is written (loc.\ cit.)
\[\xymatrix{
  k^\ast f_!^\bullet f^!_\bullet \sL \ar[r] \ar[d]^-{(1)}
    & k^\ast \sL \ar@{=}[d] \\
  f_!^\bullet f^!_\bullet k^\ast \sL \ar[r]
    & k^\ast \sL \text{.}
}\]
It is checked component by component, reducing to a compatibility for $\sL$ a sheaf and for each pair of adjoint functors $(f_i^!,f_!^i)$, denoted simply by $f^!$ and $f_!$. Arrow (1) is defined from the isomorphism $(2):k^\ast f_!=f_!k^\ast$ and a morphism $(3):k^\ast f^!\to f^!k^\ast$ deduced from it. In loc.\ cit., one first introduces a morphism $(4):k_!f_!\to f_!k_!$, obtained from (2) by adjunction for $k^\ast=k^!$ and $k_!$:
\[
  k_!f_!\to k_!f_!k^\ast k_!\to k_!k^\ast f_!k_!\to f_!k_!,
\]
and then deduces (3) from (4) by adjunction. Equivalently, (3) may be deduced from (2) by adjunction for $f_!$ and $f^!$:
\[
  k^\ast\to f^!f_!k^\ast f^!\to f^!k^\ast f_!f^!\to f^!k^\ast.
\]

Writing $k$ for $k^\ast$ and omitting $\sL$, the desired commutativity is that of the outer boundary of
\[\xymatrix@=.5cm{
  k f_! f^! \ar[r] \ar@{=}[d]
    & f_! k f^! \ar[r]
    & f_! f^! f_! k f^! \ar[r]
    & f_! f^! k f_! f^! \ar@{=}[d] \\
  k f_! f^! \ar[d]
    & & & \ar[lll] f_! f^! k f_! f^! \ar[d] \\
  k
    & & & \ar[lll] f_! f^! k \text{.}
}\]
In the first row, one finds a homomorphism $k f_!f^!\to f_!f^!k f_!f^!$, defined because $k f_!f^!$ is of the form $f_!X$ (with $X=kf^!$); the adjunction morphism in the second row is a retraction of it, and commutativity of the lower square is the functoriality of this adjunction morphism.
% END INLINED FILE: sga4.5-5.tex
\clearpage
\clearpage
% BEGIN INLINED FILE: sga4.5-6.tex
\chapter{Applications of the trace formula to trigonometric sums}\label{VI}

In this exposé, I explain how the trace formula can be used to calculate or study various trigonometric sums and how, together with the Weil conjecture, it can be used to bound them.

The first two sections provide a ``user's guide'' to these tools. Section 3 presents Weil's results on one-variable sums in cohomological language. Sections 4 through 6 give a detailed study of Gauss and Jacobi sums, including Weil's older and more recent results on Hecke characters defined by Jacobi sums. In Section 7 we study a multivariable generalization of Kloosterman sums. Finally, Section 8 gives some indications of other uses that have been, or can be, made of these methods.

\section*{Notations}\label{VI:0}

\subsection{}\label{VI:0-1}

We use the notation of \hyperref[II]{Report}, Section \ref{II:1}. We shall often consider a finite extension $\dF_{q^n}\subset\dF$ of $\dF_q$. A subscript $0$ denotes an object over $\dF_q$, and a subscript $1$ an object over $\dF_{q^n}$. Replacing a subscript $0$ by a subscript $1$ (respectively, omitting the subscript) means extending scalars to $\dF_{q^n}$ (respectively, to $\dF$).

\subsection{}\label{VI:0-2}

Let $\ell$ denote a prime number different from $p$. We freely use the language of $\dQ_\ell$-sheaves, as well as that of $E_\lambda$-sheaves, where $E_\lambda$ is a finite extension of $\dQ_\ell$ (cf. \hyperref[II]{Report}, Section \ref{II:2}, especially \ref{II:2-11}).

\subsection{}\label{VI:0-3}

Let $H^\bullet$ be a graded vector space. If $T$ is an endomorphism of $H^\bullet$, put (cf. \hyperref[IV]{Cycles}, \ref{IV:eq:1-3-5})
\[
  \tr(T,H^\bullet) = \sum (-1)^i \tr(T,H^i) \text{.}
\]

\section{Principles}\label{VI:1}

\subsection{}\label{VI:1-1}

Let $X_0$ be a separated scheme of finite type over $\dF_q$, let $E_\lambda$ be a finite extension of $\dQ_\ell$, and let $\sF_0$ be an $E_\lambda$-sheaf on $X_0$. The trace formula says that
\begin{equation*}\tag{1.1.1}\label{VI:eq:1-1-1}
  \sum_{x\in X^F} \tr(F_x^\ast,\sF) = \sum (-1)^i \tr\left(F^\ast,\h_c^i(X,\sF)\right) \text{.}
\end{equation*}

With the notation of \ref{VI:0-3}, the right-hand side is simply $\tr(F^\ast,\h_c^\bullet(X,\sF))$.

This trace formula for $E_\lambda$-sheaves may either be deduced from \hyperref[II]{Report}, \ref{II:4-10}, by passage to the limit (cf. \hyperref[II]{Report}, \ref{II:4-11}--\ref{II:4-13}), or from the trace formula for $\dQ_\ell$-sheaves (\hyperref[II]{Report}, \ref{II:3-2}) by the method of (\hyperref[III]{$L$-functions mod.\ $\ell^n$}, \ref{III:4-3}).

We shall interpret various trigonometric sums as the left-hand side of \eqref{VI:eq:1-1-1}, for suitable $X_0$ and $\sF_0$.

% NOTE: in this section, following the source (but not convention in the
% LaTeX version, torsors are denoted in roman (not script) letters
\subsection{}\label{VI:1-2}

Let $A$ be a finite group. An \emph{$A$-torsor} on a scheme $X$ is a sheaf $T$ on $X$, equipped with a right action of $A$, which locally on $X$ for the étale topology is isomorphic to the constant sheaf $A$ on which $A$ acts by right translation.

If $\tau:A\to B$ is a homomorphism and $T$ is an $A$-torsor, there is, up to unique isomorphism, a unique $B$-torsor $\tau(T)$ equipped with a morphism of sheaves $\tau:T\to\tau(T)$ such that
\begin{equation*}\tag{1.2.1}\label{VI:eq:1-2-1}
  \tau(t a) = \tau(t) \tau(a)\text{.}
\end{equation*}

If $T$ is an $A$-torsor on $X$, and $\rho$ is a linear representation of $A$, $\rho:A\to\gl(V)$, with $V$ a finite-dimensional vector space over $E_\lambda$, there is, up to unique isomorphism, a unique lisse $E_\lambda$-sheaf $\sF$ of rank $\dim(V)$ equipped with a morphism of sheaves
\begin{equation*}\tag{1.2.2}\label{VI:eq:1-2-2}
  \rho:T\to \underline{\operatorname{Isom}}(V,\sF) \text{.}
\end{equation*}
such that $\rho(ta)=\rho(t)\rho(a)$. It is denoted by $\rho(T)$. For $\tau:A\to B$ and $\rho$ a representation of $B$, there is a canonical isomorphism $\rho(\tau(T))=(\rho\tau)(T)$.

In this section (except for the appendix), we consider only the case in which $A$ is commutative and $V=E_\lambda$: thus $\rho$ is a character $A\to E_\lambda^\times$, and $\rho(T)$ is a lisse $E_\lambda$-sheaf of rank one. The morphism \eqref{VI:eq:1-2-2} is interpreted as an everywhere nonzero morphism
\begin{equation*}\tag{1.2.3}\label{VI:eq:1-2-3}
  \rho:T\to \rho(T)
\end{equation*}
such that $\rho(ta)=\rho(a)\rho(t)$.

\subsection{}\label{VI:1-3}

Let $S$ be a scheme, $G$ a commutative group scheme over $S$, and $G'$ an extension of commutative group schemes over $S$:
\[\xymatrix{
  0 \ar[r]
    & A \ar[r]
    & G' \ar[r]^-\pi
    & G \ar[r]
    & 0 \text{.}
}\]
The sheaf $T$ of local sections of $\pi$ is then an $A$-torsor on $G$.

If $X$ is a scheme over $S$, and $f,g$ are two $S$-morphisms from $X$ to $G$, the fact that $T$ is defined by an extension gives a canonical isomorphism
\begin{equation*}\tag{1.3.1}\label{VI:eq:1-3-1}
  (f+g)^\ast T = f^\ast T + g^\ast T \text{.}
\end{equation*}
On the left, $f+g$ is the sum of $f$ and $g$ for the group law of $G$; on the right, $+$ denotes the sum of torsors.

If $f:X\to G$ factors through $G'$, then $f^\ast T$ is trivial (and a factorization gives a trivialization). Up to nonunique isomorphism, the $A$-torsor $f^\ast T$ depends only on the image of $f$ in $\hom_S(X,G)/\pi\hom_S(X,G')$: it is given by the morphism $\partial$ in the exact sequence
\[\xymatrix{
  \hom(X,G') \ar[r]
    & \hom(X,G) \ar[r]^-\partial
    & \h^1(X,A) \text{.}
}\]

We shall apply these constructions to the Kummer exact sequences $0\to\dmu_n\to\dG_m\to\dG_m\to0$, and to Lang torsors.

\subsection{The Lang torsor}\label{VI:1-4}

Let $G_0$ be a connected commutative algebraic group over $\dF_q$ (for the noncommutative case, see the appendix). The group law is written \emph{multiplicatively}. The Lang isogeny
\[
\fL:G_0 \to G_0: x\mapsto F x\cdot x^{-1}
\]
is étale; its image, being an open subgroup of $G_0$, can only be $G_0$ itself, and its kernel is the finite group $G_0(\dF_q)$. The \emph{Lang torsor} $L$ is the $G_0(\dF_q)$-torsor on $G_0$ defined by the exact sequence
\begin{equation*}\tag{1.4.1}\label{VI:eq:1-4-1}
\xymatrix{
  0 \ar[r]
    & G_0(\dF_q) \ar[r]
    & G_0 \ar[r]^-\fL
    & G_0 \ar[r]
    & 0 \text{.}
}
\end{equation*}

\paragraph{Notation}
We again denote by $L$ and $\fL$ (or by $L_{(q)}$ and $\fL_{(q)}$) the objects deduced from $L$ and $\fL$ by extension of the base field. When necessary, a subscript $0$ or $1$ will be specified.

\paragraph{Examples}
For $G_0=\dG_a$ or $\dG_m$, the sequences \eqref{VI:eq:1-4-1} become
\begin{align*}\tag{1.4.2}\label{VI:eq:1-4-2}
&\xymatrix{
  0 \ar[r]
    & \dF_q \ar[r]
    & \dG_a \ar[r]^-{x^q-x}
    & \dG_a \ar[r]
    & 0
} \\
\tag{1.4.3}\label{VI:eq:1-4-3}
&\xymatrix{
  0 \ar[r]
    & \dmu_{q-1} \ar[r]
    & \dG_m \ar[r]^-{x^{q-1}}
    & \dG_m \ar[r]
    & 0
}
\end{align*}

\subsection{}\label{VI:1-5}

Let us calculate the endomorphism $F^\ast$ of the fiber of the Lang torsor at $\gamma\in G^F=G_0(\dF_q)$. If $g\in\fL^{-1}(\gamma)\subset G$, then $Fg=Fg\cdot g^{-1}\cdot g=\gamma g$. Hence (\hyperref[II]{Report}, \ref{II:1-2})
\begin{equation*}\tag{1.5.1}\label{VI:eq:1-5-1}
  \text{on $L_0(G_0)_\gamma\simeq\fL^{-1}(\gamma)$,
  $F^\ast$ is $g\mapsto g\gamma^{-1}$.}
\end{equation*}

\subsection{}\label{VI:1-6}

Over $\dF_{q^n}$, the identity $F_{(q^n)}=F_{(q)}^n$ between endomorphisms of $G_1$ implies $\fL_{(q^n)}=\fL_{(q)}\circ\prod_{i=0}^{n-1}F_{(q)}^i$. On $G_0(\dF_{q^n})$, $F_{(q)}^i$ acts as the element $x\mapsto x^{q^i}$ of $\gal(\dF_{q^n}/\dF_q)$. Hence on $G_0(\dF_{q^n})$, $\prod F_{(q)}^i$ is the composite of the norm $G_0(\dF_{q^n})\to G_0(\dF_q)$ and the inclusion $G_0(\dF_q)\subset G_0(\dF_{q^n})$: the diagram
\begin{equation*}\tag{1.6.1}\label{VI:eq:1-6-1}
\xymatrix{
  0 \ar[r]
    & G_0(\dF_{q^n}) \ar[r] \ar[d]^-N
    & G_1 \ar[r]^-{\fL_{(q^n)}} \ar[d]^-{\prod_{i=0}^{n-1} F_{(q)}^i}
    & G_1 \ar[r] \ar@{=}[d]
    & 0 \\
  0 \ar[r]
    & G_0(\dF_q) \ar[r]
    & G_1 \ar[r]^-{\fL_{(q)}}
    & G_1 \ar[r]
    & 0
}
\end{equation*}
is commutative. In the language of torsors, this gives a canonical isomorphism of $G_0(\dF_q)$-torsors on $G_1$:
\begin{equation*}\tag{1.6.2}\label{VI:eq:1-6-2}
  N L_{(q^n)} = L_{(q)} \text{.}
\end{equation*}

\begin{definition_}\label{VI:1-7}
Let $f_0:X_0\to G_0$ be a morphism and $\chi:G_0(\dF_q)\to E_\lambda^\times$ a character. Put $\sF(\chi,f_0) = \chi^{-1}(f_0^\ast (L_0(G_0))) = f_0^\ast\chi^{-1}(L_0(G_0))$.
\end{definition_}

The following functoriality properties hold:
\begin{enumerate}[a)]
  \item $\sF(\chi,f_0)$ is multiplicative in both $\chi$ and $f_0$:
    \begin{align*}\tag{1.7.1}\label{VI:eq:1-7-1}
      \sF(\chi,f_0'\cdot f_0'') &= \sF(\chi,f_0')\otimes \sF(\chi,f_0'') \text{,} \\
      \tag{1.7.2}\label{VI:eq:1-7-2}
      \sF(\chi'\chi'',f_0) &= \sF(\chi',f_0)\otimes \sF(\chi'',f_0) \text{.}
    \end{align*}
    This follows from the definition of $\chi(T)$ for a torsor $T$, together with \eqref{VI:eq:1-3-1} for \eqref{VI:eq:1-7-1}.
  \item For a morphism $g_0:Y_0\to X_0$, one has
    \begin{equation*}\tag{1.7.3}\label{VI:eq:1-7-3}
      \sF(\chi,f_0\circ g_0) = g_0^\ast \sF(\chi,f_0) \text{.}
    \end{equation*}
    In particular, $\sF(\chi,f_0)=f_0^\ast\sF(\chi)$, where $\sF(\chi)$ denotes $\sF(\chi,\operatorname{id}_{G_0})$.
  \item For a morphism $u_0:G_0\to H_0$ and a character $\chi:H_0(\dF_q)\to E_\lambda^\times$, one has
    \begin{equation*}\tag{1.7.4}\label{VI:eq:1-7-4}
      \sF(\chi,u_0 f_0) = \sF(\chi u_0,f_0) \text{.}
    \end{equation*}
  \item For $G_0=\prod_{i\in I}G_0^i$, with $\chi$ having components $\chi_i$ ($i\in I$) and $f_0$ having components $f_0^i$, it follows from \eqref{VI:eq:1-7-1}--\eqref{VI:eq:1-7-4} that
    \begin{equation*}\tag{1.7.5}\label{VI:eq:1-7-5}
      \sF(\chi,f_0) = \bigotimes_{i\in I} \sF(\chi_i,f_0^i) \text{.}
    \end{equation*}
\end{enumerate}

Note the $\chi^{-1}$ in Definition \ref{VI:1-7}. It ensures that, for $x\in X^F$, on $\sF(\chi,f_0)_x$ one has
\begin{equation*}\tag{1.7.6}\label{VI:eq:1-7-6}
  F_x^\ast = \chi f_0(x)
\end{equation*}
(use the fact that the stalk at $x$ of the morphism \eqref{VI:eq:1-2-2}, for $\rho=\chi^{-1}$, commutes with $F_x^\ast$, and use \eqref{VI:eq:1-5-1}).

If $\sF(\chi,f_0)_1$ is the sheaf deduced from $\sF(\chi,f_0)$ by extending the base field to $\dF_{q^n}$, then \eqref{VI:eq:1-6-2} gives
\begin{equation*}\tag{1.7.7}\label{VI:eq:1-7-7}
  \sF(\chi,f_0)_1 = \sF(\chi\circ N,f_1) \text{.}
\end{equation*}

\subsection{Abuse of notation}\label{VI:1-8}

(i) If $\Xi$ denotes the composite map $\chi f_0:X_0(\dF_q)\to E_\lambda^\times$, we shall sometimes write $\sF(\Xi)$ in place of $\sF(\chi,f_0)$. In view of \eqref{VI:eq:1-7-1}--\eqref{VI:eq:1-7-5}, there is little risk of ambiguity. For example:
\begin{enumerate}[a)]
  \item we write $\sF(\chi)$ for $\sF(\chi,\operatorname{id}_{G_0})$ (notation already used in \ref{VI:1-7}b);
  \item we write $\sF(\chi f_0)$ for $\sF(\chi,f_0)$;
  \item with the notation of \ref{VI:1-7}d), we write $\sF(\prod\chi_i f_0^i)$ for $\sF(\chi,f_0)$.
\end{enumerate}

With this notation, \eqref{VI:eq:1-7-4} says that $\sF(\chi u_0f_0)$ is unambiguous; \eqref{VI:eq:1-7-1}, \eqref{VI:eq:1-7-2}, and \eqref{VI:eq:1-7-5} express multiplicativity of $\sF(\Xi)$ in $\Xi$; and \eqref{VI:eq:1-7-6} becomes $F_x^\ast=\Xi(x)$ on $\sF(\Xi)_x$.

(ii) If $X$ is a scheme over an extension $k$ of $\dF_q$, and $f$ is a morphism from $X$ to $G_0\otimes_{\dF_q}k$, we again denote by $\sF(\chi,f)$, $\sF(\chi f)$, or $\sF(\chi)$ (when $f$ is an inclusion) the inverse image under $f$ of the $E_\lambda$-sheaf deduced from $\chi^{-1}(L_0(G))$ by extending scalars from $\dF_q$ to $k$. If $f_0$ is the composite $X\to G_0\otimes_{\dF_q}k\to G_0$, this is again the $\sF(\chi,f_0)$ of \ref{VI:1-7}. If $k$ is a finite extension of $\dF_q$, then
\[
  \sF(\chi f)=\sF(\chi\circ N_{k/\dF_q},f) \text{.}
\]

Applying \eqref{VI:eq:1-1-1} to \eqref{VI:eq:1-7-6} and \eqref{VI:eq:1-7-7}, we obtain:

% NOTE: originally a 'scholium'
\begin{theorem_}\label{VI:1-9}
Let $S_0$ be a separated scheme of finite type over $\dF_q$, let $G_0$ be a connected commutative algebraic group over $\dF_q$, let $f_0:S_0\to G_0$ be a morphism, and let $\chi:G_0(\dF_q)\to E_\lambda^\times$ be a character. Then
\begin{equation*}\tag{1.9.1}\label{VI:eq:1-9-1}
  \sum_{s\in S_0(\dF_q)} \chi f_0(s) = \tr(F^\ast,\h_c^\bullet(S,\sF(\chi,f_0)))
\end{equation*}
and, for every integer $n\geqslant1$,
\begin{equation*}\tag{1.9.2}\label{VI:eq:1-9-2}
  \sum_{s\in S_0(\dF_{q^n})} \chi\circ N_{\dF_{q^n}/\dF_q} f_0(s) = \tr({F^\ast}^n, \h_c^\bullet\left(S,\sF(\chi,f_0))\right) \text{.}
\end{equation*}
\end{theorem_}

\addtocounter{subsubsection}{2}
\subsubsection{Remark}\label{VI:1-9-3}

Take $G_0$ to be a product. With the notation of \ref{VI:1-7}d) and \ref{VI:1-8}c), formula \eqref{VI:eq:1-9-2} becomes
\[
  \sum_{s\in S_0(\dF_{q^n})} \prod_i \chi_i\circ N_{\dF_{q^n}/\dF_q} (f_0^i(s)) = \tr\Big({F^\ast}^n, \h_c^\bullet\Big(S,\sF\Big(\prod_i \chi_i f_0^i\Big)\Big)\Big) \text{.}
\]

\subsubsection{Remark}\label{VI:1-9-4}

Take $G_0=\{e\}$. Formula \eqref{VI:eq:1-9-1} becomes
\[
  \left|S_0(\dF_q)\right| = \sum_{s\in S_0(\dF_q)} 1 = \tr(F^\ast,\h_c^\bullet(S,\dQ_\ell)) \text{.}
\]

It is rare that the right-hand side of \eqref{VI:eq:1-9-1} can be calculated explicitly. Here is an amusing example with $G_0=\{e\}$.

\subsection{Example}\label{VI:1-10}

A nonsingular projective quadric of odd dimension over \texorpdfstring{$\dF_q$}{Fq} has the same number of rational points as projective space of the same dimension over \texorpdfstring{$\dF_q$}{Fq}.

If, over $\dC$, $X$ is a nonsingular quadric hypersurface in projective space $\dP^{2N}$ and $Y$ is a hyperplane, then for $i\leqslant2\dim(X)=2\dim(Y)$ the inclusions $X\hookrightarrow\dP^{2N}\hookleftarrow Y$ induce isomorphisms in ordinary cohomology:
\[\xymatrix{
  \h^i(X,\dQ)
    & \ar[l]_-\sim \h^i(\dP^{2 N},\dQ) \ar[r]^-\sim
    & \h^i(Y,\dQ) \text{.}
}\]

By specialization, it follows that if $X_0'$ is a nonsingular quadric hypersurface in projective space $\dP_0^{2N}$ over $\dF_q$, and $Y_0'$ is a hyperplane, then the inclusions $X_0'\hookrightarrow\dP_0^{2N}\hookleftarrow Y_0'$ induce isomorphisms
\[\xymatrix{
  \h^i(X',\dQ_\ell)
    & \ar[l]_-\sim \h^i(\dP^{2 N},\dQ_\ell) \ar[r]^-\sim
    & \h^i(Y',\dQ_\ell) \text{.}
}\]
These isomorphisms commute with $F^\ast$, and one applies the trace formula.

\subsection{}\label{VI:1-11}

The trace formula can also be used to understand the classical formulas giving the number of rational points of linear groups over finite fields, and of certain homogeneous spaces (see Section 8).

\subsection{}\label{VI:1-12}

There is a dictionary that translates various kinds of classical manipulations of trigonometric sums into cohomological terms. This dictionary will be given in Section 2. Since the cohomological statement is more ``geometric,'' it can sometimes apply in situations where the classical argument applies only after extending the finite base field. Here is an example, a special case of a theorem of Kazhdan.

\begin{theorem_}[Kazhdan]\label{VI:1-13}
Let $X_0$ be a scheme over $\dF_q$. Suppose there is an action $\rho$ of $\dG_a$ on $X$, and a morphism $f:X\to Y$ of schemes over $\dF$, making $X$ a $\dG_a$-torsor over $Y$. Then the number of rational points of $X_0$ is divisible by $q$.
\end{theorem_}

First suppose that $\rho$, $Y$, and $f$ are defined over $\dF_q$.

\paragraph{Classical argument}
Every fiber of $f:X_0(\dF_q)\to Y_0(\dF_q)$ has $q$ elements: $|X_0(\dF_q)|=q\cdot|Y_0(\dF_q)|$, proving the divisibility.

\paragraph{Cohomological translation}
The higher direct image sheaves $\eR^if_!\dQ_\ell$ are
\[
  \eR^i f_!\dQ_\ell = \begin{cases}
                       \dQ_\ell(-1) & \text{if $i=2$} \\
                       0            & \text{if $i\ne2$}
                     \end{cases}
\]
The Leray spectral sequence for $f$ therefore degenerates to an isomorphism
\[
  \h_c^i(X,\dQ_\ell) = \h_c^{i-2}(Y,\dQ_\ell)(-1) \text{.}
\]
To understand the effect of a Tate twist on the eigenvalues of Frobenius, it is most convenient to use the Galois viewpoint (\hyperref[II]{Report}, \ref{II:1-8}) and note that the Frobenius substitution $\varphi$ acts on $\dQ_\ell(1)$ by multiplication by $q$. The eigenvalues of $F^\ast$ on $\h_c^i(X,\dQ_\ell)$ are therefore $q$ times the eigenvalues of $F^\ast$ on $\h_c^{i-2}(Y,\dQ_\ell)$. The latter are known to be algebraic integers \cite[XXI 5.2.2]{sga7}; hence

\begin{equation*}\tag{$*$}\label{VI:eq:*}
  \substack{\displaystyle\text{the eigenvalues of $F^\ast$ acting on
  $\h_c^i(X,\dQ_\ell)$ are}\\
  \displaystyle\text{algebraic integers divisible by $q$.}}
\end{equation*}

\paragraph{Descent}
Return to the hypotheses of the theorem. For a suitable $n$, we may assume that $\rho$, $Y$, and $f$ are defined over $\dF_{q^n}$. Since the Frobenius morphism relative to $\dF_{q^n}$ is the $n$th power of the one relative to $\dF_q$, statement \eqref{VI:eq:*} for the scheme over $\dF_{q^n}$ obtained from $X_0$ by extension of scalars gives:
\begin{equation*}\tag{$**$}\label{VI:eq:**}
  \substack{\displaystyle\text{the $n$th powers of the eigenvalues of
  $F^\ast$ acting on}\\
  \displaystyle\text{$\h_c^i(X,\dQ_\ell)$ are algebraic integers divisible
  by $q^n$.}}
\end{equation*}

This assertion implies \eqref{VI:eq:*}. The number of rational points of $X_0$, namely $\sum(-1)^i\tr(F^\ast,\h_c^i(X,\dQ_\ell))$, is therefore an algebraic integer divisible by $q$. Since it is a rational integer, it is divisible by $q$ in $\dZ$.

\subsection{}\label{VI:1-14}

Formula \ref{VI:1-9-3} controls the dependence on $n$ of the trigonometric sum on the left. It implies identities between a trigonometric sum and those deduced from it by ``extension of scalars.'' The most famous of these is the Hasse--Davenport identity. Let $\chi:\dF_q^\times\to\dC^\times$ and $\psi:\dF_q\to\dC^\times$ be characters, with $\psi$ nontrivial, and define the Gauss sum $\tau(\chi,\psi)$ by
\begin{equation*}\tag{1.14.1}\label{VI:eq:1-14-1}
  \tau(\chi,\psi) = - \sum_{x\in \dF_q^\times} \psi(x) \chi^{-1}(x) \text{.}
\end{equation*}

\begin{theorem_}[Hasse--Davenport]\label{VI:1-15}
One has
\begin{equation*}\tag{1.15.1}\label{VI:eq:1-15-1}
  \tau\left(\chi\circ N_{\dF_{q^n}/\dF_q},\psi\circ \tr_{\dF_{q^n}/\dF_q}\right) = \tau(\chi,\psi)^n \text{.}
\end{equation*}
\end{theorem_}

Let $E\subset\dC$ be a number field containing the values of $\chi$ and $\psi$, and let $E_\lambda$ be the completion of $E$ at a place of residual characteristic $\ell\ne p$. Identity \eqref{VI:eq:1-15-1} is an identity in $E$. It is equivalent to prove it in $\dC$ or in $E_\lambda$. We prove it in $E_\lambda$, regarding $\chi$ and $\psi$ as taking values in $E_\lambda^\times$.

Apply \ref{VI:1-9-3} with $X_0=\dG_m$ over $\dF_q$ and $G_0=\dG_m\times\dG_a$. The groups $\h_c^i(\dG_m,\sF(\chi^{-1}\psi))$ vanish for $i\ne1$, and $\h_c^1$ has dimension $1$ (\ref{VI:4-2}). By \ref{VI:1-9-3}, $\tau(\chi\circ N_{\dF_{q^n}/\dF_q}, \psi\circ\tr_{\dF_{q^n}/\dF_q})$ is the unique eigenvalue of $(F^\ast)^n$ on $\h_c^1$, and \eqref{VI:eq:1-15-1} follows.

Analogous examples are given by Weil \cite[App.\ V]{we74}.

\subsection{}\label{VI:1-16}

An $E_\lambda$-sheaf $\sF_0$ on $X_0$ is said to be \emph{pointwise of weight $n$} if, for every closed point $x\in|X_0|$, the eigenvalues of $F_x^\ast$ (\hyperref[II]{Report}, \ref{II:1-2}) are algebraic numbers all of whose complex conjugates have absolute value $q_x^{n/2}$, where $q_x$ is the cardinality of $k(x)$. The following theorem will be proved in \cite{de80}.

\begin{theorem_}\label{VI:1-17}
If $\sF_0$ is pointwise of weight $n$, then for every eigenvalue $\alpha$ of the endomorphism $F^\ast$ of $\h_c^i(X,\sF)$, there is an integer $m\leqslant n+i$ such that all complex conjugates of $\alpha$ have absolute value $q^{m/2}$.
\end{theorem_}

The sheaves considered in \ref{VI:1-9} have weight $0$. The theorem therefore gives, for the sum \ref{VI:eq:1-9-1} (or rather, for all its complex conjugates), the bound
\begin{equation*}\tag{1.17.1}\label{VI:eq:1-17-1}
  \left|\sum_{s\in S_0(\dF_q)} \prod_i \chi_i(f^i(s)) \right| \leqslant \sum_i \dim \h_c^i\left(S,\sF\left(\prod \chi_i f^i\right)\right)\cdot q^{i/2} \text{.}
\end{equation*}

\subsection{Remarks}\label{VI:1-18}

\begin{enumerate}[a)]
  \item If $\dim S=n$, then for every $E_\lambda$-sheaf $\sF$ on $S$,
    \[
      \h_c^i(S,\sF)=0 \qquad \text{for $i\notin[0,2n]$.}
    \]
  \item The group $\h_c^0(S,\sF)$ is the group of global sections of $\sF$ on $S$ whose support is proper. If $S$ is connected and nonproper and $\sF$ is lisse, or if $S$ is connected and $\sF$ is a nonconstant lisse sheaf of rank one, this group is zero.
  \item If $S$ is smooth and purely of dimension $n$, and $\sF$ is lisse, Poincaré duality says that the vector spaces $\h_c^i(S,\sF)$ and $\h^{2n-i}(S,\sF^\vee(n))$ are dual to one another (for the effect of a Tate twist on Frobenius eigenvalues, cf. the second part of the proof of \ref{VI:1-13}).
  \item If $\dim S=n$, Poincaré duality permits the following computation of $\h_c^{2n}(S,\sF)$. Choose an open $U$ of $S_\text{red}$ whose complement has dimension $<n$, such that $U$ is smooth and purely of dimension $n$ and $\sF$ is lisse on $U$. Then $\h_c^{2n}(U,\sF)\iso\h_c^{2n}(S,\sF)$, since in the long exact cohomology sequence $\h_c^i(S\setminus U,\sF)=0$ for $i=2n,2n-1$. By Poincaré duality,
    \[
      \h_c^{2n}(S,\sF)
      \iso\left(\h^0(U,\sF^\vee)(n)\right)^\vee \text{.}
    \]

    Suppose that $U$ is connected, and let $u$ be a geometric point of $U$. The ``local system'' $\sF$ corresponds to a representation of $\pi_1(U,u)$ on $\sF_u$, and $\h^0(U,\sF^\vee)$ is $(\sF_u^\vee)^{\pi_1(U,u)} =((\sF_u)_{\pi_1(U,u)})^\vee$, the dual of the coinvariants. Hence
    \[
      \h_c^{2n}(S,\sF)\simeq(\sF_u)_{\pi_1(U,u)}(-n) \text{.}
    \]
  \item These remarks often permit the calculation of $b_0$ and $b_{2n}$. Calculating the other $b_i$ can be difficult. If $b_{2n}=0$ and $b_{2n-1}=0$, the bound \eqref{VI:eq:1-17-1} is of Lang--Weil type, with, for large $q$, a gain of $\sqrt q$ over the trivial $O(q^n)$ bound. One should nevertheless note that \eqref{VI:eq:1-17-1} does not formally imply the trivial bound
    \[
      \left|\sum_{s\in S_0(\dF_q)} \prod_i \chi_i(f^i(s))\right| \leqslant \left| S_0(\dF_q)\right| \text{,}
    \]
    so it can be useful to take the latter into account (cf. the proof of \ref{VI:3-8}).
\end{enumerate}

I explain below a method for proving that $b_i=0$ for $i\ne n$. When it applies, it gives an $O(q^{n/2})$ estimate. The implicit constant is $b_n=(-1)^n\chi(S,\sF(\prod\chi_if^i))$, and it can be difficult to calculate.

\begin{proposition_}\label{VI:1-19}
Let $\bar X$ be a proper scheme over an algebraically closed field $k$, let $j:X\hookrightarrow\bar X$ be an open subscheme of $\bar X$, and let $\sF$ be an $E_\lambda$-sheaf on $X$. Suppose that
\begin{enumerate}[\indent a)]
  \item $(j_\ast\sF)_x=0$ for $x\in\bar X\setminus X$, i.e. $j_!\sF\iso j_*\sF$;
  \item $\eR^ij_\ast\sF=0$ for $i>0$.
\end{enumerate}
Then the maps
\[
  \h_c^i(X,\sF) \to \h^i(X,\sF)
\]
are isomorphisms.
\end{proposition_}

Hypotheses a) and b) say that $j_!\sF\iso\eR j_*\sF$, whence
\[
  \h_c^i(X,\sF) = \h^i(\bar X,j_!\sF) \iso \hh^i(\bar X,\eR j_\ast \sF) = \h^i(X,\sF) \text{.}
\]
In other words, the Leray spectral sequence for $j$,
\[
  E_2^{pq} = \h^p(\bar X,\eR^q j_\ast \sF) \Rightarrow \h^{p+q}(X,\sF)
\]
degenerates, by b), to isomorphisms
\[
  \h^i(\bar X,j_\ast\sF) \iso \h^i(X,\sF) \text{,}
\]
while, by a),
\[
  \h_c^i(X,\sF) = \h^i(\bar X,j_!\sF) \iso \h^i(\bar X,j_\ast\sF) \text{.}
\]

\subsubsection{Example}\label{VI:1-19-1}

If $X$ is a curve, or if $\bar X$ is smooth, $X$ is the complement of a normal-crossings divisor, and $\sF$ is tamely ramified, then hypothesis a) of \ref{VI:1-19} (no invariants under local monodromy) implies hypothesis b).

\begin{proposition_}\label{VI:1-20}
Let $X_0$ be a smooth affine scheme purely of dimension $n$ over $\dF_q$, and let $\sF_0$ be a lisse $E_\lambda$-sheaf pointwise of weight $m$ on $X_0$. Suppose that the maps $\h_c^i(X,\sF)\to\h^i(X,\sF)$ are isomorphisms (cf. \ref{VI:1-19}). Then $\h_c^i(X,\sF)=0$ for $i\ne n$, and the complex conjugates of the eigenvalues of $F^\ast$ on $\h_c^n(X,\sF)$ have absolute value $q^{(m+n)/2}$.
\end{proposition_}

Poincaré duality pairs $\h_c^i(X,\sF)$ (respectively, $\h_c^i(X,\sF^\vee)$) with $\h^{2n-i}(X,\sF^\vee(n))$ (respectively, $\h^{2n-i}(X,\sF(n))$). By \cite[XIV 3.2]{sga4}, the cohomology groups without support vanish for $i>n$. By duality, the $\h_c^i$ vanish for $i<n$; since $\h_c^i(X,\sF)\iso\h^i(X,\sF)$, these groups vanish for $i\ne n$. Apply \ref{VI:1-17} to $\sF_0$ and to $\sF_0^\vee(n)$, which has weight $-m-2n$. One finds that the complex conjugates $\alpha$ of the eigenvalues of $F^\ast$ on $\h_c^n(X,\sF)=\h_c^n(X,\sF^\vee(n))^\vee$ satisfy
\begin{align*}
  |\alpha| &\leqslant q^{(m+n)/2} && \text{et} \\
  |\alpha^{-1}| &\leqslant q^{((-m-2n)+n)/2} = q^{-(m+n)/2} \text{,}
\end{align*}
which proves the assertion.

\subsection{Remark}\label{VI:1-21}

The last argument shows that, if $X_0$ is a separated smooth scheme over $\dF_q$, $\sF_0$ is a lisse $E_\lambda$-sheaf pointwise of weight $m$ on $X_0$, and $\h_c^i(X,\sF)\hookrightarrow\h^i(X,\sF)$, then the complex conjugates of the eigenvalues of $F^\ast$ on $\h_c^i(X,\sF)$ have absolute value $q^{(m+i)/2}$.

\subsection*{Appendix. The Lang isogeny in the noncommutative case}

\subsection{}\label{VI:1-22}

Let $G_0$ be a connected algebraic group over $\dF_q$, and let $\fL$ be the morphism $x\mapsto Fx\cdot x^{-1}$ from $G_0$ to itself. It is a morphism of homogeneous $G_0$-spaces if $G_0$ acts on the source by left translation, $x\mapsto(y\mapsto xy)$, and on the target by $x\mapsto(y\mapsto Fx\cdot y\cdot x^{-1})$. One has $\fL(xy)=\fL(x)$ for $y\in G_0(\dF_q)$, and $\fL$ induces an isomorphism $G_0/G_0(\dF_q)\iso G_0$.

As in the commutative case, $\fL$ therefore makes $G_0$ into a $G_0(\dF_q)$-torsor over $G_0$, the \emph{Lang torsor} $L$. If $\rho:G_0(\dF_q)\to\gl(n,E_\lambda)$ is a linear representation of $G_0(\dF_q)$, we denote the $E_\lambda$-sheaf $\rho(L)$ by $\sF(\rho)$ (\ref{VI:1-2}).

\subsection{}\label{VI:1-23}

Let $\gamma\in G^F=G_0(\dF_q)$, and calculate the endomorphism $F^\ast$ of $\sF(\rho)_\gamma$. Choose $g\in G$ such that $\fL(g)=\gamma$, giving a frame $\rho(g)\in\operatorname{Isom}(E_\lambda^n,\sF(\rho)_\gamma)$. For $e\in E_\lambda^n$, one has
\[
  F(\rho(g)(e)) = \rho(F g)(e) = \rho(\gamma g)(e) = \rho(g\cdot g^{-1}\gamma g)(e) \text{.}
\]
We shall see that $g^{-1}\gamma g\in G_0(\dF_q)$, whence
\[
  F(\rho(g)(e)) = \rho(g) \rho(g^{-1}\gamma g)(e)
\]
et
\begin{equation*}\tag{1.23.1}\label{VI:eq:1-23-1}
  \rho(g)^{-1} (F_\gamma^\ast)^{-1} \rho(g) = \rho(g^{-1} \gamma g) \text{:}
\end{equation*}
Thus $\rho(g)$ identifies the inverse of $F^\ast$ at $\gamma$ with the automorphism $\rho(g^{-1}\gamma g)$ of $E_\lambda^n$. It remains to understand $g^{-1}\gamma g$.

\begin{lemma_}\label{VI:1-24}
\begin{enumerate}[(i)]
  \item If $Fg\cdot g^{-1}\in G_0(\dF_q)$, then $g^{-1} (F g\cdot g^{-1}) g = g^{-1} F g\in G_0(\dF_q)$.
  \item Let $\gamma\in G_0(\dF_q)$, and let $g$ satisfy $Fg\cdot g^{-1}=\gamma$. The conjugacy class of $\gamma'=g^{-1}Fg$ in $G_0(\dF_q)$ depends only on that of $\gamma$.
  \item The map induced by $\gamma\mapsto\gamma'$ on the set $G_0(\dF_q)^\natural$ of conjugacy classes of $G_0(\dF_q)$ in itself is bijective. Its inverse is the corresponding map for the opposite group.
\end{enumerate}
\end{lemma_}

(i) If $\gamma=Fg\cdot g^{-1}\in G_0(\dF_q)$, the identities $Fg=\gamma g$ and $F\gamma=\gamma$ give $F(g^{-1}Fg)=(Fg)^{-1}FFg =g^{-1}\gamma^{-1}F(\gamma g) =g^{-1}\gamma^{-1}F\gamma Fg=g^{-1}Fg$, so $g^{-1}Fg\in G_0(\dF_q)$.

(ii) For $\alpha,\gamma\in G_0(\dF_q)$, the elements $g$ such that $Fg\cdot g^{-1}=\gamma$ form a right coset under $G_0(\dF_q)$; and if $Fg\cdot g^{-1}=\gamma$, then $F(\alpha g\alpha^{-1})(\alpha g\alpha^{-1})^{-1} =\alpha\gamma\alpha^{-1}$. As $\gamma$ runs through a conjugacy class in $G_0(\dF_q)$, the corresponding $g$ therefore run through a double coset under $G_0(\dF_q)$, and $\gamma'=g^{-1}Fg$ runs through a conjugacy class.

(iii) Finally, the inverse of $Fg\cdot g^{-1}\mapsto g^{-1}Fg$ is $g^{-1}Fg\mapsto Fg\cdot g^{-1}$.

\subsection{}\label{VI:1-25}

Formula \eqref{VI:eq:1-23-1} may be read briefly as follows: $F^\ast$ at $\gamma$ is $\rho(\gamma')^{-1}$. If $\gamma$ belongs to the identity component of its centralizer, one may take $g$ in $Z(\gamma)$, and then $\gamma=\gamma'$. If this condition is not satisfied, $\gamma$ and $\gamma'$ need not be conjugate.

\subsection{An example}\label{VI:1-26}

For $G_0=\operatorname{SL}(2)$, with $q$ odd, $\alpha\in\dF_q^\times\setminus{\dF_q^\times}^2$, and $\bar\gamma$ the conjugacy class of $-1(1+N)$ with $N^2=0$, the class $\bar\gamma'$ is the conjugacy class of $-1\cdot(1+\alpha N)$; hence $\bar\gamma'\ne\bar\gamma$ if $N\ne0$.

\section{Dictionary}\label{VI:2}

An elementary manipulation of trigonometric sums can often be viewed, via the trace formula, as the reflection of a cohomological statement. In this section I list such statements and their reflections; they carry the same number, with a $\ast$ appended to the cohomological statement.

\subsection{}\label{VI:2-1}

``A sum of algebraic integers is an algebraic integer'' has the following cohomological analogue, proved in \cite[XXI 5.2.2]{sga7} (cf. the use of \ref{VI:2-1}$^\ast$ in \ref{VI:1-13}).

\begin{theorem*}[2.1*]\label{VI:2-1*}
Let $X_0$ be a separated scheme of finite type over $\dF_q$, and let $\sF_0$ be an $E_\lambda$-sheaf on $X_0$. If, for every $x\in|X_0|$, the eigenvalues of $F_x^\ast$ on $\sF_0$ are algebraic integers, then the eigenvalues of $F^\ast$ on $\h_c^i(X,\sF)$ are algebraic integers.
\end{theorem*}

\subsection{}\label{VI:2-2}

Other integrality results are proved in \cite[XXI, 5.2.2 and 5.4]{sga7}: if $\dim(X_0)\leqslant n$, and $\alpha$ is an eigenvalue of $F^\ast$ on $\h_c^i(X,\sF)$, then:
\begin{enumerate}[a)]
  \item under the hypotheses of \ref{VI:2-1}$^\ast$, if $i>n$, then $q^{i-n}$ divides $\alpha$;
  \item if the inverses of the eigenvalues of $F_x^\ast$ are algebraic integers, then $\alpha^{-1}$ is integral away from $p$. More precisely, $q^{\inf(n,i)}\alpha^{-1}$ is an algebraic integer.
\end{enumerate}

For the sheaves considered in \ref{VI:1-9}, the eigenvalues of the $F_x^\ast$ are roots of unity, so the hypotheses of the theorem and of b) above are satisfied.

\subsection{}\label{VI:2-3}

Let $f:A\to B$ be a map between finite sets, and let $\varepsilon$ be a function on $A$. Then
\begin{equation*}\tag{2.3.1}\label{VI:eq:2-3-1}
  \sum_{a\in A} \varepsilon(a) = \sum_{b\in B} \sum_{f(a)=b} \varepsilon(a) \text{.}
\end{equation*}

\subsection*{2.3*}\label{VI:2-3*}

Let $f:X\to Y$ be a morphism of separated schemes of finite type over an algebraically closed field $k$, and let $\sF$ be an $E_\lambda$-sheaf on $X$. The Leray spectral sequence for $f$ in cohomology with proper support is
\begin{align*}\tag{2.3.1*}\label{VI:eq:2-3-1*}
  E_2^{p q} = \h_c^p(Y,\eR^q f_! \sF) \Rightarrow \h_c^{p+q}(X,\sF) \text{.}
\end{align*}

Let $f_0:X_0\to Y_0$ be a morphism of separated schemes of finite type over $\dF_q$, and let $\sF_0$ be an $E_\lambda$-sheaf on $X_0$. The following compatibility holds between \eqref{VI:eq:2-3-1} and \eqref{VI:eq:2-3-1*}.

\begin{enumerate}[a)]
  \item The sheaf $\eR^qf_!\sF$ is obtained from the sheaf $\eR^qf_{0!}\sF_0$ on $Y_0$ by extending scalars from $\dF_q$ to $\dF$. It is therefore equipped with a Frobenius correspondence $F^\ast:F^\ast\eR^qf_!\sF\to\eR^qf_!\sF$. By functoriality of $\eR f_!$, it is deduced from the Frobenius correspondence of $\sF$; it is the composite
    \[\xymatrix{
      F^\ast \eR^q f_! \sF \ar[r]^-\sim
        & \eR^q f_! F^\ast \sF \ar[r]^-\sim
        & \eR^q f_! \sF \text{.}
    }\]
  \item The trace formula for $\sF_0$ reads
    \begin{align*}\tag{1}\label{VI:eq:1}
      \sum_{x\in X_0(\dF_q)} \tr(F_x^\ast,\sF_0)
      = \tr(F^\ast,\h_c^\bullet(X,\sF)) \text{;}
    \end{align*}
    that for $\eR^qf_!\sF$ reads
    \begin{align*}\tag{2$_q$}\label{VI:eq:2_q}
      \sum_{y\in Y_0(\dF_q)} \tr(F_y^\ast,\eR^q f_! \sF)
      = \tr(F^\ast,\h_c^\bullet(Y,\eR^q f_! \sF)) \text{.}
    \end{align*}
\end{enumerate}

For $y\in Y(\dF)$ one has $(\eR^qf_!\sF)_y=\h_c^q(f^{-1}(y),\sF)$, and if $y\in Y_0(\dF_q)$,
\[
  \tr(F_y^\ast,\eR^q f_! \sF) = \tr(F^\ast,\h_c^q(f^{-1}(y),\sF)) \text{,}
\]
so the trace formula on the fiber gives
\[
  \sum_{\substack{x\in X_0(\dF_q) \\ f(x) = y}} \tr(F_x^\ast,\sF_0) = \sum (-1)^q \tr(F_y^\ast,\eR^q f_! \sF) \text{.}
\]
Taking the alternating sum of the identities \eqref{VI:eq:2_q}, one gets
\begin{align*}\tag{2}\label{VI:eq:2}
  \sum_{y\in Y_0(\dF_q)} \sum_{\substack{x \in X_0(\dF_q) \\ f(x) = y}} \tr(F_x^\ast,\sF_0) = \tr(F^\ast,\h^\bullet(Y,\eR^\bullet f_! \sF)) \text{.}
\end{align*}

Equality of the left-hand sides of \eqref{VI:eq:1} and \eqref{VI:eq:2} follows from \eqref{VI:eq:2-3-1}; equality of the right-hand sides follows from \eqref{VI:eq:2-3-1*}, taking into account that $F^\ast$ is an endomorphism of the entire spectral sequence.

\subsection{}\label{VI:2-4}

Let $(A_i)_{i\in I}$ be a finite family of finite sets, let $A=\prod_{i\in I}A_i$, let $\varepsilon_i$ be a function on $A_i$, and let $\varepsilon$ be the function $a\mapsto\prod\varepsilon_i(a_i)$ on $A$. Then
\begin{align*}\tag{2.4.1}\label{VI:eq:2-4-1}
  \sum_{a\in A} \varepsilon(a) = \prod_{i\in I} \sum_{a_i\in A_i} \varepsilon_i(a_i) \text{.}
\end{align*}

\subsection*{2.4*}\label{VI:2-4*}

Let $(X_i)_{i\in I}$ be a finite family of separated schemes of finite type over an algebraically closed field $k$, let $X=\prod_{i\in I}X_i$, let $\sF_i$ be an $E_\lambda$-sheaf on $X_i$, and let $\sF$ be their external tensor product $\bigboxtimes_{i\in I}\sF_i=\bigotimes_{i\in I}\pr_i^\ast(\sF_i)$. Then the K\"unneth formula is
\begin{align*}\tag{2.4.1*}\label{VI:eq:2-4-1*}
  \h_c^\bullet(X,\sF) = \bigotimes_{i\in I} \h_c^\bullet(X_i,\sF_i) \text{.}
\end{align*}

To make \eqref{VI:eq:2-4-1*} canonical and independent of a choice of total ordering on $I$, the right-hand side must be the graded tensor product in the sense of the Koszul rule (\hyperref[IV]{Cycles}, \ref{IV:1-3}).

\subsection{}\label{VI:2-5}

Let $A$ be a finite set, $B$ a subset of $A$, and $\varepsilon$ a function on $A$. Then
\begin{align*}\tag{2.5.1}\label{VI:2-5-1}
  \sum_{a\in A} \varepsilon(a) = \sum_{a\in A\setminus B}\varepsilon(a) + \sum_{a\in B} \varepsilon(a) \text{.}
\end{align*}

\subsection*{2.5*}\label{VI:2-5*}

Let $X$ be a separated scheme of finite type over an algebraically closed field $k$, let $Y$ be a closed subscheme, and let $\sF$ be an $E_\lambda$-sheaf on $X$. There is a long exact cohomology sequence
\begin{align*}\tag{2.5.1*}\label{VI:eq:2-5-1*}
\xymatrix{
  \cdots \ar[r]^-\partial
    & \h_c^i(X\setminus Y,\sF) \ar[r]
    & \h_c^i(X,\sF) \ar[r]
    & \h_c^i(Y,\sF) \ar[r]^-\partial
    & \cdots
}
\end{align*}

More generally, starting with a finite filtration of $X$ by closed parts $X_p$ ($p\in\dZ$), where $X_p\supset X_{p+1}$, $X_p=X$ for $p$ sufficiently small, and $X_p=\varnothing$ for $p$ sufficiently large, there is a spectral sequence
\begin{align*}\tag{2.5.2*}\label{VI:eq:2-5-2*}
  E_1^{p q} = \h_c^{p+q}(X_p\setminus X_{p+1},\sF)
  \Rightarrow \h_c^{p+q}(X,\sF) \text{.}
\end{align*}

\subsection{}\label{VI:2-6}

Let $A$ be a finite set, let $(A_i)_{i\in I}$ be a finite covering of $A$, and let $\varepsilon$ be a function on $A$. For $J\subset I$, let $A_J$ be the intersection of the $A_j$ for $j\in J$. Then
\begin{align*}\tag{2.6.1}\label{VI:eq:2-6-1}
  \sum_{J\subset I} (-1)^{|J|} \sum_{a\in A_J} \varepsilon(a) = 0 \text{.}
\end{align*}

\subsection*{2.6*}\label{VI:2-6*}

Let $X$ be a separated scheme of finite type over an algebraically closed field $k$, let $(X_i)_{i\in I}$ be a finite closed (respectively, open) covering of $X$ by subschemes, and let $\sF$ be an $E_\lambda$-sheaf on $X$. For $J\subset I$, let $X_J$ be the intersection of the $X_j$ for $j\in J$. There are respective Leray spectral sequences
\begin{align*}\tag{2.6.1*}\label{VI:eq:2-6-1*}
  E_1^{p q} &= \bigoplus_{|J|=p+1>0} \h_c^q(X_J,\sF) \otimes_\dZ \textstyle\bigwedge^{|J|} \dZ^J \Rightarrow \h_c^{p+q}(X,\sF) \\
  \tag{2.6.2*}\label{VI:eq:2-6-2*}
  E_1^{p q} &= \bigoplus_{|J|=1-p>0} \h_c^q(X_J,\sF) \otimes_\dZ \textstyle\bigwedge^{|J|} \dZ^J \Rightarrow \h_c^{p+q}(X,\sF)
\end{align*}
If $J=\varnothing$ were not excluded, there would similarly be spectral sequences converging to $0$. The factors $\bigwedge^{|J|}\dZ^J$ avoid the need to choose a total ordering on $I$.

\subsection{}\label{VI:2-7}

Let $A$ be a finite commutative group, and let $\chi:A\to E_\lambda^\times$ be a nontrivial character. Then
\begin{align*}\tag{2.7.1}\label{VI:eq:2-7-1}
  \sum_{a\in A} \chi(a) = 0 \text{.}
\end{align*}
We prove the cohomological analogue of \eqref{VI:eq:2-7-1} by transposing the following proof: if $x\in A$, then $a\mapsto xa$ is a permutation of $A$, and
\begin{align*}
  \sum_{a\in A}\chi(a) = \sum_{a\in A} \chi(x a)
    &= \chi(x) \sum_{a\in A} \chi(a) && \text{, whence} \\
  (\chi(x)-1)\sum_{a\in A} &= 0 \text{,}
\end{align*}
by hypothesis, there is an $x$ such that $\chi(x)-1\ne0$.

\begin{theorem*}[2.7*]\label{VI:2-7*}
Let $G_0$ be a connected commutative algebraic group over $\dF_q$, and let $\chi:G_0(\dF_q)\to E_\lambda^\times$ be a nontrivial character. Then
\begin{align*}\tag{2.7.1*}\label{VI:eq:2-7-1*}
  \h_c^\bullet(G,\sF(\chi)) = 0 \text{.}
\end{align*}
\end{theorem*}

For a rational point $x$ of $G$, let $t_x$ denote the translation $t_x(g)=xg$ of $G$. The formula $\fL t_x=t_{\fL(x)}\fL$ says that $(t_x,t_{\fL(x)})$ is an automorphism of the diagram $G\xrightarrow\fL G$, with $G$ equipped with the Lang torsor. Let $\rho(g)$ be the induced automorphism of $(G,\sF(\chi))$. For $g\in G_0(\dF_q)$, that is, for $\fL(g)=e$, it is the identity on $G$ and multiplication by $\chi(g)^{-1}$ on $\sF(\chi)$.

Let $\rho_H(g)$ denote the automorphism of $\h_c^\bullet(G,\sF(\chi))$ induced by $\rho(g)$. A homotopy argument (Lemma \ref{VI:2-8} below) shows that $\rho_H(g)=\rho_H(e)$, hence that it is the identity. On the other hand, for $g\in G_0(\dF_q)$, $\rho_H(g)$ is multiplication by $\chi^{-1}(g)$. Thus multiplication by $\chi^{-1}(g)-1$ on $\h_c^\bullet(G,\sF(\chi))$ is zero. Taking $g$ such that $\chi(g)\ne1$ proves \ref{VI:2-7*}.

\begin{lemma_}\label{VI:2-8}
Let $X$ and $Y$ be schemes over an algebraically closed field $k$, with $X$ separated of finite type and $Y$ connected. Let $\sF$ be a sheaf on $X$, and let $(\rho,\varepsilon)$ be a family of endomorphisms of $(X,\sF)$ parametrized by $Y$:
\begin{align*}
  \rho:Y\times_k X &\to Y\times_k X && \text{is a $Y$-morphism, and} \\
  \varepsilon:\rho^\ast \pr_2^\ast \sF &\to \pr_2^\ast\sF
    && \text{is a morphism of sheaves.}
\end{align*}
Assume that $\rho$ is proper. For $y\in Y(k)$, let $\rho_H(y)^\ast$ be the endomorphism of $\h_c^\bullet(X,\sF)$ induced by $\rho_y:X\to X$ and $\varepsilon_y:\rho_y^\ast\sF\to\sF$. Then $\rho_H(y)^\ast$ is independent of $y$.
\end{lemma_}

Indeed, $\eR^p\pr_{1!}\pr_2^\ast\sF$ is the constant sheaf on $Y$ with value $\h_c^p(X,\sF)$, and $\rho_H(y)^\ast$ is the stalk at $y$ of the endomorphism
\[\xymatrix{
  \eR^p \pr_{1!} \pr_2^\ast \sF \ar[r]^-{\rho^\ast}
    & \eR^p \pr_{1!} \rho^\ast \pr_2^\ast \sF \ar[r]^-\varepsilon
    & \eR^p \pr_{1!} \pr_2^\ast \sF
}\]
of this sheaf.

\subsection{Remark}\label{VI:2-9}

Let $G_0$ be a connected algebraic group over $\dF_q$, and let $\rho:G_0(\dF_q)\to\operatorname{GL}(V)$ be a linear representation such that $V^{G_0(\dF_q)}=0$. Then again $\h_c^\bullet(G,\rho(L))=0$.

\section{One-variable sums}\label{VI:3}

\subsection{}\label{VI:3-1}

Weil was the first to apply ``cohomological'' methods to the study of one-variable trigonometric sums. Since he had proved the analogue of the Riemann hypothesis for Artin $L$-functions over function fields, these methods gave him excellent estimates, of order the square root of the number of terms, as well as the behavior of these sums under extension of the base field.

The main purpose of this section is to present his results in cohomological language.

\subsection{}\label{VI:3-2}

Here the cohomological methods lead to the calculation of the groups $\h_c^i(X,\sF)$, where $\sF_0$ is an $E_\lambda$-sheaf on a curve $X_0$ over $\dF_q$. These groups vanish for $i\ne0,1,2$, and for $i=0,2$ they have a simple interpretation (\ref{VI:1-18}a,b,c). Moreover, the Euler--Poincaré characteristic $\chi_c(\sF)=\sum(-1)^i\dim\h_c^i(X,\sF)$, which is also equal to $\chi(\sF)=\sum(-1)^i\dim\h^i(X,\sF)$, can be calculated in terms of $X$, the rank of $\sF$ at the generic points of $X$, and the ramification properties of $\sF$. The essential result is the following.

Let $\bar X$ be a smooth connected projective curve of genus $g$ over an algebraically closed field $k$, let $X$ be the dense open complement in $\bar X$ of a finite set $S$ of points, and let $\sF$ be a lisse $E_\lambda$-sheaf on $X$. Let $\chi(X)=2-2g-\#S$ be the Euler--Poincaré characteristic of $X$, and let $\operatorname{rg}(\sF)$ be the rank of $\sF$. For each $s\in S$, one defines an integer $\swan_s(\sF)$, the Swan conductor, measuring the wild ramification of $\sF$, and
\begin{align*}\tag{3.2.1}\label{VI:eq:3-2-1}
  \chi_c(\sF) = \operatorname{rg}(\sF) \cdot \chi(X) - \sum_{s\in S} \swan_s(\sF)
\end{align*}
\cite{ra65}.

\subsection{}\label{VI:3-3}

Let $\bar X$, $j:X\hookrightarrow\bar X$, and $\sF$ be as in \ref{VI:3-2}. Poincaré duality takes the particularly convenient form that $\h^i(\bar X,j_*\sF)$ and $\h^{2-i}(\bar X,j_*(\sF^\vee(1)))$ are dual to one another (\hyperref[V]{Duality}, \ref{V:1-3}).

\subsection{}\label{VI:3-4}

Take $k=\dF$, and suppose that $X$, $\bar X$, and $\sF$ come from $X_0$, $\bar X_0$, and $\sF_0$ over $\dF_q$. If $\sF_0$ becomes trivial on a finite covering of $X_0$, Weil proved that the complex conjugates of the eigenvalues of $F^\ast$ on $\h^i(\bar X,j_*\sF)$ have absolute value $q^{i/2}$. Such sheaves $\sF_0$ correspond to Artin $L$-functions over the function field of $\bar X_0$.

\subsection*{3.5}\label{VI:3-5_} % originally one of two 3.5s

Under the same hypotheses, or more generally if $E_\lambda$ is the completion at a place $\lambda$ of a number field $E$ and $\sF_0$ belongs to an infinite compatible system of $v$-adic representations, with $v$ ranging over the places of $E$, one can calculate
\[
  \prod_i \det(-F^\ast,\h^i(X,\sF))^{(-1)^i}
\]
from \emph{local} information on $\sF_0$. This is explained in \cite{de73}. When $\sF_0$ is trivial on a finite covering of $X_0$, this expresses the constant in the functional equation of an Artin $L$-function as a product of local constants.

\subsection{Example}\label{VI:3-5}

Let $\psi$ be the character $\exp\left(\frac{2\pi i}{p}\tr_{\dF_q/\dF_p}\right)$ of $\dF_q$; then $\psi(a^p-a)=1$. Let $X_0$ be an absolutely irreducible smooth projective curve of genus $g$ over $\dF_q$, and let $f$ be a rational function on $X_0$, that is, a morphism $f:X_0\to\dP^1$ not identically equal to $\infty$. We consider the sum
\[
  S_f' = \sum_{\substack{x\in X_0(\dF_q) \\ f(x)\ne \infty}} \psi(f(x))
  \text{.}
\]
This sum vanishes when $f$ is of the form $g^p-g$. This suggests modifying $S_f'$ as follows:
\begin{enumerate}[\indent a)]
  \item for every closed point $x$ of $X_0$, let $v_x(f)$ be the order of the pole of $f$ at $x$ if $f(x)=\infty$, and set $v_x(f)=0$ otherwise; set $v_x^\ast(f)=\inf_g v_x(f+g^p-g)$;
  \item if $v_x^\ast(f)=0$, $x\in X_0(\dF_q)$, and $f+g^p-g$ is regular at $x$, set $\psi(f(x))=\psi((f+g^p-g)(x))$. Finally, put
    \begin{equation*}\tag{3.5.1}\label{VI:eq:3-5-1}
      S_f = \sum'_{x\in X_0(\dF_q)} \psi(f(x)) \text{,}
    \end{equation*}
    where the prime indicates that the sum is over the $x$ such that $v_x^\ast(f)=0$.
\end{enumerate}

One has $S_f=S_{f+g^p-g}$. We shall verify that if $f$ is not of the form $g^p-g+c$, with $c$ constant, then
\begin{equation*}\tag{3.5.2}\label{VI:eq:3-5-2}
  |S_f|\leqslant \left(2 g-2+\sum_{v_x^\ast(f)\ne 0} [k(x):\dF_q] (1+v_x^\ast(f))\right) q^{1/2} \text{.}
\end{equation*}

First pass from complex to $\ell$-adic coefficients, replacing $\psi$ by a character of the form $\psi_0\circ\tr_{\dF_q/\dF_p}$, with $\psi_0:\dF_p\hookrightarrow E_\lambda^\times$. If $j:U_0\hookrightarrow X_0$ is the inclusion of the open subset on which $f\ne\infty$, one then proves that
\begin{equation*}\tag{3.5.3}\label{VI:eq:3-5-3}
  S_f = \sum (-1)^i \tr(F^\ast,\h^i(X,j_\ast \sF(\psi f))) \text{.}
\end{equation*}

The trace formula reduces this assertion to the following statements, for $x\in X_0(\dF_q)$:
\begin{enumerate}[\indent a)]
  \item if $v_x^\ast(f)=0$, then $F_x^\ast$ on $j_*\sF(\psi f)$ is $\psi(f(x))$;
  \item if $v_x^\ast(f)\ne0$, then $j_*\sF(\psi f)$ is ramified at $x$: $(j_*\sF(\psi f))_{\bar x}=0$.
\end{enumerate}

Statement a) follows from \eqref{VI:eq:1-7-6} if $v_x(f)=0$, and one reduces to this case by replacing $f$ by $f+g^p-g$: up to isomorphism, this does not change $\sF(\psi f)_0$ on the open set where $f$ and $g$ are regular (\ref{VI:1-3}).

Statement b) follows from
\begin{equation*}\tag{3.5.4}\label{VI:eq:3-5-4}
  \swan_x \sF(f,\psi) = v_x^\ast(f) \text{.}
\end{equation*}
After reduction to the case $v_x(f)=v_x^\ast(f)$, whence $p\nmid v_x(f)$, and extension of scalars to $\dF$, this formula is \cite[4.4]{se61}.

Finally, the sheaf $\sF(\psi f)$ on $U$ is nonconstant if $f$ is not of the form $g^p-g+c$ with $c$ constant. Indeed, since $\psi_0$ is injective, triviality of $\sF(\psi f)=\sF(\psi_0f)$ (\ref{VI:1-8}.ii) would imply that the corresponding $\dF_p$-torsor is pulled back from a torsor over $\spec(\dF_q)$ with equation $T^p-T-\lambda=0$. The torsor on $U_0$ with equation $T^p-T-(f-\lambda)=0$ would then be trivial, so $f=g^p-g+\lambda$. Thus the groups $\h^i(\bar X,j_*\sF(\psi f))$ vanish for $i\ne1$, and \eqref{VI:eq:3-5-3} reduces to
\begin{equation*}\tag{3.5.5}\label{VI:eq:3-5-5}
  S_f = -\tr(F^\ast,\h^1(\bar X,j_\ast \sF(\psi f)))
\end{equation*}
By \eqref{VI:eq:3-2-1} and \eqref{VI:eq:3-5-4}, this $\h^1$ has dimension $2g-2+\sum_{v_x^\ast(f)>0}[k(x):\dF_q](1+v_x^\ast(f))$, and $-S_f$ is the sum of that many eigenvalues of $F^\ast$, each of whose complex conjugates has absolute value $q^{1/2}$.

\subsection{}\label{VI:3-6}

Suppose that, for an automorphism $\sigma$ of $X_0$, one has
\[
  f(\sigma x) = - f(x) \text{.}
\]
Then the sum $S_f$ is real. If $\sigma$ is an involution, Poincaré duality gives a little more. Put $\sF_0=j_*\sF(\psi f)$ and $\sG_0=j_*\sF(\psi(-f))$. Then:
\begin{enumerate}[\indent a)]
  \item $\sF_0$ and $\sG_0$ are dual;
  \item $\sigma^\ast\sF_0\iso\sG_0$ and $\sigma^\ast\sG_0\iso\sF_0$, by natural isomorphisms such that the composite $\sF_0=(\sigma^2)^\ast\sF_0\iso\sigma^\ast\sG_0\iso\sF_0$ is the identity;
  \item the pairing $\sF_0\otimes\sG_0\to E_\lambda$ satisfies $\sigma^\ast(f\cdot g) = \sigma^\ast(f) \cdot \sigma^\ast (g)$.
\end{enumerate}

Passing to cohomology, one finds that $\h^1(X,\sF)$ and $\h^1(X,\sG)$ are in perfect duality with values in $E_\lambda(1)$, and that the bilinear form $\alpha\cdot\sigma^\ast\beta$ on $\h^1(X,\sF)$ is \emph{alternating}: $\sigma^\ast$ acts trivially on $E_\lambda(1)$, and
\[
  \alpha\cdot \sigma^\ast \beta = \sigma^\ast(\alpha\cdot \sigma^\ast \beta) = \sigma^\ast \alpha\cdot \beta = -\beta\cdot \sigma^\ast \alpha \text{.}
\]
It follows that $\h^1(X,\sF)$ has even dimension (one also readily checks that every nonzero $v_x^\ast(f)$ is odd) and that the eigenvalues of $F^\ast$ on $\h^1(X,\sF)$ occur in pairs $\alpha$ and $q/\alpha$.

\subsection{Example}\label{VI:3-7}

This applies to the Kloosterman sums $\sum_{x\in\dF_q^\times}\psi(x+a/x)$: take $X_0=\dP^1$, $f=x+a/x$, and $\sigma x=-x$. The poles of $f$ are $0$ and $\infty$, and at each one $v_x^\ast(f)=1$; thus $\h^1$ has dimension $-2+2+2=2$. Consequently, for $a\ne0$,
\[
  \sum_{\substack{x y=a \\ x,y\in \dF_{q^n}}}
  \exp\left(\frac{2\pi i}{p}
  \tr_{\dF_{q^n}/\dF_p}(x+y)\right)
  =-(\alpha^n+\bar\alpha^n)
  \qquad\text{with }\alpha\bar\alpha=q \text{.}
\]
This result is due to L.\ Carlitz \cite{ca69}.

Here is an application of a Lang--Weil method.

\begin{proposition_}\label{VI:3-8}
Let $P\in\dF_q[X_1,\dotsc,X_n]$ be a polynomial in $n$ variables of degree $d$, assumed not to be of the form $Q^p-Q+C$ with $C$ constant. Again put $\psi(x)=\exp\left(\frac{2\pi i}{p}\tr_{\dF_q/\dF_p}(x)\right)$. Then
\[
  \left|\sum_{x_i\in \dF_q} \psi P(x_1,\dots,x_n)\right| \leqslant (d+1) q^{n-\frac 1 2} \text{.}
\]
\end{proposition_}

The left-hand side has the trivial bound $q^n$. It is therefore enough to prove \ref{VI:3-8} when $d-1<q^{1/2}$. Assume merely that $d<q+1$, and let $P_d$ be the homogeneous part of degree $d$ of $P$. Recall the following.

\begin{lemma_}\label{VI:3-9}
A hypersurface of degree $d$ in $\dP^r$, with $r\geqslant1$, cannot pass through all the $\dF_q$-rational points unless $d\geqslant q+1$.
\end{lemma_}

Proceed by induction. If there is a rational hyperplane not wholly contained in the hypersurface (which is automatic for $r=1$), apply the induction hypothesis to the trace of the hypersurface on that hyperplane. Otherwise, the degree $d$ is at least the number of rational hyperplanes, hence at least $q+1$.

We now distinguish two cases.

\paragraph{Case 1: $p\nmid d$.}
Applying the lemma to $P_d$, and making a linear change of variables if necessary, we may assume that $P_d(1,0,\ldots,0)\ne0$, that is, that the coefficient of $X_1^d$ in $P$ is nonzero. A polynomial in one variable
\[
  S(X) = \sum_{i=0}^d a_i x^i \text{,}
\]
with $a_d\ne0$ and $p\nmid d$ is never of the form $Q^p-Q+C$ with $C$ constant, and estimate \eqref{VI:eq:3-5-2} becomes
\[
  \left| \sum \psi(S(x)) \right| \leqslant (d-1) q^{1/2} \text{.}
\]
Applying this estimate to the partial sums obtained by varying only $x_1$ gives
\[
  \left| \sum \psi P(x_1,\dots,x_n) \right| = \left| \sum_{x_2,\dots,x_n} \sum_{x_1} \psi P(x_1,\dots,x_n)\right| \leqslant q^{n-1} (d-1) q^{1/2}
\]
as claimed.

\paragraph{Case 2: $p\mid d$ ($d>0$).}
If $P_d$ is a $p$th power, replacing $P$ by $P-(P_d-P_d^{1/p})$ does not change the sum and lowers the degree; this case is eliminated by induction on $d$. Otherwise, the differential of $P_d$ is not identically zero. Applying \ref{VI:3-9}, and making a linear change of variables if necessary, we may assume that it is nonzero at $(1,0,\dotsc,0)$. Write
\[
  P_d = \sum_{i=0}^d X_1^{d-i} S_i(X_2,\dots,X_n) \text{,}
\]
This means that the linear form $S_1$ is not identically zero.

The form $S_0$ is constant, and by replacing $P$ with $P-(S_0X_1^d-S_0^{1/p}X_1^{d/p})$, we may assume it is zero. Finally, let $-\lambda$ be the coefficient of $X_1^{d-1}$ in $P$. For fixed $x_2,\dotsc,x_n$, one has
\[
  P(X_1,x_2,\dotsc,x_n)
  =(S_1(x_2,\dotsc,x_n)-\lambda)X_1^{d-1}
  +\text{terms of lower degree in $X_1$}
\]
and if $S_1(x_2,\dotsc,x_n)\ne\lambda$, then
\[
  \left| \sum_{x_1} \psi P(x_1,\dots,x_n)\right| \leqslant (d-2) q^{1/2} \text{.}
\]
Altogether,
\begin{align*}
  \left|\sum \psi P(x_1,\dots,x_n)\right|
    &\leqslant \left|\sum_{S=\lambda} \psi P(x_1,\dots,x_n)\right| + \sum_{S(X_2,\dots,X_n)\ne\lambda} \left| \sum_{x_1} \psi P(x_1,\dots,x_n)\right| \\
    &\leqslant q^{n-1} + (q^{n-1}-q^{n-2})(d-2) q^{1/2} \\
    &< (d-2) q^{n-1/2} + q^{n-1} \\
    &< (d-1) q^{n-1/2} \text{.}
\end{align*}

Another result of this kind is due to R.\ A.\ Smith \cite{sm70}.

\section{Gauss sums and Jacobi sums}\label{VI:4}

\subsection{}\label{VI:4-1}

Let $k$ be a finite field of characteristic $p$, let $\chi$ be a character of $k^\times$, and let $\psi$ be a nontrivial character of the additive group of $k$. We take as our definition of the Gauss sums:
\begin{equation*}\tag{4.1.1}\label{VI:eq:4-1-1}
  \tau(\chi,\psi) = -\sum_{x\in k^\times} \psi(x) \chi^{-1}(x)
\end{equation*}
(note the sign). Classically, $\chi$ and $\psi$ are complex-valued. We shall take them to have values in $E_\lambda^\times$ (cf.\ the proof of \ref{VI:1-15}). Regard $-\tau(\chi,\psi)$ as a sum over the rational points of the scheme $\dG_m$ over $k$. By Section~\ref{VI:1},
\[
  -\tau(\chi,\psi)
  =\tr\bigl(F^\ast,\h_c^\ast(\dG_m,\sF(\psi\chi^{-1}))\bigr).
\]
Moreover:

\begin{proposition_}\label{VI:4-2}
The cohomology of $\dG_m$ with coefficients in $\sF(\psi\chi^{-1})$ satisfies:
\begin{enumerate}[(i)]
  \item $\h_c^i=0$ for $i\ne 1$, and $\dim \h_c^1=1$.
  \item On $\h_c^1$, $F^\ast$ acts as multiplication by $\tau(\chi,\psi)$.
  \item If $\chi$ is nontrivial, then $\h_c^\bullet \iso \h^\bullet$.
\end{enumerate}
\end{proposition_}
\begin{proof}
For every $n$ prime to $p$, let $K_n$ denote the $\dmu_n$-torsor on $\dG_m$ defined by the Kummer exact sequence $0 \to \dmu_n \to \dG_m \xrightarrow{x^n} \dG_m \to 0$. If $k$ has $q$ elements, then over $k$ one has $\dmu_{q-1}=k^\times$, and the Lang torsor on $\dG_m/k$ is $K_{q-1}$. Hence $\sF(\chi^{-1})=\chi(K_{q-1})$. Assertion~\ref{VI:4-2}(ii) follows from (i) and (iii), which themselves are contained in the following geometric statement. Here $\sF(\psi)$ denotes the sheaf on $\dG_m/\dF$ obtained from $\sF(\psi)$ on $\dG_m/k$ by extension of scalars from $k$ to $\dF$ (\ref{VI:1-8}.b).
\end{proof}

\begin{proposition_}\label{VI:4-3}
Let $\chi:\dmu_n \to E_\lambda^\times$. The cohomology of $\dG_m$ with coefficients in $\sF(\psi)\otimes \chi(K_n)$ satisfies:
\begin{enumerate}[\indent (i)]
  \item $\h_c^i=0$ for $i\ne 1$, and $\dim \h_c^1{} = 1$.
  \item If $\chi$ is nontrivial, then $\h_c^\bullet{}\iso \h^\bullet{}$.
\end{enumerate}
\end{proposition_}
\begin{proof}
The sheaf $\sF(\psi)$ is the restriction to $\dG_m$ of a locally constant sheaf on $\dG_a$, wildly ramified at infinity, with Swan conductor $1$. The sheaf $\chi(K_n)$ is constant if $\chi=1$; if $\chi\ne1$, it is tamely ramified at $0$ and at infinity.

Thus $\sF(\psi)\otimes\chi(K_n)$ is ramified at infinity. Applying \ref{VI:1-18}(b,c), we find that $\h_c^i(\dG_m,\sF(\psi)\otimes\chi(K_n))=0$ for $i\ne1$. If $\chi\ne1$, it is ramified at $0$ and at infinity, and (ii) follows from \ref{VI:1-19}.a.

The Swan conductors are $0$ at $0$ and $1$ at infinity. By \eqref{VI:eq:3-2-1}, the Euler--Poincaré characteristic is therefore $-1$, which completes the proof.
\end{proof}

\subsection{Remark}\label{VI:4-4}

If $\chi$ is nontrivial, \ref{VI:4-2}(iii) and Poincaré duality show that $\h_c^1(\dG_m,\sF(\psi\chi^{-1}))$ and $\h_c^1(\dG_m,\sF(\psi^{-1}\chi))$ are dual (with values in $E_\lambda(-1)$). Consequently $\tau(\chi,\psi)\tau(\chi^{-1},\psi^{-1})=q$, i.e. $|\tau(\chi,\psi)|=q^{1/2}$.

\subsection{}\label{VI:4-5}

If $k$ is an extension of degree $N$ of $\dF_q$, we may also regard \eqref{VI:eq:4-1-1} as a sum in $N$ variables over $\dF_q$. More generally, let $k$ be an étale algebra over $\dF_q$, of degree $N$ over $\dF_q$. It is a product of fields $k_i$ of degrees $N_i$, and we set
\begin{equation*}\tag{4.5.1}\label{VI:eq:4-5-1}
  \varepsilon(k) = (-1)^{\sum (N_i+1)} \text{.}
\end{equation*}
This is the sign of the permutation of $S=\hom_{\dF_q}(k,\dF)$ induced by the Frobenius substitution $\varphi\in\gal(\dF/\dF_q)$.

Let $\psi:\dF_q\to E_\lambda^\times$ be a nontrivial character, and let $\chi:k^\times\to E_\lambda^\times$. We set
\begin{equation*}\tag{4.5.2}\label{VI:eq:4-5-2}
  \tau_{\dF_q}(\chi,\psi) = (-1)^N \sum_{x\in k^\times} \psi \tr_{k/\dF_q}(x) \cdot \chi^{-1}(x) \text{.}
\end{equation*}

If the components of $\chi$ are the characters $\chi_i:k_i^\times\to E_\lambda^\times$, then one has the immediate identity
\begin{equation*}\tag{4.5.3}\label{VI:eq:4-5-3}
  \tau_{\dF_q}(\chi,\psi) = \varepsilon(k) \prod \tau(\chi_i,\psi\circ \tr_{k_i/\dF_q}) \text{.}
\end{equation*}

After some preliminaries, in \ref{VI:4-10} we shall give a cohomological interpretation of the sums \eqref{VI:eq:4-5-2}. In \ref{VI:4-12}, we shall interpret the Hasse--Davenport identity as a twisted form (cf.\ \ref{VI:1-12}) of the following special case of \eqref{VI:eq:4-5-3}: for $k=\dF_q^N$ and $\chi$ a character of $\dF_q^\times$, one has $\tau_{\dF_q}(\chi\circ N_{k/\dF_q},\psi)=\tau(\chi,\psi)^N$.

\subsection{}\label{VI:4-6}

Recall that if $M$ is a finite-type projective module over a ring $A$, the functor $\spec(B)\mapsto M\otimes_A B$, from affine schemes over $\spec(A)$ to $\mathsf{Sets}$, is represented by the affine scheme $\dV(M^\vee)$ (EGA notation), the spectrum of $\operatorname{Sym}_A^\bullet(M^\vee)$. If $M=A^n$, this is the standard affine space of dimension $n$.

Suppose that $M$ is a unital $A$-algebra, and set $V=\dV(M^\vee)$. By definition, for every $A$-algebra $B$, the set $V(B)$ of $B$-valued points of $V$ is $M\otimes_A B$. Since the functor $B\mapsto M\otimes_A B$ takes values in unital rings, $V$ is a unital ring scheme over $\spec(A)$. The norm and trace maps $M\otimes_A B\to B$ are functorial in $B$; they therefore correspond to scheme morphisms $N$ and $T$ from $V$ to $\dG_a$. We shall need the following schemes derived from $V$:
\begin{enumerate}[a)]
  \item $V^\ast$ is the open subscheme of invertible elements of $V$ (a group scheme for multiplication);
  \item $W$ is the hyperplane defined by $T=0$, and $W^\ast=W\cap V^\ast$;
  \item $P=(V\setminus\{0\})/\dG_m$ is the projective space at infinity of the affine space $V$ over $\spec(A)$. If $Q$ is the hyperplane at infinity of $W$, then $W^\ast/\dG_m$ and $V^\ast/\dG_m$ are open subschemes of the projective spaces $Q$ and $P$ over $\spec(A)$.
    \begin{equation*}\tag{4.6.1}\label{VI:eq:4-6-1}
    \xymatrix{
      W^\ast \ar@{^{(}->}[r] \ar[d]^-\pi
        & V^\ast \ar[d]^-\pi \\
      W^\ast/\dG_m \ar@{^{(}->}[r]
        & V^\ast/\dG_m
    }\qquad\subset\qquad
    \xymatrix{
      W\setminus \{0\} \ar@{^{(}->}[r] \ar[d]^-\pi
        & V\setminus \{0\} \ar[d]^-\pi \\
      Q \ar@{^{(}->}[r]
        & P \text{.}
    }
    \end{equation*}
\end{enumerate}

If $M=A^I$, where $I$ is a finite set, then $V\simeq\dG_a^I$, $V^\ast=\dG_m^I$, and $V^\ast/\dG_m$ is the torus $\dG_m^I/(\text{diagonal }\dG_m)$; the maps $N$ and $T$ are respectively $\prod x_i$ and $\sum x_i$, and $Q$ is therefore the projective hyperplane in $P$ defined by $\sum x_i=0$.

The case of interest to us is that in which $M$ is finite étale over $A$. The preceding description then holds locally on $\spec(A)$ for the étale topology, and one may write $V^\ast=\dG_m^I$, where $I$ is a locally constant sheaf of finite sets on $\spec(A)$. For example, if $A$ is a field with algebraic closure $\bar A$, and $M$ is a separable $A$-algebra, set $I=\hom_A(M,\bar A)$. Over $\bar A$ one then has $V\simeq\dG_a^I$. Under this isomorphism, $N$ and $T$ are respectively $\prod x_i$ and $\sum x_i$.

\subsection{Kummer torsor}\label{VI:4-7}

Let $S$ be a scheme, let $n$ be an integer invertible on $S$, and let $G$ be a commutative group scheme with connected fibers over $S$, written multiplicatively. The sequence
\[\xymatrix{
  0 \ar[r]
    & G_n \ar[r]
    & G \ar[r]^-{x^n}
    & G \ar[r]
    & 0
}\]
is exact. It defines a $G_n$-torsor $K_n(G)$ (or simply $K_n$) on $G$. For a homomorphism $\chi:G_n\to E_\lambda^\times$ from the étale $S$-sheaf $G_n$ to the constant sheaf $E_\lambda^\times$, denote by $\sK_n(\chi)$ the $E_\lambda$-sheaf $\chi^{-1}(K_n)$.

The torsors $K_n$ form an inverse system of torsors under the inverse system of groups $G_n$ (transition maps $x\mapsto x^d:G_{nd}\to G_n$). This expresses the commutativity of the diagrams
\[\xymatrix{
  0 \ar[r]
    & G_{n d} \ar[r] \ar[d]^-{x^d}
    & G \ar[r]^-{x^{n d}} \ar[d]^-{x^d}
    & G \ar[r] \ar@{=}[d]
    & 0 \\
  0 \ar[r]
    & G_n \ar[r]
    & G \ar[r]^-{x^n}
    & G \ar[r]
    & 0 \text{.}
}\]
Thus $\sK_n(\chi)=\sK_{nd}(\chi\circ x^d)$.

\subsection{}\label{VI:4-8}

Apply the construction of \ref{VI:4-6} with $A=\dF_q$ and $M=k$, an étale algebra over $\dF_q$. In accordance with our general conventions, the $\dF_q$-schemes denoted by $V,V^\ast,\ldots$ in \ref{VI:4-6} will here be written with a subscript $0$. The same letters without a subscript denote the schemes over $\dF$ obtained from them by extension of scalars.

Set $I=\hom(k,\dF)$. Thus $V\simeq\dG_a^I$, and $T$ is $(x_i)\mapsto\sum x_i$.
\begin{equation*}\tag{4.8.1}\label{VI:eq:4-8-1}
  \sF(\psi\circ \tr_{k/\dF_q}) = T^\ast \sF(\psi) = \bigotimes_i \pr_i^\ast \sF(\psi) \text{.}
\end{equation*}
Since $V^\ast\simeq\dG_m^I$, a character $\chi$ of $V_n^\ast\simeq\dmu_n^I$ with values in $E_\lambda^\times$ may be written as an $I$-indexed family $(\chi_i)_{i\in I}$ of characters of $\dmu_n$. It is defined over $\dF_q$ if, for every $\sigma\in\gal(\dF/\dF_q)$, one has $\chi_{\sigma i}=\chi_i\circ\sigma^{-1}$. It then defines an $E_\lambda$-sheaf $\sK_n((\chi_i)_{i\in I})$ on $V_0^\ast$. If the product of the $\chi_i$ is trivial, $\chi$ factors through a character of $(V_0^\ast/\dG_m)_n$, and $\sK_n((\chi_i)_{i\in I})$ is the inverse image of a sheaf, denoted by the same symbol, on $V_0^\ast/\dG_m$. This construction is related to the Lang torsor as follows.

\begin{lemma_}\label{VI:4-9}
\begin{enumerate}[(i)]
  \item For $n$ sufficiently divisible, there is a commutative diagram
    \[\xymatrix{
      0 \ar[r]
        & (V_0^\ast)_n \ar[r] \ar[d]^-\tau
        & V_0^\ast \ar[r]^-{x^n} \ar[d]^-\tau
        & V_0^\ast \ar[r] \ar@{=}[d]
        & 0 \\
      0 \ar[r]
        & k^\times \ar[r]
        & V_0^\ast \ar[r]^-\fL
        & V_0^\ast \ar[r]
        & 0 \text{;}
    }\]
    if $\chi$ is a character $\chi:k^\times\to E_\lambda^\times$, then $\sF(\chi)=\sK_n(\chi\circ\tau)$.
  \item For $n$ sufficiently divisible, $\tau$ identifies $k^\times$ with the coinvariants of $\gal(\dF/\dF_q)$ acting on $V_n^\ast\simeq\dmu_n^I$.
  \item If $k$ is a field, $N=[k:\dF_q]$, $\omega\in I$ is an embedding of $k$ into $\dF$, and $n=q^N-1$, then the component of $\chi\circ\tau$ indexed by $\omega$ is $\chi\circ\omega^{-1}$.
  \item If $\chi$ is nontrivial on every factor of $k$, then all the $(\chi\circ\tau)_i$ are nontrivial.
\end{enumerate}
\end{lemma_}

It is enough to prove the lemma when $k$ is a field of degree $N$ over $\dF_q$. In this case, $n$ is ``sufficiently divisible'' if $q^N-1\mid n$. Choose an embedding $\omega$ of $k$ into $\dF$; then $I$ is identified with $\dZ/N$, with $i\in\dZ/N$ corresponding to $\omega_i=\omega^{q^i}$. Under the isomorphism $V_0^\ast(\dF)={\dF^\times}^I$, one has
\begin{enumerate}[\indent a)]
  \item $x\in k^\times$ corresponds to $(\omega_i(x))\in{\dF^\times}^I$;
  \item $F((x_i)_{i\in \dZ/N})=(x_{i-1}^q)_{i\in \dZ/N}$.
\end{enumerate}

A character $\chi=(\chi_i)$ of $V_n^\ast\simeq\dmu_n^I$ is defined over $\dF_q$ if $\chi_{-i}(x)=\chi_0(x^{q^i})$ for $i\in\dZ$. If $q^N-1\mid n$, there are $q^N-1$ such characters: $\chi_0$ factors through $\dmu_{q^N-1}$ and determines the $\chi_i$. To prove (ii), it is therefore enough to verify (i) and (iii) (or its consequence (iv)), which ensures that $\chi\mapsto\chi\circ\tau$ is injective.

We prove (i) and (iii) for $n=q^N-1$. Write the group of endomorphisms of $V^\ast\simeq\dG_m^I$ additively; in particular, write $n$ for the operator $x\mapsto x^n$. If $\alpha$ is the cyclic permutation operator $(x_i)\mapsto(x_{i-1})$, then $q^N-1=(q\alpha)^N-1=(q\alpha-1)((q\alpha)^{N-1}+\cdots+1)$. This determines $\tau$:
\[
  \tau((x_i))=\prod_{0\leqslant j<N}x_{i-j}^{q^j}.
\]
The induced map from $(V_0^\ast)_n$ to $k^\times$ is, in the notation of a) above, $(x_i)\mapsto\prod x_{-i}^{q^i}$; (iii) follows.

\begin{proposition_}\label{VI:4-10}
The cohomology of $V^\ast$ with coefficients in $\sF(\chi^{-1}\cdot\psi\tr_{k/\dF_q})$ satisfies:
\begin{enumerate}[\indent (i)]
  \item $\h_c^i=0$ for $i\ne N$, and $\dim\h_c^N=1$.
  \item On $\h_c^N$, $F^\ast$ acts as multiplication by $\tau_{\dF_q}(\chi,\psi)$.
  \item If $\chi$ is nontrivial on every factor of $k$, then $\h_c^\bullet\iso\h^\bullet$.
\end{enumerate}
\end{proposition_}

Apply \ref{VI:4-9}(i): over $\dF$, if $\chi\circ\tau=(\chi_i)_{i\in I}$, then $\sF(\chi)=\sK_n((\chi_i)_{i\in I})$. Assertions (i) and (iii) therefore follow from the following more geometric statement (for (iii), apply \ref{VI:4-9}(iv)), and (ii) then follows from the trace formula.

\begin{proposition_}\label{VI:4-11}
Let $(\chi_i)_{i\in I}$ be a family of characters of $\dmu_n(\dF)$. The cohomology of $V^\ast\simeq\dG_m^I$ with coefficients in $\sK_n((\chi_i)_{i\in I})\otimes\sF(\psi\tr_{k/\dF_q})$ satisfies:
\begin{enumerate}[\indent (i)]
  \item $\h_c^i{}=0$ for $i\ne N$, and $\dim\h_c^N{}=1$.
  \item If all the $\chi_i$ are nontrivial, then $\h_c^\bullet{}\iso\h^\bullet$.
\end{enumerate}
\end{proposition_}

One has $V^\ast\simeq\dG_m^I$, $\sK((\chi_i))=\bigotimes_i\pr_i^\ast\chi_i(K_n(\dG_m))$, and $\sF(\psi\circ\tr_{k/\dF_q})=T^\ast\sF(\psi) =\bigotimes_i\pr_i^\ast\sF(\psi)$. The Künneth formula therefore applies, and \ref{VI:4-11} follows from \ref{VI:4-3}.

\subsection{}\label{VI:4-12}

It also follows from \ref{VI:2-4*}* and (\hyperref[VI]{Cycle}, \ref{VI:1-3}, Example~2) that the group of permutations $\sigma$ of $I$ such that $\chi_i=\chi_{\sigma i}$ for $i\in I$ acts on $\h_c^N$ by multiplication by the sign $\varepsilon(\sigma)$. We shall deduce from this a second proof of the Hasse--Davenport identity.

If $\chi$ is a character of $\dF_q^\times$, the fact that over $\dF$ the norm $N$ is $(x_i)\mapsto\prod x_i$ gives $\sF(\chi\circ N)=N^\ast\sF(\chi) =\bigotimes_i\pr_i^\ast\sF(\chi)$, and the Künneth formula gives
\[
  \h_c^\bullet\bigl(V^\ast,
    \sF(\chi^{-1}\circ N)\otimes\sF(\psi\circ\tr)\bigr)
  \simeq
  \h_c^\bullet(\dG_m,\sF(\chi^{-1}\psi))^{\otimes I},
\]
where the tensor product on the right is taken in the sense of \ref{VI:2-4*}*. Since $\h_c^1(\dG_m,\sF(\psi\chi^{-1}))$ is one-dimensional, its ordinary $I$-indexed tensor power depends only on the cardinality $N$ of $I$. Applying (\hyperref[IV]{Cycle}, \ref{IV:1-3}, Example~2), one obtains a canonical isomorphism
\begin{equation*}\tag{4.12.1}\label{VI:eq:4-12-1}
  \h_c^N\left(V^\ast,\sF(\chi^{-1}\circ N)\otimes T^\ast\sF(\psi)\right)
  \simeq
  \h_c^1(\dG_m,\sF(\psi\chi^{-1}))^{\otimes N}
  \otimes_\dZ \textstyle\bigwedge^N\dZ^I.
\end{equation*}
To compute the action of $F^\ast$, it is most convenient to adopt the Galois viewpoint: since the isomorphism \eqref{VI:eq:4-12-1} is canonical, it is compatible with the action of $\gal(\dF/\dF_q)$ by transport of structure. If $\varepsilon(k)$ is the sign of the permutation $\varphi$ of $I$, one finds
\begin{align*}
  \tau_{\dF_q}(\chi\circ N_{k/\dF_q},\psi)
    &= \tr\left(\varphi^{-1},\h_c^N\left(V^\ast,N^\ast\sF(\chi^{-1})\otimes T^\ast \sF(\psi)\right)\right) \\
    &= \varepsilon(k) \tr\left(\varphi^{-1},\h_c^1(\dG_m,\sF(\chi^{-1}\psi))\right)^N \\
    &= \varepsilon(k) \tau(\chi,\psi)^N \text{,}
\end{align*}
that is,
\begin{equation*}\tag{4.12.2}\label{VI:eq:4-12-2}
  \tau_{\dF_q}\left(\chi\circ N_{k/\dF_q},\psi\right) = \varepsilon(k) \tau(\chi,\psi)^N \text{.}
\end{equation*}

If $k$ is a field, $\varphi$ is a cyclic permutation of $I$, $\varepsilon(k)=(-1)^{N+1}$, and one recovers the Hasse--Davenport identity:
\[
  \tau\left(\chi\circ N_{k/\dF_q},\psi\circ \tr_{k/\dF_q}\right) = (-1)^{N+1} \tau_{\dF_q}\left(\chi\circ N_{k/\dF_q},\psi\right) = \tau(\chi,\psi)^N \text{.}
\]

\begin{lemma_}\label{VI:4-13}
For $\chi\circ\tau=(\chi_i)_{i\in I}$ as in \ref{VI:4-9}(i), the following conditions are equivalent:
\begin{enumerate}[\indent (i)]
  \item $\chi|_{\dF_q^\times}$ is trivial;
  \item the product of the $\chi_i$ is trivial;
  \item $\sF(\chi)=\sK_n((\chi_i))$ is the inverse image of a unique sheaf on $V^\ast/\dG_m$;
  \item the inverse image of $\sF(\chi)$ on $\dG_m$ (mapped into $V^\ast$ by the structural morphism $\dF_q\to k$) is trivial.
\end{enumerate}
\end{lemma_}

(i)$\Rightarrow$(ii). Regard $\chi$ as a character of $(V^\ast/\dG_m)(\dF_q)=k^\times/\dF_q^\times$, and take $\sF(\chi)$ on $V^\ast/\dG_m$.

(ii)$\Rightarrow$(iii). Similarly, regard $(\chi_i)$ as a character of $(V^\ast/\dG_m)_n$. Uniqueness in (iii) follows because $V^\ast$ is a bundle with connected fibers over $V^\ast/\dG_m$, and (iii)$\Rightarrow$(iv) is immediate.

(iv)$\Rightarrow$(i), (ii). This inverse image is both $\sF(\chi|_{\dF_q^\times})$ and $\sK_n(\prod\chi_i)$.

\subsection{}\label{VI:4-14}

Let $k$ be an étale algebra over $\dF_q$ of dimension $N+1$, and let $\chi$ be a nontrivial character of $k^\times$ whose restriction to $\dF_q^\times$ is trivial. The Jacobi sum $J(\chi)$ is defined by
\begin{equation*}\tag{4.14.1}\label{VI:eq:4-14-1}
  J(\chi) = (-1)^{N-1} \sum_{\substack{x\in k^\times/\dF_q^\times \\ \tr(x)=0}} \chi^{-1}(x) \text{.}
\end{equation*}

Gauss sums and Jacobi sums are related by the following identity.

\begin{proposition_}\label{VI:4-15}
For $\chi$ as above and $\psi$ a nontrivial additive character of $\dF_q$, one has
\[
  q J(\chi) = \tau_{\dF_q}(\chi,\psi) \text{.}
\]
\end{proposition_}

In the special case $k=\dF_q^{N+1}$, this formula may be rewritten, using \eqref{VI:eq:4-5-3}, as
\begin{equation*}\tag{4.15.1}\label{VI:eq:4-15-1}
  q J(\chi) = \prod_i \tau(\chi_i,\psi) \text{,}
\end{equation*}
or, in the more customary form,
\begin{align*}
  \chi_0(-1) J(\chi) &= (-1)^{N-1} \sum_{\substack{x_1,\dots,x_N\in \dF_q^\times \\ \sum x_i = 1}} \prod_{i=1}^N \chi_i^{-1}(x_i) = \tau(\chi_0^{-1},\psi)^{-1} \prod_{i=1}^N \tau(\chi_i,\psi) \\
  \chi_0^{-1} &= \prod_{i=1}^N \chi_i
\end{align*}

\begin{proof}[Proof of \ref{VI:4-15}]
\[
  \tau_{\dF_q}(\chi,\psi) = (-1)^{N+1} \sum_{x\in k^\times} \chi(x)^{-1} \psi \tr(x) = (-1)^{N+1} \sum_{x\in k^\times/\dF_q^\times} \chi(x)^{-1} \sum_{\lambda\in \dF_q^\times} \psi(\lambda \tr x) \text{.}
\]
The sum $\sum_{\lambda\in\dF_q^\times}\psi(\lambda\tr x)$ is $q-1$ if $\tr(x)=0$, and $-1$ if $\tr(x)\ne0$. Hence
\[
  \tau_{\dF_q}(\chi,\psi)
  =(-1)^{N+1}\left(
    q\sum_{\substack{x\in k^\times/\dF_q^\times\\ \tr(x)=0}}\chi(x)^{-1}
    -\sum_{x\in k^\times/\dF_q^\times}\chi(x)^{-1}\right).
\]
The second term on the right is zero, since it is the sum of the values of a nontrivial character, and \ref{VI:4-15} follows.
\end{proof}

Proposition~\ref{VI:4-15} has the following cohomological counterpart.

\begin{proposition_}\label{VI:4-16}
The cohomology of $W^\ast/\dG_m$ (\ref{VI:4-6}, \ref{VI:4-8}) with coefficients in $\sF(\chi^{-1})$ satisfies:
\begin{enumerate}[\indent (i)]
  \item $\h_c^i=0$ for $i\ne N-1$, and $\dim\h_c^{N-1}{}=1$.
  \item On $\h_c^{N-1}$, $F^\ast$ acts as multiplication by $J(\chi)$.
  \item $\h_c^{N-1}$ is canonically isomorphic to $\h_c^{N+1}\left(V^\ast,\sF(\chi^{-1},\psi\circ \tr_{k/\dF_q})\right)(1)$.
\end{enumerate}
\end{proposition_}

Assertion (ii) follows from (i) and the trace formula; (i) and (iii) follow from \ref{VI:4-11} and the following more geometric statement.

\begin{proposition_}\label{VI:4-17}
Let $(\chi_i)_{i\in I}$ be a family of nontrivial characters of $\dmu_n(\dF)$ whose product is $1$. Then $\h_c^{i-1}\left(W^\ast/\dG_m,\sK_n((\chi_i)_{i\in I})\right)$ is canonically isomorphic to \newline $\h_c^{i+1}\left(V^\ast,\sK_n((\chi_i)_{i\in I})\otimes \sF(\psi\tr_{k/\dF_q})\right)(1)$.
\end{proposition_}

The first line in the proof of \ref{VI:4-15} becomes, with $\pi$ as in \eqref{VI:eq:4-6-1},
\begin{equation*}\tag{4.17.1}\label{VI:eq:4-17-1}
  \eR^i \pi_!\left(\sK_n((\chi_i)_{i\in I})\otimes \sF(\psi \tr_{k/\dF_q})\right) = \sK_n((\chi_i)_{i\in I}) \otimes \eR^i \pi_! \sF(\psi\tr_{k/\dF_q})
\end{equation*}
(on $V^\ast/\dG_m$). Let us compute the sheaves $\eR^i\pi_!\sF(\psi\tr_{k/\dF_q}) =\eR^i\pi_!T^\ast\sF(\psi)$. Let $v_0:\widetilde V_0\to V_0$ be the blowup of $V_0$ at $\{0\}$. In the diagram
\begin{equation*}\tag{4.17.2}\label{VI:eq:4-17-2}
\xymatrix{
  V_0\setminus \{0\} \ar@{^{(}->}[r] \ar[dr]_-\pi
    & \widetilde V_0 \ar[r]^-{v_0} \ar[d]^-{\bar\pi}
    & V_0 \\
  & P_0
}
\end{equation*}
$\pi$ is a fibration whose fibers are punctured lines, $\widetilde V_0$ is the corresponding line bundle, and $Z_0=v_0^{-1}(0)$ is its zero section. The sheaf $T^\ast\sF(\psi)$ on $V_0\setminus\{0\}$ extends to $(Tv_0)^\ast\sF(\psi)$ on $\widetilde V_0$, and the restriction of this sheaf to $Z_0$ is the constant sheaf $E_\lambda$, since $T_0v_0$ vanishes on $Z_0$. The long exact sequence of cohomology with proper support is therefore
\begin{equation*}\tag{4.17.3}\label{VI:eq:4-17-3}
\xymatrix{
  \ar[r]
    & \eR^i \pi_! T^\ast \sF(\psi) \ar[r]
    & \eR^i \bar\pi_! (T v)^\ast \sF(\psi) \ar[r]
    & (E_\lambda\text{ for $i=0$, $0$ otherwise}) \ar[r]
    &
}
\end{equation*}

The inverse image $\bar\pi^{-1}(Q_0)$ is the blowup $\widetilde W_0$ of $W_0$ at $0$; $T_0v_0$ vanishes on $\widetilde W_0$. Thus $(Tv)^\ast\sF(\psi)=E_\lambda$ on $\bar\pi^{-1}(Q_0)$, and since $\bar\pi$ is a line bundle,
\begin{equation*}\tag{4.17.4}\label{VI:eq:4-17-4}
  \eR^i \bar\pi_! (T v)^\ast \sF(\psi)|Q =
    \begin{cases}
      0 & \text{for $i\ne 2$}, \\
      E_\lambda(-1) & \text{for $i=2$.}
    \end{cases}
\end{equation*}

If $x\in P$ and $x\notin Q$, the line $D=\bar\pi^{-1}(x)$ is mapped isomorphically onto $\dG_a$ by $Tv$; therefore, by \ref{VI:2-7}*,
\begin{align} \notag
  \h_c^\bullet(D,(T v)^\ast \sF(\psi)) &= \h^\bullet(\dG_a,\sF(\psi)) = 0 && \text{et} \\
  \tag{4.17.5}\label{VI:eq:4-17-5}
  \eR^i \bar\pi_! (T v)^\ast \sF(\psi)
    &= \begin{cases}
         0 & \text{for $i\ne 2$}, \\
         E_\lambda(-1)_Q & \text{for $i=2$}
       \end{cases}
\end{align}

Combining \eqref{VI:eq:4-17-3} and \eqref{VI:eq:4-17-5}, we finally obtain
\begin{equation*}\tag{4.17.6}\label{VI:eq:4-17-6}
  \eR^i \pi_! T^\ast\sF(\psi) =
    \begin{cases}
      0 & \text{for $i\ne 1,2$}, \\
      E_\lambda & \text{for $i=1$}, \\
      E_\lambda(-1)_Q & \text{for $i=2$.}
    \end{cases}
\end{equation*}

\subsection{}\label{VI:4-18}

Compute the cohomology with proper support of $\sK_n((\chi_i))\otimes T^\ast\sF(\psi)$ on $V^\ast$ using the Leray spectral sequence for $\pi:V^\ast\to V^\ast/\dG_m$. Applying \eqref{VI:eq:4-17-1} and \eqref{VI:eq:4-17-6}, the initial terms are $E_2^{p,1}=\h_c^p(V^\ast/\dG_m,\sK_n((\chi_i)))=0$ (because $(\chi_i)\ne0$) and $E_2^{p,2}=\h_c^p(W^\ast/\dG_m,\sK_n((\chi_i)))(-1)$.

The spectral sequence reduces to an isomorphism, and \ref{VI:4-17} follows.

\subsection{}\label{VI:4-19}

The sheaves $\sK_n((\chi_i)_{i\in I})$ make sense over any field, and even over any base scheme. This will allow us to generalize \ref{VI:4-16}(i). With the notation of \ref{VI:4-6}, suppose that $M$ is everywhere finite étale of rank $N$ over $A$, and let $n$ be an integer invertible in $A$. Denote by $a$ the projection from $W^\ast/\dG_m$ to $\spec(A)$. Let $\chi$ also be a character (a morphism of sheaves) $\chi:(V^\ast/\dG_m)_n\to E_\lambda^\times$, and let $\sK_n(\chi)$ be the corresponding sheaf. Locally for the étale topology, $\chi$ may be regarded as a family $(\chi_i)_{i\in I}$ of characters of $\dmu_n$ whose product is trivial.

\begin{proposition_}\label{VI:4-20}
\begin{enumerate}[(i)]
  \item If $\chi$ is nontrivial at every point, then $\eR^ia_!(\sK_n(\chi))=0$ for $i\ne N-1$, and $\eR^{N-1}a_!(\sK_n(\chi))$ is lisse of rank $1$.
  \item An automorphism $\sigma$ of $M$ that preserves $\chi$ acts on this $\eR^{N-1}a_!$ by multiplication by the sign $\varepsilon(\sigma)$ of $\sigma$, viewed as a permutation of $I$.
  \item If all the $\chi_i$ are nontrivial, then $\eR a_!\iso\eR a_\ast$.
\end{enumerate}
\end{proposition_}
\begin{proof}
The scheme $W^\ast/\dG_m$ is the complement of a relative normal-crossings divisor in $Q$, which is proper and smooth over $\spec(A)$, and $\sK_n(\chi)$ is locally constant on $W^\ast/\dG_m$ and tamely ramified at infinity. It follows that the $\eR^ia_!$ and $\eR^ia_\ast$ are lisse and that their formation commutes with every base change. A standard argument then reduces us to the case in which $A$ is a finite field, and (i) follows from \ref{VI:4-17} and \ref{VI:4-11}(i). For (ii), use the compatibility of the action of $\sigma$ with the isomorphism of \ref{VI:4-17}, together with \ref{VI:4-12}. For (iii), observe that if all the $\chi_i$ are nontrivial, then $\sK_n(\chi)$ on $W^\ast/\dG_m\subset Q$ is ramified along every divisor at infinity, and apply \ref{VI:1-19-1}.
\end{proof}

\section{Hecke characters}\label{VI:5}

\subsection{}\label{VI:5-1}

Let $F$ be a number field (of finite degree over $\dQ$), and let $k$ be a field of characteristic $0$. The following are several equivalent ways of specifying what it means for a homomorphism $\alpha:F^\times\to k^\times$ to be \emph{algebraic}.
\begin{enumerate}[(i)]
  \item Let $\{e_i\}$ be a basis of $F$ over $\dQ$. The homomorphism $\alpha$ is algebraic if it is given by a formula $\alpha(\sum x_i e_i)=A(x_i)$, with $A\in k(X_i)$.
\end{enumerate}

This means that on $F^\times$, $\alpha$ agrees with a rational map defined over $k$ from $\res_{F/\dQ}(\dG_m)$ to $\dG_m$. By the Zariski density of $F^\times$, and the fact that a rational homomorphism is everywhere defined, such a map is a homomorphism of group schemes:
\begin{enumerate}[(i)]
\setcounter{enumi}{1}
  \item $\alpha$ is induced by a homomorphism of $k$-group schemes
    \[
      \res_{F/\dQ}(\dG_m)\otimes_\dQ k \to \dG_m \text{.}
    \]
\end{enumerate}

If $k$ is a number field, the adjunction property of restriction of scalars shows that this is equivalent to:
\begin{enumerate}[(i)]
\setcounter{enumi}{2}
  \item (for $k$ a number field) $\alpha$ is induced by a homomorphism of $\dQ$-group schemes $\res_{F/\dQ}(\dG_m)\to\res_{k/\dQ}(\dG_m)$.
\end{enumerate}

Let $\bar k$ be an algebraic closure of $k$, and let $I=\hom(F,\bar k)$. Over $\bar k$, the character group of $\res_{F/\dQ}(\dG_m)$ is $\dZ^I$, with basis the embeddings of $F$ into $\bar k$. The characters defined over $k$ are those invariant under $\gal(\bar k/k)$. They may be described either as those mapping $F^\times$ into $k^\times\subset\bar k$, or in terms of the orbits of $\gal(\bar k/k)$ on $I$, which correspond to the factors of $F\otimes k$.
\begin{enumerate}[(i)]
\setcounter{enumi}{3}
  \item $\alpha$ has the form $\alpha=\prod_{\omega\in I}\omega^{n_\omega}$. The permitted families of exponents $(n_\omega)$ are those for which $n_\omega=n_{\omega'}$ whenever $\omega$ and $\omega'$ lie in the same $\gal(\bar k/k)$-orbit. Equivalently, they are those for which $\prod_\omega\omega(x)^{n_\omega}\in k$ for every $x\in F$.
  \item Write $F\otimes k=\prod_{j\in J}F_j$, with the $F_j$ fields. Then $\alpha$ has the form $\alpha=\prod N_{F_j/k}^{m_j}$.
\end{enumerate}

In the special case in which $k$ contains a normal closure of $F$, these expressions become $\alpha=\prod\omega^{n_\omega}$, where $\omega$ runs through the $[F:\dQ]$ embeddings of $F$ into $k$.

\subsection{}\label{VI:5-2}

Suppose that $k$ is a number field, and let $\alpha$ be an algebraic homomorphism from $F^\times$ to $k^\times$. There is then a unique homomorphism, again denoted $\alpha$, from the group $I(F)$ of fractional ideals of $F$ to that of $k$, such that $\alpha((x))=(\alpha(x))$. Uniqueness follows because every ideal has a power that is principal and $I(k)$ is torsion-free. For existence, one may use \ref{VI:5-1}(v): one has $\alpha(\fa)=\prod N_{F_j/k}^{m_j}(\fa)$.

\subsection{}\label{VI:5-3}

Recall the definition of algebraic Hecke characters, called by Weil Hecke characters (or Grössencharaktere) of type $A_0$. Let $F$ be a number field, let $\fm$ be an ideal of $F$ (that is, of its ring of integers), let $I_\fm$ be the group of fractional ideals of $F$ prime to $\fm$, and let $k$ be a field of characteristic $0$. A homomorphism $\chi:I_\fm\to k^\times$ is an \emph{algebraic Hecke character} (of conductor $\leqslant\fm$) if there is an algebraic homomorphism $\chi_{\mathrm{alg}}:F^\times\to k^\times$ satisfying
\begin{equation*}\tag{$*$}\label{VI:eq:5-3-1}
\parbox{0.86\textwidth}{For $x\in F^\times$ prime to $\fm$, totally
positive, and congruent to $1$ modulo $\fm$, one has
$\chi((x))=\chi_{\mathrm{alg}}(x)$.}
\end{equation*}

By density of the set of $x$ occurring in \eqref{VI:eq:5-3-1}, $\chi_{\mathrm{alg}}$ is completely determined by $\chi$. It is the \emph{algebraic part} of $\chi$. If $\chi((x))=\chi_{\mathrm{alg}}(x)$ for totally positive $x$ congruent to $1$ modulo $\fm'$, then $\chi$ extends to a Hecke character of conductor $\leqslant\inf(\fm,\fm')$: $I_{\fm+\fm'}\to k^\times$. We identify Hecke characters that agree on their common domain of definition, and call the \emph{conductor} of $\chi$ the smallest $\fm'$ for which $\chi$ has conductor $\leqslant\fm'$.

\subsection{Remark}\label{VI:5-4}

If an algebraic Hecke character $\chi$ takes its values in a subfield $k'$ of $k$, it is already a Hecke character with values in $k'$. It is enough to see that $\chi_{\mathrm{alg}}:\res_{F/\dQ}(\dG_m)\to\dG_m$ is already defined over $k'$, and this follows from the Zariski density in $\res_{F/\dQ}(\dG_m)$ of the totally positive elements $x$ congruent to $1$ modulo $\fm$. One may always take $k'$ to be a subfield of $k$ of finite degree over $\dQ$.

\subsection{}\label{VI:5-5}

If $\varepsilon$ is a totally positive unit congruent to $1$ modulo $\fm$, then $\chi_{\mathrm{alg}}(\varepsilon)=\chi((\varepsilon))=1$. Thus $\chi_{\mathrm{alg}}$ factors through the quotient $T_\fm$ of $\res_{F/\dQ}(\dG_m)$ by the Zariski closure of the group $E_\fm\subset F^\times$ of totally positive units congruent to $1$ modulo $\fm$.

According to Serre \cite[II.3]{se68}, if $k$ is an algebraic closure of $\dQ$ and $\fm$ is sufficiently large, the characters $\prod\omega^{n_\omega}$ of $\res_{F/\dQ}(\dG_m)$ that occur as $\chi_{\mathrm{alg}}$ are characterized as follows: there must be an integer $N$, the \emph{weight} of $\chi$ (or of $\chi_{\mathrm{alg}}$), such that for every element $\sigma$ of $\gal(k/\dQ)$ conjugate to complex conjugation, one has $n_\omega+n_{\sigma\omega}=N$. If $F_1$ is the largest subfield of $F$ that is a totally imaginary quadratic extension of a totally real field $F_1'$, this is equivalent to saying that $\chi_{\mathrm{alg}}$ has the form $\chi_1\circ N_{F/F_1}$ and that $\chi_1|_{F_1'^\times}=(N_{F_1'/\dQ})^N$.

If $\chi$ has weight $N$, then for every ideal $\fa$ of $F$, $\chi(\fa)$ is an algebraic number all of whose conjugates have absolute value $N(\fa)^{N/2}$ \cite[II.3, Prop.~2]{se68}.

\subsection{}\label{VI:5-6}

For a fixed number field $k$, further conditions are imposed on $\chi_{\mathrm{alg}}$. Indeed, for every ideal $\fa$ of $F$ prime to the conductor, one has
\begin{equation*}\tag{5.6.1}\label{VI:eq:5-6-1}
  \chi_\text{alg}(\fa) = (\chi(\fa)) \text{,}
\end{equation*}
so that $\chi_{\mathrm{alg}}(\fa)$ is principal. Since the group of fractional ideals of $k$ is torsion-free, it suffices to prove a suitable nonzero $n$th power of \eqref{VI:eq:5-6-1}. This allows us to replace $\fa$ by $\fa^n$, and hence to suppose that $\fa=(x)$ with $x$ totally positive and congruent to $1$ modulo $\fm$. The assertion then follows directly from the definitions.

Together with \ref{VI:5-5}, this shows that $\chi_{\mathrm{alg}}$ determines the absolute value of $\chi(\fa)$ at every place of $k$.

\subsection{}\label{VI:5-7}

Serre's group $S_\fm$ may be characterized as the group of multiplicative type whose character group (over any field) is the group of algebraic Hecke characters of conductor $\leqslant\fm$ (cf.\ \cite[II.2.1 and II.2.2]{se68}). Its relation with $\ell$-adic representations is explained in \cite[II.2.3]{se68}.

\begin{theorem_}[{\cite{se68}}]\label{VI:5-8}
Let $\chi$ be an algebraic Hecke character of $F$ with values in $E_\lambda$, where $E_\lambda$ is a finite extension of $\dQ_\ell$. There is then a unique homomorphism $\chi_\lambda:\gal(\bar F/F)^\textnormal{ab}\to E_\lambda^\times$ such that:
\begin{enumerate}[\indent (i)]
  \item $\chi_\lambda$ is unramified outside the conductor $\ff$ of $\chi$ and the primes above $\ell$;
  \item for a prime ideal $\fp$ of $F$ prime to $\ff$ and $\ell$, and $F_\fp\in\gal(\bar F/F)^\textnormal{ab}$ the geometric Frobenius at $\fp$, one has
    \[
      \chi_\lambda(F_\fp) = \chi(\fp) \text{.}
    \]
\end{enumerate}
\end{theorem_}

\subsection{}\label{VI:5-9}

In the remainder of this section we shall give a criterion for a $\lambda$-adic representation to arise in this way from a Hecke character, and in the next section we shall apply it to Jacobi sums. In this way we shall recover, in slightly generalized form, results of Weil \cite{we52,we74-2}, replacing the second part of Weil's proof by an argument using $\ell$-adic cohomology.

\begin{theorem_}\label{VI:5-10}
Let $F$ and $k$ be number fields, let $\lambda$ be a place of $k$ of residue characteristic $\ell$, and let $\chi_\lambda:\gal(\bar F/F)^\textnormal{ab}\to k_\lambda^\times$ be a homomorphism unramified outside $\ell$ and a finite set $S$ of places. Suppose that there exist a set $T$ of places, containing $S$ and the places above $\ell$, of density $0$, and an algebraic homomorphism $\chi_0:F^\times\to k^\times$, such that:
\begin{enumerate}[(i)]
  \item for an algebraic closure $\bar k$ of $k$, the homomorphism $\chi_0:F^\times\to k^\times\hookrightarrow\bar k^\times$ is of the type considered in \ref{VI:5-5}, of weight $N$;
  \item for every prime $\fp\notin T$, $\chi_\lambda(F_\fp)$ lies in $k^\times$, one has $(\chi_\lambda(F_\fp))=\chi_0(\fp)$, and every complex conjugate of $\chi_\lambda(F_\fp)$ has absolute value $(N\fp)^{N/2}$.
\end{enumerate}

Then $\chi_\lambda$ is defined, in the sense of \ref{VI:5-8}, by an algebraic Hecke character of $F$ with values in $k$ and algebraic part $\chi_0$.
\end{theorem_}

Let $\chi'$ be an algebraic Hecke character with values in a finite Galois extension $k'$ of $k$, and with algebraic part $\chi_0$ (\ref{VI:5-5}). After enlarging $T$, we may suppose that the conductor of $\chi'$ is supported in $T$. Let $\lambda'$ be a place of $k'$ above $\lambda$, let $\chi_{\lambda'}$ be the composite
\[
  \gal(\bar F/F)^\mathrm{ab}\longrightarrow k_\lambda^\times
  \longrightarrow {k'}_{\lambda'}^\times,
\]
and let $\chi_{\lambda'}':\gal(\bar F/F)^\mathrm{ab}\to {k'}_{\lambda'}^\times$ be the character defined by $\chi'$ (\ref{VI:5-8}). Set $\varepsilon=\chi_{\lambda'}({\chi'}_{\lambda'})^{-1}$.

Hypothesis (ii), together with \ref{VI:5-5}, \ref{VI:5-6}, and \ref{VI:5-8}, ensures that for $\fp\notin T$ one has $\varepsilon(F_\fp)\in{k'}^\times$, and that this number has absolute value $1$ in every completion of $k'$. It follows that $\varepsilon(F_\fp)$ belongs to the finite group $\mu$ of roots of unity in $k'$. By the Chebotarev density theorem and the density hypothesis on $T$, the elements $F_\fp$ for $\fp\notin T$ are dense in $\gal(\bar F/F)^\mathrm{ab}$. Thus $\varepsilon(\sigma)\in\mu$ for every $\sigma\in\gal(\bar F/F)^\mathrm{ab}$: $\varepsilon$ is a finite-order character $\gal(\bar F/F)^\mathrm{ab}\to\mu$. By class field theory, it corresponds to a finite-order character of the idèle class group of $F$, that is, to an algebraic Hecke character with trivial algebraic part. Correcting $\chi'$ by this character, we reduce to the case $\chi_{\lambda'}=\chi_{\lambda'}'$. It remains to show that $\chi'$, which a priori takes values in $k'$, actually takes values in $k$. If $\sigma\in\gal(k'/k)$, the characters $\chi'$ and ${\chi'}^\sigma$ agree on every prime ideal $\fp\notin T$, since
\[
  \chi'(\fp)=\chi'_{\lambda'}(F_\fp)
  =\chi_{\lambda'}(F_\fp)\in k^\times.
\]
The characters $\chi_{\lambda'}'$ and $(\chi_{\lambda'}')^\sigma$ agree on a dense subset of $\gal(\bar F/F)^\mathrm{ab}$, hence are equal. Thus $\chi'=(\chi')^\sigma$, so $\chi'$ takes values in $k$.

\section{Hecke characters defined by Jacobi sums}\label{VI:6}

The idea of this section is as follows. Consider a Gauss sum $\tau_{\dF_q}(\chi,\psi)$. Both it and the corresponding sheaf $\sF(\chi^{-1}\psi)$ are tied in two ways to the characteristic $p$ and to the base field $\dF_q$:
\begin{enumerate}[\indent a)]
  \item the additive character $\psi$ and the sheaf $\sF(\psi)$ occur in their definition;
  \item $\sF(\chi)$ is defined from the Lang exact sequence
    \[\xymatrix{
      0 \ar[r]
        & k^\times \ar[r]
        & V^\ast \ar[r]
        & V^\ast \ar[r]
        & 0 \text{.}
    }\]
\end{enumerate}
These give two difficulties if one wants to lift the corresponding cohomological constructions to a number field $F$ and find $\ell$-adic representations of $\gal(\bar F/F)$ whose Frobenius eigenvalues are Gauss sums.

If $\chi$ is trivial on $\dF_q^\times$, the sum $\tau_{\dF_q}(\chi,\psi)$ is independent of $\psi$ and is $q$ times a Jacobi sum. The latter is related to the cohomology of $W^\ast/\dG_m$ with values in a sheaf derived from $\chi$: $\psi$, and hence difficulty a), has disappeared. To resolve b), it is enough to describe the sheaves used by means of Kummer exact sequences, which make sense in every characteristic. This is possible by \ref{VI:4-9}.

Once the $\ell$-adic representations have been found, we use \ref{VI:5-10} and the Stickelberger theorem to show that they arise from algebraic Hecke characters.

\subsection{}\label{VI:6-1}

If $E$ is a field with algebraic closure $\bar E$, and $\lambda$ is a character of finite order dividing $d$, then for every multiple $n$ of $d$ we denote by ${}_n\lambda$ the character of $\mu_n(\bar E)$ that induces the same character as $\lambda$ on $\widehat\dZ(1)_{\bar E}=\varprojlim\mu_n(\bar E)$. We use the same notation when $\lambda$ is itself a finite-order character of $\widehat\dZ(1)_{\bar E}$.

Let $F$ and $k$ be number fields, let $\bar F$ be an algebraic closure of $F$, let $I$ be a finite set endowed with an action of $\gal(\bar F/F)$, and let $\lambda=(\lambda_i)_{i\in I}$ be a family of characters $\lambda_i:\widehat\dZ(1)_{\bar F}\to k^\times$. Suppose that $\lambda_i\ne1$ for every $i\in I$, that the product of the $\lambda_i$ is trivial, and that the family $\lambda$ is defined over $F$, that is, for every $\sigma\in\gal(\bar F/F)$, $\lambda_{\sigma(i)}=\lambda_i\circ\sigma^{-1}$.

The Galois set $I$ is the set of homomorphisms into $\bar F$ of a separable $F$-algebra $E$. If $C$ is the set of $\gal(\bar F/F)$-orbits in $I$, then $E$ decomposes as a product of fields $E=\prod_{\alpha\in C}E_\alpha$, where $\alpha$ is identified with the set of embeddings of $E_\alpha$ into $\bar F$. If the characters $\lambda_i$ for $i\in\alpha$ have order $d_\alpha$, then $E_\alpha$ contains the $d_\alpha$th roots of unity and there is a character $\lambda_\alpha:\mu_{d_\alpha}(E_\alpha)\hookrightarrow k^\times$ such that, if $i\in\alpha$ corresponds to the embedding $\sigma_i:E_\alpha\hookrightarrow\bar F$, then ${}_{d_\alpha}\lambda_i=\lambda_\alpha\circ\sigma_i^{-1}$.

\subsection{}\label{VI:6-2}

Let $n>0$ be a multiple of the $d_\alpha$, let $\Sigma$ be the set of places of $F$ at which $E/F$ ramifies or which divide $n$, let $\mu$ be a place of $k$ of residue characteristic $\ell$, and let $A$ be the ring of elements of $F$ integral outside $\Sigma$ and $\ell$. The integral closure $M$ of $A$ in $E$ is finite étale of rank $N=|I|$ over $A$. With the notation of \ref{VI:4-6} and \ref{VI:4-8}, the $_n\lambda_i$ define
\[
  _n\lambda:(V^\ast/\dG_m)_n \to k^\times \subset k_\mu^\times \text{.}
\]

Let $a$ be the projection from $W^\ast/\dG_m$ to $\spec(A)$. By \ref{VI:4-20}, the sheaf $\eR^{N-2}a_!\sK_n({}_n\lambda)$ is lisse of rank one; it defines an $\ell$-adic representation
\[
  j_\mu[\lambda]:\gal(\bar F/F)^\text{ab} \to k_\mu^\times
\]
(or simply $j_\mu$), unramified outside $\Sigma$ and $\ell$.

\subsection{}\label{VI:6-3}

Let us compute $j_\mu(F_\fp)$. By \ref{VI:4-20}, it is enough to do so after reduction modulo $\fp$. Set $f=A/\fp$ and $e=M/\fp M$, and let $\bar f$ be the algebraic closure of $f$ determined by a place of $\bar F$ above $\fp$. We still have $I=\hom_f(e,\bar f)$. After reduction, ${}_n\lambda$ again has a description as in \ref{VI:6-1}:
\begin{enumerate}[a)]
  \item Let $D$ be the set of $\gal(\bar f/f)$-orbits in $I$. The decomposition of $e$ as a product of fields is $e=\prod_{\beta\in D}e_\beta$, where the $\sigma_i$ for $i\in\beta$ are identified with the embeddings of $e_\beta$ into $\bar f$. For $\beta\in D$, let $\alpha(\beta)$ be the element of $C$ containing $\beta$, and put $d_\beta=d_{\alpha(\beta)}$. Define
    \[\xymatrix{
      \lambda_\beta : \mu_{d_\beta}(e_\beta)
        & \ar[l]_-\sim \mu_{d_\beta}(E_{\alpha(\beta)})
          \ar[r]^-{\lambda_{\alpha(\beta)}}
        & k^\times \text{.}
    }\]
  \item If $i\in I$ corresponds to the embedding $\sigma_i$ of $e_\beta$ into $\bar f$, then the reduction ${}_{d_\beta}\bar\lambda_i:\mu_{d_\beta}(\bar f)\to k^\times$ of ${}_{d_\beta}\lambda_i$ is $\lambda_\beta\circ\sigma_i^{-1}$.
\end{enumerate}

Let $q_\beta$ be the number of elements of $e_\beta$, set $\chi_\beta={}_{q_\beta-1}\lambda_\beta$, and let $\chi:e^\times\to k^\times$ have components $\chi_\beta$. On the affine space over $A/\fp$ defined by $e$ (\ref{VI:4-6}), or more precisely on the torus $e_\beta^\times$, there are isomorphisms
\[
\begin{split}
  \sK_n(({}_n\bar\lambda_i)_{i\in\beta})
  &=\sK_{d_\beta}
      ((\lambda_\beta\circ\sigma_i^{-1})_{i\in\beta})\\
  &=\sK_{q_\beta-1}
      ((\chi_\beta\circ\sigma_i^{-1})_{i\in\beta})
   =\sF(\chi_\beta)
\end{split}
\]
by \ref{VI:4-9}. Thus on $e^\times$ one has $\sK_n(({}_n\bar\lambda_i)_{i\in I})=\sF(\chi)$, and this isomorphism descends to the reduction modulo $\fp$ of $V^\ast/\dG_m$. Consequently:

\begin{proposition_}\label{VI:6-4}
With the notation of \ref{VI:6-2} and \ref{VI:6-3}, $j_\mu(F_\fp)$ is the Jacobi sum $J(\chi)$.
\end{proposition_}

Let us rewrite the formulas relating $(\lambda_i)_{i\in I}$ to $(\chi_\beta)_{\beta\in D}$. For $i\in\beta\subset\alpha$,
\begin{align*}
  {}_{d_\alpha}\lambda_i
    &=\lambda_\alpha\circ\sigma_i^{-1}
    &&\lambda_\alpha:\mu_{d_\alpha}(E_\alpha)\to k^\times,\\
  \lambda_\beta
    &=\text{the reduction of $\lambda_\alpha$}
    &&:\mu_{d_\alpha}(e_\beta)\to k^\times,\\
  \chi_\beta
    &={}_{q_\beta-1}\lambda_\beta:
      x\longmapsto\lambda_\beta(x^{(q_\beta-1)/d_\alpha})
    &&:e_\beta^\times\to k^\times.
\end{align*}

\begin{theorem_}\label{VI:6-5}
The representation $j_\mu$ is defined by an algebraic Hecke character of $F$ with values in $k^\times$.
\end{theorem_}

We shall verify this using \ref{VI:5-10}. At this point we rejoin Weil's proof \cite{we52,we74-2}, through the use of the Stickelberger theorem, and I shall give only a few indications.

\subsection{}\label{VI:6-6}

Given $j:\gal(\bar F/F)\to k_\lambda^\times$ and a finite extension $F'/F$, one checks using \ref{VI:5-10} that $j$ is defined by a Hecke character if and only if its restriction to $\gal(\bar F/F')$ is. This observation allows us in \ref{VI:6-5} to reduce to the case in which Galois acts trivially on $I$ and $F$ contains the $n$th roots of unity. We may then reduce to $F=\dQ(\sqrt[n]{1})$ and take $k=F$. In this case each ${}_n\lambda_i$ is specified by an element $a_i\in\dZ/n$: it is the map $x\mapsto x^{a_i}:\mu_n\to k^\times$. One has $\sum a_i=0$ in $\dZ/n$, and $a_i\ne0$.

\subsection{}\label{VI:6-7}

In a number field, the places of degree greater than $1$ always form a set of density $0$. In applying \ref{VI:5-10}, we may therefore disregard them and use Stickelberger only over a prime field. Let $p$ be prime, let $0<a<p-1$, let $\widetilde{x}:\dF_p^\times\to\dQ_p^\times$ be the Teichmüller (``multiplicative lift'') character, and set
\[
  g=-\sum_{x\in\dF_p^\times}\widetilde{x}^{-a}\zeta^x
    \in\dQ_p(\zeta),
\]
where $\zeta$ is a $p$th root of unity. Also write $x$ for the integer in $[0,p-1]$ lifting $x$, put $\zeta=1+\pi$, and expand $\zeta^x=(1+\pi)^x$ by the binomial formula. We obtain $g=-\sum_{i=0}^{p-1}A_i\pi^i$, where $A_i\in\dZ_p$ and the reduction of $A_i$ modulo $p$ is $\sum_{x\in\dF_p^\times}x^{-a}\binom{x}{i}$. If $b$ is not divisible by $p-1$, then $\sum_{x\in\dF_p^\times}x^b=0$. If $i<a$, $\binom{x}{i}$ is a polynomial in $x$ of degree $i<a$, whence $A_i\equiv0\pmod p$. Likewise $A_a\equiv(p-1)/a!\pmod p$, and $g\sim\pi^a/a!$. Thus the valuation of $g$ is given by
\[
  \frac{v(g)}{v(p)}=\frac{a}{p-1}.
\]

\subsection{}\label{VI:6-8}

Let $\sigma_i$, for $i\in(\dZ/n)^\times$, be the element of $\gal(\dQ(\sqrt[n]{1})/\dQ)$ such that $\sigma_i(\zeta)=\zeta^i$ for $\zeta^n=1$, and write the group of algebraic homomorphisms $\dQ(\sqrt[n]{1})^\times\to\dQ(\sqrt[n]{1})^\times$ additively. For $\lambda$ as in \ref{VI:6-6}, defined by a family $(a_j)_{j\in I}$ of nonzero integers modulo $n$ whose sum is zero, \ref{VI:5-10}, \eqref{VI:eq:4-15-1}, and \ref{VI:6-7} show that $j_\mu[\lambda]$ is defined by an algebraic Hecke character $j[a]$ and that, if $N$ denotes the norm, its algebraic part is given by
\begin{equation*}\tag{6.8.1}\label{VI:eq:6-8-1}
  (N j[a])_{\mathrm{alg}}
  =\sum_{i\in(\dZ/n)^\times}\sum_{j\in I}
    \left\langle\frac{i a_j}{n}\right\rangle\sigma_i^{-1},
\end{equation*}
where $\langle\,\cdot\,\rangle$ denotes the fractional part.

\subsection{}\label{VI:6-9}

The most interesting case is that in which $F\subset\dQ(\sqrt[n]{1})$ and the family of ${}_n\lambda_i$ is constructed as follows:
\begin{enumerate}[a)]
  \item Take a finite family $(A_j)_{j\in J}$ of orbits in $\dZ/n$ under $\gal(\dQ(\sqrt[n]{1})/F)=H\subset(\dZ/n)^\times$. Suppose that $A_j\ne\{0\}$ and $\sum_{j\in J}\sum_{a\in A_j}a=0$ in $\dZ/n$.
  \item Take $I=\coprod_{j\in J}A_j$; Galois acts on each $A_j$.
  \item Take $\bar F\supset\dQ(\sqrt[n]{1})$, $k=\dQ(\sqrt[n]{1})$, and, if $i\in I$ corresponds to $a\in A_j\subset\dZ/n$, set ${}_n\lambda_i(x)=x^a$.
\end{enumerate}

Using \ref{VI:4-20}, one checks that a general $j_\mu$ relative to $F_1$ is obtained from a $j_\mu$ as above, for $F\subset\dQ(\sqrt[n]{1})$ and $F_1/F$ finite, by restriction to $\gal(\bar F/F_1)$ and multiplication by a character of order two.

An element $a\in\gal(k/\dQ)=(\dZ/n)^\times$ transforms the family $(A_j)$ into $(aA_j)$. If $a\in H$, the family $(A_j)$ is therefore preserved, and the values $j_\mu(F_\fp)$ are invariant under $H$.

\begin{proposition_}\label{VI:6-10}
Under the hypotheses of \ref{VI:6-9}, $j_\mu$ is defined by an algebraic Hecke character of $F$ with values in $F$.
\end{proposition_}

Return to $F=\dQ(\sqrt[n]{1})$, and consider the Hecke characters $j[a]$ defined by a family $(a_i)_{i\in I}$ with $a_i\in\dZ/n$, $a_i\ne0$, and $\sum a_i=0$. The following theorem was obtained independently by Vishik.

\begin{theorem_}\label{VI:6-11}
Every algebraic Hecke character of $\dQ(\sqrt[n]{1})$ has a nonzero power in the group generated by the $j[a]$ and the norm.
\end{theorem_}

Since we are working up to torsion, only the algebraic part of the Hecke characters under consideration matters. Applying \eqref{VI:eq:6-8-1} to the $j[a]$, with $a$ an element of $\dZ/n$ repeated $n$ times, reduces the assertion to the following lemma.

\begin{lemma_}\label{VI:6-12}
Write the elements of the group algebra $\dQ[(\dZ/n)^\times]$ as $\sum x_j\sigma_j$, and let $X$ be the subspace of those for which $x_j+x_{-j}$ is constant. Then $X$ is generated over $\dQ$ by $N=\sum_j\sigma_j$ and by
\[
  g_a=\sum_{i\in(\dZ/n)^\times}
      \left\langle\frac{ia}{n}\right\rangle\sigma_i^{-1}
  \qquad(a\in\dZ/n,\ a\ne0).
\]
\end{lemma_}

Both $X$ and $Y=\langle N,(g_a)\rangle$ are ideals of the group algebra, and it is enough to verify that they are annihilated by the same complex characters $\chi$ of $(\dZ/n)^\times$. This amounts to saying that an odd character, for which $\chi(-j)=-\chi(j)$, never annihilates all the $g_a$. If $\chi$ is primitive, then
\[
  \chi(g_1)=\sum_{i\in(\dZ/n)^\times}\frac{i}{n}\chi^{-1}(i)
  =-L(\chi,0)\ne0.
\]
In general, if $\chi$ is induced by a primitive character modulo $d$, then $\chi(g_{n/d})$ reduces to the analogous sum over $(\dZ/d)^\times$. This proves the lemma.

\section{Generalized Kloosterman sums}\label{VI:7}

\subsection{}\label{VI:7-1}

Let $\dF_q$ be a finite field with $q$ elements, and let $\psi$ be a nontrivial additive character. In this section we make a cohomological study of the generalized Kloosterman sums
\begin{equation*}\tag{7.1.1}\label{VI;eq:7-1-1}
  \kloos_{n,a} = \sum_{x_1\dotsm x_n = a} \psi(x_1+\cdots + x_n) \qquad (n\geqslant 1) \text{.}
\end{equation*}
For $a=0$, this sum is elementary (cf.\ \ref{VI:7-7}):
\begin{equation*}\tag{7.1.2}\label{VI:eq:7-1-2}
  \kloos_{n,0} = (-1)^{n-1} \text{.}
\end{equation*}
The interesting case is $a\ne0$. The $x_i$ then belong to $\dF_q^\times$. In this case we shall prove the bound
\begin{equation*}\tag{7.1.3}\label{VI:eq:7-1-3}
  \left|\kloos_{n,a}\right|\leqslant n q^{\frac{n-1}{2}}.
\end{equation*}

Our essential tools will be cohomological statements reflecting the following identities:
\begin{align*}\tag{7.1.4}\label{VI:eq:7-1-4}
  \kloos_{n,a} &= \sum_{x_1} \psi(x_1) \sum_{x_2\dotsm x_n = a/x_1} \psi(x_2+\cdots + x_n) \\ \notag
    &= \sum_{x\in \dF_q^\times} \psi(x) \kloos_{n-1,a/x} \qquad (a\ne 0,n\geqslant 2) \\ \tag{7.1.5}\label{VI:eq:7-1-5}
  \sum_a \kloos_{n,a} &= \sum_a \sum_{x_1 \dotsm x_n = a} \psi\left(\sum x_i\right) = \sum_{x_1,\dots,x_n} \psi\left(\sum x_i\right) = 0
\end{align*}
For $\chi$ a nontrivial character of $\dF_q^\times$,
\begin{equation*}\tag{7.1.6}\label{VI:eq:7-1-6}
  \sum \chi(a) \kloos_{n,a} = \sum_{x_1,\dots,x_n} \chi\left(\prod x_i\right) \cdot \psi\left(\sum x_i\right) = (-\tau(\chi^{-1},\psi))^n
\end{equation*}
(a Gauss sum).

\subsection{}\label{VI:7-2}

Let $k$ be an étale algebra of degree $n$ over $\dF_q$, and let $\varepsilon(k)$ be as in \eqref{VI:eq:4-5-1}. Set
\[
  \kloos_{k,a} = \sum_{N_{k/\dF_q}(x)=a} \psi \tr(x) \text{.}
\]
One again has
\begin{align*}\tag{7.2.1}\label{VI:eq:7-2-1}
  \kloos_{k,0} &= (-1)^{n-1} \cdot \varepsilon(k) \\ \tag{7.2.2}\label{VI:eq:7-2-2}
  \sum_a \kloos_{k,a} &= 0 \\ \tag{7.2.3}\label{VI:eq:7-2-3}
  \sum_a \chi(a) \kloos_{k,a} &= (-1)^n \tau_{\dF_q}\left(\chi^{-1}\circ N_{k/\dF_q},\psi\right) \text{.}
\end{align*}

Fourier inversion on $\dF_q^\times$ gives an expression for $\kloos_{k,a}$ in terms of Gauss sums:
\begin{align*}
  \kloos_{k,a} &= \frac{1}{q-1} \sum_\chi \chi(a) \sum_{b\in \dF_q^\times} \chi^{-1}(b) \kloos_{k,b} \\ \tag{7.2.4}\label{VI:eq:7-2-4}
  (-1)^n \kloos_{k,a} &= \frac{1}{q-1} \left(\varepsilon(k) + \sum_{\chi\ne 1} \chi(a) \tau_{\dF_q}\left(\chi\circ N_{k/\dF_q},\psi\right) \right) \text{.}
\end{align*}

The Hasse--Davenport identity now allows us to compare $\kloos_{k,a}$ with $\kloos_{n,a}=\kloos_{\dF_q^n,a}$:
\begin{equation*}\tag{7.2.5}\label{VI:eq:7-2-5}
  \kloos_{k,a} = \varepsilon(k) \kloos_{n,a} \text{.}
\end{equation*}

I know no cohomological interpretation of \eqref{VI:eq:7-2-4}, but I shall give one for \eqref{VI:eq:7-2-5}. Let us return to the sums $\kloos_{n,a}$.

\subsection{}\label{VI:7-3}

Let $\dF$ be an algebraic closure of $\dF_q$, and for $a\in\dF$ let $V_a^{n-1}$, or simply $V_a$, be the hypersurface in $\dA^n$ defined by $x_1\cdots x_n=a$. As usual, we regard $\psi$ as $\ell$-adic-valued rather than complex-valued. With the notation of \ref{VI:1-8}(ii), the bound \eqref{VI:eq:7-1-3} follows from the following theorem.

\begin{theorem_}\label{VI:7-4}
The cohomology of $V_a^{n-1}$ with coefficients in $\sF(\psi(\sum x_i))$ satisfies:
\begin{enumerate}[\indent (i)]
  \item $\h_c^i=0$ for $i\ne n-1$;
  \item $\h_c^\bullet{} \iso \h^\bullet{}$;
  \item for $a\ne0$, $\dim\h_c^{n-1}{}=n$;
  \item for $a=0$, $\h_c^{n-1}{}$ is canonically isomorphic to $E_\lambda$.
\end{enumerate}
\end{theorem_}

\subsection{}\label{VI:7-5}

As usual, this theorem also controls the dependence of $\kloos_{n,a}$ on $q$: for each $a\in\dF_q^\times$, there are $n$ Frobenius eigenvalues $\alpha_1,\dots,\alpha_n$, each of complex absolute value $q^{(n-1)/2}$, such that
\[
  \sum_{\substack{x_1\dotsm x_n = a \\ x_i\in \dF_{q^m}}} \psi\circ \tr\left(\sum x_i\right) = (-1)^{n-1} \left(\alpha_1^m + \cdots + \alpha_n^m \right) \text{.}
\]

For even $n$, $x\mapsto-x$ is an involution of $V_a$ that transforms $\sF(\psi\sigma)$ into its dual. Arguing as in \ref{VI:3-6}, we obtain a canonical nondegenerate alternating form on $\h_c^{n-1}(V_a,\sF(\psi\sigma))$ with values in $\dQ_\ell(-(n-1))$, and hence the following corollary.

\begin{corollary_}\label{VI:7-6}
With the notation of \ref{VI:7-5}, if $n$ is even, the $\alpha_i$ fall into $n/2$ pairs $\alpha,\bar\alpha$ whose product is $q^{n-1}$.
\end{corollary_}

\subsection{}\label{VI:7-7}

Let us verify the case $a=0$ of \ref{VI:7-4}. The hypersurface $V_0$ is the union of the coordinate hyperplanes $x_i=0$. Consider the Leray spectral sequence \eqref{VI:eq:2-6-1*} for this covering of $V_0$, for cohomology either with or without support and with coefficients in $\sF(\psi\sigma)$. By \ref{VI:2-7*}*, all the initial terms vanish except the cohomology of the intersection $\{0\}$ of all the coordinate hyperplanes. What remains is an isomorphism $\h_c^0\iso\h^0=E_\lambda$, which proves this case of \ref{VI:7-4}.

We shall prove \ref{VI:7-4} and \ref{VI:7-8} below by simultaneous induction on $n$, using the spectral sequence reflected by \eqref{VI:eq:7-1-4}. This requires control of the dependence of $\kloos_{n,a}$ on $a$.

Let $\pi:\dA^n\to\dA^1$ be the coordinate-product map, and let $\sigma$ be the sum of the coordinates. The hypersurfaces $V_a$ are the fibers of $\pi$.

\begin{theorem_}\label{VI:7-8}
\begin{enumerate}[(i)]
  \item The sheaf $\eR^{n-1}\pi_!\sF(\psi\sigma)$ is lisse of rank $n$ on $\dA^1\setminus\{0\}$.
  \item Its extension by zero to $\dP^1$ equals the direct image of its restriction to $\dA^1\setminus\{0\}$.
  \item At $0$, the monodromy is unipotent with a single Jordan block.
  \item At $\infty$, wild inertia has no nonzero fixed vector, and the Swan conductor is $1$.
  \item One has $\eR^i\pi_!\sF(\psi\sigma)\iso\eR^i\pi_\ast\sF(\psi\sigma)$; these sheaves vanish for $i\ne n-1$.
\end{enumerate}
\end{theorem_}
\begin{proof}
We first prove that \ref{VI:7-4} for a given value of $n$ implies \ref{VI:7-8} for the same value of $n$.

The stalk of $\eR^i\pi_!\sF(\psi\sigma)$ at $a$ is $\h_c^i(V_a,\sF(\psi\sigma))$. On the dense open subset on which the $\eR^i\pi_!\sF(\psi^{-1}\sigma)$ are locally constant, the stalk of $\eR^i\pi_\ast\sF(\psi\sigma)$ is $\h^i(V_a,\sF(\psi\sigma))$ (\hyperref[VII]{Finiteness Theorem}, \ref{VII:2-1}). Hypothesis \ref{VI:7-4}($n$) therefore gives the following lemma.

\begin{lemma_}\label{VI:7-9}
\begin{enumerate}[(i)]
  \item $\eR^i\pi_!\sF(\psi\sigma)=0$ for $i\ne n-1$.
  \item For $i=n-1$, the stalk of this sheaf at every point $a\ne0$ has constant rank $n$. At $0$, it has rank $1$.
  \item On a dense open subset $U$, one has $\eR^i\pi_!\sF(\psi\sigma)\iso \eR^i\pi_\ast\sF(\psi\sigma)$.
\end{enumerate}
\end{lemma_}

\subsection{}\label{VI:7-10}

We now use the Leray spectral sequence for $\pi$, for cohomology with or without support and with coefficients in $\sF(\psi\sigma)$. Its abutment is known: $\h_c^\bullet(\dA^n,\sF(\psi\sigma)) =\h^\bullet(\dA^n,\sF(\psi\sigma))=0$ (\ref{VI:2-7*}*). For cohomology with proper support, \ref{VI:7-9}(i) ensures that $E_2^{pq}=0$ for $q\ne n-1$. All the $E_2$-terms therefore vanish:
\begin{equation*}\tag{7.10.1}\label{VI:eq:7-10-1}
  \h_c^p(\dA^1,\eR^{n-1} \pi_!\sF(\psi\sigma)) = 0 \text{.}
\end{equation*}

For $p=0$, this means that $\eR^{n-1}\pi_!\sF(\psi\sigma)$ has no section with point support. In view of the constancy of rank in \ref{VI:7-9}(ii), it follows that
\begin{equation*}\tag{7.10.2}\label{VI:eq:7-10-2}
  \begin{array}{c}
    \text{The sheaf $\eR^{n-1}\pi_!\sF(\psi\sigma)$ on $\dA^1$ is lisse on
    $\dA^1\setminus\{0\}$ and is}\\
    \text{a subsheaf of the direct image $\sG$ of its restriction to
    $\dA^1\setminus\{0\}$.}
  \end{array}
\end{equation*}

The long exact cohomology sequence of
\[\xymatrix{
  0 \ar[r]
    & \eR^{n-1} \pi_! \sF(\psi\sigma) \ar[r]
    & \sG \ar[r]
    & \sQ \ar[r]
    & 0 \text{.}
}\]
where the support of $\sQ$ is concentrated at $0$, gives
\[\xymatrix{
  \h_c^0(\dA^1,\sG) \ar[r]
    & \sQ_0 \ar[r]
    & \h_c^1(\dA^1,\eR^{n-1} \pi_! \sF(\psi\sigma)) \text{.}
}\]
The two outer terms vanish, so $\sQ=0$, and
\begin{equation*}\tag{7.10.3}\label{VI:eq:7-10-3}
  \text{$\eR^{n-1}\pi_!\sF(\psi\sigma)$ is the direct image of its
  restriction to $\dA^1\setminus\{0\}$.}
\end{equation*}

Let $j$ be the inclusion of $\dA^1$ into $\dP^1$. The derived-category form of the Leray spectral sequences is
\begin{align*}
  \hh^\bullet(\dP^1,j_! \eR\pi_! \sF(\psi\sigma)) &= \hh_c^\bullet(\dA^1,\eR \pi_! \sF(\psi\sigma)) = \h_c^\bullet(\dA^n,\sF(\psi\sigma)) = 0 \\
  \hh^\bullet(\dP^1,\eR j_\ast \eR \pi_\ast \sF(\psi\sigma)) &= \hh^\bullet(\dA^1,\eR \pi_\ast \sF(\psi\sigma)) = \h^\bullet(\dA^n,\sF(\psi\sigma)) = 0 \text{.}
\end{align*}
Let $\Delta$ be the mapping cone of $j_!\eR\pi_!\sF(\psi\sigma)\to \eR j_\ast\eR\pi_\ast\sF(\psi\sigma)$. By \ref{VI:7-9}(iii), the cohomology sheaves of this complex have finite support. On the other hand, the long exact cohomology sequence of the triangle $(j_!\eR\pi_!\sF(\psi\sigma), \eR j_\ast\eR\pi_\ast\sF(\psi\sigma),\Delta)$ shows that $\hh^\bullet(\dP^1,\Delta)=0$. Since $\hh^\bullet(\dP^1,\Delta) =\h^0(\dP^1,\sH^\bullet(\Delta))$, it follows that $\Delta=0$:
\[\xymatrix{
  j_! \eR\pi_! \sF(\psi\sigma) \ar[r]^-\sim
    & \eR j_\ast \eR \pi_\ast \sF(\psi\sigma) \text{.}
}\]

In particular, $\eR\pi_!\sF(\psi\sigma)\iso\eR\pi_\ast\sF(\psi\sigma)$; and since these complexes have only one nonzero cohomology sheaf,
\begin{equation*}\tag{7.10.4}\label{VI:eq:7-10-4}
\xymatrix{
  j_! \eR \pi_! \sF(\psi\sigma) \ar[r]^-\sim
    & j_\ast \eR \pi_! \sF(\psi\sigma) \text{.}
}
\end{equation*}
This completes the proof of (i), (ii), and (v).

\subsection{}\label{VI:7-11}

Let $\swan_0$ and $\swan_\infty$ be the Swan conductors at $0$ and $\infty$ of $\eR^{n-1}\pi_!\sF(\psi\sigma)$. By \eqref{VI:eq:7-10-1}, $\chi_c(\eR^{n-1}\pi_!\sF(\psi\sigma))=0$, and the Euler--Poincaré formula \eqref{VI:eq:3-2-1} reduces to
\begin{equation*}\tag{7.11.1}\label{VI:eq:7-11-1}
  \swan_0 + \swan_\infty =1 \text{:}
\end{equation*}
one of these conductors is $0$ (tame ramification), and the other is $1$.

\subsection{}\label{VI:7-12}

Let $\sG$ be a nonconstant rank-one sheaf on $\dG_m$ of Kummer type \ref{VI:4-7}, say $\sG=\sK_n(\chi)$ with $\chi\ne1$. The sheaf $\sF(\psi\sigma)\otimes\pi^\ast\sG$ on $V^\ast=\pi^{-1}(\dG_m)=\dA^n\setminus\pi^{-1}(0)$ is then of the type encountered in the study of Gauss sums, and \ref{VI:4-11}(ii) gives
\begin{equation*}\tag{7.12.1}\label{VI:eq:7-12-1}
\xymatrix{
  \h_c^\bullet(V^\ast,\pi^\ast\sG\otimes \sF(\psi\sigma)) \ar[r]^-\sim
    & \h^\bullet(V^\ast,\pi^\ast \sG\otimes \sF(\psi\sigma)) \text{.}
}
\end{equation*}
The Leray spectral sequence for $\pi$ transforms this statement into
\begin{equation*}\tag{7.12.2}\label{VI:eq:7-12-2}
\xymatrix{
  \h_c^\bullet(\dG_m,\sG\otimes \eR^{n-1} \pi_! \sF(\psi\sigma)) \ar[r]^-\sim
    & \h^\bullet(\dG_m,\sG\otimes \eR^{n-1} \pi_\ast \sF(\psi\sigma)) \text{.}
}
\end{equation*}

Let $i$ be the inclusion of $\dG_m$ into $\dP^1$, and let $\Delta$ be the mapping cone of
\[
  \delta:i_!\bigl(\sG\otimes\eR^{n-1}\pi_!\sF(\psi\sigma)\bigr)
  \longrightarrow
  \eR i_\ast\bigl(\sG\otimes\eR^{n-1}\pi_!\sF(\psi\sigma)\bigr),
\]
where $\eR^{n-1}\pi_!=\eR^{n-1}\pi_\ast$. Arguing as in \ref{VI:7-10}, we deduce from \eqref{VI:eq:7-12-2} that $\Delta=0$. In particular,
\begin{equation*}\tag{7.12.3}\label{VI:eq:7-12-3}
  \text{At $0$ and $\infty$, inertia has no nonzero fixed vector on
  $\sG\otimes\eR^{n-1}\pi_!\sF(\psi\sigma)$.}
\end{equation*}

At $\infty$, this holds even when $\sG$ is trivial, by \eqref{VI:eq:7-10-4}; varying the tame Kummer twist shows that wild inertia has no fixed vector. In particular the ramification is wild. By \ref{VI:7-11}, one has $\swan_\infty=1$, and the ramification at $0$ is tame.

At $0$, \eqref{VI:eq:7-12-3} forces the tame ramification of $\eR^{n-1}\pi_!\sF(\psi\sigma)$ to be unipotent. By \eqref{VI:eq:7-10-3} and \ref{VI:7-9}(ii), there is a single Jordan block. This completes the proof of \ref{VI:7-8}($n$).

\subsection{}\label{VI:7-13}

For $n=1$, \ref{VI:7-4} is immediate. To complete the proof of \ref{VI:7-4} and \ref{VI:7-8}, it remains to prove that \ref{VI:7-8} for a given $n$ implies \ref{VI:7-4} for $n+1$. The case $a=0$ was treated in \ref{VI:7-7}, so assume $a\ne0$. Write $x_0,\dots,x_n$ for the coordinates of $\dA^{n+1}$. Let $g=x_0:V_a\to\dG_m$, let $\tau$ be the involution $x\mapsto ax^{-1}$ of $\dG_m$, and, as above, let $\pi:\dA^n\to\dA^1$ be the product of the coordinates $x_1,\dots,x_n$, with $\sigma$ their sum. We use $\sF(\psi\sigma)$ to denote the sheaf $\sF(\psi(\sum x_i))$ on both $\dA^{n+1}$ and $\dA^n$, as the context dictates.

\subsection{}\label{VI:7-14}

Write down the Leray spectral sequence for $g$. Since $\sF(\psi\sigma)$ on $\dA^{n+1}$ is the external tensor product of $\sF(\psi)$ on $\dA^1$ and $\sF(\psi\sigma)$ on $\dA^n$, one obtains (cf.\ \eqref{VI:eq:7-1-4})
\begin{align*}
  'E_2^{p q} &= \h_c^p(\dG_m,\sF(\psi)\otimes \tau^\ast \eR^q \pi_! \sF(\psi\sigma)) \Rightarrow \h_c^{p+q}(V_a,\sF(\psi\sigma)) \\
  ''E_2^{p q} &= \h^p(\dG_m,\sF(\psi)\otimes \tau^\ast \eR^q \pi_\ast \sF(\psi\sigma)) \Rightarrow \h^{p+q}(V_a,\sF(\psi\sigma)) \text{.}
\end{align*}
By hypothesis, $\eR^q\pi_!\sF(\psi\sigma)\iso\eR^q\pi_\ast\sF(\psi\sigma)$, and this sheaf vanishes for $q\ne n-1$. The sheaf $\sF(\psi)$ on $\dG_m$ is wildly ramified at $\infty$ and unramified at $0$, whereas $\tau^\ast\eR^q\pi_!\sF(\psi\sigma)$ is wildly ramified at $0$, with no wild-inertia invariants and Swan conductor $1$, and is tamely ramified at $\infty$. Their tensor product is therefore:
\begin{enumerate}[\indent a)]
  \item wildly ramified at $0$ and $\infty$, with no wild-inertia invariants;
  \item of Swan conductor $1$ at $0$ and $n$, the rank of $\eR^q\pi_!\sF(\psi\sigma)$, at $\infty$.
\end{enumerate}
Thus $'E_2^{pq}\iso{}''E_2^{pq}$, and $'E_2^{pq}=0$ except for $p=1$ and $q=n-1$. Finally, the Euler--Poincaré formula gives
\[
  \dim{}'E_2^{1,n-1}=1+n.
\]
This completes the proof.
\end{proof}

\subsection{Remark}\label{VI:7-15}

The isomorphism obtained in \ref{VI:7-14},
\begin{equation*}\tag{7.15.1}\label{VI:eq:7-15-1}
  \h_c^n(V_a^n,\sF(\psi\sigma)) = \h_c^1(\dG_m,\sF(\psi)\otimes \tau^\ast \eR^{n-1} \pi_! \sF(\psi\sigma))
\end{equation*}
allows one, with the notation of \ref{VI:7-5}, to compute the product $\det(F)$ of the Frobenius eigenvalues $\alpha_i$. The formalism of \cite{de73} applies, since $\sF(\psi)$ and $\eR^{n-1}\pi_!\sF(\psi\sigma)$ belong to infinite compatible systems of $\ell$-adic representations. This reduces the calculation to local questions, which may be simplified by observing that (cf.\ \cite[9.5]{de73}):
\begin{enumerate}[\indent a)]
  \item the global constants for $\sF(\psi)$ and $\tau^\ast\eR^{n-1}\pi_!\sF(\psi\sigma)$ are easy to compute, since the cohomology of these sheaves is of a trivial nature;
  \item at every point, for either sheaf, the corresponding $\ell$-adic representation of the decomposition group has unramified semisimplification.
\end{enumerate}

For the sum in $n$ variables, the result is
\begin{equation*}\tag{7.15.2}\label{VI:eq:7-15-2}
  \det\left(F,\h_c^{n-1}(V_a^{n-1},\sF(\psi))\right) = q^{\frac{n(n-1)}{2}} \text{.}
\end{equation*}

\subsection{Remark}\label{VI:7-16}

For $p=2$, the sheaf $\sF(\psi)$ is orthogonal. For $a\ne0$ and odd $n$, \ref{VI:7-4}(ii) and Poincaré duality therefore give a symmetric bilinear form on $\h_c^{n-1}(V_a^{n-1},\sF(\psi\sigma))$ with values in $\bar\dQ_\ell(1-n)$. Relative to this form, $F$ is an orthogonal similitude with multiplier $q^{n-1}$, and $q^{(1-n)/2}F$ belongs to the special orthogonal group by \eqref{VI:eq:7-15-2}. One of the eigenvalues of $F$ is therefore $q^{(n-1)/2}$, and the others fall into pairs $\alpha,\bar\alpha$ whose product is $q^{n-1}$.

For $n=3$, L.\ Carlitz \cite{ca69-2} obtained a more precise result, equivalent to the following proposition.

\begin{proposition_}\label{VI:7-17}
For $p=2$ and $\psi(x)=(-1)^{\tr_{\dF_q/\dF_2}(x)}$, there is an isomorphism
\[
  \h_c^2(V_a^2,\sF(\psi\sigma)) = \operatorname{Sym}^2 \h_c^1(V_a^1,\sF(\psi\sigma)) \text{.}
\]
\end{proposition_}

The tensor square of $\h_c^1(V_a^1,\sF(\psi\sigma))$ is $\h_c^2(V_a^1\times V_a^1, \sF(\psi\sigma)\boxtimes\sF(\psi\sigma))$ by Künneth. The symmetry $x\otimes y\mapsto y\otimes x$ is represented geometrically by $-\tau^\ast$, where $\tau$ is the automorphism $(x,y)\mapsto(y,x)$ of $V_a^1\times V_a^1$. Let $X'$ be the quotient of $V_a^1\times V_a^1$ by $\{\operatorname{id},\tau\}$, and let $\pi$ be the projection onto $X'$. The group $\operatorname{Sym}^2\h_c^1(V_a^1,\sF(\psi\sigma))$ is therefore identified with the cohomology of $X'$ with coefficients in the $\tau$-anti-invariant part of $\pi_\ast(\sF(\psi\sigma)\boxtimes\sF(\psi\sigma))$. Let $s$ be the function on $X'$ such that $s\pi=\sigma\pr_1+\sigma\pr_2$. One has
\[
  \sF(\psi\sigma)\boxtimes\sF(\psi\sigma)
  =\sF(\psi(\sigma\pr_1+\sigma\pr_2))
  =\sF(\psi s\pi)=\pi^\ast\sF(\psi s),
\]
compatibly with $\tau$. On the image $\pi(\Delta)$ of the diagonal, the desired anti-invariant sheaf is therefore zero. On the complement $X=X'\setminus\pi(\Delta)$, it is $\sF(\psi s)\otimes\varepsilon(P)$, where $\varepsilon$ is the unique character of order $2$ of $\{1,\tau\}$ and $P$ is the $\{1,\tau\}$-torsor on $X$ given by the induced double cover $V_a^1\times V_a^1\to X'$. Thus
\begin{equation*}\tag{7.17.1}\label{VI:eq:7-17-1}
  \operatorname{Sym}^2 \h_c^1(V_a^1,\sF(\psi\sigma)) = \h_c^2(X,\sF(\psi\sigma)\otimes \varepsilon(P)) \text{.}
\end{equation*}

Let us compute $X$, $s$, and $\varepsilon(P)$. On $V_a^1\times V_a^1$, write $x_1,x_2$ for the coordinates pulled back from the first factor and $x_1',x_2'$ for those pulled back from the second. Thus $x_1x_2=x_1'x_2'=a$. Identify functions on $X'$ with $\tau$-invariant functions on $V_a^1\times V_a^1$. Put $s_1=x_1+x_1'$, $s_2=x_2+x_2'$, $p_1=x_1x_1'$, and $p_2=x_2x_2'$. One has $s_1p_2=as_2$ and $s_2p_1=as_1$, while $X$ is defined in $X'$ by $s_1\ne0$ and $s_2\ne0$. Hence on $X$, $p_1=as_1s_2^{-1}$, and $(s_1,s_2)$ identifies $X$ with $\dG_m\times\dG_m$. On $X$, the cover $P$ has equation $T^2-s_1T+as_1s_2^{-1}=0$, with $T=x_1$. On putting $T=s_1t$, this becomes $t^2-t+as_1^{-1}s_2^{-1}=0$. Consequently $\varepsilon(P)=\sF(\psi(as_1^{-1}s_2^{-1}))$. Since $s=s_1+s_2$, one has
\[
  \sF(\psi s)\otimes\varepsilon(P)
  =\sF\bigl(\psi(s_1+s_2+as_1^{-1}s_2^{-1})\bigr),
\]
and the right-hand side of \eqref{VI:eq:7-17-1} is identified with the left-hand side of \ref{VI:7-17}.

\paragraph{Remark}

$\h_c^1(V_a^1,\sF(\psi\sigma))$ is also the first cohomology group of the elliptic curve obtained by completing the affine curve with equations
\begin{align*}
  y^2-y &= x_1+x_2 \\
  x_1 x_2 &= a \text{;}
\end{align*}
its modular invariant is $j=a^{-2}$.

\subsection{Remark}\label{VI:7-18}

For a multiplicative character $\chi$ of $\dF_q^\times$ and $a\ne0$, the method of \ref{VI:7-14} applies to the study of the sum
\[
  S=\sum_{x_0\dotsm x_n=a}\chi(x_0)\psi(x_0+\cdots+x_n)
\]
(replace the sheaf $\sF(\psi)$ on $\dG_m$ by $\sF(\chi\cdot\psi)$). One again obtains a cohomology group of dimension $n+1$, and $|S|\leqslant(n+1)q^{n/2}$. Such sums were considered by Salié and Mordell \cite{mo73}.

One may hope that our methods can be used to study the sums
\[
  S_{k,\chi,a} = \sum_{N(x)=a} \chi(x) \psi \tr(x)
\]
where $k$ is an étale algebra over $\dF_q$, $a\in\dF_q^\times$, and $\chi$ is a character of $k^\times$. The analogue of \ref{VI:7-4} for $a\ne0$ should hold.

\subsection{}\label{VI:7-19}

The symmetric group $S_n$ acts on $\dA^n$, preserving $V_a$ and the map $\sigma$, and hence also the sheaf $\sF(\psi\sigma)$, the inverse image under $\sigma$ of $\sF(\psi)$ on $\dG_a$. It therefore acts on $\h_c^{n-1}(V_a,\sF(\psi\sigma))$. Formula \eqref{VI:eq:7-2-5} has the following cohomological interpretation.

\begin{proposition_}\label{VI:7-20}
The action of $S_n$ on $\h_c^{n-1}(V_a,\sF(\psi\sigma))$ is multiplication by the sign character.
\end{proposition_}

It is enough to show that a transposition $\tau$ acts as multiplication by $-1$. Since $\tau^2=1$, its only possible eigenvalues are $\pm1$. Let $V_a/\tau$ be the quotient of $V_a$ by $\{\operatorname{id},\tau\}$. The map $\sigma$, and hence the sheaf $\sF(\psi\sigma)$, descends to the quotient. The $\tau$-fixed subspace of $\h_c^{n-1}(V_a,\sF(\psi\sigma))$ is $\h_c^{n-1}(V_a/\tau,\sF(\psi\sigma))$. We must show that this cohomology group vanishes.

Assume $a\ne0$. For $\tau=(1,2)$, the quotient $V_a/\tau$ is identified with $\dG_a\times\dG_m^{n-2}$ by
\[
  (x_1,x_2,x_3,\dots,x_n)
  \longmapsto(x_1+x_2,x_3,\dots,x_n).
\]
Indeed, $x_1$ and $x_2$ are determined up to order by $x_1+x_2$ and $x_1x_2=a/(x_3\cdots x_n)$. In the resulting coordinates $(t,x_3,\dots,x_n)$, the sheaf is
\[
  \sF(\psi(t+x_3+\cdots+x_n))
  =\sF(\psi(t))\otimes\sF(\psi(x_3+\cdots+x_n)).
\]
The Künneth formula and $\h_c^\bullet(\dG_a,\sF(\psi))=0$ give the required vanishing.

For $a=0$, Proposition~\ref{VI:7-20} follows from the calculation in \ref{VI:7-7}.

\subsection{}\label{VI:7-21}

The sums $\kloos_{n,a}$ in this section were studied by S.\ Sperber, in his thesis and in \cite{sp77}, by $p$-adic methods inspired by Dwork and Bombieri's paper \cite{bo66}. For $p\ne2$, these methods give him the dependence on $q$ described in \ref{VI:7-5}, the functional equation $\alpha_i=q^{n-1}\beta_i^{-1}$ between the Frobenius eigenvalues for $\kloos_{n,a}$ and $\kloos_{n,b}$ when $b=(-1)^na$ (cf.\ \ref{VI:7-6}), and the value in \ref{VI:7-15} of the product of the roots. For $p\geqslant n+3$, he also determines their $p$-adic valuations. The result is striking: if the roots are ordered suitably as $\alpha_0,\dots,\alpha_{n-1}$, then $\alpha_iq^{-i}$ is a unit. The restrictions on $p$ are probably not essential.

\section{Other applications}\label{VI:8}

\subsection{}\label{VI:8-1}

Identity~\ref{VI:1-9-4} can sometimes be used to calculate the number of rational points of an algebraic variety over $\dF_q$. In most cases in which an exact result has been obtained, the cohomology can be expressed in terms of algebraic cycles. This is so for rational surfaces, quadrics, smooth intersections of two quadrics of odd dimension, flag varieties, and so forth. It is not always so, however: algebraic cycles do not appear in Lusztig's beautiful paper \cite{lu76}.

\subsection{}\label{VI:8-2}

For reductive groups, the cohomology may be calculated by lifting to characteristic $0$. If $B$ is a Borel subgroup of $G$, then $G\to G/B$ is a $B$-torsor. The cohomology of the smooth projective variety $G/B$ is invariant under specialization, and the Leray spectral sequence for $G\to G/B$ is used to prove the same result for $G$. This cohomology is conveniently expressed in terms of a maximal torus $T$ of $G$ and the action of its Weyl group $W$ on $T$ (for the definition of the maximal torus, cf.\ \cite[VIII, \S2, Rem.~2]{bo05} or \cite[p.~105]{dl76}). It is the exterior algebra on its primitive part, which, up to a shift, coincides with the quotient $I(\h^\bullet(BG,\dQ_\ell))$ of the cohomology of $BG$ formed by the indecomposable elements; moreover $\h^\bullet(BG,\dQ_\ell)=\h^\bullet(BT,\dQ_\ell)^W$. If $X$ is the character group of $T$, altogether one has
\[
  \h^\bullet(G,\dQ_\ell) = \bigwedge \left(I\left(\operatorname{Sym}^\bullet(X\otimes \dQ_\ell(-1))^W\right)[-1]\right) \text{.}
\]
If $G$ is defined over a finite field $\dF_q$, then $\gal(\dF/\dF_q)$ acts on $X$, and this isomorphism is Galois compatible. Passing to cohomology with proper support gives the classical formulas for the number of rational points of $G$. If $F^\ast:X\to X$ is defined by the Frobenius morphism $F:T\to T$, and $\deg_G(F)$ is the degree $q^{\dim G}$ of $F:G\to G$, then
\begin{equation*}\tag{8.2.1}\label{VI:eq:8-2-1}
  \left|G_0(\dF_q)\right| = \deg_G(F) \cdot \det\left(1-{F^\ast}^{-1},I(\operatorname{Sym}^\bullet(X\otimes \dQ)^W)\right) \text{.}
\end{equation*}

The same formula holds for the Ree and Suzuki groups. These groups are of the form $G^F$, with $G$ reductive and $F^2$ a Frobenius, and it is enough to verify that the Lefschetz trace formula holds for $F$ acting on $G$. If $B$ is an $F$-stable Borel subgroup, this is checked directly for $F$ acting on $B$; for the action on the proper smooth variety $G/B$, the general theorems apply.

\subsection{}\label{VI:8-3}

The trace formula was used by Deligne--Lusztig \cite{dl76}, Kazhdan \cite{ka77}, and Springer \cite{sp76} in the study of the complex representations of the finite groups $G_0(\dF_q)$, for $G_0$ reductive over $\dF_q$. The work of Kazhdan and Springer contains admirable examples of how the trace formula permits ``analytic continuation'' from a split situation to a nonsplit one (cf.\ \ref{VI:1-13}).
% END INLINED FILE: sga4.5-6.tex
\clearpage
\clearpage
% BEGIN INLINED FILE: sga4.5-7.tex
\chapter{Finiteness theorems in \texorpdfstring{$\ell$}{l}-adic cohomology}\label{VII}

\section{Statements of the theorems}\label{VII:1}

Throughout this exposé, $S$ is a noetherian scheme and $A$ is a left noetherian torsion ring annihilated by an integer invertible on $S$. Our main result is the following.

\begin{theorem_}\label{VII:1-1}
Suppose that $S$ is regular of dimension $0$ or $1$. Let $f:X\to Y$ be a morphism of finite-type $S$-schemes, and let $\sF$ be a constructible sheaf of left $A$-modules on $X$. Then the sheaves $\eR^if_\ast\sF$ are constructible.
\end{theorem_}

The proof will be given in Section~\ref{VII:2} and in \ref{VII:3-10}.

\subsection{Remark}\label{VII:1-2}

In characteristic $0$, this result is less general than \cite[XIX, \S5]{sga4}, which proves the conclusion of the theorem for every finite-type morphism of excellent schemes of characteristic $0$. The proof in \cite[XIX]{sga4} uses both resolution of singularities and the fact that the schemes are equicharacteristic, in order to deduce the purity theorem from resolution.

\subsection{Remark}\label{VII:1-3}

We call a complex $K\in\ob\eD(X,A)$ \emph{constructible} if its cohomology sheaves are constructible, and use a subscript $c$ for the subcategory of $\eD(X,A)$ (or of $\eD^+$, $\eD^-$, or $\eD^b$) consisting of constructible complexes. In the language of derived categories, which we shall use freely, \ref{VII:1-1} says that $\eR f_\ast:\eD^+(X,A)\to\eD^+(Y,A)$ maps $\eD_c^+$ into $\eD_c^+$. More explicitly:
\begin{enumerate}[\indent a)]
  \item if $K$ is a sheaf $\sF$ concentrated in degree $0$, then $\sH^i\eR f_\ast K=\eR^if_\ast\sF$; thus the derived statement implies the theorem;
  \item conversely, one uses the spectral sequence
    \[
      E_1^{p q} = \eR^q f_\ast \sH^p(K) \Rightarrow \sH^{p+q} \eR f_\ast K \text{.}
    \]
\end{enumerate}

The advantage of the language of derived categories is that the Leray spectral sequence $E_2^{pq}=\eR^pf_\ast\eR^qg_\ast\Rightarrow \eR^{p+q}(fg)_\ast$ is replaced by the simple transitivity formula $\eR(fg)_\ast=\eR f_\ast\eR g_\ast$.

Under the hypotheses of \ref{VII:1-1}, it follows from \cite{sga4} that $\eR f_\ast$ has finite cohomological dimension. Thus $\eD^+$ above may be replaced by $\eD$, $\eD^-$, or $\eD^b$.

\subsection{Ideas of the proof}\label{VII:1-4}
\begin{enumerate}[\indent a)]
  \item Poincaré duality handles the case in which $X$ is smooth, $\sF$ is locally constant, and $Y=S$. The dimension hypothesis enters through the calculation of $\rHom$ on $S$.
  \item Factor $f$ as $gj$, with $g$ proper and $j$ an open immersion. One has $\eR f_\ast=\eR g_\ast\eR j_\ast$, and the finiteness theorem for proper morphisms controls $\eR g_\ast$. By dévissage, one reduces to the case in which $X$ is smooth, $\sF$ is locally constant, $f$ is an open immersion, and $Y$ is proper over $S$.
  \item Induction on $\dim X$, allowing a change of $S$, makes it possible, roughly speaking, to suppose the theorem true away from a part of $Y$ finite over $S$. If $a:X\to S$ and $b:Y\to S$ are the structural morphisms, one then shows that if the theorem were false for $\sF$ and $\eR f_\ast$, it would also be false for $\sF$ and $\eR b_\ast\eR f_\ast=\eR a_\ast$, a contradiction.
\end{enumerate}

\begin{corollary_}\label{VII:1-5}
Under the hypotheses of the theorem, the four operations $\eR f_\ast$, $\eR f_!$, $f^\ast$, and $\eR f^!$ preserve the constructible derived categories $\eD_c(X,A)$ and $\eD_c(Y,A)$.
\end{corollary_}

The case of $\eR f_\ast$ was treated above, that of $\eR f_!$ in \cite[XVII, 5.3]{sga4}, and the case of $f^\ast$ is clear. It remains to consider $\eR f^!$. The question is local, so we may suppose that $f$ factors as $X\xrightarrow{i}Z\xrightarrow{g}Y$, where $i$ is a closed immersion and $g$ is smooth and purely of relative dimension $n$. Poincaré duality, $\eR g^!K=K(n)[2n]$, and transitivity, $\eR f^!=\eR i^!\eR g^!$, reduce us to proving \ref{VII:1-5} for $i$. For every $L\in\eD(Z,A)$, if $j$ is the inclusion in $Z$ of $U=Z\setminus X$, then $i_\ast\eR i^!L$ is the mapping cone, shifted by $-1$, of $L\to\eR j_\ast j^\ast L$. Thus \ref{VII:1-5} for $i$ follows from \ref{VII:1-1} for $j$.

\begin{corollary_}\label{VII:1-6}
Let $X$ be as in the theorem, and suppose that $A$ is commutative. If $\sF$ and $\sG$ are constructible sheaves of $A$-modules on $X$, then the sheaves $\Ext_A^i(\sF,\sG)$ are constructible; more generally, $\rHom$ maps $\eD_c^-(X,A)\times\eD_c^+(X,A)$ into $\eD_c^+(X,A)$.
\end{corollary_}

By dévissage of $\sF$, one reduces to the case $\sF=j_!\sF_1$, where $j:Y\hookrightarrow X$ is a locally closed immersion and $\sF_1$ is locally constant on $Y$. One then has
\[
  \rHom(j_! \sF_1,\sG) = \eR j_\ast \Hom(\sF_1,\eR j^! \sG) \text{:}
\]
By \ref{VII:1-3}, we are reduced to the case in which $\sF$ is locally constant, and even constant since the question is local. For constructible locally constant $\sF$, one has $\Hom(\sF,\sG)_x=\hom(\sF_x,\sG_x)$, and similarly for $\rHom$. When $\sF$ is constant, the $\Ext^i$ may therefore be computed using a finite-type projective resolution of its constant value. They are constructible when $\sG$ is, which proves the corollary.

\subsection{Remark}\label{VII:1-7}

Since $\eR f_\ast$ has finite cohomological dimension, it maps complexes of finite Tor-dimension (respectively, Tor-dimension $\leqslant d$) to complexes with the same property \cite[XVII, 5.2.11]{sga4}. \cite[XVII, 5.2.10]{sga4} and the proof of \ref{VII:1-5} then show that finite Tor-dimension is stable under the four operations $\eR f_\ast$, $\eR f_!$, $f^\ast$, and $\eR f^!$. A variant of the proof of \ref{VII:1-6}---use a partition of $X$ to reduce to locally constant cohomology sheaves, then work locally to replace $K$ by a finite complex of finite-type locally free sheaves---shows that, for commutative $A$, $\rHom$ induces
\[
  \rHom:\eD_\text{ctf}^b(X,A)\times \eD_\text{tf}^b(X,A) \to \eD_\text{tf}^+(X,A)
\]
(t.f.\ means finite Tor-dimension).

\subsection{}\label{VII:1-8}

Under the hypotheses of the theorem, the same method allows us to prove a local biduality theorem $K\iso DDK$ (Section~\ref{VII:4}). We also prove a finiteness theorem for vanishing-cycle sheaves. For arbitrary $S$, one still obtains a generic theorem:

\begin{theorem_}\label{VII:1-9}
Let $f:X\to Y$ be a morphism of finite-type $S$-schemes, and let $\sF$ be a constructible sheaf of $A$-modules on $X$. There exists a dense open subset $U$ of $S$ such that:
\begin{enumerate}[\indent (i)]
  \item over $U$, the $\eR^if_\ast\sF$ are constructible, and all but finitely many vanish;
  \item the formation of the $\eR^if_\ast\sF$ commutes with every base change $S'\to U\subset S$.
\end{enumerate}
\end{theorem_}

If $S$ is the spectrum of a field, one may take $U=S$.

When $S$ is the spectrum of a field, a limit argument gives compatibility of the $\eR^if_\ast$ with $S$-base change for every quasi-compact, quasi-separated morphism of $S$-schemes and every sheaf $\sF$.

\begin{corollary_}\label{VII:1-10}
Let $X$ be a finite-type scheme over a separably closed field $k$, and let $\sF$ be a constructible sheaf of $A$-modules. Then the $\h^i(X,\sF)$ are of finite type.
\end{corollary_}

This is the special case $S=\spec(k)$ and $Y=S$.

\begin{corollary_}\label{VII:1-11}
Let $X$ and $Y$ be finite-type schemes over a separably closed field $k$, and let $K\in\ob\eD^-(X,A^\circ)$ and $L\in\ob\eD^-(Y,A)$ be complexes of sheaves of right and left modules, respectively. The Künneth map
\[
  \eR \Gamma(X,K) \lotimes \eR \Gamma(Y,L) \to \eR\Gamma(X\times Y,\pr_1^\ast K \lotimes \pr_2^\ast L)
\]
is an isomorphism.
\end{corollary_}

Proceed as in \cite[XVII, 5.4.3]{sga4}:
\[\xymatrix{
  X\times Y \ar[r]^-{\text{pr}_1} \ar[d]^-{\text{pr}_2}
    & X \ar[d]^-a \\
  Y \ar[r]^-b
    & \spec(k)
}\]
base change by $b$ shows that $\eR\pr_{2,\ast}\pr_1^\ast K=b^\ast\eR\Gamma(X,K)$. Hence (cf.\ \cite[XVII, 5.2.11]{sga4} and its proof)

\begin{align*}
  \eR\Gamma(X\times Y,\pr_1^\ast K\lotimes \pr_2^\ast L)
    &= \eR\Gamma(Y,\eR \pr_{2,\ast}(\pr_1^\ast K\lotimes \pr_2^\ast L)) \\
    &= \eR\Gamma(Y,(\eR\pr_{2,\ast}\pr_1^\ast K)\lotimes L) \\
    &= \eR\Gamma(Y,b^\ast \eR\Gamma(X,K)\lotimes L) \\
    &= \eR\Gamma(X,K)\lotimes \eR\Gamma(Y,L)
\end{align*}

\subsection{}\label{VII:1-12}

Finally, as a local counterpart of \ref{VII:1-9}(ii), we shall prove a generic local acyclicity theorem (\ref{VII:2-13}).

\section{Generic theorems}\label{VII:2}

In this section we prove \ref{VII:1-9} and the generic local acyclicity theorem \ref{VII:2-13}. To prove \ref{VII:1-9}, we immediately reduce to the case in which $S$ is integral. Let $\eta$ be its generic point.

\subsection{}\label{VII:2-1}

We first prove \ref{VII:1-9} under the following additional hypotheses: $X$ is smooth over $S$, purely of relative dimension $n$, $A=\dZ/m$, $\sF$ is locally constant, and $Y=S$. Put $\sF'=\Hom(\sF,\dZ/m)$. We shall show that if the $\eR^if_!\sF'$ are locally constant, then the conclusions of \ref{VII:1-9} hold with $U=S$. Recall that if $\sL$ is a constructible locally constant sheaf, then the $\Ext^i(\sL,\sG)$ are computed stalk by stalk and are constructible if $\sG$ is. Also recall that $\dZ/m$ is an injective $\dZ/m$-module and is the dualizing module, so that $\sF=\rHom(\sF',\dZ/m)$. Poincaré duality,
\[
  \eR f_\ast \rHom(K,\eR f^! L) = \rHom(\eR f_! K,L) \text{,}
\]
with $K=\sF'$, $L=\dZ/m$, and $\eR f^!L=\dZ/m[2n](n)$, therefore gives
\[
  \eR^{2n-i} f_\ast \sF = \Hom(\eR^i f_! \sF',\dZ/m)(-n) \text{.}
\]

The sheaf on the right is locally constant and its formation commutes with every base change, proving the assertion.

\subsection{}\label{VII:2-2}

We prove \ref{VII:1-9} under the additional hypotheses that $X$ is smooth over $S$, $\sF$ is locally constant, and $Y=S$.
\begin{enumerate}[a)]
  \item Decomposing $X$ into connected components, we reduce to the case in which it is purely of relative dimension $n$ over $S$.
  \item Decomposing $A$ as a product, we reduce to the case $\ell^mA=0$, where $\ell$ is a prime invertible on $S$. Filter $\sF$ by the $\ell^k\sF$; the corresponding spectral sequence reduces us to $\ell\sF=0$. We may then replace $A$ by $A/\ell A$ and suppose $\ell A=0$.
  \item Replacing $S$ by a suitable $U$, we may suppose that there is a finite étale Galois covering $X_1/X$, with Galois group $G$, such that the inverse image $\sF_1$ of $\sF$ on $X_1$ is a constant sheaf with value $F$. If $f_1$ denotes the projection from $X_1$ to $S$, then
    \[
      \eR^i {f_1}_\ast \sF_1 = (\eR^i {f_1}_\ast \dZ/\ell)\otimes_{\dZ/\ell} F \text{.}
    \]
    There is also the Hochschild--Serre spectral sequence (equivalently, the Leray spectral sequence for the covering $X_1/X$)
    \[
      \underline\h^p(G,\eR^q {f_1}_\ast \sF_1) \Rightarrow \eR^{p+q} f_\ast \sF \text{.}
    \]
    By \ref{VII:2-1}, after shrinking $U$ we may suppose that the $\eR^i{f_1}_\ast\dZ/\ell$ are locally constant and that their formation commutes with every base change. The same property then holds for the $\eR^if_\ast\sF$. Finally, if $\bar\eta$ is a geometric point over the generic point $\eta$ of $S$, all but finitely many $\eR^if_\ast\sF$ vanish, because the same is true of their stalks $(\eR^if_\ast\sF)_{\bar\eta}=\h^i(X_{\bar\eta},\sF)$.
\end{enumerate}

\subsection{}\label{VII:2-3}

We prove by induction on $n$ that
\begin{equation*}\tag{$\ast_n$}\label{VII:eq:2-3-1}
  \begin{array}{c}
    \text{The conclusions of \ref{VII:1-9} hold when
    $\dim X_\eta\leqslant n$}\\
    \text{and $f$ is an open immersion with dense image.}
  \end{array}
\end{equation*}

For $n=0$, after shrinking $S$ one has $X=Y$, so $(\ast_0)$ is clear. Assume $(\ast_{n-1})$ and prove \eqref{VII:eq:2-3-1}. In $(\ast_{n-1})$, ``open immersion'' may be replaced by ``immersion'', as is seen by factoring an immersion into an open immersion followed by a closed immersion.

\begin{lemma_}\label{VII:2-4}
After shrinking $S$, the conclusions of \ref{VII:1-9} hold over an open subset $Y'\subset Y$ whose complement $Y_1$ is finite over $S$.
\end{lemma_}

The assertion is local on $Y$, which we may suppose affine, with $Y\subset\dA_S^n$. The induction hypothesis $(\ast_{n-1})$ applies to
\[\xymatrix{
  X \ar[r]^-f \ar[dr]
    & Y \ar[d]^-{\pr_i} \\
  & \dA_S^1
}\]
For each $i$ there is therefore a dense open subset $U_i$ of $\dA_S^1$ such that the conclusions of \ref{VII:1-9} hold over $\pr_i^{-1}(U_i)$. They hold over the union of the $\pr_i^{-1}(U_i)$, and \ref{VII:2-4} follows.

\addtocounter{subsection}{1} % a section is skipped in the original
\subsection{}\label{VII:2-6}

We prove \eqref{VII:eq:2-3-1} when $X$ is smooth over $S$ and $\sF$ is locally constant on $X$. Since the question is local on $Y$, we may suppose that $Y$ is affine and then make it projective by replacing it with its closure in a projective space. Let $i:Y_1\hookrightarrow Y$ and $j:Y'\to Y$ be as in \ref{VII:2-4}.
\[\xymatrix{
  X \ar@{^{(}->}[r] \ar[dr]_-a
    & Y \ar[d]^-b
    & \ar[l]_-i \ar[dl]^-{b_1} Y_1 \\
  & S
}\]
After shrinking $S$, we know that $j^\ast\eR f_\ast\sF$ is constructible and its formation commutes with every base change on $S$. We also know that $\eR a_\ast\sF=\eR b_\ast\eR f_\ast\sF$ is constructible and has the same base-change property.

Apply $\eR b_\ast$ to the triangle defined by the exact sequence
\begin{equation*}\tag{1}\label{VII:eq:2-6-1}
\xymatrix{
  0 \ar[r]
    & j_! j^\ast \eR f_\ast \sF \ar[r]
    & \eR f_\ast \sF \ar[r]
    & i_! i^\ast \eR f_\ast \sF \ar[r]
    & 0 \text{:}
}
\end{equation*}
One obtains a triangle
\begin{equation*}\tag{2}\label{VII:eq:2-6-2}
\xymatrix{
  \ar[r]
    & \eR b_\ast j_! j^\ast \eR f_\ast \sF \ar[r]
    & \eR a_\ast \sF \ar[r]
    & {b_1}_\ast i^\ast \eR f_\ast \sF \ar[r]
    &
}
\end{equation*}
in which the first two terms are constructible and their formation commutes with every base change on $S$; for the first, this follows from the finiteness theorem for the proper morphism $b$. The same is therefore true of the third. Since $b_1$ is finite, it follows that $i^\ast\eR f_\ast\sF$ is constructible and has the base-change property. By \eqref{VII:eq:2-6-1}, the same holds for $\eR f_\ast\sF$.

\subsection{}\label{VII:2-7}

We prove \eqref{VII:eq:2-3-1} in general. First reduce to the case in which $X$ contains a dense open subset $V$ smooth over $S$. If $S$ is the spectrum of a perfect field, it is enough to replace $X$ by $X_{\mathrm{red}}$ and $Y$ by $Y_{\mathrm{red}}$. In general, one must shrink $S$, make a finite, radicial, surjective base change $S'\to S$, and replace $X$ and $Y$ by their reductions. Since the étale topology is insensitive to finite, radicial, surjective morphisms, this causes no problem. After shrinking $V$, we may suppose that $\sF$ is locally constant on $V$:
\[\xymatrix{
  V \ar@{^{(}->}[r]^j
    & X \ar@{^{(}->}[r]^-f
    & Y \text{.}
}\]
Define $\Delta$ by the triangle
\begin{equation*}\tag{1}\label{VII:eq:2-7-1}
\xymatrix{
  \ar[r]
    & \sF \ar[r]
    & \eR j_\ast j^\ast \sF \ar[r]
    & \Delta \ar[r]
    & \text{.}
}
\end{equation*}
The cohomology sheaves of $\Delta$ are supported on $X\setminus V$, and $\dim(X\setminus V)_\eta<n$. The induction hypothesis therefore allows us to suppose that $\eR f_\ast\Delta$ is constructible. Applying $\eR f_\ast$ to the triangle \eqref{VII:eq:2-7-1} gives a triangle
\[\xymatrix{
  \ar[r]
    & \eR f_\ast \sF \ar[r]
    & \eR(f j)_\ast j^\ast \sF \ar[r]
    & \eR f_\ast \Delta \ar[r]
    &
}\]
in which two of the vertices are constructible and their formation commutes with every base change on $S$. The same is therefore true of the third.

\subsection{}\label{VII:2-8}

We prove \ref{VII:1-9}. The question is local on $Y$, which we may suppose affine. Taking an affine covering of $X$ and using the Leray spectral sequence reduces us to the case in which $X$ is also affine. This ensures that $f$ can be factored as an open immersion followed by a proper morphism: $f=gj$, whence $\eR f_\ast=\eR g_\ast\eR j_\ast$. The open immersion is covered by \eqref{VII:eq:2-3-1}, and the proper morphism by the finiteness theorem.

The following corollaries are proved as in Section~\ref{VII:1}.

\begin{corollary_}\label{VII:2-9}
Under the hypotheses of the theorem, for $K$ in $\eD_c^b(X,A)$ or $\eD_c^b(Y,A)$ as appropriate, there exists a nonempty open subset $U$ of $S$ over which $\eR f_\ast K$, $\eR f_!K$, $f^\ast K$, and $\eR f^!K$ belong to $\eD_c^b(Y,A)$ or $\eD_c^b(X,A)$ as appropriate, and their formation commutes with every base change $S'\to U\subset S$.
\end{corollary_}

\begin{corollary_}\label{VII:2-10}
Let $X$ be as in the theorem, and suppose $A$ commutative. If $\sF$ and $\sG$ are constructible sheaves of $A$-modules on $X$, then over some dense open subset $U$ of $S$, the sheaves $\Ext_A^i(\sF,\sG)$ are constructible and their formation commutes with every base change $S'\to U\subset S$.
\end{corollary_}

\subsection{}\label{VII:2-11}

If $x$ is a geometric point of a scheme $X$, denote by $X_x$ the strict henselization of $X$ at $x$. For $f:X\to S$ and a geometric point $t$ of $S$, denote by $X_t$ the geometric fiber of $X$ at $t$. Finally, for a geometric point $x$ of $X$ and a geometric point $t$ of $S_{f(x)}$, $(X_x)_t$ is the fiber at $t$ of $X_x\to S_{f(x)}$.

\begin{definition_}\label{VII:2-12}
Let $f:X\to S$ and $K\in\ob\eD^+(X,A)$. We say that $f$ is \emph{locally acyclic at $x$ relative to $K$} if, for every geometric point $t$ of $S_{f(x)}$, one has $K_x=\eR\Gamma(X_x,K)\iso\eR\Gamma((X_x)_t,K)$. We say that $f$ is \emph{locally acyclic relative to $K$} if this holds for every $x$, and \emph{universally locally acyclic relative to $K$} if it remains true after every base change $S'\to S$.
\end{definition_}

\begin{theorem_}\label{VII:2-13}
Let $f:X\to S$ be a finite-type morphism and let $\sF$ be constructible on $X$. There exists a dense open subset $U$ of $S$ over which $f$ is universally locally acyclic relative to $\sF$.
\end{theorem_}

We use the following result, proved in the appendix.

\begin{lemma_}\label{VII:2-14}
Let $X\xrightarrow{f}S_1\xrightarrow{g}S_2$. If $g$ is smooth and $f$ is universally locally acyclic relative to $K$, then $gf$ is also.
\end{lemma_}

We reduce to the case in which $S$ is integral, with generic point $\eta$, and proceed by induction on $\dim X_\eta$. First, the induction hypothesis implies that
\begin{equation*}\tag{A}\label{VII:eq:2-14-1}
  \begin{array}{c}
    \text{After shrinking $S$, there exists $T\subset X$, finite over
    $S$, such that $f$}\\
    \text{is universally locally acyclic outside $T$.}
  \end{array}
\end{equation*}

Since the question is local, work locally on $X$ and factor $f$ as $X\xrightarrow{u}\dA_S^1\to S$. The induction hypothesis applies to $u$, and \ref{VII:2-14}, together with the usual arguments, gives \eqref{VII:eq:2-14-1}.

To prove the theorem, we may suppose $X$ proper, since the question is local. Shrink $S$ so that the $\eR^if_\ast\sF$ are locally constant and \eqref{VII:eq:2-14-1} applies, and use the following lemma.

\begin{lemma_}\label{VII:2-15}
Let $f:X\to S$ be proper and let $\sF$ be constructible. Suppose that the $\eR^if_\ast\sF$ are locally constant and that, outside a subset $T\subset X$ finite over $S$, $f$ is locally acyclic relative to $\sF$. Then $f$ is locally acyclic relative to $\sF$.
\end{lemma_}

Let $s$ be a geometric point of $S$, let $S_s$ be the strict localization of $S$ at $s$, and let $t$ be a geometric point of $S_s$. Let $X_{(s)}$ be the inverse image of $X$ over $S_s$, let $\bar j:X_t\to X_{(s)}$, and let $i:X_s\hookrightarrow X_{(s)}$. Still denoting by $\sF$ its inverse image on $X_{(s)}$ or $X_t$, we must prove that $i^\ast\sF\iso i^\ast\eR\bar j_\ast\sF$. Let $\Delta$ be the mapping cone of this map. Its cohomology sheaves are supported on $T_s$. Thus, to prove $\Delta=0$, it is enough to prove $\eR\Gamma(X_s,\Delta)=0$. This is the mapping cone of
\[
\begin{split}
  (\eR f_\ast\sF)_s=\eR\Gamma(X_s,\sF)
  &\longrightarrow
  \eR\Gamma(X_s,i^\ast\eR\bar j_\ast\sF)\\
  &=\eR\Gamma(X_{(s)},\eR\bar j_\ast\sF)
   =\eR\Gamma(X_t,\sF)=(\eR f_\ast\sF)_t.
\end{split}
\]
By hypothesis this specialization morphism is an isomorphism, so $\Delta=0$.

\begin{corollary_}\label{VII:2-16}
If $S$ is the spectrum of a field, every $S$-scheme $X$ is universally locally acyclic relative to every $K\in\ob\eD^+(X,A)$.
\end{corollary_}

This follows from \ref{VII:2-13} by passage to the limit, which is possible because the $U$ of \ref{VII:2-13} is always equal to $S$.

\section{Proof of \texorpdfstring{\ref{VII:1-1}}{1.1} and constructibility of vanishing-cycle sheaves}\label{VII:3}

\subsection{}\label{VII:3-1}

Let $S$ be a strictly local trait, the spectrum of a henselian discrete valuation ring with separably closed residue field. Denote its closed and generic points by $s$ and $\eta$, and let $\bar\eta$ be a geometric point over $\eta$, corresponding to a separable closure of $k(\eta)$. Let $X$ be an $S$-scheme and $\sF$ a sheaf on $X_\eta$. Recall the definition of the vanishing-cycle sheaves $\eR^i\Psi_\eta(\sF)$: these are sheaves on $X_s$ endowed with an action of $\gal(\bar\eta/\eta)$. Let $i:X_s\hookrightarrow X$, $j:X_\eta\hookrightarrow X$, and let $\bar j$ be the composite $X_{\bar\eta}\to X_\eta\hookrightarrow X$. Then
\[
  \eR^i \Psi_\eta(\sF) = i^\ast \eR^i \bar j_\ast (\bar j^\ast \sF) \text{.}
\]

The main theorem of this section is the following. It improves \cite[XIII, 2.3.1 and 2.4.2]{sga7}.

\begin{theorem_}\label{VII:3-2}
Suppose that $X$ is of finite type over $S$ and $\sF$ is constructible. Then the vanishing-cycle sheaves $\eR^i\Psi_\eta(\sF)$ are constructible.
\end{theorem_}

We proceed by induction on the dimension $n$ of $X_\eta$. We may and do suppose that $X_\eta$ is dense in $X$: replacing $X$ by the closure $\overline{X_\eta}$ does not change the $\eR^i\Psi$, which are supported on $\overline{X_\eta}$.

\begin{lemma_}\label{VII:3-3}
If \ref{VII:3-2} holds when $\dim X_\eta<n$, and if $\dim X_\eta=n$, then there exist constructible subsheaves $\sG_i$ of $\eR^i\Psi_\eta(\sF)$ such that the local sections of $\eR^i\Psi_\eta(\sF)/\sG_i$ have finite support.
\end{lemma_}

Let $s'$ be a geometric generic point of the affine line over $s$, and let $S'$ be the strict localization at $s'$ of the affine line $\dA_S^1$. This is again a strictly local trait, and a uniformizer for $S$ is also a uniformizer for $S'$.
\[\xymatrix{
  & \spec(k') \ar[rr] \ar[dd] \ar[dr]
    & & (S',\eta',s') \ar[dr] \ar[dd] \\
  \bar\eta' \ar[ur] \ar[dr]
    & & \dA_{\bar S}^1 \ar[d]
    & & \dA_S^1 \ar[dl] \\
  & \bar\eta\ar@{^{(}->}[r]
    & \bar S \ar[r]
    & (S,\eta,s)
}\]
Put $k'=k(\bar\eta)\otimes_{k(\eta)}k(\eta')$. If $\bar S$ is the normalization of $S$ in $\bar\eta$, then $k'$ is the fraction field of the strict localization of $\dA_{\bar S}^1$ at $s'$. It is therefore a field, and $\gal(\bar\eta/\eta)=\gal(k'/\eta')$. Let $\bar\eta'$ be the spectrum of an algebraic closure of $k'$. The Galois group $P=\gal(\bar\eta'/k')$ is a pro-$p$ group, where $p$ is the characteristic exponent of $k(s)$.

\begin{lemma_}\label{VII:3-4}
For every scheme $X'$ over $S'$ and every sheaf $\sF$ on $X_\eta'=X'_{\eta'}$, the vanishing-cycle sheaves for $X'/S'$ and for $X'/S$ are related by $\eR^i\Psi_\eta(\sF)=\eR^i\Psi_{\eta'}(\sF)^P$.
\end{lemma_}

The following diagram compares the morphisms used to define $\eR\Psi$ for $X'/S'$ and $X'/S$:
\[\xymatrix{
  X_{\bar\eta'} \ar[rr] \ar[d]
    & & X_{\eta'}' \ar@{=}[dd] \ar@{^{(}->}[dr]
    & & S' \\
  X_{\eta'}\times_{\eta'}\spec(k') \ar@{=}[d] \ar[urr]
    & & & X \ar[ur] \ar[dr] \\
  X_\eta' \ar[rr]
    & & X_\eta' \ar@{^{(}->}[ur]
    & & S
}\]
The lemma follows from the Hochschild--Serre spectral sequence, since $P$ is a pro-$p$ group and $p$ is invertible in $A$.

We prove \ref{VII:3-3}. The question is local, so we may suppose $X$ affine, with $X\subset\dA_S^n$. Let $f$ be one of the coordinate projections $X\subset\dA_S^n\to\dA_S^1$, let $X'=X\times_{\dA_S^1}S'$ be the corresponding ``localization'', and let $\sF'$ be the inverse image of $\sF$ on $X'$:
\[\xymatrix{
  \sF'\text{ on }X' \ar[r] \ar[d]^-\lambda
    & S' \ar[d]^-\lambda \\
  \sF \text{ on }X \ar[r]^-f
    & \dA_S^1 \ar[r]
    & S \text{.}
}\]
On $X_s'=X_{s'}'$, one has
\[
  \lambda^\ast \eR^i\Psi_\eta(\sF) = \eR^i \Psi_\eta(\sF') = \eR^i\Psi_{\eta'}(\sF')^P \text{:}
\]
This sheaf is constructible because the induction hypothesis applies to $X'/S'$. We then use the following lemma, applied to $X_s\subset\dA_s^n$ and to the vanishing-cycle sheaves.

\begin{lemma_}\label{VII:3-5}
Let $\dA^n$ be standard affine $n$-space over a field $k$, let $X\subset\dA^n$, let $\sF$ be a sheaf of $A$-modules on $X$, and let $\bar\eta$ be a geometric generic point of $\dA^1$. Let $X_{\bar\eta,i}$ be the geometric generic fiber of $X\subset\dA^n\xrightarrow{\pr_i}\dA^1$, and let $\sF_{\bar\eta,i}$ be the restriction of $\sF$ to $X_{\bar\eta,i}$. If the $\sF_{\bar\eta,i}$ are constructible, then there is a constructible subsheaf $\sF'\subset\sF$ such that the local sections of $\sF/\sF'$ have finite support.
\end{lemma_}

By passage to the limit, for each $i$:
\begin{enumerate}[a)]
  \item there is an étale neighborhood $U$ of $\bar\eta$ and a constructible sheaf $\sH$ on $X_{U,i}=X\times_{\dA^1,\pr_i}U$ whose restriction to $X_{\bar\eta,i}$ is $\sF_{\bar\eta,i}$;
  \item for suitable $U$, the isomorphism $\sH_{\bar\eta,i}\iso\sF_{\bar\eta,i}$ comes from a morphism $u:\sH\to\sF_{U,i}$, where $\sF_{U,i}$ is the inverse image of $\sF$ on $X_{U,i}$. Let $\varphi:X_{U,i}\to X$. The morphism $u$ defines $v_i:\varphi_!\sH\to\sF$; let its image be $\sF_i'$. Then $(\sF/\sF_i')_{\bar\eta,i}=0$, and one takes for $\sF'$ the sum of the $\sF_i'$.
\end{enumerate}

\subsection{}\label{VII:3-6}

We prove \ref{VII:3-2} in dimension $n$. We may suppose $X$ affine, since the question is local on $X_s$, and then projective over $S$ by taking its closure in a projective embedding. There is then a spectral sequence \cite[I, 2.2.3]{sga7}: the Leray spectral sequence for $\bar j$, together with the base-change theorem for the proper morphism $X\to S$,
\begin{equation*}\tag{1}\label{VII:eq:3-6-1}
  E_2^{pq}=\h^p(X_s,\eR^q\Psi_\eta(\sF))
  \Rightarrow\h^{p+q}(X_{\bar\eta},\sF).
\end{equation*}
Let $\sG^q\subset\eR^q\Psi_\eta(\sF)$ be as in \ref{VII:3-3}, and let $\sH^q$ be the quotient sheaf. One has $\h^i(X_s,\sH^q)=0$ for $i\ne0$. For $\sH^q$, and hence $\eR^q\Psi_\eta(\sF)$, to be constructible, it is enough that $\h^0(X_s,\sH^q)$ be of finite type.

Compute \eqref{VII:eq:3-6-1} modulo finite-type modules, that is, in the quotient of the category of $A$-modules by the thick subcategory of finite-type $A$-modules. The long exact cohomology sequence associated with $0\to\sG^q\to\eR^q\Psi_\eta(\sF)\to\sH^q\to0$, together with \ref{VII:1-10} applied to $\sG^q$, gives $E_2^{pq}\sim\h^p(X_s,\sH^q)$. In particular, $E_2^{pq}\sim0$ for $p\ne0$, while $E_2^{0q}\sim\h^q(X_{\bar\eta},\sF)\sim0$. This completes the proof.

Parallel arguments give the following result, improving \cite[XIII, 2.1.12 and 2.4.2]{sga7}.

\begin{proposition_}\label{VII:3-7}
The formation of vanishing-cycle sheaves commutes with change of traits.
\end{proposition_}

Let $g:(S',\eta',s')\to(S,\eta,s)$ be a surjective morphism of strictly local traits, let $X/S$, let $\sF$ on $X_\eta$ be of torsion prime to the residue characteristic, and let $(X',\sF')$ be their inverse images over $S'$. Still writing $g$ for the morphisms parallel to $g$, such as $X_{s'}'\to X_s$, we must prove that
\begin{equation*}\tag{3.7.1}\label{VII:eq:3-7-1}
\xymatrix{
  g^\ast \eR\Psi_\eta(\sF) \ar[r]^-\sim
    & \eR\Psi_{\eta'}(\sF') \text{.}
}
\end{equation*}
This morphism is defined by the commutative diagram
\[\xymatrix{
  X_{s'}' \ar@{^{(}->}[r]^-{i'} \ar[d]^-g
    & X' \ar[d]^-g
    & X_{\eta'}' \ar[l]_-{\bar j'} \ar[d]^-g \\
  X_s \ar@{^{(}->}[r]^-i
    & X
    & \ar[l]_-{\bar j} X_{\bar\eta}
}\]

A limit argument reduces us to the case in which $X$ is of finite type over $S$, and we proceed by induction on $\dim X_\eta$. Since the question is local, we may suppose $X$ affine, and then projective. Suppose \eqref{VII:eq:3-7-1} known for $\dim X_\eta<n$. For $\dim X_\eta=n$, it follows from the following two assertions, where $\Delta$ is the mapping cone of \eqref{VII:eq:3-7-1}.

\begin{enumerate}[(A)]
  \item The local sections of the cohomology sheaves of $\Delta$ have finite support. \phantomsection\hypertarget{VII:3-A}{}
  \item $\eR\Gamma(X_s,\Delta)=0$, hence $\Delta=0$ by (\hyperlink{VII:3-A}{A}). \phantomsection\hypertarget{VII:3-B}{}
\end{enumerate}

\subsection{}\label{VII:3-8}

The proof of (\hyperlink{VII:3-A}{A}) is parallel to that of \ref{VII:3-3}. Work locally on $X$ and consider maps $X\to\dA_S^1$. If $S(x)$ denotes the strict henselization of $\dA_S^1$ at the generic point of the special fiber, apply the induction hypothesis to $X_1/S(x)$, where $X_1=X\times_{\dA^1}S(x)$, and to $S'(x)$. Taking invariants under a pro-$p$ group, one finds that $\Delta$ vanishes on the geometric generic fiber of $X_{s'}\to\dA_s^1$. Since this holds locally on $X$ and for every projection, (\hyperlink{VII:3-A}{A}) follows.

\subsection{}\label{VII:3-9}

For (\hyperlink{VII:3-B}{B}), properness of $X$ over $S$ ensures that
\[
  \eR\Gamma(X_{s'}',\eR\Psi_{\eta'}(\sF')) = \eR\Gamma(X_{\bar\eta}',\sF')
\]
et
\[
  \eR\Gamma(X_s,\eR\Psi_\eta(\sF)) = \eR\Gamma(X_{\bar\eta},\sF) \text{.}
\]
One has $\eR\Gamma(X_s,\eR\Psi_\eta(\sF))\iso \eR\Gamma(X_{s'}',g^\ast\eR\Psi_\eta(\sF))$ by invariance under change of separably closed base field. Thus $\eR\Gamma(X_{s'}',\Delta)$ is the mapping cone of $\eR\Gamma(X_{\bar\eta},\sF)\to \eR\Gamma(X_{\bar\eta}',\sF')$, which is an isomorphism by the same invariance theorem.

\subsection{}\label{VII:3-10}

We prove \ref{VII:1-1}. We may suppose $S$ connected, hence integral. If $\dim S=0$, so that $S$ is the spectrum of a field, apply \ref{VII:1-9} and \ref{VII:1-10}. If $\dim S=1$, Theorem~\ref{VII:1-9} shows that the $\eR^if_\ast\sF$ are constructible over the complement of a finite set $T$ of points of $S$. It remains to show that their restrictions to $Y_t$, for $t\in T$, are constructible. It is enough to do this after localization at $t$, so we may suppose that $S$ is a strictly local trait. The essential case is the following.

\begin{lemma_}\label{VII:3-11}
Let $X$ be of finite type over $S$, let $j:X_\eta\hookrightarrow X$ and $i:X_s\hookrightarrow X$, and let $\sF$ be constructible on $X_\eta$. Then $i^\ast\eR j_\ast\sF$ is constructible.
\end{lemma_}

Let $I=\gal(\bar\eta/\eta)$ be the inertia group. It is an extension of $\widehat\dZ_{p'}(1)$ by a pro-$p$ group $P$, where $p$ is the characteristic exponent of the residue field. The Hochschild--Serre spectral sequence gives
\[
  E_2^{p q} = \sH^p(I,\eR^q\Psi_\eta(\sF))\Rightarrow i^\ast \eR^{p+q}j_\ast\sF \text{.}
\]
Let $R_t^q=\eR^q\Psi_\eta(\sF)^P$, where $t$ stands for tame. If $\sigma$ is a generator of $\widehat\dZ_{p'}(1)$, then $E_2^{0 q} = \ker(\sigma-1,R_t^q)$, $E_2^{1 q} = \operatorname{coker}(\sigma-1,R_t^q)$ and $E_2^{p q} = 0$ for $p\ne0,1$. These sheaves are constructible because the vanishing-cycle sheaves are, and \ref{VII:3-11} follows.

\subsection{}\label{VII:3-12}

We treat the general case, with $S$ still a strictly henselian trait. If $X=X_\eta$, then $f$ is the composite $X\to Y_\eta\to Y$, and we apply \ref{VII:1-9} and \ref{VII:3-11}. In general, let $j:X_\eta\hookrightarrow X$, and let $\Delta$ be the mapping cone of $\sF\to\eR j_\ast j^\ast\sF$. Its cohomology sheaves are supported on $X_s$. By \ref{VII:1-9}, applied to $X_s\to Y_s$, $\eR f_\ast\Delta$ is therefore constructible. The complex $\eR f_\ast\eR j_\ast j^\ast\sF =\eR(fj)_\ast j^\ast\sF$ is constructible as well, and the conclusion follows from the triangle
\[\xymatrix{
  \ar[r]
    & \eR f_\ast \sF \ar[r]
    & \eR f_\ast \eR j_\ast j^\ast \sF \ar[r]
    & \eR f_\ast \Delta \ar[r]
    & \text{.}
}\]

\section{Local biduality}\label{VII:4}

\subsection{}\label{VII:4-1}

Let $S$ be regular of dimension $0$ or $1$, and put $K_S=A$, the constant sheaf with value $A$. For $K\in\ob\eD_{\mathrm{ctf}}^b(S,A)$, set $DK=\rHom(K,K_S)\in\ob\eD^+(S,A^\circ)$, a complex of sheaves of right $A$-modules. One again has $DK\in\ob\eD_{\mathrm{ctf}}^b(S,A^\circ)$, and an explicit calculation, possible because $S$ has dimension at most $1$, shows that $K\iso DDK$; for $A=\dZ/n$, see \hyperref[V]{Duality}, \ref{V:1-4}.

\subsection{}\label{VII:4-2}

We suppose $A$ commutative, an assumption that is probably unnecessary. For $a:X\to S$ of finite type over $S$, put $K_X=\eR a^!K_S$. For $K\in\ob\eD_{\mathrm{ctf}}^b(X,A)$, set $DK=\rHom(K,K_X)\in\ob\eD_{\mathrm{ctf}}^b(X,A)$. Our main result is:

\begin{theorem_}\label{VII:4-3}
One has $K\iso DDK$ for $K\in\ob\eD_\textnormal{ctf}^b(X,A)$.
\end{theorem_}

\begin{lemma_}\label{VII:4-4}
If $a$ is proper, then $\eR a_\ast K\iso\eR a_\ast DDK$.
\end{lemma_}

Poincaré duality,
\[
  \eR a_! \rHom(K,\eR a^! K_S) = \rHom(\eR a_! K,K_S)
\]
gives an isomorphism $\eR a_\ast D=D\eR a_!=D\eR a_\ast$. The lemma follows from the commutativity of the diagram (\hyperref[V]{Duality}, \ref{V:1-2})
\[\xymatrix{
  \eR a_\ast K \ar[r] \ar[d]^-\wr
    & \eR a_\ast D D K \ar[d]^-\wr \\
  D D \eR a_\ast K \ar[r]^-\sim
    & D \eR a_\ast D K \text{.}
}\]

\subsection{}\label{VII:4-5}

By localization, it is enough to treat the cases in which $S$ is the spectrum of a field or a trait. When $S$ is the spectrum of a field, proceed by induction on $\dim X$. Since the question is local, we may suppose $X$ proper. The induction hypothesis and the usual arguments show that the mapping cone $\Delta$ of $K\to DDK$ has skyscraper cohomology sheaves. The lemma then gives $\eR a_\ast\Delta=0$, hence $\Delta=0$.

\subsection{}\label{VII:4-6}

When $S$ is a trait, \ref{VII:4-3} is already known on the generic fiber by \ref{VII:4-5}. We may again suppose $X$ proper over $S$ and proceed by induction on $\dim X_s$. The mapping cone $\Delta$ again has skyscraper cohomology sheaves, now on $X_s$, and we conclude as above.

\subsection{}\label{VII:4-7}

Suppose $A=\dZ/m$. There is then a variant of the preceding theory for arbitrary $K\in\eD_c^b(X,A)$. The complex $\eR a^!\dZ/m$ has finite injective dimension \cite[XVIII, 3.1.7]{sga4}, which ensures that $DK$ still belongs to $\eD_c^b(X,A)$. To prove local biduality on $S$, use the fact that $\dZ/m$ is the dualizing $\dZ/m$-module, and then proceed as above.

\section{Appendix, by L.\ Illusie}\label{VII:5}

In this appendix, which draws on parts of the former \cite[II]{sga5} and on a letter from P.\ Deligne to the author, we prove \ref{VII:2-14}, as well as several generalizations and variants of the specialization theorems \cite[XVI, 2.1]{sga4} and (\hyperref[I]{Étale cohomology}, \ref{I:5-1-7}).

All schemes and morphisms considered are assumed quasi-compact and quasi-separated. If $X$ is a scheme and $x$ a geometric point of $X$, write $X(x)$ for the strict localization of $X$ at $x$.

Fix a base scheme $S$ and a sheaf of rings $A$ on $S$, not necessarily commutative or noetherian. If $X$ is an $S$-scheme, write $\eD(X)$ for $\eD(X,A_X)$.

\subsection{Cohomological properness}\label{VII:5-1}

\begin{proposition}\label{VII:5-1-1}
Let $f:X\to Y$ be an $S$-morphism and $E\in\ob\eD^+(X)$. The following conditions are equivalent:
\begin{enumerate}[\indent (i)]
  \item The formation of $\eR f_\ast E$ commutes with every finite base change $S'\to S$.
  \item The formation of $\eR f_\ast E$ commutes with every quasi-finite base change, or inverse limit of quasi-finite morphisms, $S'\to S$.
  \item For every specialization map $i:t\to S(s)$ (\hyperref[I]{Étale cohomology}, \ref{I:5-1-2}), let $\bar S$ be the normalization in $k(\bar t)$ of the integral scheme given by the closure of $i(t)$ in $S(s)$. By \cite[IV, 6.15.5, 18.6.11, 18.8.16]{ega}, $\bar S$ is thus integral and local, with closed point radicial over $s$. Let $\bar f:\bar X\to\bar Y$ be obtained from $f$ by the base change $\bar S\to S$, and let $f_s:X_s\to Y_s$ be the fiber of $f$ at $s$. Then the base-change map
    \[
      (\eR \bar f_\ast (E|\bar X))|Y_s \to \eR {f_s}_\ast(E|X_s)
    \]
    is an isomorphism.
\end{enumerate}
\end{proposition}

The equivalence of (i) and (ii) follows from Zariski's Main Theorem, and it is clear that (ii) implies (iii). We prove that (iii) implies (ii). For every $S'\to S$ that is an inverse limit of quasi-finite morphisms, and every geometric point $s$ of $S'$, let $C(S',s)$ be the cone of the base-change map
\[
  (\eR f_\ast'(E|X'))|Y_s \to \eR{f_s}_\ast (E|X_s) \text{,}
\]
where $f':X'\to Y'$ is obtained from $f$ by the base change $S'\to S$, and $f_s:X_s\to Y_s$ is the fiber of $f$ at $s$. We must prove that $C(S',s)$ is acyclic for every such $(S',s)$. We show by induction on $N$ that $\h^iC(S',s)=0$ for $i\leqslant N$. By the Main Theorem, passage to the limit, and elimination of radicial morphisms, we may suppose that $S$ and $S'$ are strictly local, that $S'$ is integral with closed point $s$, and that $S'\to S$ is finite. Let $t$ be a geometric generic point of $S'$, and let $\bar S$ be the normalization of $S'$ in $k(t)$. Let $\bar S_\bullet$ be the simplicial coskeleton of $\bar S\to S'$, so that $\bar S_0=\bar S,\ldots,\bar S_n=(\bar S/S')^{n+1},\ldots$, and put $X_\bullet=X\times_S\bar S_\bullet$. Let $g_\bullet:\bar X_\bullet\to X'$ be the canonical projection, and write $E'$ for the inverse image of $E$ on $X'$. The bicomplex ${g_\bullet}_\ast g_\bullet^\ast E'$, whose columns are ${g_n}_\ast g_n^\ast E'$, is a resolution of $E'$ \cite[VIII, 8]{sga4}. Thus $C(S',s)$ can be represented by a bicomplex whose columns are the $C(\bar S_n,s)$. By (iii), $C(\bar S_0,s)=C(\bar S,s)$ is acyclic. For $n>0$, the induction hypothesis gives $\h^iC(\bar S_n,s)=0$ for $i\leqslant N$. It follows that $\h^iC(S',s)=0$ for $i\leqslant N+1$, completing the proof.

\subsubsection{Remark}\label{VII:5-1-2}

Suppose $(f,E)$ satisfies condition (i) of \ref{VII:5-1-1}.
\begin{enumerate}[a)]
  \item If $E\in\ob\eD^b(X)$ and $f_\ast$ has finite cohomological dimension, then for every $F\in\ob\eD^-(S,A^\circ)$, the canonical map
    \begin{equation*}\tag{$*$}\label{VII:eq:5-1-1-1}
      \eR f_\ast(E)\lotimes_A q^\ast F \to \eR f_\ast(E\lotimes_A p^\ast F) \text{,}
    \end{equation*}
    where $p:X\to S$ and $q:Y\to S$ are the structural morphisms, is an isomorphism. The same conclusion holds with $E\in\ob\eD^-(X)$.
\end{enumerate}

By dévissage, it is enough to verify the assertion for $F=i_!A_{S'}$, with $i:S'\to S$ quasi-finite. By the Main Theorem, one reduces to the case in which $i$ is finite; the conclusion is then equivalent to commutation of $\eR f_\ast(E)$ with base change by $i$.
\begin{enumerate}[a)]
\setcounter{enumi}{1}
  \item Suppose $A$ is noetherian and constructible. Then for every constructible $F\in\ob\eD^b(S,A^\circ)$ of finite Tor-dimension, the map \eqref{VII:eq:5-1-1-1} above is an isomorphism.
\end{enumerate}

Indeed, dévissage reduces to the case $F=i_!M$, where $i:S'\to S$ is locally closed, $A$ is locally constant on $S'$, and $M$ is constructible, of finite Tor-dimension, with locally constant cohomology on $S'$. After base change from $S$ to $S'$, we are reduced to the case in which $F$ is perfect \cite[I]{sga6}, and finally, by localization and dévissage, to $F=A$, for which the conclusion is immediate.

\subsubsection{}\label{VII:5-1-3}

Let $f:X\to Y$ be an $S$-morphism and $E\in\ob\eD^+(X)$. We say that $(f,E)$ is \emph{cohomologically proper relative to $S$} if the formation of $\eR f_\ast E$ commutes with every base change $S'\to S$. Equivalently, after every base change $S'\to S$, the formation of $\eR f'_\ast E'$, where $f'=f\times_SS'$ and $E'$ is the inverse image of $E$ on $X\times_SS'$, commutes with every finite base change $S''\to S'$ (cf.\ \cite[XII, 6.1]{sga4}). Here are some examples of cohomological properness. For the analogous notion for sheaves of sets or noncommutative groups, see \cite[XIII]{sga1}.
\begin{enumerate} % originally 1.3.1, 1.3.2, ...
  \item Suppose $A$ is torsion and $f$ is proper. Then for every $E\in\ob\eD^+(X)$, $(f,E)$ is cohomologically proper relative to $S$ by the proper base-change theorem \cite[XII, 5.1]{sga4}.
  \item Suppose $S$ is finite, $X$ and $Y$ are of finite type over $S$, and $A$ is annihilated by an integer invertible on $S$. Then for every $E\in\ob\eD^+(X)$, $(f,E)$ is cohomologically proper relative to $S$ by \ref{VII:1-9}.
  \item Suppose $A$ is constant and noetherian, annihilated by an integer $n$ invertible on $S$, and $f:X=Y\setminus D\to Y$ is the inclusion of the complement of a divisor $D$ on $Y$ with normal crossings relative to $S$ \cite[XIII, 2.1]{sga1}. Let $E$ be a sheaf of $A$-modules on $X$ satisfying one of the following conditions:
    \begin{enumerate}[(i)]
      \item $E$ is constructible locally constant and tamely ramified along $D$;
      \item étale-locally on $Y$ and on $S$, $E$ is the inverse image of a sheaf of $A$-modules on $S$.
    \end{enumerate}
    Then $(f,E)$ is cohomologically proper relative to $S$, and $\eR f_\ast E$ is constructible (compare \cite[XIII, 2.4]{sga1}).
\end{enumerate}

We sketch the proof in case (i). The question is local near a point $y$ of $Y$. We may suppose that $D$ is a sum of smooth divisors, $D=\sum_{i=1}^mD_i$. By the relative Abhyankar lemma \cite[XIII, 5.5]{sga1}, near $y$ one can trivialize $E$ by a finite covering $\widetilde Y\to Y$ of the form
\[
  \widetilde Y
  =Y[T_1,\dots,T_r]/(T_1^{n_1}-t_1,\dots,T_r^{n_r}-t_r),
\]
where the $t_i$ are local equations for the smooth divisors passing through $y$, and the $n_i$ are prime to the characteristic of $k(y)$. The inverse image $\widetilde D$ of $D$ in $\widetilde Y$ still has normal crossings relative to $S$, and if $g:\widetilde X=\widetilde Y\setminus\widetilde D\to X$ is the projection, then $g^\ast E$ is constant. Since $E$ injects into $g_\ast g^\ast E$ and the quotient is tamely ramified, an easy dévissage reduces us to the case in which $E$ is constant. Let $p:X\to S$ and $q:Y\to S$ be the structural morphisms, and for $1\leqslant i\leqslant m$ let $f_i:Y\setminus D_i\to Y$ be the canonical inclusion. The relative purity theorem \cite[XVI, 3.7]{sga4} readily implies, by induction on $m$, that for every locally constant $A$-module $M$ on $S$, the canonical map
\begin{equation*}\tag{5.1.3.1}\label{VII:eq:5-1-3-3-1} % originally 1.3.3.1
  \eR f_\ast (\dZ/n)\lotimes q^\ast M\to \eR f_\ast (p^\ast M)
\end{equation*}
is an isomorphism, and moreover:
\begin{align*}\tag{5.1.3.2}\label{VII:eq:5-1-3-3-2}
  {f_i}_\ast(\dZ/n) &= \dZ/n \text{,} \\
  \eR^1 {f_i}_\ast (\dZ/n) &= (\dZ/n)(-1)_{D_i} \text{,}
\end{align*}
a canonical isomorphism being given by the fundamental class $\cl(D_i)$ (\hyperref[IV]{Cycle}, \ref{IV:2-1-4}), and
\begin{align*}
  \eR^q {f_i}_\ast(\dZ/n) &= 0\quad\text{for $q>1$}, \\
  \eR f_\ast (\dZ/n) &= \bigotimes_i^\eL \eR {f_i}_\ast (\dZ/n),
    \quad\text{i.e.}, \\
  \eR^1 f_\ast(\dZ/n) &= \bigoplus_{i=1}^m (\dZ/n)(-1)_{D_i} \text{,} \\
  \textstyle\bigwedge^\bullet \eR^1 f_\ast(\dZ/n) &\iso \eR^\bullet f_\ast (\dZ/n) \text{.}
\end{align*}
The conclusion of Example~3 in case (i) follows from \eqref{VII:eq:5-1-3-3-1} and \eqref{VII:eq:5-1-3-3-2}. In case (ii), passage to the limit reduces to $E=p^\ast M$ with $M$ constructible; dévissage and base change, using (i), then reduce to the already treated case in which $M$ is constant.

The preceding proof shows that \eqref{VII:eq:5-1-3-3-1} is an isomorphism for every $M\in\ob\eD^+(S)$.

Compare \eqref{VII:eq:5-1-3-3-2} with the analogous formulas in integral cohomology when $S=\spec(\dC)$ (see, for example, \cite[3.1]{de71-2}). Those formulas follow from the fact that, locally for the classical topology, the complement of $D$ has the homotopy type of a torus.

Passing to the limit in \eqref{VII:eq:5-1-3-3-2} gives similar formulas in $\ell$-adic cohomology, with $\dZ/n$ replaced by $\dZ_\ell$, where $\ell$ is a prime invertible on $S$. Suppose $q:Y\to S$ is proper and smooth. The Weil conjectures \cite{de74,de80} then imply that the Leray spectral sequence
\[
  E_2^{i j} = \eR^i q_\ast \eR^j f_\ast \dZ_\ell \Rightarrow \eR^\bullet p_\ast \dZ_\ell
\]
degenerates at $E_3$ modulo torsion \cite[\S6]{de71-1}. The same phenomenon occurs in integral cohomology when $S=\spec(\dC)$; see the reference just cited and \cite[3.2.13]{de71-2}.

\subsection{Specialization and cospecialization}\label{VII:5-2}

Let $f:X\to S$ and $E\in\ob\eD^+(X)$.

\subsubsection{}\label{VII:5-2-1}

Suppose that the formation of $\eR f_\ast E$ commutes with every finite base change. For every specialization $u:t\to S(s)$ of geometric points of $S$, define a map, called the \emph{specialization map},
\begin{equation*}\tag{5.2.1.1}\label{VII:eq:5-2-1-1}
  \spe(u):\eR\Gamma(X_s,E|X_s) \to \eR\Gamma(X_t,E|X_t) \text{,}
\end{equation*}
as follows. As in \ref{VII:5-1-1}(iii), let $\bar S$ be the normalization in $k(t)$ of the closure of $u(t)$ in $S(s)$, and let $\bar f:\bar X\to\bar S$ be obtained from $f$ by the base change $\bar S\to S$. By the immediate implication (i)$\Rightarrow$(iii) of \ref{VII:5-1-1}, the base-change map
\begin{equation*}\tag{5.2.1.2}\label{VII:eq:5-2-1-2}
  \eR\Gamma(\bar X,E|\bar X)=\eR\bar f_\ast (E|\bar X)_s \to \eR\Gamma(X_s,E|X_s)
\end{equation*}
is an isomorphism. Define \eqref{VII:eq:5-2-1-1} as the composite of the inverse of \eqref{VII:eq:5-2-1-2} with the restriction map
\begin{equation*}\tag{5.2.1.3}\label{VII:eq:5-2-1-3}
  \eR\Gamma(\bar X,E|\bar X) \to \eR\Gamma(X_t,E|X_t) \text{.}
\end{equation*}

\subsubsection{}\label{VII:5-2-2}

Suppose that $f$ is locally acyclic relative to $E$ (\ref{VII:2-12}). For every specialization $u:t\to S(s)$ of geometric points of $S$, define a map, called the \emph{cospecialization map},
\begin{equation*}\tag{5.2.2.1}\label{VII:eq:5-2-2-1}
  \cosp(u):\eR\Gamma(X_t,E|X_t) \to \eR\Gamma(X_s,E|X_s) \text{,}
\end{equation*}
as follows. Let $g:X_t\to X$ and $g':\bar X\to X$ be the canonical maps. The local acyclicity hypothesis implies that the canonical map
\begin{equation*}\tag{5.2.2.2}\label{VII:eq:5-2-2-2}
  \eR g_\ast'(E|\bar X)\to \eR g_\ast(E|X_t)
\end{equation*}
is an isomorphism. Indeed, we may suppose $S$ strictly local and integral, with closed point $s$, and $t$ a geometric generic point. Compute the stalk of \eqref{VII:eq:5-2-2-2} at a geometric point $x$ of $X$, with image $y$ in $S$. Replacing $S$ by its strict localization at $y$, we may suppose $y=s$. Since $\bar S\to S$ is an inverse limit of finite morphisms and is radicial at $s$, one finds $\eR g'_\ast(E|\bar X)_x=E_x$, and the stalk of \eqref{VII:eq:5-2-2-2} at $x$ is identified with the restriction map
\[
  E_x=\eR\Gamma(X(x),E)\longrightarrow
  \eR\Gamma(X(x)_t,E|X(x)_t),
\]
which is an isomorphism by hypothesis. Applying $\eR\Gamma(X,-)$ to \eqref{VII:eq:5-2-2-2}, we deduce that the restriction map \eqref{VII:eq:5-2-1-3} is an isomorphism. Define \eqref{VII:eq:5-2-2-1} as the composite of the inverse of \eqref{VII:eq:5-2-1-3} and the base-change map \eqref{VII:eq:5-2-1-2}.

\subsubsection{}\label{VII:5-2-3}

If $s'' \xrightarrow{u'} s' \xrightarrow u s$ are specialization maps of geometric points of $S$, one checks that
\begin{align*}
  \spe(u')\spe(u) &= \spe(u u'), \\
  \cosp(u) \cosp(u') &= \cosp(u u'),
\end{align*}
whenever both sides are defined.

\subsubsection{}\label{VII:5-2-4}

Suppose that the formation of $\eR f_\ast E$ commutes with every finite base change and that $f$ is locally acyclic relative to $E$. It follows from the definitions that, for every specialization $u:t\to S(s)$ of geometric points of $S$, the maps $\spe(u)$ and $\cosp(u)$ are isomorphisms inverse to one another.

These hypotheses hold in particular when $A$ is noetherian and annihilated by an integer $n$ invertible on $S$, when $f$ is proper and smooth, and when $E$ has locally constant cohomology (cohomological properness of proper morphisms (\ref{VII:5-1-3}) and local acyclicity of smooth morphisms \cite[XV]{sga4}). This recovers the ``specialization theorem'' \cite[XVI 2.1]{sga4}.

On the other hand, \ref{VII:5-1-1} has the following consequence.

\begin{proposition}\label{VII:5-2-5}
Suppose that $f$ is locally acyclic relative to $E$, and that for every specialization $u:t\to S(s)$ the cospecialization map $\cosp(u)$ is an isomorphism. Then the formation of $\eR f_\ast E$ commutes with every finite base change.
\end{proposition}

In view of \ref{VII:5-2-4}, \ref{VII:5-2-5} generalizes (\hyperref[I]{Étale cohomology} \ref{I:5-1-7}).

\begin{corollary}\label{VII:5-2-6}
Suppose that $f$ is locally acyclic relative to $E$. Let $x$ be a geometric point of $X$, with image $s$ in $S$, and let $f_{(x)}:X(x)\to S(s)$ be the morphism induced by $f$. Then the formation of $\eR {f_{(x)}}_\ast(E|X(x))$ commutes with every finite base change $S'\to S(s)$.
\end{corollary}

By \ref{VII:5-2-5}, it suffices to check that the cospecialization maps for $(f_{(x)},E|X(x))$ are isomorphisms. By transitivity of cospecialization maps (\ref{VII:5-2-3}), it is enough to show that, if $u:t\to S(s)$ is a specialization, then $\cosp(u):\eR\Gamma(X(x)_t,E|X(x)_t)\to \eR\Gamma(X(x)_s,E|X(x)_s)$ is an isomorphism. Since $s$ is closed in $S(s)$, however, $X(x)_s$ is strictly local with closed point $x$, $\eR\Gamma(X(x)_s,E|X(x)_s)=E_x$, and $\cosp(u)$ is an isomorphism, inverse to the restriction isomorphism $E_x\to \eR\Gamma(X(x)_t,E|X(x)_t)$.

\begin{corollary}\label{VII:5-2-7}
Let $f:X\to Y$ be an $S$-morphism, let $p:X\to S$ and $q:Y\to S$ be the projections, and let $E\in\ob\eD^+(X)$. Suppose that $A$ is locally constant, that $f$ is locally acyclic relative to $E$, and that $q$ is locally acyclic at every geometric point $y$ of $Y$ relative to every sheaf that is locally constant in a neighborhood of $y$. Then $p$ is locally acyclic relative to $E$.
\end{corollary}

Let $x$ be a geometric point of $X$, with images $y$ and $s$ in $Y$ and $S$, let $t\to S(s)$ be a specialization, and let $f_{(x)}:X(x)\to Y(y)$ be the morphism induced by $f$. We must show that the restriction map
\begin{equation*}\tag{$*$}\label{VII:eq:5-2-7-1}
  E_x \to \eR\Gamma(X(x)_t,E|X(x)_t)
\end{equation*}
is an isomorphism. We may suppose that $A$ is constant. By \ref{VII:5-2-4} and \ref{VII:5-2-5}, for every specialization $z\to Y(y)$, the specialization map
\[
  E_x = \bigl(\eR {f_{(x)}}_\ast(E|X(x))\bigr)_y
  \longrightarrow
  \bigl(\eR {f_{(x)}}_\ast(E|X(x))\bigr)_z
  = \eR\Gamma(X(x)_z,E|X(x)_z)
\]
is an isomorphism. In other words, $\eR{f_{(x)}}_\ast(E|X(x))$ is constant with value $E_x$. The acyclicity hypothesis on $q$ therefore implies that the restriction map
\[
  E_x \longrightarrow \eR\Gamma(Y(y)_t,E_x)
  = \eR\Gamma\left(
      Y(y)_t,\eR {f_{(x)}}_\ast(E|X(x))|Y(y)_t
    \right)
\]
is an isomorphism. This map identifies with \eqref{VII:eq:5-2-7-1}, which proves the assertion.

The hypothesis of \ref{VII:5-2-7} on $q$ holds in particular when $q$ is smooth and $A$ is annihilated by an integer $n$ invertible on $S$ \cite[XV]{sga4}. Moreover, if the hypotheses of \ref{VII:5-2-7} hold after every base change $S'\to S$, then $p$ is universally locally acyclic relative to $E$. In particular, assertion \ref{VII:2-14} follows from these remarks.

\begin{corollary}\label{VII:5-2-8}
Let $f:X\to S$ and $E\in\ob\eD^+(X)$. Suppose that $f$ is universally locally acyclic relative to $E$, that the formation of $\eR f_\ast E$ commutes with every finite base change, and that for every morphism $s'\to s$, where $s$ is a geometric point of $S$ and $s'$ is the spectrum of a separably closed field, the base-change map $\eR\Gamma(X_s,E|X_s)\to \eR\Gamma(X_{s'},E|X_{s'})$ is an isomorphism. Then $(f,E)$ is cohomologically proper (\ref{VII:5-1-3}).
\end{corollary}

Let $f':X'\to S'$ be obtained from $f$ by a base change $S'\to S$. We must show that the formation of $\eR f'_\ast(E|X')$ commutes with every finite base change. By \ref{VII:5-2-5}, it suffices to show that, for every specialization $u':t'\to S'(s')$, the cospecialization map
\[
  \cosp(u'):
  \eR\Gamma(X_{t'}',E|X_{t'}')
  \longrightarrow
  \eR\Gamma(X_{s'}',E|X_{s'}')
\]
is an isomorphism. Let $s$ (respectively $t$) be the geometric point that is the image of $s'$ (respectively $t'$) in $S$, and let $u:t\to S(s)$ be the specialization defined by $u'$. There is plainly a commutative square
\[\xymatrix{
  \eR\Gamma(X_t,E|X_t) \ar[r]^-{\cosp(u)} \ar[d]
    & \eR\Gamma(X_s,E|X_s) \ar[d] \\
  \eR\Gamma(X_{t'}',E|X_{t'}') \ar[r]^-{\cosp(u')}
    & \eR\Gamma(X_{s'}',E|X_{s'}')
}\]
whose vertical arrows are the base-change maps. By hypothesis they are isomorphisms. Moreover, by \ref{VII:5-2-4}, $\cosp(u)$ is an isomorphism; hence $\cosp(u')$ is an isomorphism, which completes the proof.

Here is an application of \ref{VII:5-2-8}. Suppose that $A$ is annihilated by an integer $n$ invertible on $S$, and that $f$ is locally of finite type and universally locally acyclic. Let $x$ be a geometric point of $X$, with image $s$ in $S$, and let $f_{(x)}:X(x)\to S(s)$ be the morphism induced by $f$. It follows from \ref{VII:5-2-6} and \ref{VII:1-9} (by passage to the limit) that the hypotheses of \ref{VII:5-2-8} hold for the pair $(f_{(x)},E|X(x))$, which is therefore cohomologically proper.

\subsubsection{}\label{VII:5-2-9}

Suppose that $A$ is constant. Let $f:X\to S$, and let $E\in\ob\eD^+(X)$ be a complex of finite Tor-dimension. We shall say that $f$ is \emph{strongly locally acyclic} relative to $E$ if, for every geometric point $x$ of $X$, with image $s$ in $S$, every specialization $t\to S(s)$, and every $A^\circ$-module $M$, the restriction map
\[
  E_x\lotimes_A M \to \eR\Gamma(X(x)_t,(E|X(x)_t)\lotimes_A M)
\]
is an isomorphism. I do not know whether ``locally acyclic'' implies ``strongly locally acyclic.''

\begin{proposition}\label{VII:5-2-10}
Under the hypotheses of \ref{VII:5-2-9}, suppose that $A$ is noetherian and torsion, that $S$ is noetherian, and that $f$ is strongly locally acyclic relative to $E$. Then, for every diagram with cartesian squares
\[\xymatrix{
  X \ar[d]^-f
    & \ar[l]_-j X' \ar[d]^-{f'}
    & \ar[l]_-{j'} X'' \ar[d]^-{f''} \\
  S
    & \ar[l]_-i S'
    & \ar[l]_-{i'} S'' \text{,}
}\]
in which $i$ is quasi-finite and $i'$ is an open immersion, and for every $F\in\ob\eD^+(S'',A^\circ)$, the canonical map
\begin{equation*}\tag{$*$}\label{VII:eq:5-2-10-1}
  j^\ast E \lotimes {f'}^\ast \eR i_\ast' F \to \eR j_\ast' ((j j')^\ast E\lotimes {f''}^\ast F)
\end{equation*}
is an isomorphism.
\end{proposition}

Compare this statement with \cite[XV 1.17]{sga4}.

We prove \ref{VII:5-2-10}. It is easy to reduce to the case $S=S'$, with $S$ strictly local and $F$ constructible and concentrated in degree $0$. Resolve $F$ on the right by finite sums of sheaves of the form $u_\ast M$, where $u:t\to S$ is a geometric point and $M$ is an $A^\circ$-module. The hypothesis implies that \eqref{VII:eq:5-2-10-1} is an isomorphism for $F=u_\ast M$, and the result follows.
% END INLINED FILE: sga4.5-7.tex
\clearpage
\clearpage
% BEGIN INLINED FILE: sga4.5-8.tex
\chapter{Derived categories}\label{VIII}
\blfootnote{by Jean-Louis Verdier}

\section{Triangulated categories: definitions and examples}\label{VIII:1}

\subsection{Definition of triangulated categories}\label{VIII:1-1}

\addtocounter{subsubsection}{-1}
\subsubsection{}\label{VIII:1-1-0}

Let $\cA$ be an additive category and $T$ an additive automorphism of $\cA$. For two objects $X$ and $Y$ of $\cA$, set
\[
  \hom^i(X,Y) = \hom_\cA(X,T^i(Y)) \qquad i\in \dZ
\]
Let $X$, $Y$, and $Z$ be three objects of $\cA$, with $\alpha\in\hom^i(X,Y)$ and $\beta\in\hom^j(Y,Z)$. Define their composite $\beta\circ\alpha\in\hom^{i+j}(X,Z)$ by
\[
  \beta\circ\alpha=T^i(\beta)\circ\alpha .
\]
One checks that this gives a new category whose morphism groups are graded. By abuse of language, this new category will be called the category $\cA$ graded by the translation functor $T$.

Let $\cA$ be a category graded by a translation functor $T$. Whenever a morphism is named or denoted without specifying its degree, it will always mean a morphism of degree zero.

Let $\cA$ and $\cA'$ be two additive categories graded by the functors $T$ and $T'$, and let $F:\cA\to\cA'$ be an additive functor. We say that $F$ \emph{is graded} if an isomorphism of functors has been given\footnote{Text modified in 1976: the original text required an equality $F\circ T=T'\circ F$. The case of functors in several variables is discussed in \cite[XVIII 0.2]{sga4}; it presents a sign problem.}
\[
  F\circ T = T'\circ F
\]

The definition of morphisms of graded functors is given in \ref{VIII:4-1-1}.

Let $\cA$ be an additive category graded by the functor $T$. A \emph{triangle} consists of three objects $X,Y,Z$ and three morphisms $u:X\to Y$, $v:Y\to Z$, and $w:Z\to T(X)$ ($\deg w=1$). We denote a triangle by $(X,Y,Z,u,v,w)$ or by the diagram
\[\xymatrix{
  & Z \ar[dl]_-w \\
  X \ar[rr]^-u
    & & Y \ar[lu]_-v
}\qquad \deg(w)=1
\]
For two triangles $(X,Y,Z,u,v,w)$ and $(X',Y',Z',u',v',w')$ of $\cA$, a morphism from the first triangle to the second is a triple of morphisms $f:X\to X'$, $g:Y\to Y'$, and $h:Z\to Z'$ such that the following diagram commutes:
\[\xymatrix{
  X \ar[r]^-u \ar[d]^-f
    & Y \ar[r]^-v \ar[d]^-g
    & Z \ar[r]^-w \ar[d]^-h
    & T(X) \ar[d]^-{T(f)} \\
  X' \ar[r]^-{u'}
    & Y' \ar[r]^-{v'}
    & Z' \ar[r]^-{w'}
    & T(X')
}\]

\subsubsection{}\label{VIII:1-1-1}

A \emph{triangulated category} is an additive category graded by a translation functor, together with a family of triangles called the family of distinguished triangles. This family must satisfy the following axioms.

\paragraph{TR1}
\phantomsection\hypertarget{VIII:TR1}{}

Every triangle isomorphic to a distinguished triangle is distinguished. Every morphism $u:X\to Y$ belongs to a distinguished triangle $(X,Y,Z,u,v,w)$. The triangle $(X,X,0,\operatorname{id}_{X},0,0)$ is distinguished.

\paragraph{TR2}
\phantomsection\hypertarget{VIII:TR2}{}

The triangle $(X,Y,Z,u,v,w)$ is distinguished if and only if $(Y,Z,T(X),v,w,-T(u))$ is distinguished.

\paragraph{TR3}
\phantomsection\hypertarget{VIII:TR3}{}

If $(X,Y,Z,u,v,w)$ and $(X',Y',Z',u',v',w')$ are distinguished, then for every morphism $(f,g):u\to u'$ there is a morphism $h:Z\to Z'$ such that $(f,g,h)$ is a morphism of triangles.

\paragraph{TR4}
\phantomsection\hypertarget{VIII:TR4}{}

\[\xymatrix{
  & Z' \ar[dl]|{\deg=1} \\
  X \ar[rr]^-u
    & & Y \ar[ul]_-i
}\quad
\xymatrix{
  & X' \ar[dl]|{\deg(j)=1} \\
  Y \ar[rr]^-v
    & & Z \ar[ul]
}\quad
\xymatrix{
  & Y' \ar[dl]|{\deg=1} \\
  X \ar[rr]^-w
    & & Z \ar[ul]
}\]
If these are three distinguished triangles such that $w=v\circ u$, there are two morphisms $f:Z'\to Y'$ and $g:Y'\to X'$ such that
\begin{enumerate}
  \item $(\operatorname{id}_X,v,f)$ is a morphism of triangles;
  \item $(u,\operatorname{id}_Z,g)$ is a morphism of triangles;
  \item $(Z',Y',X',f,g,T(i)j)$ is a distinguished triangle.
\end{enumerate}

For two triangulated categories $\cC$ and $\cC'$, an \emph{exact functor} from $\cC$ to $\cC'$ is an additive graded functor that takes distinguished triangles to distinguished triangles. Morphisms of exact functors are morphisms of graded functors.

Axioms \hyperlink{VIII:TR1}{TR1}, \hyperlink{VIII:TR2}{TR2}, and \hyperlink{VIII:TR3}{TR3} immediately imply the following property.

\begin{proposition}\label{VIII:1-1-2}
Let $(X,Y,Z,u,v,w)$ be a distinguished triangle and $M$ an object of a triangulated category. Then the following sequences are exact:
\[\xymatrix{
  \cdots \ar[r]
    & \hom^i(M,X) \ar[r]^-{\hom(M,u)}
    & \hom^i(M,Y) \ar[r]^-{\hom(M,v)}
    & \hom^i(M,Z) \ar[r]^-{\hom(M,w)}
    & \hom^{i+1}(M,X) \ar[r]
    & \cdots
}\]
\[\xymatrix{
  \cdots \ar[r]
    & \hom^i(Z,M) \ar[r]^-{\hom(v,M)}
    & \hom^i(Y,M) \ar[r]^-{\hom(u,M)}
    & \hom^i(X,M) \ar[r]^-{\hom(w,M)}
    & \hom^{i+1}(Z,M) \ar[r]
    & \cdots
}\]
\end{proposition}

This proposition then yields statements such as the following:
\begin{itemize}
  \item If $(f,g,h)$ is a morphism of distinguished triangles and $f$ and $g$ are isomorphisms, then $h$ is an isomorphism.
  \item The distinguished triangles constructed on a given morphism (\hyperlink{VIII:TR1}{TR1}) are all isomorphic (\emph{but the isomorphisms are not in general uniquely determined}).
  \item If, in a distinguished triangle, one of the morphisms has a kernel or a cokernel, the triangle splits; that is, it has the form
    \[\xymatrix{
      & Z+T(X) \ar[dl] \\
      X+Y \ar[rr]
        & & Y+Z \ar[ul]
    }\]
  \item The direct sum of two distinguished triangles is distinguished. (One uses axiom \hyperlink{VIII:TR4}{TR4}.)
\end{itemize}

\subsection{Examples}\label{VIII:1-2}

\subsubsection{}\label{VIII:1-2-1}

We first fix some notation. For an additive category $\cA$, let $\eC(\cA)$ denote the following category:
\begin{itemize}
  \item the objects of $\eC(\cA)$ are complexes in $\cA$, unbounded in either direction, whose differential has degree $+1$;
  \item the morphisms of $\eC(\cA)$ are degree-preserving morphisms of complexes (that is, morphisms commuting with the differential).
\end{itemize}

Let $T$ be the following functor: for every object $X^\bullet$ of $\eC(\cA)$,
\begin{align*}
  T(X^\bullet)_i &= (X^\bullet)_{i+1} \\
  d_i(T(X^\bullet)) &= -d_{i+1}(X^\bullet)  \text{.}
\end{align*}
On morphisms, $T$ acts as follows: for every morphism $f:X^\bullet\to Y^\bullet$, the morphism $T(f):T(X^\bullet)\to T(Y^\bullet)$ is defined by
\[
  T(f)_i = f_{i+1} \text{.}
\]
It is immediate that $T$ is an additive automorphism of $\eC(\cA)$. Morphisms of degree $n$, as defined by this translation functor, are then precisely the morphisms that increase degree by $n$ and commute or anticommute with the differential according to the parity of $n$. Thus we recover the usual definition.

Let $\sC$ denote the following family of triangles. A triangle $(X^\bullet,Y^\bullet,Z^\bullet,u,v,w)$ belongs to this family if\footnote{Text modified in 1976 to ensure the validity of \ref{VIII:1-2-4}.}
\begin{itemize}
  \item $Z^\bullet$ is the total complex associated with the double complex
    \[\xymatrix{
      \cdots \ar[r]
        & 0 \ar[r]
        & X^\bullet \ar[r]^-u
        & Y^\bullet \ar[r]
        & 0 \ar[r]
        & \cdots
    }\]
    in which the objects of $Y^\bullet$ have first degree zero and those of $X^\bullet$ have first degree $-1$;
  \item $v$ is the image, under the functor from double complexes to total complexes, of the following morphism of double complexes:
    \[\xymatrix{
      \cdots \ar[r]
        & 0 \ar[r] \ar[d]
        & 0^\bullet \ar[r] \ar[d]
        & Y^\bullet \ar[r] \ar[d]
        & 0 \ar[r] \ar[d]
        & \cdots \\
      \cdots \ar[r]
        & 0 \ar[r]
        & X^\bullet \ar[r]
        & Y^\bullet \ar[r]
        & 0 \ar[r]
        & \cdots
    }\]
  \item $w$ is the negative of the image, under the same functor, of the following morphism of double complexes:
    \[\xymatrix{
      \cdots \ar[r]
        & 0 \ar[r]
        & X^\bullet \ar[r]
        & 0 \ar[r]
        & 0 \ar[r]
        & \cdots \\
      \cdots  \ar[r]
        & 0 \ar[r] \ar[u]
        & X^\bullet \ar[r] \ar[u]
        & Y^\bullet \ar[r] \ar[u]
        & 0 \ar[r] \ar[u]
        & \cdots
    }\]
\end{itemize}

\subsubsection{}\label{VIII:1-2-2}

The null-homotopic morphisms form a two-sided ideal (if $f$ and $g$ are composable and either $f$ or $g$ is null-homotopic, their composite is null-homotopic; if $f$ and $g$ are null-homotopic, then so is $f\oplus g$). For each pair of objects $X^\bullet,Y^\bullet$, let $\operatorname{Htp}(X^\bullet,Y^\bullet)$ denote the subgroup of $\hom_{\eC(\cA)}(X^\bullet,Y^\bullet)$ consisting of null-homotopic morphisms. Define $\eK(\cA)$ to be the category whose objects are those of $\eC(\cA)$ and in which the group of morphisms from $X^\bullet$ to $Y^\bullet$ is
\[
  \hom_{\eC(\cA)}(X^\bullet,Y^\bullet) / \operatorname{Htp}(X^\bullet,Y^\bullet) \text{.}
\]
Composition in $\eK(\cA)$ is induced from composition in $\eC(\cA)$ by passing to the quotient. The category $\eK(\cA)$ is additive. The translation functor $T$ on $\eC(\cA)$ passes to the quotient (the translate of a null-homotopic morphism is null-homotopic). The resulting additive automorphism of $\eK(\cA)$ will again be denoted by $T$. By abuse of notation, $\eK(\cA)$ will henceforth denote the category $\eK(\cA)$ graded by $T$.

A distinguished triangle in $\eK(\cA)$ is any triangle \emph{isomorphic} to one arising from $\sC$. The family of distinguished triangles satisfies axioms \hyperlink{VIII:TR1}{TR1}, \hyperlink{VIII:TR2}{TR2}, \hyperlink{VIII:TR3}{TR3}, and \hyperlink{VIII:TR4}{TR4}. By abuse of notation, $\eK(\cA)$ will also denote $\eK(\cA)$ equipped with this triangulated structure.

\subsubsection{}\label{VIII:1-2-3}

We shall also use the following collection of subcategories:
\begin{itemize}
  \item A complex is \emph{bounded below} if all but finitely many of its objects in negative degrees vanish. Bounded-above and bounded complexes are defined similarly. Let $\eK^+(\cA)$, $\eK^-(\cA)$, and $\eK^b(\cA)$ denote the full subcategories of $\eK(\cA)$ on these respective families of objects. These are ``\emph{triangulated subcategories}'' of $\eK(\cA)$: every distinguished triangle two of whose objects belong to the subcategory is isomorphic to a triangle all three of whose objects belong to it.
  \item Suppose that $\cA$ is abelian. A complex is \emph{cohomologically bounded below} if all but finitely many of its cohomology objects in negative degrees vanish. Cohomologically bounded-above, cohomologically bounded, and acyclic complexes are defined similarly. Let $\eK^{\infty,+}(\cA)$, $\eK^{\infty,-}(\cA)$, $\eK^{\infty,b}(\cA)$, and $\eK^{\infty,\varnothing}(\cA)$ denote the full subcategories on these families of objects. They are triangulated subcategories.

    We likewise have triangulated subcategories $\eK^{+,b}(\cA)$, $\eK^{+,\varnothing}(\cA)$, and so on (the first superscript describes the objects of the complexes, and the second their cohomology objects).
  \item Without assuming that $\cA$ is abelian, let $O$ be a collection of objects of $\cA$ stable under isomorphisms and direct sums. Let $O(\cA)$ denote the full subcategory of $\eC(\cA)$ consisting of complexes whose object in every degree belongs to $O$, and let $\underline O(\cA)$ denote the corresponding triangulated subcategory of $\eK(\cA)$. When $\cA$ is abelian, one may take for $O$ the collection $I$ of injective objects or the collection $P$ of projective objects.
\end{itemize}

\subsubsection{}\label{VIII:1-2-4}

The category $\eK(\cA)$ is uniquely determined by $\eC(\cA)$ and the family $\sC$. More precisely, every additive graded functor from $\eC(\cA)$ to a triangulated category that sends each triangle in $\sC$ to a distinguished triangle factors uniquely through $\eK(\cA)$, and the resulting functor is exact.

For every complex $X^\bullet$ of $\eC(\cA)$, let $\bar X^\bullet$ denote the complex obtained from $X^\bullet$ by setting its differential equal to zero. A three-term sequence
\[\xymatrix{
  0 \ar[r]
    & X^\bullet \ar[r]
    & Y^\bullet \ar[r]
    & Z^\bullet \ar[r]
    & 0
}\]
is called \emph{degreewise split} if the sequence
\[\xymatrix{
  0 \ar[r]
    & \bar X^\bullet \ar[r]
    & \bar Y^\bullet \ar[r]
    & \bar Z^\bullet \ar[r]
    & 0
}\]
is split; that is, if it is isomorphic to the sequence
\[\xymatrix{
  0 \ar[r]
    & \bar X^\bullet \ar[r]
    & \bar X^\bullet + \bar Z^\bullet \ar[r]
    & \bar Z^\bullet \ar[r]
    & 0 \text{.}
}\]

Let $0\to X^\bullet\to Y^\bullet\to Z^\bullet\to0$ be a degreewise split sequence. Choose a splitting, that is, for each $n$ an isomorphism
\[
  (Y^\bullet)^n \simeq (X^\bullet)^n + (Z^\bullet)^n \text{.}
\]
With respect to the direct-sum decomposition above, the matrix of the differential of $Y^\bullet$ has the form
\[
  \begin{pmatrix}
    d_{X^\bullet} & \varphi_n \\
    0 & d_{Z^\bullet}
  \end{pmatrix}
\]
where $\varphi_n$ is a morphism from $(Z^\bullet)^n$ to $(X^\bullet)^{n+1}$. The $\varphi_n$ determine a morphism of complexes
\[
  \varphi:Z^\bullet \to T(X^\bullet) \text{.}
\]
Changing the splitting does not change the homotopy class of $\varphi$. Moreover, if $\underline u$, $\underline v$, and $\underline\varphi$ denote the images in $\eK(\cA)$ of $u$, $v$, and $\varphi$, then the triangle $(X^\bullet,Y^\bullet,Z^\bullet,\underline u,\underline v, \underline\varphi)$ is distinguished. Let $\mathsf{SSS}(\cA)$ denote the category of degreewise split sequences in $\eC(\cA)$, and let $\mathsf{TrK}(\cA)$ denote the category of distinguished triangles in $\eK(\cA)$. The preceding construction defines a functor
\[
  \rho:\mathsf{SSS}(\cA) \to \mathsf{TrK}(\cA) \text{.}
\]
This functor is essentially surjective.

\subsection{Examples of cohomological and exact functors}\label{VIII:1-3}

\begin{definition}
Let $\cA$ be a triangulated category and $\cB$ an abelian category. An additive functor $F:\cA\to\cB$ is called a \emph{cohomological functor} if, for every distinguished triangle $(X,Y,Z,u,v,w)$, the sequence
\[\xymatrix{
  F(X) \ar[r]^-{F(u)}
    & F(Y) \ar[r]^-{F(v)}
    & F(Z)
}\]
is exact.
\end{definition}

The functor $F\circ T^i$ will often be denoted by $F^i$. By axiom \hyperlink{VIII:TR2}{TR2} for triangulated categories, there is an unbounded exact sequence
\[\xymatrix{
  \cdots \ar[r]
    & F^i(X) \ar[r]
    & F^i(Y) \ar[r]
    & F^i(Z) \ar[r]
    & F^{i+1}(X) \ar[r]
    & \cdots
}\]

\subsubsection{Examples}\label{VIII:1-3-2}

\begin{enumerate}
  \item Let $\cA$ be a triangulated category. For every object $X$ of $\cA$, the functor $\hom_\cA(X,-)$ is cohomological. The functor $\hom_\cA(-,X)$ is a cohomological functor $\cA^\circ\to\mathsf{Ab}$.
  \item Let $\cA$ be an abelian category and $\eC(\cA)$ its category of complexes. Define a functor $\h^0:\eC(\cA)\to\cA$ as follows. For an object $X^\bullet$ of $\eC(\cA)$, set
    \[
      \h^0(X^\bullet) = \ker d^0 / \im{d^{-1}} \text{.}
    \]
    The functor $\h^0$ annihilates null-homotopic morphisms and hence factors uniquely through $\eK(\cA)$. The resulting functor $\eK(\cA)\to\cA$ will again be denoted by $\h^0$.
\end{enumerate}

\subsubsection{An example of an exact bifunctor}\label{VIII:1-3-3}

Let $\cA$, $\cA'$, and $\cA''$ be three additive categories, and let
\[
  P:\cA\times \cA'\to \cA''
\]
be a bilinear functor (that is, additive in each argument separately). It induces a bilinear functor
\[
  P^\ast:\eC(\cA)\times \eC(\cA') \to \eC(\cA'')
\]
as follows.

For $X^\bullet\in\eC(\cA)$ and $Y^\bullet\in\eC(\cA')$, $P(X^\bullet,Y^\bullet)$ is a double complex in $\cA''$. Define $P^\ast(X^\bullet,Y^\bullet)$ to be its associated total complex.

If $f$ is a null-homotopic morphism in $\eC(\cA)$ (respectively $\eC(\cA')$) and $Z^\bullet$ is an object of $\eC(\cA')$ (respectively $\eC(\cA)$), then $P^\ast(f,Z^\bullet)$ (respectively $P^\ast(Z^\bullet,f)$) is null-homotopic. It follows that $P^\ast$ uniquely defines a functor
\[
  P^\bullet:\eK(\cA) \times \eK(\cA') \to \eK(\cA'') \text{.}
\]
$P^\bullet$ is an exact bifunctor.

In particular, for an additive category $\cA$, one may take for $P$ the functor
\begin{align*}
  \cA^\circ\times \cA &\to \mathsf{Ab} \\
  (X,Y) &\mapsto \hom(X,Y)
\end{align*}

The preceding construction then gives a functor
\[
  \hom^\bullet:\eK(\cA)^\circ \times \eK(\cA) \to \eK(\mathsf{Ab})
\]
which, when composed with $\h^0:\eK(\mathsf{Ab})\to\mathsf{Ab}$, plainly recovers the functor $\hom_{\eK(\cA)}$.

\section{Quotient categories}\label{VIII:2}

\subsection{Thick subcategories. Multiplicative systems of morphisms}\label{VIII:2-1}

We begin with two definitions.

\begin{definition}\label{VIII:2-1-1}
A subcategory $\cB$ of a triangulated category $\cA$ is \emph{thick} if it is a full triangulated subcategory of $\cA$ and also has the following property: whenever a morphism $f:X\to Y$ factors through an object of $\cB$ and belongs to a distinguished triangle $(X,Y,Z,f,g,h)$ with $Z$ in $\cB$, both the source and the target of $f$ belong to $\cB$.
\end{definition}

The set of thick subcategories of $\cA$ will be denoted by $\sN$.

The intersection of any family of thick subcategories is thick.

\begin{definition}
Let $\cA$ be a triangulated category. A set $S$ of morphisms of $\cA$ is called a \emph{multiplicative system of morphisms} if it has the following five properties.

\paragraph{FR1}
\phantomsection\hypertarget{VIII:FR1}{} If $f,g\in S$ are composable, then $f\circ g\in S$. For every object $X$ of $\cA$, the identity morphism of $X$ belongs to $S$.

\paragraph{FR2}
\phantomsection\hypertarget{VIII:FR2}{} In $\cA$, every diagram
\[\xymatrix{
  & Y \ar[d]^-{s\in S} \\
  Z \ar[r]^-f
    & X
}\]
can be completed to a commutative diagram
\[\xymatrix{
  P \ar[r]^-g \ar[d]^-{t\in S}
    & Y \ar[d]^-{s\in S} \\
  Z \ar[r]^-f
    & X
}\]
The symmetric property also holds.

\paragraph{FR3}
\phantomsection\hypertarget{VIII:FR3}{} For morphisms $f$ and $g$, the following conditions are equivalent:
\begin{enumerate}[\indent i)]
  \item there is an $s\in S$ such that $s\circ f=s\circ g$;
  \item there is a $t\in S$ such that $f\circ t=g\circ t$.
\end{enumerate}

\paragraph{FR4}
\phantomsection\hypertarget{VIII:FR4}{} Let $T$ be the translation functor of $\cA$. For every $s\in S$, one has $T(s)\in S$.

\paragraph{FR5}
\phantomsection\hypertarget{VIII:FR5}{} If $(X,Y,Z,u,v,w)$ and $(X',Y',Z',u',v',w')$ are two distinguished triangles and $f,g\in S$ are such that $(f,g)$ is a morphism from $u$ to $u'$, then there is a morphism $h\in S$ such that $(f,g,h)$ is a morphism of triangles.
\end{definition}

A multiplicative system of morphisms is \emph{saturated} if it has the following property:

a morphism $f$ belongs to $S$ if and only if there are morphisms $g$ and $g'$ such that $g\circ f\in S$ and $f\circ g'\in S$.

The intersection of any family of saturated multiplicative systems is a saturated multiplicative system.

The set of saturated multiplicative systems will be denoted by $\sS$.

For a thick subcategory $\cB$, let $\varphi(\cB)$ be the set of morphisms $f$ belonging to a distinguished triangle $(X,Y,Z,f,g,h)$ in which $Z$ is an object of $\cB$. Then $\varphi(\cB)$ is a saturated multiplicative system.

For a saturated multiplicative system $S$, let $\psi(S)$ be the full subcategory on the objects $Z$ belonging to a distinguished triangle $(X,Y,Z,f,g,h)$ in which $f\in S$. The subcategory $\psi(S)$ is thick.

The main result of this subsection is that $\varphi$ is an order-preserving isomorphism from $\sN$ onto $\sS$, where the order is defined by inclusion. Its inverse is $\psi$.

\subsection{Universal problems}\label{VIII:2-2}

Let $\cA$ and $\cA'$ be two triangulated categories and $F:\cA\to\cA'$ an exact functor. Let $S(F)$ be the set of morphisms of $\cA$ that $F$ takes to isomorphisms, and let $\cB(F)$ be the full subcategory on the objects of $\cA$ that $F$ takes to zero objects of $\cA'$. Then $S(F)$ is a saturated multiplicative system, $\cB(F)$ is a thick subcategory, and $S(F)=\varphi(\cB(F))$. Now let $\cA$ be a triangulated category, $\cB$ a thick subcategory, and $S=\varphi(\cB)$ the corresponding saturated multiplicative system. The following two universal problems are equivalent.

\paragraph{Problem 1}
\phantomsection\hypertarget{VIII:prob1}{} Find a triangulated category $\cA/\cB$ and an exact functor $Q:\cA\to\cA/\cB$ such that every exact functor from $\cA$ to a triangulated category $\cA'$ that takes the objects of $\cB$ to zero objects of $\cA'$ admits a unique factorization of the form $G\circ Q$, where $G$ is exact.

\paragraph{Problem 2}
\phantomsection\hypertarget{VIII:prob2}{} Find a triangulated category $\cA_S$ and an exact functor $Q:\cA\to\cA_S$ such that every exact functor from $\cA$ to a triangulated category $\cA'$ that takes the morphisms of $S$ to isomorphisms of $\cA'$ admits a unique factorization of the form $G\circ Q$, where $G$ is exact.

(The editor asks the reader's indulgence for this antiquated and reactionary language.) In the next subsection we shall show that \hyperlink{VIII:prob2}{Problem 2} has a solution. The corresponding \hyperlink{VIII:prob1}{Problem 1} therefore has the same solution.

\subsection{Calculus of fractions}\label{VIII:2-3}

\subsubsection{}\label{VIII:2-3-1}

Let $\cA$ be a triangulated category and $S$ a saturated multiplicative system. Let $\cA[S^{-1}]$ denote the category of fractions of $\cA$ with respect to $S$, and let $Q:\cA\to\cA[S^{-1}]$ be the canonical functor. The pair $(\cA[S^{-1}],Q)$ solves the following problem: every functor from $\cA$ to an arbitrary category (not necessarily additive) that takes the morphisms of $S$ to isomorphisms factors uniquely through $(\cA[S^{-1}],Q)$. Such a pair always exists, with no hypothesis on the set $S$ of morphisms \cite[1.1]{gz67}\footnote{The original text references [C.G.G] here. Most likely this is the lost work \emph{Catégories et foncteurs} by Chevalley, Gabriel, and Grothendieck, referenced in \cite{sc72}.}. Because $S$ is a saturated multiplicative system, however, it satisfies \hyperlink{VIII:FR1}{FR1}, \hyperlink{VIII:FR2}{FR2}, and \hyperlink{VIII:FR3}{FR3}. Consequently, by \cite{gz67}, the category $\cA[S^{-1}]$ can be obtained by a right or left calculus of fractions. We recall the right calculus of fractions. (The left calculus is obtained by reversing the arrows.)

\subsubsection{}\label{VIII:2-3-2}

The objects of $\cA[S^{-1}]$ are the objects of $\cA$.
\begin{itemize}
  \item For each object $X$ of $\cA$, let $S_X$ denote the category of arrows in $S$ with target $X$. To every object $Y$ of $\cA$ we associate, functorially in $Y$, the following set-valued functor on $S_X^\circ$ (where $S_X^\circ$ denotes the opposite category):
    \[
      H_Y:s\mapsto \hom(\operatorname{source}(s),Y) \text{.}
    \]
    Set
    \[
      \hom_{\cA[S^{-1}]}(X,Y) = \varinjlim_{S_X^\circ} H_Y \text{.}
    \]
    For all facts about direct limits, we refer to \cite{ar62}.

    Using \hyperlink{VIII:FR1}{FR1}, \hyperlink{VIII:FR2}{FR2}, and \hyperlink{VIII:FR3}{FR3}, one checks with pleasure that $S_X^\circ$ has properties L1, L2, and L3. Thus these direct limits have all the good properties of direct limits over filtered ordered sets.
  \item Let $X,Y,Z$ be three objects of $\cA$, let $a\in\hom_{\cA[S^{-1}]}(X,Y)$, and let $b\in\hom_{\cA[S^{-1}]}(Y,Z)$. Choose an object $s$ of $S_X$ and an element $\underline a\in\hom(\operatorname{source}(s),Y)$ whose image is $a$, and choose an object $t$ of $S_Y$ and an element $\underline b\in\hom(\operatorname{source}(t),Z)$ whose image is $b$. We then have a diagram
    \begin{equation*}\tag{1}\label{VIII:eq:2-3-2-1}
    \xymatrix{
      & \bullet \ar[dl]_-s \ar[dr]^-{\underline a}
        &
        & \bullet \ar[dl]_-t \ar[dr]^-{\underline b} \\
      X
        &
        & Y
        &
        & Z
    }
    \end{equation*}
    which, by \hyperlink{VIII:FR2}{FR2}, can be completed to a diagram
    \[\xymatrix{
      & & \bullet \ar[dl]_-{t'} \ar[dr]^-{\underline c} \\
      & \bullet \ar[dl]_-s \ar[dr]^-{\underline a}
        & & \bullet \ar[dl]_-t \ar[dr]^-{\underline b} \\
      X
        & & Y
        & & Z
    }\]
    Let $b\circ a$ be the image in $\hom_{\cA[S^{-1}]}(X,Z)$ of the morphism $\underline{b\circ c}$. Properties \hyperlink{VIII:FR1}{FR1}, \hyperlink{VIII:FR2}{FR2}, and \hyperlink{VIII:FR3}{FR3} show that $b\circ a$ is independent of the chosen representatives $\underline a$ and $\underline b$ and of the way in which \eqref{VIII:eq:2-3-2-1} is completed. They also show that this defines a category $\cA[S^{-1}]$.
\end{itemize}

Let $Q:\cA\to\cA[S^{-1}]$ denote the evident functor.

Since $\cA$ is additive, $\cA[S^{-1}]$ is additive. Using \hyperlink{VIII:FR4}{FR4}, one also proves that there is a unique translation functor on $\cA[S^{-1}]$, denoted by $T$, satisfying
\[
  Q\circ T = T\circ Q \text{.}
\]
Finally, using \hyperlink{VIII:FR5}{FR5}, one proves that $\cA[S^{-1}]$ has a unique triangulated structure for which $Q$ is exact. Its distinguished triangles are precisely the triangles isomorphic to the images under $Q$ of distinguished triangles of $\cA$. If $\cA_S$ denotes $\cA[S^{-1}]$ with this triangulated structure, then it is straightforward to prove that $(\cA_S,Q)$ is a solution to \hyperlink{VIII:prob2}{Problem 2}.

\begin{definition}\label{VIII:2-3-3}
Let $\cA$ be a triangulated category, $\cB$ a thick subcategory of $\cA$, and $(Q,\cA/\cB)$ the solution to \hyperlink{VIII:prob1}{Problem 1} (\S \ref{VIII:2-2}). The category $\cA/\cB$ is called the \emph{quotient category} of $\cA$ by $\cB$, and $Q$ is called the canonical quotient functor. More generally, let $\cN$ be a triangulated subcategory of $\cA$, and let $\underline\cN$ be the smallest thick subcategory containing $\cN$. The quotient category of $\cA$ by $\cN$ is, by definition, $\cA/\cN=\cA/\underline\cN$.
\end{definition}

\subsection{Properties of quotient categories}\label{VIII:2-4}

\subsubsection{}\label{VIII:2-4-1}

Let $\cA$ be a triangulated category and $H$ a cohomological functor with values in an abelian category $\cG$. Let $\cB_H$ be the full subcategory on the objects all of whose translates are taken by $H$ to zero objects of $\cG$, and let $S_H$ be the set of morphisms all of whose translates are taken by $H$ to isomorphisms in $\cG$. Then $\cB_H$ is a thick subcategory of $\cA$, $S_H$ is a saturated multiplicative system, and $S_H=\varphi(\cB_H)$ (\S \ref{VIII:2-1}). The functor $H$ factors uniquely through $(Q,\cA/\cB_H)$.

\begin{theorem}\label{VIII:2-4-2}
Let $\cA$ be a triangulated category, $\cB$ a triangulated subcategory, $\cN$ a thick subcategory of $\cA$, and $S=\varphi(\cN)$ the corresponding saturated multiplicative system.
\begin{enumerate}[(a)]
  \item The category $\cN\cap\cB$ is a thick subcategory of $\cB$. Its corresponding multiplicative system is $S\cap\cB$.
  \item The following two properties are equivalent:
    \begin{enumerate}[(i)]
      \item For every object $X$ of $\cB$ and every morphism $s:R\to X$ with $s\in S$, there is a morphism $s':R'\to R$ such that $s\circ s'\in S\cap\cB$ and $R'\in\operatorname{Ob}\cB$.
      \item Every morphism from an object $X$ of $\cB$ to an object $Y$ of $\cN$ factors through an object of $\cN\cap\cB$.
    \end{enumerate}
  \item The properties (i)' and (ii)' obtained from (i) and (ii) by reversing the arrows are equivalent.
  \item If (i) and (ii) hold, or if (i)' and (ii)' hold, the canonical functor
    \[\xymatrix{
      \cB/\cN\cap \cB \ar[r] & \cA/\cN
    }\]
    is faithful.

    Since this functor is injective on objects, it identifies $\cB/(\cN\cap\cB)$ with a subcategory of $\cA/\cN$.
  \item If, in addition, $\cB$ is a full subcategory of $\cA$, then $\cB/(\cN\cap\cB)$ is a full subcategory of $\cA/\cN$.
\end{enumerate}
\end{theorem}

\begin{corollary}\label{VIII:2-4-3}
Let $\cA\subset\cB\subset\cC$ be three triangulated categories, with $\cA$ a thick subcategory of $\cB$ and $\cB$ a thick subcategory of $\cC$. Then $\cA$ is a thick subcategory of $\cC$, and $\cB/\cA$ is a thick subcategory of $\cC/\cA$. The canonical functor $\cC/\cB\to(\cC/\cA)/(\cB/\cA)$ is an isomorphism.
\end{corollary}

\subsection{Properties of the quotient functor}\label{VIII:2-5}

Let $\cA$ be a triangulated category, $\cB$ a thick subcategory, $S$ the corresponding multiplicative system, and $Q:\cA\to\cA/\cB$ the canonical quotient functor.

\begin{proposition}\label{VIII:2-5-1}
For objects $X$ and $Y$ of $\cA$, the following assertions are equivalent:
\begin{enumerate}[(a)]
  \item Every morphism $f:X\to Y$ for which there is a morphism $s:Z\to X$ in $S$ satisfying $f\circ s=0$ is zero.
  \item Every morphism $f:X\to Y$ for which there is a morphism $s:Y\to Z$ in $S$ satisfying $s\circ f=0$ is zero.
  \item Every morphism $f:X\to Y$ that factors through an object of $\cB$ is zero.
  \item The canonical map
    \[
      \hom_\cA(X,Y) \to \hom_{\cA/\cB}(Q(X),Q(Y))
    \]
    is injective.
\end{enumerate}
\end{proposition}

\begin{proposition}\label{VIII:2-5-2}
Let $X$ and $Y$ be two objects of $\cA$. The following assertions are equivalent:
\begin{enumerate}[(a)]
  \item Every diagram ($s\in S$)
    \[\xymatrix{
      & Z \ar[dl]_-s \ar[dr]^-f \\
      X
        & & Y
    }\]
    can be completed to a diagram ($s,t\in S$)
    \[\xymatrix{
      & Z' \ar[d]_-t \\
      & Z \ar[dl]_-s \ar[dr]^-f \\
      X
        & & Y
    }\]
    in which $f\circ t$ factors through $(X,s\circ t)$.
  \item The assertion obtained by reversing the arrows in (a) and interchanging $X$ and $Y$.
  \item Every diagram
    \[\xymatrix{
      X \ar[r]^-f
        & N \ar[r]^-g
        & Y
    }\]
    in which $g\circ f=0$ and $N$ is an object of $\cB$ can be completed to a diagram
    \[\xymatrix{
      & N' \ar[d]^-i \\
      X \ar[ur]^-h \ar[r]^-f
        & N \ar[r]^-g
        & Y
    }\]
    in which $f=i\circ h$ and $g\circ i=0$.
  \item The assertion obtained by reversing the arrows in (c) and interchanging $X$ and $Y$.
  \item The canonical map
    \[\xymatrix{
      \hom_\cA(X,Y) \ar[r]^-Q
        & \hom_{\cA/\cB}(Q(X),Q(Y))
    }\]
    is surjective.
\end{enumerate}
\end{proposition}

\begin{proposition}\label{VIII:2-5-3}
Let $X$ be an object of $\cA$. The following assertions are equivalent:
\begin{enumerate}[(a)]
  \item For every object $Y$ of $\cA$, the canonical map
    \[\xymatrix{
      \hom_\cA(X,Y) \ar[r]^-Q
        & \hom_{\cA/\cB}(Q(X),Q(Y))
    }\]
    is an isomorphism.
  \item Every morphism from $X$ to an object of $\cB$ is zero.
\end{enumerate}
Likewise, the following two assertions are equivalent:
\begin{enumerate}[(a)']
  \item For every object $Y$ of $\cA$, the canonical map
    \[\xymatrix{
      \hom_\cA(Y,X) \ar[r]^-Q
        & \hom_{\cA/\cB}(Q(Y),Q(X))
    }\]
    is an isomorphism.
  \item Every morphism from an object of $\cB$ to $X$ is zero.
\end{enumerate}
\end{proposition}

\begin{definition}\label{VIII:2-5-4}
An object satisfying the equivalent conditions (a) and (b) of the preceding proposition is called an \emph{object left-free over $\cA/\cB$}, or a \emph{left $Q$-free object}. Similarly, an object satisfying the equivalent conditions (a)' and (b)' is called an \emph{object right-free over $\cA/\cB$}, or a \emph{right $Q$-free object}. It is determined up to isomorphism by $Q(X)\in\operatorname{Ob}(\cA/\cB)$.
\end{definition}

\subsection{Functors adjoint to the quotient functor}\label{VIII:2-6}

\begin{definition}\label{VIII:2-6-1}
Two triangulated subcategories $\cN$ and $\cN'$ of a triangulated category $\cA$ are called \emph{orthogonal} if, for every object $X$ of $\cN$ and every object $Y$ of $\cN'$, one has
\[
  \hom_\cA(X,Y) = 0
\]
In this case $\cN'$ is said to be right orthogonal to $\cN$, and $\cN$ left orthogonal to $\cN'$.
\end{definition}

\begin{proposition}\label{VIII:2-6-2}
Let $\cN$ be a triangulated subcategory of a triangulated category $\cA$. The full category $\cN^\bot$ (respectively $^\bot\cN$) on the objects $X$ of $\cA$ such that $\hom_\cA(Y,X)=0$ (respectively $\hom_\cA(X,Y)=0$) for every object $Y$ of $\cN$ is a thick subcategory of $\cA$. By abuse of language, $\cN^\bot$ is called the right orthogonal of $\cN$.
\end{proposition}

This proposition lets us reformulate Proposition \ref{VIII:2-5-3}.

\begin{proposition}\label{VIII:2-6-3}
Let $\cB\subset\cA$ be a thick subcategory of a triangulated category $\cA$. The full category on the objects right-free over $\cA/\cB$ (respectively left-free) is precisely the right orthogonal (respectively left orthogonal) of $\cB$.
\end{proposition}

\begin{proposition}\label{VIII:2-6-4}
Let $\cA$ be a triangulated category, $\cB$ a thick subcategory of $\cA$, $\cB^\bot$ the right orthogonal of $\cB$, $S(\cB)$ and $S(\cB^\bot)$ the corresponding multiplicative systems, and $Q$ and $Q^\bot$ the canonical quotient functors. For an object $X$ of $\cA$, consider the following properties:
\begin{enumerate}[(i)]
  \item There is a morphism in $S(\cB)$ from $X$ to an object of $\cB^\bot$.
  \item There is a morphism in $S(\cB^\bot)$ from an object of $\cB$ to $X$.
  \item The category $S_X(\cB)$ of arrows in $S(\cB)$ with source $X$ has a final object.
  \item The category $S^X(\cB^\bot)$ of arrows in $S(\cB^\bot)$ with target $X$ has an initial object.
  \item The category $\cB/X$ of objects of $\cB$ over $X$ has a final object.
  \item The category $\cB^\bot/X$ of objects of $\cB^\bot$ under $X$ has an initial object.
\end{enumerate}
Then
\[\xymatrix@=0.5cm{
  & & & & (iv) \\
  (i) \ar@{<=>}[r]
    & (ii) \ar@{<=>}[r]
    & (iii) \ar@{<=>}[r]
    & (v) \ar@{=>}[ur] \ar@{=>}[dr] \\
  & & & & (vi)
}\]

If, in addition, $^\bot(\cB^\bot)=\cB$, then all these properties are equivalent.
\end{proposition}

Let $\cB$ be a thick subcategory of a triangulated category $\cA$, let $i:\cB\to\cA$ be the inclusion functor, let $Q:\cA\to\cA/\cB$ be the quotient functor, let $\cB^\bot$ be the right orthogonal of $\cB$, and let $i^\bot$ and $Q^\bot$ be the corresponding functors. Proposition \ref{VIII:2-6-4} immediately gives the following propositions.

\begin{proposition}\label{VIII:2-6-5}
The following two properties are equivalent:
\begin{enumerate}[(i)]
  \item The functor $i$ has a right adjoint.
  \item The functor $Q$ has a right adjoint.
\end{enumerate}
\end{proposition}

The proposition remains true when ``right'' is replaced by ``left.''

\begin{proposition}\label{VIII:2-6-6}
The following two properties are equivalent:
\begin{enumerate}[(i)]
  \item The functor $i$ has a right adjoint.
  \item The functor $i^\bot$ has a left adjoint, and the left orthogonal of $\cB^\bot$ is equal to $\cB$.
\end{enumerate}
\end{proposition}

\begin{proposition}\label{VIII:2-6-7}
Suppose that the properties of Propositions \ref{VIII:2-6-5} and \ref{VIII:2-6-6} hold. Let $i^\ast$ and $Q^\ast$ be the right adjoints of $i$ and $Q$, and let $^\ast i^\bot$ and $^\ast Q^\bot$ be the left adjoints of $i^\bot$ and $Q^\bot$. All these functors are exact. Moreover:

there is a functorial isomorphism $Q^\ast\circ Q\iso i^\bot\circ{}^\ast i^\bot$ such that the following diagram commutes:
\[\xymatrix{
  & \operatorname{id} \ar[dl] \ar[dr] \\
  Q^\ast \circ Q \ar[rr]^\sim
    & & i^\bot\circ ^\ast i^\bot
}\]
(the oblique arrows are defined by the adjunction properties).

Likewise, there is an isomorphism $^\ast Q^\bot\circ Q^\bot\iso i\circ i^\ast$ such that the following diagram commutes:
\[\xymatrix{
  ^\ast Q^\bot\circ Q^\bot \ar[rr]^-\sim \ar[dr]
    & & i\circ i^\ast \ar[dl] \\
  & \operatorname{id}{}
}\]

Finally, there is a functorial morphism of degree $1$, $\partial:i^\bot\circ{}^\ast i^\bot\to i\circ i^\ast$, such that the following triangle is distinguished ($\deg\partial=1$):
\[\xymatrix{
  i^\bot \circ ^\ast i^\bot \ar[rr]^-\partial
    & & i\circ i^\ast \ar[dl] \\
  & \operatorname{id}{} \ar[ul]
}\]

Such a morphism $\partial$ is unique and satisfies $\partial(TX)=-T\partial(X)$ ($X$ an object of $\cA$, $T$ the translation).
\end{proposition}

\section{Derived categories of an abelian category}\label{VIII:3}

\subsection{The derived categories}\label{VIII:3-1}

We use the notation of \ref{VIII:1-2-3}, together with the following.

\subsubsection{Notation}\label{VIII:3-1-1}

Let $\cA$ be an abelian category. Set
\begin{align*}
  \eD(\cA) &= \eK^{\infty,\infty}(\cA)/\eK^{\infty,\varnothing}(\cA) \\
  \eD^b(\cA) &= \eK^{b,b}(\cA)/\eK^{b,\varnothing}(\cA) \\
  \eD^+(\cA) &= \eK^{+,+}(\cA) / \eK^{+,\varnothing}(\cA) \\
  \eD^-(\cA) &= \eK^{-,-}(\cA) / \eK^{-,\varnothing}(\cA) \text{.}
\end{align*}

\begin{proposition}\label{VIII:3-1-2}
The natural functors in the following diagram
\[\xymatrix@=0.5cm{
  & \eD(\cA) \\
  \eD^+(\cA) \ar[ur]
    & & \eD^-(\cA) \ar[ul] \\
  & \eD^b(\cA) \ar[ul] \ar[ur]
}\]
are fully faithful.

The natural functors
\begin{align*}
  \eD^b(\cA) &= \eK^{b,b}(\cA)/\eK^{b,\varnothing}(\cA) \to \eK^{+,b}(\cA)/\eK^{+,\varnothing}(\cA) \\
  \eD^b(\cA) &= \eK^{b,b}(\cA)/\eK^{b,\varnothing}(\cA) \to \eK^{-,b}(\cA) / \eK^{-,\varnothing}(\cA)
\end{align*}
are equivalences of categories.

Since the functors in the diagram
\[\xymatrix@=0.5cm{
  & \eD(\cA) \\
  \eD^+(\cA) \ar[ur]
    & & \eD^-(\cA) \ar[ul] \\
  & \eD^b(\cA) \ar[ul] \ar[ur]
}\]
are injective on objects, they realize $\eD^b(\cA)$, $\eD^+(\cA)$, and $\eD^-(\cA)$ as full subcategories of $\eD(\cA)$.
\end{proposition}

The proof is contained in Theorem \ref{VIII:2-4-2}.

\begin{definition}\label{VIII:3-1-3}
The category $\eD(\cA)$ is called the \emph{derived category} of $\cA$. Objects of $\eD(\cA)$ isomorphic to objects of $\eD^b(\cA)$ are called \emph{bounded}; those isomorphic to objects of $\eD^+(\cA)$ are called \emph{bounded below}; and those isomorphic to objects of $\eD^-(\cA)$ are called \emph{bounded above}.
\end{definition}

Let $D$ denote the canonical functor
\[
  D : \eC(\cA) \to \eD(\cA) \text{.}
\]
The ``restriction'' of $D$ to the full subcategory $\cA$ of $\eC(\cA)$ will again be denoted by $D$. The restricted functor $D|_\cA$ is fully faithful, and $D$ factors uniquely through $\eK(\cA)$. Let $X$ and $Y$ be two objects of $\cA$.

It is easy to check that, for every $m>0$, the groups
\[
  \hom_{\eD(\cA)}(D(X),T^{-m}D(Y))
\]
vanish (where $T$ is the translation functor). The cohomological dimension of $\cA$ is the least integer $n$ such that, for every $m>n$,
\[
  \hom_{\eD(\cA)}(D(X),T^m D(Y)) = 0
  \qquad\text{for every pair of objects $X,Y$ of $\cA$.}
\]
If no such integer exists, the cohomological dimension of $\cA$ is said to be infinite.

Let $\eI^+(\cA)$ (respectively $\eI(\cA)$) be the full triangulated subcategory of $\eK^+(\cA)$ (respectively $\eK(\cA)$) consisting of complexes whose objects are injective in every degree.

\begin{proposition}\label{VIII:3-1-4}
Suppose that $\cA$ has enough injectives. The natural functor
\[
  Q^+:\eK^+(\cA) \to \eD^+(\cA)
\]
induces an equivalence of categories
\[
  \eI^+(\cA) \to \eD^+(\cA)  \text{.}
\]
The quasi-inverse is a right adjoint of $Q^+$.
\end{proposition}

If, in addition, $\cA$ has finite cohomological dimension, the preceding assertion remains true after deleting the superscripts $+$.

There is an analogous statement for projectives.

The proof of Proposition \ref{VIII:3-1-4} rests on \ref{VIII:2-6-4}.

The next proposition describes the triangulated structure of $\eD(\cA)$ more precisely.

Let $\eE(\cA)$ be the category of short exact sequences in $\eC(\cA)$. The category $\mathsf{SSS}(\cA)$ of degreewise split exact sequences (\S\ref{VIII:1-2-4}) is a full subcategory of $\eE(\cA)$. In the cited subsection we defined a functor
\[
  \rho:\mathsf{SSS}(\cA) \to \mathsf{TrK}(\cA)
\]
where $\mathsf{TrK}(\cA)$ denotes the category of distinguished triangles of $\eK(\cA)$. It induces a functor
\[
  \sigma : \mathsf{SSS}(\cA) \to \mathsf{TrD}(\cA)
\]
where $\mathsf{TrD}(\cA)$ denotes the category of distinguished triangles of $\eD(\cA)$.

\begin{proposition}\label{VIII:3-1-5}
There is a unique functor
\[
  \Pi:\eE(\cA) \to \mathsf{TrD}(\cA)
\]
of the following form ($\deg\delta(S)=1$):
\[
(S=0 \to X^\bullet \to Y^\bullet \to Z^\bullet \to 0)
\mapsto
\xymatrix{
  & D(Z^\bullet) \ar[dl]_-{\delta(S)} \\
  D(X^\bullet) \ar[rr]^-{D(u)}
    & & D(Y^\bullet) \ar[ul]_-{D(v)}
}
\]
whose restriction to $\mathsf{SSS}(\cA)$ is $\sigma$. \emph{This functor is essentially surjective.}
\end{proposition}

\begin{proposition}\label{VIII:3-1-6}
The canonical cohomological functor
\[
  \h^0:\eK(\cA) \to \cA
\]
factors uniquely through a functor, again denoted by $\h^0$,
\[
  \h^0:\eD(\cA) \to \cA \text{.}
\]
On $\eD(\cA)$, the functor $\h^0$ and its translates $\h^i$ form a conservative family; that is, a morphism $f:X^\bullet\to Y^\bullet$ is an isomorphism if and only if, for every integer $i$,
\[
  \h^i(f):\h^i(X) \to \h^i(Y)
\]
is an isomorphism.
\end{proposition}

\begin{definition}\label{VIII:3-1-7}
A morphism $f:X^\bullet\to Y^\bullet$ of $\eC(\cA)$ (respectively of $\eK(\cA)$) is called a \emph{quasi-isomorphism} if, for every integer $i$,
\[
  \h^i(f):\h^i(X) \to \h^i(Y)
\]
is an isomorphism.
\end{definition}

By the preceding proposition, the quasi-isomorphisms are precisely the morphisms that become isomorphisms in $\eD(\cA)$.

\begin{definition}\label{VIII:3-1-8}
Let $\sigma$ be a collection of objects of $\cA$, and let $X^\bullet$ be an object of $\eC(\cA)$ (respectively $\eK(\cA)$). A right (respectively left) resolution of $X^\bullet$ of type $\sigma$ is a quasi-isomorphism $X^\bullet\to V^\bullet$ (respectively $V^\bullet\to X^\bullet$), where every object of the complex $V^\bullet$ belongs to $\sigma$. We say \emph{injective resolution} (respectively projective resolution) for a right (respectively left) resolution of injective (respectively projective) type.
\end{definition}

In general, the words ``right'' and ``left'' will be omitted when no ambiguity can result.

\begin{proposition}\label{VIII:3-1-9}
\begin{enumerate}
  \item Suppose that every object of $\cA$ embeds into an object of $\sigma$ (respectively is a quotient of an object of $\sigma$). Then every object of $\eC^+(\cA)$ (respectively $\eC^-(\cA)$) has a right (respectively left) resolution of type $\sigma$ in $\eC^+(\cA)$ (respectively $\eC^-(\cA)$).
  \item Suppose in addition that there is an integer $n$ such that, for every exact sequence
    \begin{align*}
      Y^0 &\to Y^1 \to \cdots \to Y^{n-1} \to Y^n \to 0 \\
      \text{(resp. } 0 &\to Y^n \to Y^{n-1} \to \cdots \to Y^1 \to Y^0 \text{)}
    \end{align*}
    of objects of $\cA$ in which $Y^i$ belongs to $\sigma$ for every $0\leqslant i\leqslant n-1$, the object $Y^n$ also belongs to $\sigma$. Then every object of $\eC(\cA)$ has a right (respectively left) resolution of type $\sigma$.
\end{enumerate}
\end{proposition}

\subsection{Study of $\ext$}\label{VIII:3-2}

Let $X^\bullet$ and $Y^\bullet$ be two objects of $\eC(\cA)$ (respectively $\eK(\cA)$), and let $D:\eC(\cA)\to\eD(\cA)$ (respectively $Q:\eK(\cA)\to\eD(\cA)$) be the canonical functor.

\begin{definition}\label{VIII:3-2-1}
The \emph{$i$-th hyper-Ext}, denoted by
\[
  (X^\bullet,Y^\bullet)\mapsto \ext^i(X^\bullet,Y^\bullet)
\]
is the functor $\hom_{\eD(\cA)}(D(X^\bullet),T^iD(Y^\bullet))$ (respectively $\hom_{\eD(\cA)}(Q(X^\bullet),T^iQ(Y^\bullet))$).
\end{definition}

\begin{proposition}\label{VIII:3-2-2}
Let $0\to X^\bullet\to Y^\bullet\to Z^\bullet\to0$ be an exact sequence in $\eC(\cA)$, and let $V^\bullet$ be an object of $\eC(\cA)$. Then there are unbounded exact sequences
\[\xymatrix@=0.5cm{
  \cdots \ar[r]
    & \ext^i(V^\bullet,X^\bullet) \ar[r]
    & \ext^i(V^\bullet,Y^\bullet) \ar[r]
    & \ext^i(V^\bullet,Z^\bullet) \ar[r]^-\delta
    & \ext^{i+1}(V^\bullet,X^\bullet) \ar[r]
    & \cdots \\
  \cdots \ar[r]
    & \ext^i(Z^\bullet,V^\bullet) \ar[r]
    & \ext^i(Y^\bullet,V^\bullet) \ar[r]
    & \ext^i(X^\bullet,V^\bullet) \ar[r]^-\delta
    & \ext^{i+1}(Z^\bullet,V^\bullet) \ar[r]
    & \cdots
}\]
\end{proposition}

Let $X^\bullet$ and $Y^\bullet$ be two objects of $\eK(\cA)$. Let $\mathsf{Qis}/X^\bullet$ (respectively $\mathsf{Qis}^+/X^\bullet$, $\mathsf{Qis}^-/X^\bullet$, or $\mathsf{Qis}^b/X^\bullet$) denote the category of quasi-isomorphisms with target $X^\bullet$ and source in $\eK(\cA)$ (respectively $\eK^+(\cA)$, $\eK^-(\cA)$, or $\eK^b(\cA)$). Similarly define the categories $Y^\bullet/\mathsf{Qis}$, $Y^\bullet/\mathsf{Qis}^+$, and so on (quasi-isomorphisms with source $Y^\bullet$).

The following proposition summarizes the various equivalent descriptions of $\ext^i$.

\begin{proposition}\label{VIII:3-2-3}
There are isomorphisms of functors
\[\xymatrix{
  \ext^0(X^\bullet,Y^\bullet) = \hom_{\eD(\cA)}(Q(X^\bullet),Q(Y^\bullet)) \ar[r]^-\sim
    & \varinjlim_{(\mathsf{Qis}/X^\bullet)^\circ} \hom_{\eK(\cA)}(-,Y^\bullet) \ar[r]^-\sim
    & \\
  \varinjlim_{Y^\bullet/\mathsf{Qis}} \hom_{\eK(\cA)}(X^\bullet,-) \ar[r]^-\sim
    & \varinjlim_{(\mathsf{Qis}/X^\bullet)^\circ \times Y^\bullet/\mathsf{Qis}} \hom_{\eK(\cA)}(-,-) \text{.}
}\]
If $X^\bullet$ is an object of $\eK^-(\cA)$, then
\[
  \ext^0(X^\bullet,Y^\bullet) \iso \varinjlim_{(\mathsf{Qis}^-/X^\bullet)^\circ} \hom_{\eK(\cA)}(-,Y^\bullet)
\]
If $Y^\bullet$ is an object of $\eK^+(\cA)$, then
\[
  \ext^0(X^\bullet,Y^\bullet) \iso
  \varinjlim_{Y^\bullet/\mathsf{Qis}^+}
  \hom_{\eK(\cA)}(X^\bullet,-)
\]
If $X^\bullet$ is an object of $\eK^-(\cA)$ and $Y^\bullet$ an object of $\eK^b(\cA)$, then
\[
  \ext^0(X^\bullet,Y^\bullet) \iso \varinjlim_{Y^\bullet/\mathsf{Qis}^b} \hom_{\eK(\cA)}(X^\bullet,-)
\]
If $X^\bullet$ is an object of $\eK^b(\cA)$ and $Y^\bullet$ an object of $\eK^+(\cA)$, then
\[
  \ext^0(X^\bullet,Y^\bullet) \iso \varinjlim_{(\mathsf{Qis}^b/X^\bullet)^\circ} \hom_{\eK(\cA)}(-,Y^\bullet)
\]
If $\cA$ has enough injectives and $Y^\bullet$ is an object of $\eK^+(\cA)$, then $Y^\bullet$ has an injective resolution, unique up to isomorphism in $\eK^+(\cA)$,
\[
  Y^\bullet \to I(Y^\bullet)
\]
and
\[
  \ext^0(X^\bullet,Y^\bullet) \iso \hom_{\eK(\cA)}(X^\bullet,I(Y^\bullet)) \text{.}
\]
There is an analogous statement for projectives. If $\cA$ has enough injectives and has \emph{finite cohomological dimension}, then every object $Y^\bullet$ of $\eK(\cA)$ has a unique injective resolution
\[
  Y^\bullet \to I(Y^\bullet)
\]
and
\[
  \ext^0(X^\bullet,Y^\bullet) \iso \hom_{\eK(\cA)}(X^\bullet,I(Y^\bullet))
\]
There is an analogous statement for projectives.
\end{proposition}

\paragraph{Remark.}
Writing out the isomorphism in Proposition \ref{VIII:3-2-3}(3) recovers Yoneda's definition.

\section{Derived functors}\label{VIII:4}

\subsection{Definition of derived functors}\label{VIII:4-1}

\begin{definition}\label{VIII:4-1-1}
Let $\cC$ and $\cC'$ be two graded categories (write $T$ for the translation functor of both $\cC$ and $\cC'$), and let $F$ and $G$ be two graded functors from $\cC$ to $\cC'$. A \emph{morphism of graded functors} is a morphism of functors
\[
  u:F\to G
\]
with the following property:

for every object $X$ of $\cC$, the following diagram commutes:
\[\xymatrix{
  F(T X) \ar[r]^-{u(T X)}
    & G(T X) \\
  T F(X) \ar[r]^-{T u(X)} \ar[u]_-\wr
    & T G(X) \ar[u]_-\wr
}\]
\end{definition}

Let $\cC$ and $\cC'$ be two triangulated categories. Denote by $\mathsf{Fex}(\cC,\cC')$ the category of exact functors from $\cC$ to $\cC'$, with morphisms of graded functors as its morphisms.

Let $\cA$ and $\cB$ be two abelian categories, and let $\Phi:\eK^\ast(\cA)\to\eK^{\ast'}(\cB)$ be an exact functor (where $\ast$ and $\ast'$ each denote one of $+$, $-$, $b$, or the empty symbol). Composition with the canonical functor $Q:\eK^\ast(\cA)\to\eD^\ast(\cA)$ gives a functor
\[\xymatrix{
  \mathsf{Fex}(\eD^\ast(\cA),\eD^{\ast'}(\cB)) \ar[r]
    & \mathsf{Fex}(\eK^\ast(\cA),\eD^{\ast'}(\cB))
}\]
Thus, if $Q'$ also denotes the canonical functor $\eK^{\ast'}(\cB)\to\eD^{\ast'}(\cB)$, we obtain a functor $\%$ (respectively $\%'$) from $\mathsf{Fex}(\eD^\ast(\cA),\eD^{\ast'}(\cB))$ to $\mathsf{Ab}$:
\begin{align*}
  \Psi  &\mapsto \hom(Q'\circ \Phi,\Psi\circ Q) \\
  (\text{resp. }\Psi &\mapsto \hom(\Psi\circ Q,Q'\circ \Phi)\text{ ).}
\end{align*}

\begin{definition}\label{VIII:4-1-2}
We say that $\Phi$ has a right (respectively left) \emph{total derived functor} if the functor $\%$ (respectively $\%'$) is representable. An object representing $\%$ (respectively $\%'$) is called the right (respectively left) total derived functor of $\Phi$ and is denoted by $\eR\Phi$ (respectively $\eL\Phi$).\footnote{A more manageable definition, equivalent in practice, is given in \cite[XVII 1.2]{sga4}.}
\end{definition}

\subsubsection{Notation}\label{VIII:4-1-3}

Let $\cA$ and $\cB$ be two abelian categories and $f:\cA\to\cB$ an additive functor. The right \emph{total derived functor} of $\eK^+(f):\eK^+(\cA)\to\eK^+(\cB)$ is denoted by $\eR^+f$. The functors $\eR^-f$, $\eR f$, $\eL^+f$, $\eL^-f$, and $\eL f$ are defined similarly.

Likewise, let $F:\eK^\ast(\cA)\to\cB$ be a cohomological functor. Composition with the canonical quotient functor $Q:\eK^\ast(\cA)\to\eD^\ast(\cA)$ defines a functor
\[\xymatrix{
  \mathsf{Foco}(\eD^\ast(\cA),\cB) \ar[r]
    & \mathsf{Foco}(\eK^\ast(\cA),\cB)
}\]
(where $\mathsf{Foco}(-,-)$ denotes the category of cohomological functors). Hence there is a functor $\%$ (respectively $\%'$) $\mathsf{Foco}(\eD^\ast(\cA),\cB)\to\mathsf{Ab}$ given by $G\mapsto\hom(F,G\circ Q)$ (respectively $G\mapsto\hom(G\circ Q,F)$).

\begin{definition}\label{VIII:4-1-4}
We say that $F$ has a right (respectively left) \emph{cohomological derived functor} if $\%$ (respectively $\%'$) is representable. An object representing $\%$ (respectively $\%'$) is called the right (respectively left) cohomological derived functor of $F$ and is denoted by $\eR F$ (respectively $\eL F$).
\end{definition}

\subsubsection{Notation}\label{VIII:4-1-5}

Let $\cA$ and $\cB$ be two abelian categories and $f:\cA\to\cB$ an additive functor. The right cohomological derived functor of
\[
  \h^0{}\circ \eK^+ f:\eK^+(\cA) \to \cB
\]
is denoted by $\eR^+f$. When no confusion can result, the functor $\eR^+f\circ T^i$ (where $T$ is the translation functor) is denoted by $\eR^if$. The right cohomological derived functor of $\h^0\circ\eK f$ is denoted by $\eR f$. When no confusion can result, $\eR f\circ T^i$ is denoted by $\eR^if$.

The functors $\eL^-f$, $\eL f$, and $\eL^if$ are defined similarly.

\subsubsection{Remarks}\label{VIII:4-1-6}

As we shall see in the next subsection, these definitions include the classical definition of derived functors when the categories have enough injectives or projectives and the complexes are suitably bounded. They are more general, but without further hypotheses they are difficult to use. For example, let $f:\cA\to\cB$ be an additive functor between abelian categories. Suppose that the total derived functors $\eR^+f$ and $\eR f$ exist. I do not know how to prove that $\eR f$ induces $\eR^+f$ on $\eD^+(\cA)$. Suppose in addition that the cohomological derived functor $\eR^+f$ exists. I do not know how to prove that $\h^0\circ\eR^+f$ is isomorphic to that $\eR^+f$. In practice, however, these properties will hold.

\subsection{Existence of derived functors}\label{VIII:4-2}

\begin{proposition}\label{VIII:4-2-1}
Retain the hypotheses of Definition \ref{VIII:4-1-4}. Suppose in addition that $\cA$ is a $\sU$-category, where $\sU$ is a universe containing the set $\dZ$ of rational integers, and that the set of objects of $\cA$ belongs to $\sU$. Suppose also that $\cB$ is the category of abelian $\sU$-groups or, more generally, the category of sheaves of abelian $\sU$-groups on an arbitrary $\sU$-site. (A $\sU$-site is a $\sU$-category whose set of objects belongs to $\sU$, equipped with a topology.) Then $F$ has a right cohomological derived functor, computed as follows.

Let $X$ be an object of $\eD^\ast(\cA)$ and $\underline X$ the object of $\eK^\ast(\cA)$ lying over $X$. There is a functorial isomorphism
\[
  \eR F(X) = \varinjlim_{\underline X/\mathsf{Qis}^\ast} F(-) \text{.}
\]
\end{proposition}

In particular, let $f:\cA\to\cB$ be an additive functor. The functors $\eR f$ and $\eR^+f$ exist and are computed as above. The restriction of $\eR f$ to $\eD^+(\cA)$ is $\eR^+f$. Since no confusion can then result, we shall use the notation $\eR^if$.

\subsubsection{Remarks}\label{VIII:4-2-2}

Proposition \ref{VIII:4-2-1} can be stated more generally for triangulated categories. The hypotheses needed on $\cB$ are the existence and exactness of direct limits indexed by pseudofiltered $\sU$-categories whose sets of objects belong to $\sU$. At least in practice, Proposition \ref{VIII:4-2-1} introduces an asymmetry between right and left derived functors.

\begin{theorem}\label{VIII:4-2-3} % labelled 4.2.2 in original
Retain the data of Definition \ref{VIII:4-1-2}. Let $\cC$ also be a full triangulated subcategory of $\eK^\ast(\cA)$, and let $\cN$ be the full triangulated subcategory of acyclic objects of $\cC$. Suppose that
\begin{enumerate}[(1)]
  \item $\Phi$ takes every object of $\cN$ to an acyclic object of $\eK^{\ast'}(\cB)$.
  \item Every object $X$ of $\eK^\ast(\cA)$ is the source (respectively target) of a quasi-isomorphism whose target (respectively source) is an object of $\cC$.
\end{enumerate}
\begin{enumerate}[a)]
  \item Then $\Phi$ has a right (respectively left) total derived functor.
  \item The functor $\eR\Phi$ (respectively $\eL\Phi$) is obtained as follows.

    The restriction to $\cC$ of $Q'\circ\Phi:\eK^\ast(\cA)\to\eD^{\ast'}(\cB)$ vanishes on the objects of $\cN$, and therefore factors through $\cC/\cN$. This gives a functor $\Phi':\cC/\cN\to\eD^{\ast'}(\cB)$. The natural functor $\cC/\cN\to\eD^\ast(\cA)$ is an equivalence of categories (Theorem \ref{VIII:2-4-2}); composing $\Phi'$ with its quasi-inverse gives $\eR\Phi$ (respectively $\eL\Phi$).
  \item Let $X$ be an object of $\eK^\ast(\cA)$, let $Y$ be an object of $\cC$, and let $X\to Y$ (respectively $Y\to X$) be a quasi-isomorphism. Then $\eR\Phi\circ Q(X)$ (respectively $\eL\Phi\circ Q(X)$) is isomorphic to $Q'\circ\Phi(Y)$.
  \item The functor $X\mapsto\eR\Phi\circ Q(X)$ (respectively $X\mapsto\eL\Phi\circ Q(X)$) is isomorphic to $X\mapsto\varinjlim_{X/\mathsf{Qis}^\ast}Q'\circ\Phi(-)$ (respectively $X\mapsto\varinjlim_{\mathsf{Qis}^\ast/X}Q'\circ\Phi(-)$).
  \item The functor $\h^0\circ\Phi:\eK^\ast(\cA)\to\cB$ has a right (respectively left) cohomological derived functor, namely $\h^0\circ\eR\Phi$ (respectively $\h^0\circ\eL\Phi$).
\end{enumerate}
\end{theorem}

\begin{corollary}\label{VIII:4-2-3-1} % originally corollary 1
In the hypotheses of the preceding theorem, suppose that $\ast=+$ (respectively $\ast=-$) and that $\cA$ has enough injectives (respectively projectives). Then Theorem \ref{VIII:4-2-3} applies with $\cC$ the category of complexes of injectives (respectively projectives). If, in addition, $\cA$ has finite cohomological dimension, the theorem applies without restriction on $\ast$, again taking $\cC$ to be the category of complexes of injectives (respectively projectives).
\end{corollary}

\begin{definition}\label{VIII:4-2-5} % originally 2.4
Let $\cA$ and $\cB$ be two abelian categories and $f:\cA\to\cB$ an additive functor.

We say that $f$ has finite right (respectively left) \emph{cohomological dimension} if $\eR^+f$ (respectively $\eL^-f$) exists and there is an integer $m\geqslant0$ such that, for every $n>m$, $\eR^nf$ (respectively $\eL^{-n}f$) vanishes on the objects coming from $\cA$ under $D$. The right (respectively left) \emph{cohomological dimension} of $f$ is the least such integer $m$. If $\eR^+f$ (respectively $\eL^-f$) exists but $f$ does not have finite cohomological dimension, then $f$ is said to have infinite cohomological dimension.

An object $X$ of $\cA$ is called right (respectively left) \emph{$f$-acyclic} if, for every integer $n>0$, $\eR^nf(D(X))=0$ (respectively $\eL^{-n}f(D(X))=0$).
\end{definition}

\begin{corollary}\label{VIII:4-2-3-2} % originally corollary 2
Retain the data of Definition \ref{VIII:4-2-5}. Suppose that there is a collection $\sigma$ of objects of $\cA$, stable under direct sums, with the following properties:
\begin{enumerate}
  \item Every acyclic complex in $\operatorname{Ob}\eK^+(\cA)$ (respectively in $\operatorname{Ob}\eK^-(\cA)$) whose objects belong to $\sigma$ is taken by $f$ to an acyclic complex.
  \item Every object $X$ of $\cA$ embeds into (respectively is a quotient of) an object of $\sigma$. The hypotheses of Theorem \ref{VIII:4-2-3} hold with $\ast=+$ (respectively $\ast=-$) and with $\cC$ the category of complexes whose objects in every degree belong to $\sigma$. The elements of $\sigma$ are right (respectively left) $f$-acyclic objects.
\end{enumerate}

If, in addition, $f$ has finite right (respectively left) cohomological dimension, the collection of right (respectively left) $f$-acyclic objects has, in addition to properties (1) and (2) above, property (2) of Proposition \ref{VIII:3-1-9}. Theorem \ref{VIII:4-2-3} then applies to $\eK f$ with $\cC$ the category of complexes whose objects in every degree are right (respectively left) $f$-acyclic. The functor $\eR f$ (respectively $\eL f$) exists and induces $\eR^+f$ (respectively $\eL^+f$) on $\eD^+(\cA)$ and $\eR^-f$ (respectively $\eL^-f$) on $\eD^-(\cA)$.
\end{corollary}

\subsection{Product. Composition}\label{VIII:4-3}

Let $\cA$ be an abelian category, let $F:\eD(\cA)\to\mathsf{Ab}$ be a cohomological functor, and let $X^\bullet$ and $Y^\bullet$ be two objects of $\eC(\cA)$. The functor $F$ defines a map
\[
  \ext^0(X^\bullet,Y^\bullet) \to \hom_{\mathsf{Ab}}(F(D(X^\bullet)),F(D(Y^\bullet)))
\]
and hence a pairing $F(D(X^\bullet))\times \ext^0(X^\bullet,Y^\bullet) \to F(D(Y^\bullet))$.

Applying translations gives pairings
\[
  F^i(D(X^\bullet))\times \ext^j(X^\bullet,Y^\bullet) \to F^{i+j}(D(Y^\bullet)) \text{.}
\]
For cohomological derived functors, these pairings are, by definition, the Yoneda products. Their properties follow immediately from this definition, and we shall not develop them here.

Let $\cA$, $\cA'$, and $\cA''$ be three abelian categories, and let $f:\cA\to\cA'$ and $g:\cA'\to\cA''$ be additive functors. Suppose that $\eR^+f$ and $\eR^+g$ exist. I do not know how to deduce that $\eR^+(g\circ f)$ exists; and even if it does, the formula
\[\xymatrix{
  \eR^+(g\circ f)\ar[r]^-\sim
    & \eR^+ g \circ \eR^+ f
}\]
is probably not true in general. We do, however, have the following proposition.

\begin{proposition}\label{VIII:4-3-1}
Suppose that $\cA'$ has enough right $g$-acyclic objects and that $\cA$ has enough right $f$-acyclic objects which $f$ takes to right $g$-acyclic objects. Then $\eR^+(g\circ f)$ exists and there is an isomorphism
\[\xymatrix{
  \eR^+(g\circ f) \ar[r]^-\sim
    & \eR^+ g \circ \eR^+ f \text{.}
}\]
\end{proposition}

There are analogous statements for left derived functors.

\section{Examples}\label{VIII:6} % originally chapter 2, section 3

\subsection{The functor \texorpdfstring{$\mathsf{Rhom}^\bullet$}{Hom*}}\label{VIII:6-1}

Let $\cA$ be an abelian category with enough injectives.

For $X^\bullet\in\eK(\cA)$ and $Y^\bullet\in\eK^+(\cA)$, the functor $Y^\bullet\mapsto\hom^\bullet(X^\bullet,Y^\bullet)$ (\ref{VIII:1-3}) has a right total derived functor (Corollary \ref{VIII:4-2-3-1}), denoted by $\rhom^\bullet(X^\bullet,-)$. For an object $Y$ of $\eD^+(\cA)$, the functor
\begin{align*}
  \eK(\cA) &\to \eD(\mathsf{Ab}) \\
  X^\bullet &\mapsto \rhom^\bullet(X^\bullet,Y)
\end{align*}
is exact and vanishes on acyclic complexes, and hence gives an exact bifunctor
\[\xymatrix{
  \eD(\cA)^\circ \times \eD^+(\cA) \ar[r]
    & \eD(\mathsf{Ab})
}\]
which we denote by $\mathsf{Rhom}^\bullet$.

When $\cA$ has finite cohomological dimension, $\mathsf{Rhom}^\bullet$ extends to $\eD(\cA)^\circ\times\eD(\cA)$ (loc.\ cit.).

If $\cA$ also has enough projectives, one can derive $\hom$ in its first argument. The two definitions agree on their common domain of validity.

\subsection{\texorpdfstring{$\Tor$}{Tor} of sheaves}\label{VIII:6-2}

Let $E$ be a ringed site (in particular, a ringed topological space), let $\sA$ be its sheaf of rings, and let $\mathsf{Mod}(E)$ be the category of $\sA$-modules. An object $X$ of $\mathsf{Mod}(E)$ is \emph{$\sA$-flat} if, for every exact sequence
\[\xymatrix{
  0 \ar[r]
    & Y' \ar[r]
    & Y \ar[r]
    & Y'' \ar[r]
    & 0
}\]
the sequence $0\to Y'\otimes_\sA X\to Y\otimes_\sA X\to Y''\otimes_\sA X\to0$ is exact. There are enough $\sA$-flat objects: direct sums of objects $\sA_U$, where $U$ is an object of $E$ and $\sA_U$ is the sheaf $\sA$ \emph{extended by zero outside $U$}. Thus Theorem \ref{VIII:4-2-3} applies and defines the functor
\begin{align*}
  \eD^-(E,\sA) \times \eD^-(E,\sA) &\to \eD^-(E,\sA) \\
  (X,Y) &\mapsto X\lotimes Y
\end{align*}
the total tensor product, which is the left derived functor of the tensor product. Passing to cohomology defines the local hyper-Tor sheaves $\Tor_i^\sA(X,Y)$.

Since $\mathsf{Mod}(E)$ has enough injectives, one can derive the functor $\hom^\bullet(X,Y)$, the complex of local homomorphisms, thereby obtaining $\mathsf{Rhom}(X,Y)$ for $X\in\eD$ and $Y\in\eD^+$.

For $X,Y\in\eD^-(E,\sA)$ and $Z\in\eD^+(E,\sA)$, there is a functorial isomorphism
\[\xymatrix{
  \hom_{\eD(E,\sA)}(X\lotimes Y,Z) \ar[r]^-\sim
    &\hom_{\eD(E,\sA)}(X,\mathsf{Rhom}(Y,Z)) \text{.}
}\]

\subsection{Left derived functor of inverse image}\label{VIII:6-3}

Let $f:E\to E'$ be a map of ringed topological spaces. The inverse-image functor
\[
  f^\ast:\mathsf{Mod}(E') \to \mathsf{Mod}(E)
\]
need not be exact. One can nevertheless define $f^\ast$-flat objects, and there are enough of them. Thus the functor $\eL^-f^\ast$ is defined. There is also the following adjunction formula.

For $X\in\eD^-(E',\sA')$ and $Y\in\eD^+(E,\sA)$, there is a functorial isomorphism
\[\xymatrix{
  \hom_{\eD(E,\sA)}(\eL^- f^\ast(X),Y) \ar[r]^-\sim
    & \hom_{\eD(E',\sA')}(X,\eR^+ f_\ast(Y)) \text{.}
}\]
% END INLINED FILE: sga4.5-8.tex
\clearpage
\clearpage
% BEGIN INLINED FILE: sga4.5-erratum.tex
\appendix
\chapter{Erratum for SGA 4, Volume 3}

\paragraph{XIV p.18 1.14}
Read (XIX 6) instead of (XX 6).

\paragraph{XVI 2.2}
One must assume that $F$ is locally constant!

\paragraph{XVI 5.2}
The proof given is incomplete. After stating the induction hypothesis, one must first reduce to the case in which $F$ is constant ($F$ becomes constant on an étale Galois covering $\pi:X'\to X$ of $X$, with Galois group $G$, and one invokes the Hochschild--Serre spectral sequence $E_2^{p q}=\h^p(G,\h^q(X',\pi^\ast F))\Rightarrow \h^{p+q}(X,F)$). The arguments that follow are then correct.

\paragraph{XVII 1.1.8}
The sign in (1.1.8.1) is incorrect when $F$ is contravariant in some variables. It should read:
\begin{equation*}\tag{1.1.8.1}
  \rho^{\underline k} = (-1)^{A(\underline k)}: \text{ automorphism of }
F\circ (G_j)\left(K_i^{k_i\varepsilon(i)\varepsilon(\psi(i))}\right)
\end{equation*}
where
\[
  A(\underline k) = \sum_{\varepsilon(i) = \varepsilon(\psi(i))=-} k_i + \sum_{\varepsilon(j)=-} \sum_{\substack{\psi(a)=\psi(b)=j \\ a<b}} k_a k_b
\]

\paragraph{XVII 2.1.3}
The proof contains obvious errors. Delete the third line (p.~34, l.~9) and replace the sixth and seventh lines (p.~34, ll.~12--13) by:

The arrows in diagram (2.1.3.2) induce maps
\[
  \hom_y(F,F) \to \hom_{f g'}({g'}^\ast F,f_\ast F) = \hom_{x'}(f'_\ast {g'}^\ast F,g^\ast f_\ast F) \text{,}
\]

\paragraph{XVIII 2.14.4}
Read XVII 6.2.7.2 instead of XVII 6.2.4.

\paragraph{XVIII p.99 1-1}
Read $u$ instead of $U$, and 3.1.16.1 instead of 3.1.11.1.
% END INLINED FILE: sga4.5-erratum.tex
\clearpage

% BEGIN INLINED BIBLIOGRAPHY
\providecommand{\bysame}{\leavevmode\hbox to3em{\hrulefill}\thinspace} \providecommand{\MR}{\relax\ifhmode\unskip\space\fi MR }
% \MRhref is called by the amsart/book/proc definition of \MR.
\providecommand{\MRhref}[2]{%
  \href{http://www.ams.org/mathscinet-getitem?mr=#1}{#2} } \providecommand{\href}[2]{#2}
\begin{thebibliography}{10}

\bibitem{ar62} Michael Artin, \emph{Grothendieck topologies}, Harvard University Dept. Math., 1962.

\bibitem{sga4} Michael Artin, Alexandre Grothendieck, and Jean-Louis Verdier, \emph{Th\'eorie des topos et cohomologie \'etale des sch\'emas ({SGA 4})}, Lecture Notes in Mathematics, vol. 269, 270, 305, Springer, 1972.

\bibitem{ba76} Hyman Bass, \emph{Euler characteristics and characters of discrete groups}, Invent. Math. \textbf{35} (1976), 155--196. \MR{0432781 (55 \#5764)}

\bibitem{bfm75} Paul Baum, William Fulton, and Robert MacPherson, \emph{Riemann-{R}och for singular varieties}, Inst. Hautes \'Etudes Sci. Publ. Math. (75), no.~45, 101--145.

\bibitem{sga6} Pierre Berthelot, Alexandre Grothendieck, and Luc Illusie, \emph{Th\'eorie des intersections et th\'eor\`eme de {R}iemann-{R}och ({SGA} 6)}, Lectures Notes in Mathematics, vol. 225, Springer-Verlag, 1971.

\bibitem{bo66} Enrico Bombieri, \emph{On exponential sums in finite fields}, Amer. J. Math. \textbf{88} (1966), 71--105.

\bibitem{bo05} Nicolas Bourbaki, \emph{Lie groups and {L}ie algebras. {C}hapters 7--9}, Elements of mathematics, Springer-Verlag, Berlin, 2005, Translated from the 1975 and 1982 French originals by Andrew Pressley.

\bibitem{ca69-2} L.~Carlitz, \emph{A note on exponential sums}, Pacific J. Math. \textbf{30} (1969), 35--37.

\bibitem{ca69} \bysame, \emph{Kloosterman sums and finite field extensions}, Acta Arith. \textbf{16} (1969/1970), 179--193.

\bibitem{dl76} P.~Deligne and G.~Lusztig, \emph{Representations of reductive groups over finite fields}, Ann. of Math. (2) \textbf{103} (1976), no.~1, 103--161.

\bibitem{de71-1} Pierre Deligne, \emph{Th\'eorie de {H}odge. {I}}, Actes du {C}ongr\`es {I}nternational des {M}ath\'ematiciens ({N}ice, 1970), {T}ome 1, Gauthier-Villars, Paris, 1971, pp.~425--430.

\bibitem{de71-2} \bysame, \emph{Th\'eorie de {H}odge. {II}}, Inst. Hautes \'Etudes Sci. Publ. Math. (1971), no.~40, 5--57.

\bibitem{de73} \bysame, \emph{Les constantes des \'equations fonctionnelles des fonctions {$L$}}, Modular functions of one variable, {II} ({P}roc. {I}nternat. {S}ummer {S}chool, {U}niv. {A}ntwerp, {A}ntwerp, 1972), Springer, Berlin, 1973, pp.~501--597. Lecture Notes in Math., Vol. 349.

\bibitem{de74} \bysame, \emph{La conjecture de {W}eil. {I}}, Inst. Hautes \'Etudes Sci. Publ. Math. (1974), no.~43, 273--307.

\bibitem{de80} \bysame, \emph{La conjecture de {W}eil. {II}}, Inst. Hautes \'Etudes Sci. Publ. Math. (1980), no.~52, 137--252.

\bibitem{sga7} Pierre Deligne, Alexandre Grothendieck, and Nicholas Katz, \emph{Groupes de monodromie en g\'eom\'etrie alg\'ebrique ({SGA 7})}, Lecture Notes in Mathematics, vol. 288, 340, Springer, 1972-73.

\bibitem{do61} Albrecht Dold and Dieter Puppe, \emph{Homologie nicht-additiver {F}unktoren. {A}nwendungen}, Ann. Inst. Fourier Grenoble \textbf{11} (1961), 201--312. \MR{0150183 (27 \#186)}

\bibitem{gz67} P.~Gabriel and M.~Zisman, \emph{Calculus of fractions and homotopy theory}, Ergebnisse der Mathematik und ihrer Grenzgebiete, vol.~35, Springer-Verlag, New York, 1967.

\bibitem{gi64} Jean Giraud, \emph{Analysis situs}, S\'eminaire {B}ourbaki, 1962/63. {F}asc. 3, {N}o. 256, Secr\'etariat math\'ematique, Paris, 1964, p.~11. \MR{0193122 (33 \#1343)}

\bibitem{ega} Alexandre Grothendieck, \emph{\'{E}l\'ements de g\'eom\'etrie alg\'ebrique}, Inst. Hautes \'Etudes Sci. Publ. Math., 1960-67.

\bibitem{sga5} \bysame, \emph{Cohomologie {$\ell$}-adique et fonctions {$L$} ({SGA 5})}, Lecture Notes in Mathematics, vol. 589, Springer, 1965-66.

\bibitem{gr95} \bysame, \emph{Formule de {L}efschetz et rationalit\'e des fonctions {$L$}}, S\'eminaire {B}ourbaki, {V}ol.\ 9, Soc. Math. France, Paris, 1995, pp.~Exp.\ No.\ 279, 41--55. \MR{1608788}

\bibitem{sga1} Alexandre Grothendieck and Michele Raynaud, \emph{Rev\^etements \'etales et groupe fondamentale ({SGA 1})}, Lecture Notes in Mathematics, vol. 224, Springer, 1971.

\bibitem{sga2} \bysame, \emph{Cohomologie locale des faisceaux coh\'erents et th\'eor\`emes de {L}efschetz locaux et globaux ({SGA 2})}, Documents Math\'ematiques (Paris), 4, Soci\'et\'e Math\'ematique de France, Paris, 2005, S{\'e}minaire de G{\'e}om{\'e}trie Alg{\'e}brique du Bois Marie, 1962, Augment{\'e} d'un expos{\'e} de Mich{\`e}le Raynaud, With a preface and edited by Yves Laszlo, Revised reprint of the 1968 French original.

\bibitem{ka77} D.~Kazhdan, \emph{Proof of {S}pringer's hypothesis}, Israel J. Math. \textbf{28} (1977), no.~4, 272--286.

\bibitem{la83} Serge Lang, \emph{Abelian varieties}, Springer-Verlag, New York, 1983, Reprint of the 1959 original. \MR{713430 (84g:14041)}

\bibitem{lu76} G.~Lusztig, \emph{Coxeter orbits and eigenspaces of {F}robenius}, Invent. Math. \textbf{38} (1976/77), no.~2, 101--159.

\bibitem{mo73} L.~J. Mordell, \emph{Some exponential sums}, Proceedings of the {I}nternational {C}onference on {N}umber {T}heory ({M}oscow, 1971), vol. 132, 1973, pp.~30--34, 264.

\bibitem{ra65} Michel Raynaud, \emph{Caract\'eristique d'{E}uler-{P}oincar\'e d'une faisceau et cohomologie des vari\'et\'es ab\'eliennes}, S\'eminaire {B}ourbaki. {V}ol.\ 9, vol. Exp.\ No.\ 286, Soc. Math. France, Paris, 1965, pp.~129--147.

\bibitem{ra70} \bysame, \emph{Anneaux locaux hens\'eliens}, Lecture Notes in Mathematics, Vol. 169, Springer-Verlag, Berlin, 1970. \MR{0277519 (43 \#3252)}

\bibitem{sc72} G\'erard Schiffmann, \emph{Th\'eorie \'el\'ementaire des cat\'egories}, S\'eminaire Banach (C.~Houzel, ed.), Lecture Notes in Mathematics, vol. 277, Springer, Berlin, 1972, pp.~1--33.

\bibitem{se55} Jean-Pierre Serre, \emph{G\'eom\'etrie alg\'ebrique et g\'eom\'etrie analytique}, Ann. Inst. Fourier, Grenoble \textbf{6} (1955--1956), 1--42. \MR{0082175 (18,511a)}

\bibitem{se59} \bysame, \emph{Groupes alg\'ebriques et corps de classes}, Publications de l'institut de math\'ematique de l'universit\'e de Nancago, VII. Hermann, Paris, 1959. \MR{0103191 (21 \#1973)}

\bibitem{se61} \bysame, \emph{Sur les corps locaux \`a corps r\'esiduel alg\'ebriquement clos}, Bull. Soc. Math. France \textbf{89} (1961), 105--154.

\bibitem{se68} \bysame, \emph{Abelian {$\ell$}-adic representations and elliptic curves}, McGill University lecture notes, W. A. Benjamin, Inc., New York-Amsterdam, 1968.

\bibitem{se94} \bysame, \emph{Cohomologie galoisienne}, fifth ed., Lecture Notes in Mathematics, vol.~5, Springer-Verlag, Berlin, 1994. \MR{1324577 (96b:12010)}

\bibitem{sm70} R.~A. Smith, \emph{The distribution of rational points on hypersurfaces defined over a finite field}, Mathematika \textbf{17} (1970), 328--332.

\bibitem{sp77} S.~Sperber, \emph{{$p$}-adic hypergeometric functions and their cohomology}, Duke Math. J. \textbf{44} (1977), no.~3, 535--589.

\bibitem{sp76} T.~A. Springer, \emph{Trigonometric sums, {G}reen functions of finite groups and representations of {W}eyl groups}, Invent. Math. \textbf{36} (1976), 173--207.

\bibitem{ve67} Jean-Louis Verdier, \emph{The {L}efschetz fixed point formula in etale cohomology}, Proc. {C}onf. {L}ocal {F}ields ({D}riebergen, 1966), Springer, Berlin, 1967, pp.~199--214. \MR{0242846 (39 \#4173)}

\bibitem{we48} Andr{\'e} Weil, \emph{Vari\'et\'es ab\'eliennes et courbes alg\'ebriques}, Actualit\'es Sci. Ind., no. 1064 = Publ. Inst. Math. Univ. Strasbourg 8 (1946), Hermann \& Cie., Paris, 1948. \MR{0029522 (10,621d)}

\bibitem{we52} \bysame, \emph{Jacobi sums as ``{G}r\"ossencharaktere''}, Trans. Amer. Math. Soc. \textbf{73} (1952), 487--495.

\bibitem{we74} \bysame, \emph{Basic number theory}, third ed., 1974, New York, 1974, Die Grundlehren der Mathematischen Wissenschaften, Band 144.

\bibitem{we74-2} \bysame, \emph{Sommes de {J}acobi et caract\`eres de {H}ecke}, Nachr. Akad. Wiss. G\"ottingen Math.-Phys. Kl. II (1974), no.~1, 1--14.

\end{thebibliography}
% END INLINED BIBLIOGRAPHY

\end{document}
